
%\documentclass[oneside,final,14pt]{extreport}
\documentclass[12pt, a4paper, openany]{book}



\usepackage[left=2cm,right=2cm,top=2cm,bottom=2cm,bindingoffset=0cm]{geometry}
%\usepackage[koi8-r]{inputenc}
%\usepackage[russianb]{babel}
\usepackage{vmargin}

\setpapersize{A4}
\usepackage[T2A]{fontenc}
\usepackage[utf8x]{inputenc}
\usepackage[english, russian]{babel}
\setmarginsrb{2cm}{2cm}{2cm}{2cm}{0pt}{5mm}{0pt}{0mm}
\usepackage{indentfirst}
\usepackage{nicefrac} % For comparison
%\usepackage{xfrac} % Works better with other fonts
%\usepackage[unicode, pdftex]{hyperref}
\usepackage{lettrine}
\usepackage[usenames]{color}
\usepackage{colortbl}
\usepackage{mathtext}
\usepackage{epigraph}
\usepackage{amsmath, amsfonts, amssymb, mathrsfs}
%\usepackage{mathptmx}
%\usepackage{txfonts}
\usepackage{pxfonts}
\usepackage[pagestyles]{titlesec}
\usepackage{ebgaramond}
\usepackage{awesomebox}
\usepackage{enumitem}
\usepackage{makeidx}
\usepackage[letterspace=150]{microtype}
\makeindex
\usepackage{etoolbox}
\makeatletter
\newlength\epitextskip
\pretocmd{\@epitext}{\em}{}{}
\apptocmd{\@epitext}{\em}{}{}
\patchcmd{\epigraph}{\@epitext{#1}\\}{\@epitext{#1}\\[\epitextskip]}{}{}
\makeatother
%running fraction with slash - requires math mode.
\newcommand*\rfrac[2]{{}^{#1}\!/_{#2}}

\DeclareSymbolFont{Xlargesymbols}{OMX}{cmex}{m}{n}
\DeclareMathSymbol{\Xsum}{\mathop}{Xlargesymbols}{80}

\setlength\epigraphrule{0pt}
\setlength\epitextskip{2ex}
\setlength\epigraphwidth{.8\textwidth}

\usepackage{xfrac} % Works better with other fonts
\usepackage[colorlinks=true,linkcolor=black,urlcolor=black,bookmarksopen=true]{hyperref}


\usepackage[
backend=biber,
style=alphabetic,
sorting=ynt
]{biblatex}

\addbibresource{references.bib} %Imports bibliography file

\usepackage{fancyhdr} % пакет для установки колонтитулов
\pagestyle{fancy} % смена стиля оформления страниц
\fancyhf{} % очистка текущих значений
\fancyhead[C]{\thepage} % установка верхнего колонтитула
\renewcommand{\headrulewidth}{0pt} % убрать разделительную линию


% Настройка вертикальных и горизонтальных отступов
\titlespacing{\chapter}{0pt}{5pt}{5pt}
\titlespacing{\section}{\parindent}{4mm}{4mm}
\titlespacing{\subsection}{\parindent}{3mm}{3mm}


%
%\renewcommand{\tabcolsep}{1cm} %% increase table column spacing

% Настройка задачи со зведочкой
\newcounter{namedlistcounter} % number the items
\newenvironment{withdot}
{\begin{list}
		{\arabic{namedlistcounter}*.} % labeling 
		{\usecounter{namedlistcounter} % set counter
			\setlength{\leftmargin}{3em}} % set spacing 
	}
	{\end{list}}


\newcommand{\anonsection}[1]{ \section*{#1} \addcontentsline{toc}{section}{\numberline {}#1}} 

\makeatletter %%%%% <---- Starting chapter without a pagebreak
\renewcommand\chapter{\par%
	\thispagestyle{plain}% \global\@topnum\z@
	\@afterindentfalse \secdef\@chapter\@schapter}
\makeatother %%%%% <---- Starting chapter without a pagebreak
\titleformat{\chapter}[display]
{\normalfont\bfseries}{}{0pt}{\Large}

\newpagestyle{mystyle}{
	\sethead[\thepage][][]{}{}{\thepage}	
}

\renewcommand{\rmdefault}{cmr}


\pagestyle{mystyle}
\sloppy
\begin{document}

	
	\begin{titlepage}
		
		\begin{center}
			\topskip0pt
			\vspace*{\fill}
			
			
			{\large\bf А. И. МИКОЯН\\}
			\ \\
			\ \\
			{\Huge\bf Мысли и воспоминания о ЛЕНИНЕ\\}
			\ \\
			\ \\
			\ \\
			\vspace*{\fill}
			
			\vfill
			
			
Издательство политической литературы
			
\textls{Москва} $\cdot$ \textls{1970}
			

		\end{center}
		
	\end{titlepage}
	
	\thispagestyle{empty} % выключаем отображение номера для этой страницы
	
	\newpage
	

\setcounter{secnumdepth}{0}
	
	\section[О ленинском предвидении (Вместо вступления)]{\center \textbf{О ленинском предвидении} \\ {\normalsize (Вместо вступления)}}


О Владимире Ильиче Ленине написано много исследований, монографий, статей, воспоминаний и, несомненно, еще многое будет написано.

Для меня встречи с Лениным — это немногочисленные, но запомнившиеся навсегда беседы с ним, это живое слово Ильича на партийных съездах и конференциях, на съездах Советов, на конгрессах Коминтерна...

Впервые о Ленине я узнал, будучи учащимся средней школы, когда стал читать его книги: «Развитие капитализма в России», «Что делать?», «Шаг вперед, два шага назад» и др. Поэтому встречи с Лениным — это также встречи с его книгами и статьями — бесценным богатством, бездонным кладезем мудрости партии и народа, ставшими теперь интернациональным достоянием. Это более чем полувековые размышления над его мудрыми мыслями, его изумительной прозорливостью. Это, наконец, многотрудные дела его как главы первого социалистического государства и вождя партии, проявлявшиеся в огромном многообразии свершений молодой республики Советов.

Велик и многогранен Ленин, неисчерпаемо его влияние на жизнь современников, на судьбы грядущих поколений. Рассказать о Ленине все — невозможно, посильна лишь задача поделиться личным восприятием, оттенить некоторые ленинские грани, напомнить об отдельных моментах, которые представляются важными и сегодня, выделить то, что хочется воскресить в памяти, рассказать о переживаниях и впечатлениях людей старшего поколения — современников В. И. Ленина.

\subsection*{\center * * *}

Что больше всего поражает и восхищает в Ленине — революционере, ученом и государственном деятеле?

Мне представляется, что в первую очередь это, пожалуй, дар предвидения. Предвидения того, как события общественной жизни будут развиваться, в каком направлении они пойдут, как будут эволюционировать различные классы, а с ними и партии и т. д.

Это был не просто дар гениального человека. Ленинское предвидение являлось результатом титанического, самоотверженного труда. Оно было качеством одновременно и революционера и ученого, который с позиций передового класса анализирует происходящие явления общественной жизни и умеет делать из них правильные заключения и выводы как стратегического, так и тактического порядка.

В связи с этим обращают на себя внимание замечательные слова самого Ленина, написанные им в декабре 1902 года: «Пусть нам не говорят, что мы преувеличиваем серьезность момента. Кто способен заглядывать дальше поверхности ряби, кто способен распознавать совершающийся в глубине процесс, тот не заподозрит нас в преувеличении»\footnote{\textit{В. И. Ленин}. Полн. собр. соч., т. 7, стр. 90.}.

Причин того, почему прогнозы марксистов, большевиков оказывались реальными и оправдывались, а расчеты представителей буржуазии проваливались, Ленин коснулся на IX Всероссийском съезде Советов, на котором мне довелось присутствовать в качестве делегата от Нижнего Новгорода.

В отчете ВЦИК и Совнаркома о внутренней политике республики, сделанном 23 декабря 1921 года, он, в частности, говорил:

«...наш расчет, в большом масштабе взятый, оказывается более правильным, чем их расчет. И не потому, что у них нет людей, которые умеют правильно рассчитывать, — наоборот, у них их больше, чем у нас, — а потому, что нельзя рассчитывать правильно, когда стоишь на пути к гибели»\footnote{\textit{В. И. Ленин}. Полн. собр. соч., т. 44, стр. 303.}.

И наоборот, отсюда следует, что возможность рассчитывать правильно, далеко вперед появляется, когда представляешь класс восходящий, класс, которому историей отведена роль возглавить человечество на путях общественного прогресса. Ленин, как никто другой, мастерски пользовался такой объективной возможностью.

При этом Ленин неоднократно предупреждал о пагубности недооценки возможностей классовых врагов. Он подчеркивал, что у буржуазии есть люди, которые умеют правильно рассчитывать. Но все дело в том, что идеологи капитализма не в состоянии признать за истину единственно реальную историческую перспективу — неизбежную гибель капитализма как общественной системы. В апологетике капитализма корень их фундаментальных просчетов. Однако Ленин вовсе не исключал, что в меньшем масштабе, в тактическом плане представители идущего к закату класса в тот или другой момент сделают правильные расчеты, предугадают расстановку сил и непосредственный ход событий и смогут нанести серьезные удары по тем или иным отрядам революционного движения.

Конечно, наряду с этим можно привести многочисленные примеры, когда и умудренным опытом деятелям враждебного стана не удавалось сделать правильных тактических выводов, рассчитать реальное соотношение сил. За этим, как правило, кроется недооценка сил противника, неумение трезво оценить скрытые потенции, великую революционную силу масс, поднимающихся на борьбу за свою свободу и независимость.

Нельзя не вспомнить оценок, делавшихся на Западе в трудную годину гражданской войны. Безнадежным по всем «объективным данным» им представлялось положение Петрограда осенью 1919 года. Город был почти окружен, войска Юденича заняли ближайшие к нему подступы; на суше, на Балтийском море и на озерах действовали войска и флот иностранных интервентов. В такой критической, казалось бы, безвыходной ситуации Ленин и ЦК партии, петроградский пролетариат приняли героические меры по спасению города. В дни, когда решающие бои шли и на деникинском фронте, в Петроград на это время были брошены буквально последние резервы главного командования. И враг был разбит, отброшен от славного Питера.

Горькими словами подвел итоги провалившихся «прогнозов» о неизбежности падения Петрограда тогдашний премьер-министр Франции Ж. Клемансо: «Случилось нечто такое, чего решительно никто не мог предвидеть. Россия, доведенная до крайних пределов разрухи, вся изголодавшимся, теснимая со всех сторон изнутри и извне, эта Россия вдруг точно прикоснулась к какому-то источнику живой воды и ощутила в себе такую мощь, что отразила и последний удар Юденича и Родзянко — удар, как всем казалось, смертельный». Опьяненные победой над Германией, Клемансо и иже с ним оказались не в состоянии правильно рассчитать возможности сторон, ход и исход событий даже на таком коротком отрезке истории.

Еще большую слепоту проявил за полтора года до этого посол германского кайзера в Москве граф Мирбах. 25 июня 1918 года, давая оценку положения в Советской России, он писал и Берлин: «Мы, несомненно, находимся у постели тяжелобольного, который случайно может еще показать видимость улучшения, но кончина которого предрешена».

Мирбаху не повезло. Он был лишен, в отличие от Клемансо, возможности лично убедиться в провале своего «исторического пророчества» в силу того, что через несколько дней после этого письма был сражен пулей террориста левого эсера Блюмкина.

Чем же Ленин достигал высшей степени точности своих расчетов? Прежде всего, он был вооружен идеями Маркса и Энгельса, отлично владел методом материалистической диалектики, который он довел до совершенства, применяя его к условиям современности Этот метод — лучшее оружие распознании сущности общественных явлений и тенденций, определяющих развитие жизни. «Чудесное пророчество есть сказка, говорил Ленин. — Но научное пророчество есть факт»\footnote{\textit{В. И. Ленин}. Полн. собр. соч., т. 30, стр. 472.}.

Изучая общественные процессы, Ленин озарил практику светом научной теории, одновременно проверяя и совершенствуя теорию на уроках опыта Его выводы обогащали теорию. Вот почему события, как правило, развивались «по Ленину» так, как он их предвидел.

К Ленину с полным основанием можно отнести оценку, данную им самим Энгельсу. В 1918 году в «Правде» в специальной статье под заголовком «Пророческие слова» Ленин рассказал о том, как Энгельс, еще 30 лет назад предсказал, что в Европе дело идет к мировой войне, что начнет эту войну Германия и участвовать в ней будут многие миллионы солдат. Война эта, писал Энгельс, будет продолжаться 3–4 года. Она принесет с собой разорение и всеобщее истощение, расстройство торговли и промышленности, «голод, эпидемии, всеобщее одичание как войск, так и народных масс, вызванное острой нуждой...».

Энгельс предвидел такой крах старых государств, когда короны дюжинами валяются по мостовым...». Результатом войны, несомненно, будут «всеобщее истощение и создание условий для окончательной победы рабочего класса».

Ленин восклицал: «Какое гениальное пророчество!» И далее: «Кое-что из того, что предсказал Энгельс, вышло иначе: еще бы не измениться миру и капитализму за тридцать лет бешено быстрого империалистского развития. Но удивительнее всего, что столь многое, предсказанное Энгельсом, идет, «как по писаному». Ибо Энгельс давал безупречно точный классовый анализ...»\footnote{\textit{В. И. Ленин}. Полн. собр. соч., т. 36, стр. 473.}.

Еще Маркс и Энгельс писали, что центр мирового революционного движения постепенно перемещается из Западной Европы в Россию, что скорее всего «толчок» придет из России и что «это будет ближайшим поворотным пунктом во всемирной истории». Такую точку зрения разделяли тогда и западноевропейские социал-демократы. Так, Каутский, который в то время был еще марксистом, в 1902 году прислал в ленинскую «Искру» статью, в которой писал: «Новое столетие начинается такими событиями, которые наводят на мысль, что мы идем навстречу дальнейшему передвижению революционного центра, именно: передвижению его в Россию».

Вот такой же безупречно точный классовый анализ событий своего времени давал Ленин. Именно Ленин, и только Ленин, был тем человеком, который лучше других отразил нашу эпоху — ее главные закономерности и основные направления развития, вскрыл ее противоречия и приоткрыл перспективы, подметил необычайное ускорение самого исторического процесса.

Ленин научно доказал неотвратимое приближение первой русской революции 1905 года и предвидел огромное международное значение, которое она будет иметь.

Уже тогда Ленин проявил замечательное умение сосредоточиться на главных вопросах, выбрать центральное звено, ухватившись за которое, как он любил говорить, можно вытащить всю цепь. Таким звеном на первых порах была общерусская партийная газета, которая могла бы помочь осуществлению главной цели — созданию партии нового типа. В 1902 году в книге «Что делать?» Ленин писал: «Дайте нам организацию революционеров — и мы перевернем Россию!». Ленин никогда не пользовался пустыми, пышными фразами. Он делал продуманные, научные выводы. И в этом случае он не только сказал: «Дайте нам организацию революционеров», — а вместе со своими соратниками выпестовал такую организацию — партию нового типа, партию коммунистов.

Почему наша партия стала партией нового типа? Хотя бы потому, что она не только требовала от каждого ее члена полной убежденности в коммунизме, принципиальности, беспредельной преданности, но и обязывала его работать на дело революции, ведя за партией массы трудящихся, бороться, не жалея жизни.

И вот через 15 лет, в октябре 1917 года, сбылось ленинское предвидение: Россия была перевернута, колесо истории человечества повернулось в новом направлении — к социализму.

Ленин-революционер был одновременно и выдающимся мыслителем, крупнейшим ученым. Но не во имя науки вообще. Он рассматривал науку, теорию прежде всего как могучий рычаг революционной борьбы, как маяк, освещающий ей путь. Например, до 1908 года Ленин специальных философских трудов не писал. После поражения революции 1905–1907 годов наступил период реакции. Именно в этих условиях некоторые видные деятели Большевистской партии, такие, как Богданов и Луначарский, не только образовали оппозиционную фракцию в партии, но и стали на идеалистические позиции в вопросах философии, ударились даже в богоискательство. Ленин увидел, что идеалистическая философия и махистское мировоззрение, проникая в ряды партии, могут нанести ей тяжелый ущерб. Он убедился, что именно эта область знания превратилась и важнейшее звено внутрипартийной борьбы и в России, и на международной арене. Еще активнее, чем прежде, взялся Ленин за изучение философских трудов и новейших научных открытий в области естествознания. Написанная им и этот период гениальная книга «Материализм и эмпириокритицизм» и другие его работы явились не только защитой марксизма, материализма от новейших форм идеализма, но и новым словом в философии марксизма. По своей сущности они были обращены в будущее, указывали пути выхода естественных наук из кризиса, который те тогда переживали. Ленин разбил в пух и прах получившие в то время широкое хождение идеалистические и метафизические концепции и спас партию от грозившего ей с этой стороны удара.

Тут нельзя не учитывать и другую сторону вопроса. Каутский, как известно, утверждал, что философия — это частное дело члена партии и партия не должна вмешиваться в его философские взгляды. Ленин же не допускал и мысли об отделении философии от политики.

\subsection*{\center * * *}

Знаменательно, что еще в 1905 году, в самый разгар революции, Ленин разработал теорию перерастания буржуазно-демократической революции под гегемонией пролетариата в революцию социалистическую.

В сентябре 1905 года в статье «Отношение социал-демократии к крестьянскому движению», опубликованной в газете «Пролетарий», Ленин писал: «...от революции демократической мы сейчас же начнем переходить и как раз в меру нашей силы, силы сознательного и организованного пролетариата, начнем переходить к социалистической революции. Мы стоим за непрерывную революцию. Мы не остановимся на полпути»\footnote{\textit{В. И. Ленин}. Полн. собр. соч., т. 11, стр. 222.}.

В 1915 году в небольшой статье «О лозунге Соединенных Штатов Европы» Ленин высказал гениальную идею, которая как молния озарила сознание революционеров. Победа социализма, утверждал Ленин, возможна первоначально в небольшой группе стран, может быть, даже только в одной, отдельно взятой стране. Это было откровением, новым словом в марксизме и явилось, как известно, одним из величайших предвидений Ленина.

Это вовсе не было случайной идеей, это было результатом огромной, плодотворной теоретической деятельности Ленина над проблемами эпохи империализма, гениальным ее итогом. Открыв закон небывалого усиления неравномерности политического и экономического развития капиталистических стран в эпоху империализма, Ленин пришел к смелому выводу о возможности победы социализма вначале даже в одной стране.

Не случайно и то, что к такому выводу Ленин пришел в разгар первой мировой империалистической войны. Когда разразилась война, социалисты европейских стран, заседавшие в парламентах, голосовали за военные кредиты, став тем самым на сторону своей национальной буржуазии и прекратившись в социал-шовинистов. И только большевики, депутаты царской государственной думы, руководствуясь указаниями Ленина, выступили против войны. Вскоре депутаты-большевики были арестованы и сосланы в Сибирь.

Ленин сразу же выступил против войны. Он выдвинул лозунг: превратить войну империалистическую в войну гражданскую, войну наций против наций превратить в войну класса против класса, в войну рабочего класса, за которым несомненно пойдет и трудовое крестьянство, против буржуазии собственной страны.

Если бы в каждой европейской стране или хотя бы в ряде решающих стран социалистические партии подняли тогда свой рабочий класс против собственной буржуазии, то ясно, какие социальные и международные последствия проистекли бы из такого единственно правильного, марксистского, революционного поведения пролетарских партий в империалистическую войну.

Призывая к гражданской войне, Ленин заглядывал вперед: к чему такая революционная, классовая война может привести, какие цели должны поставить перед собой коммунисты в такой войне.

Это, по-видимому, и было главным стимулом к постановке Лениным кардинального вопроса — о соединении борьбы рабочего класса с национально-освободительной борьбой. Тем самым был поставлен вопрос о взаимной помощи двух революционных сил: пролетариат помогает народам колоний в борьбе за независимость и получает помощь от народов, поднявшихся против колонизаторов.

В результате глубокого анализа Ленин отчетливо сформулировал мысль о том, что конечной целью гражданской войны станет социалистическая революция в развитых капиталистических странах, которая послужит толчком к усилению освободительной революции в угнетенных странах.

В то время, когда Ленин писал, что социалистическая революция возможна в одной стране, он, судя по всему, имел в виду Россию, хотя и не называл ее. Имении Россия оказалась ближе всего к осуществлению этой возможности. здесь завязался главный узел всех политических и экономических противоречий, сюда переместился с Запада центр революционного движения. Россию Ленин знал как никто другой. Находясь в Швейцарии, Ленин хотя и сетовал на то, что трудно из этого проклятого далека разобраться в происходящем в России, но он чувствовал пульс ее общественной жизни, видел развитие противоречий, раздиравших русское общество, яснее, чем товарищи, находившиеся на родине, многие из которых томились в тюрьмах и ссылке. Это тот редкий случай, когда издали действительно оказалось виднее. Это — свидетельство гениальности Ленина.

Когда в феврале 1917 года самодержавие было свергнуто, в сознании Ленина уже была готова мысль о том, что именно Россия должна стать первой страной социализма. Об этом в марте 1917 года, до приезда в Россию, Владимир Ильич писал в знаменитых четырех «Письмах из далека» (сохранилось и начало пятого, не дописанного им письма). В них проводилась мысль, что в России победила буржуазно-демократическая революция, но это только первый этап, а сейчас приходит черед второго этапа — этапа социалистической революции. Однако из четырех писем, посланных Лениным в Петроград для помещения в центральном партийном печатном органе, было опубликовано только одно, первое письмо, да и то с сокращениями и изменениями некоторых важных мест. Последующие письма вообще не были тогда опубликованы.

Как в центре, так и на местах в партийных кругах не было еще понимания того, что после Февральской революции возможно сразу поставить в повестку дня переход к революции социалистической. Хотя в сознании старых большевистских кадров жили ленинские идеи 1905 года о перерастании революции и гегемонии пролетариата и буржуазно демократической революции, довольно широко было распространено и такое представление, что после победы буржуазно-демократической революции должен пройти какой-то период для выполнения общедемократических задач, для созревания условий социалистической революции. В этом легко убедиться, перелистав большевистские газеты того времени — до апреля 1917 года. Я был тогда молодым коммунистом, работал на Кавказе и помню такого рода умонастроения партийных товарищей.

В этом таилась большая опасность. Нельзя забывать, что суть меньшевистского понимания хода революционного процесса заключалась в том. что они видели одну только экономическую и культурную отсталость России и утверждали, будто между буржуазно-демократической и социалистической революциями должна пролегать длительная полоса капиталистического развития.

Такая опасность непосредственно угрожала революционной практике. Недопонимание задач, вставших перед партией после свержения самодержавия, питало тенденцию к сближению с меньшевиками, в ряде городов имело место даже фактическое объединение с ними. И это продолжалось до тех пор, пока Ленин, вернувшись из эмиграции, не положил конец шатаниям в партии в области теории и практики. Ленин сразу поставил вопрос о полном размежевании с меньшевиками, с тем чтобы не оставить ничего связывающего с ними, вплоть до немедленного переименования партии — из социал-демократической в коммунистическую.

Ленин был великим стратегом классовых битв пролетариата, умевшим сочетать необыкновенную смелость мысли с ясностью лозунгов. Достаточно сказать о его гениальных Апрельских тезисах, которые дали партии курс на социалистическую революцию и России.

Революционная сознательность и классовое чутье петроградских большевиком ярко проявились в том, что петроградская организация в подавляющем большинстве сразу приняла его тезисы и взяла их на вооружение. Этот факт оказал огромное влияние на партийные организации других городов. Партии удалось быстро и относительно безболезненно воспринять новые ленинские идеи о задачах и тактике социалистической революции в России и перейти к их осуществлению.

Помню, какое исключительно большое впечатление произвело это открытие — ленинский призыв к социалистической революции — на меня и моих товарищей. В более глубоком понимании смысла поворота в генеральной линии большевиков мне помог следующий факт. Старый большевик Миха Цхакая, который ехал вместе с Лениным из швейцарской эмиграции через Германию, в начале мая, после Апрельской конференции, прибыл в Баку и привез нам комплекты журнала «Коммунист», газеты «Социал-демократ» и другую партийную литературу, издававшуюся в Швейцарии во время войны. Я с жадностью набросился на этот клад и стал читать.

Тогда-то мне и стало ясно, что вопрос о характере и темпах перерастания предстоящей русской революции в социалистическую был по-новому теоретически осмыслен Лениным еще в эмиграции. В первые недели после свержения царизма Ленин сделал четкие, конкретные практические выводы.

В 1917 году Россия была уже не та, что в 1905-м. Ленин видел, что в экономическом отношении она продолжает быть страной отсталой, и никто, наверное, не говорил столь горьких слов по этому поводу, как он сам. Но он придавал еще большее значение тому, что промышленность за последнее десятилетие получила заметное развитие, продвинулась вперед концентрация производства и известная централизация капитала. Выросли мощные очаги концентрации пролетариата: Петроград, Москва, Донбасс, Урал, Баку, Варшава, Нижний Новгород, Рига и др. Обобщая эти изменения, Ленин указывал, что Россия превратилась в страну со средним уровнем развития капитализма, и перешла в империалистическую стадию. Раздиравшие российское общество противоречия обострились до крайнего предела. Обстановка усугублялась поражениями на фронтах вследствие неспособности царизма и русской буржуазии справиться с создавшимся положением.

Ленин раньше других увидел, что господствующие классы окончательно потеряли способность управлять, вывести страну из того кризиса по всем линиям, в который она зашла. А народ не только не желал больше жить по-старому, но и нем уже созрели силы, способные взять на себя руководство государством.

Прошло менее семи месяцем после призыва Ленина к переходу к социалистическому этапу революции, и в октябре 1917 года была одержана историческая победа, оправдалось и это ленинское предвидение.

Чтобы представить себе всемирно-историческое значение выдвинутого Лениным лозунга о переходе к социалистическому этапу революции, задаешься вопросом: какая судьба ожидала русскую революцию, если бы задача перехода к социалистической революции была выдвинута, скажем, через год-два?

Нелегко, конечно, ответить точно на такой гипотетический вопрос. Но все же можно предположить, что буржуазная Учредиловка через пару лет дала бы куцую аграрную реформу, которая удовлетворила бы верхушечные слои деревни и в какой-то мере, на время успокоила большинство крестьянства. Так или иначе окончилась бы война и большая часть солдат разошлась по домам. В результате этих и иных подобных факторов исчезла бы благоприятная революционная ситуация лета и осени 1917 года, и пролетариату во много раз труднее было бы совершить социалистическую революцию, ее сроки отодвинулись бы на неопределенное время.

История зло посмеялась над расчетами и прогнозами меньшевиков. Ветеран марксизма в России Г. В. Плеханов (он был старше Ленина на 14 лет, и Ленин в ранние годы считал его своим учителем) вернулся из эмиграции почти одновременно с Лениным. В июне 1917 года, как бы в противовес ленинским Апрельским тезисам, Плеханов также сформулировал концепцию меньшевизма в отношении социалистической революции в России:

«Русская истории еще не смолола той муки, из которой будет со временем испечен пшеничный пирог социализма».

Вдумайтесь в эти слова!

Плеханов не чувствовал, не видел, что русская история уже перемолола муку и изготавливает тесто, из которого в октябре того же года приступит к выпечке социалистического пирога. Это была смерть русского меньшевизма: четыре месяца спустя после этих «вещих» слов в России победила социалистическая революция.

В таком же глупом положении оказался и лидер европейской социал-демократии Карл Каутский. Когда в Петрограде победила социалистическая революция, он утверждал, что большевикам и двух недель не удержаться у власти. Так проявилось банкротство правой и центристской социал-демократии в Европе.

Поразительным было умение Ленина предвидеть не только исторические повороты в развитии общества, не только формулировать стратегическую линию борьбы класса и партии, определять наличие революционной ситуации, но и с безукоризненной точностью предугадывать тактические повороты и зигзаги событий, действовать сообразно обстановке текущих событий, а подчас и данного конкретного, решающего дня.

Ленин всегда придавал огромное значение не только анализу объективной обстановки, но и субъективным факторам, влияющим на ход революционных событий.

Я уже говорил выше о роли, которую отводил Ленин партии, такой партии, какая была создана им в России. Небольшая партия, насчитывавшая к марту 1917 года всего лишь около 24 тыс. членов, объединила через полгода уже 350 тыс. человек, встала во главе рабочих и крестьянских масс, повела их за собой и одержала победу в стране со 150-миллионным населением.

Какое огромное впечатление произвело тогда, как сильно звучит и сегодня, по прошествии полувека, заявление, сделанное Лениным в июне 1917 года на I Всероссийском съезде Советов. Когда меньшевистские лидеры заявили на съезде, будто в России нет политической партии, которая выразила бы готовность взять власть целиком на себя, Ленин воскликнул: «Есть!»\footnote{\textit{В. И. Ленин}. Полн. собр. соч., т. 32, стр. 500.}.

Враги частенько обвиняли большевиков в антигуманизме, объявляя их сторонниками только вооруженных методов борьбы. Такое обвинение ни на чем не основано. Еще Маркс предупреждал коммунистов: «Восстание было бы безумием там, где мирная агитация привела бы к цели более быстрым и верным путем»\footnote{\textit{К. Маркс} и \textit{Ф. Энгельс}. Собр. соч., 2 изд., т. 17, стр. 635.}. Такой же точки зрения придерживался и Ленин.

В условиях царской тирании партия, естественно, ориентировалась на вооруженное восстание для победы революции. Большевики были воспитаны в этом духе, нас всегда учили, что только таким путем можно прийти к победе, что правящие классы России не уступит власти без боя. Превращение империалистическом войны в гражданскую открывало дорогу революции.

Поэтому можно себе представить, какое потрясающее впечатление производило на нас всех, когда в определенные моменты в 1917 году Ленин выступал за мирное развитие революции. В конкретной и особой ситуации двоевластия он провозгласил лозунг мирного перехода всей власти к Советам рабочих и солдатских депутатов.

Двоевластие после Февраля, как известно, выражалось в том, что было буржуазное временное правительство и были Советы. Рабочие и солдаты слушались своих Советов, у которых была фактическая сила. Но большинство в них составляли, однако, меньшевики и эсеры. Ленин стоял тогда за передачу всей полноты власти даже таким Советам, исходя из того, что в ходе развития революции большинство в них перейдет к большевикам.

Ленин — страстный революционер. В создавшейся тогда революционной ситуации он рвался к победе и все делал для ее приближения. Но Ленин одновременно являл пример трезвости и спокойствия в самые острые моменты, когда иные начинали горячиться и даже теряли голову. Еще в июне 1917 года, когда некоторые участники конференции военных организаций в Петрограде требовали немедленного приступа к вооруженному восстанию, Ленин выступил против такого способа захвата власти, так как мирные возможности развития не были исчерпаны да и армия не была еще готова поддержать рабочих.

И к самому мирному пути развития революции Ленин относился чрезвычайно серьезно, все время напоминая, что мирный путь возможен только при создании огромного перевеса революционных сил и что всегда надо быть готовым к тому, что контрреволюция в любой момент может навязать нам кровопролитие.

Весьма примечательны в этом отношении мысли Ленина, изложенные им 17(4) марта 1917 года, перед самым отъездом из Цюриха на родину, в письме к Коллонтай. Ленин писал, что сейчас на очереди организация масс, пробуждение новых слоев, в том числе создание ячеек в войсках, и «\textit{вооруженное выжидание, вооруженная подготовка} более \textit{широкой} базы для более \textit{высокого} этапа»\footnote{\textit{В. И. Ленин}. Полн. собр. соч., т. 49, стр. 402.}.

В июле того же 1917 года, после расстрела Временным правительством мирной демонстрации рабочих Петрограда, обстановка резко изменилась и создалась новая ситуация. Окончился период двоевластия, отпала возможность дальнейшего мирного развития революции, и Ленин, не колеблясь, выдвигает новый план — план подготовки вооруженного восстания.

И здесь проявился гений Ленина как стратега и тактика революции, политического руководителя, вождя партии и народа.

Однако и в новой обстановке Ленин с пристальным вниманием следит за малейшими, даже очень кратковременными возможностями, открывающимися для мирного пути пролетариата к власти.

После провала попытки генерала Корнилова захватить власть, когда не только большевики и шедшие за ними трудящиеся, но и широкие массы членов мелкобуржуазных партий поднялись на защиту революции, на короткое время опять создалась иная ситуация. Меньшевики и эсеры решили не идти в правительство вместе с кадетами, главной партией буржуазии.

1 сентября 1917 года Ленин в статье «О компромиссах» писал: «Теперь наступил такой крутой и такой оригинальный поворот русской революции, что мы можем, как партия, предложить добровольный компромисс — ...наш возврат к доиюльскому требованию: вся власть Советам, ответственное перед Советами правительство из эсеров и меньшевиков».

Такая возможность представлялась «может быть всего в течение нескольких дней или на одну — две недели», но это «могло бы обеспечить, с гигантской вероятностью, мирное движение вперед всей российской революции...»\footnote{\textit{В. И. Ленин}. Полн. собр. соч., т. 34, стр. 134, 135}.

На такой компромиссный шаг большевики, но словам Ленина, могли идти «только во имя этого мирного развития революции — возможности, крайне редкой в истории и крайне ценной...»\footnote{Там же, стр. 135.}.

Статья «О компромиссах» в силу обстоятельств не попала в редакцию своевременно, а уже через три дня, когда Ленин увидел новый зигзаг в развитии революции, он констатировал: «Пожалуй, те несколько дней, в течение которых мирное развитие было еще возможно, тоже прошли»\footnote{Там же, стр. 138–139.}.

Все это иному человеку может показаться эпизодом мимолетным, не представляющим большого теоретического и даже практического интереса. На самом деле это совсем не так: здесь проявилась примечательная черта ленинского гения. Вторично на протяжении полугода он возвращается к идее мирного развития русской революции в социалистическую Это была возможность, внезапно возникшая и так же внезапно, прежде чем успели высохнуть чернила, исчезнувшая.

Статья была послана Лениным в редакцию уже после, с просьбой озаглавить ее «Запоздалые мысли». «...Иногда, может быть, и с запоздалыми мыслями ознакомиться небезынтересно»\footnote{Там же, стр. 139.}, приписал к ней 3 сентября Ленин. Действительно, какой большой теоретический интерес и сегодня представляют собой «запоздалые» мысли великого Ленина!

Но важна еще и другая сторона дела. Ставя вопрос о вооруженном восстании, рассматривая его, вслед за Марксом, как искусство. Ленин продолжал оставаться величайшим гуманистом, проявлял постоянную заботу о том, чтобы восстание было подготовлено настолько тщательно и момент для него выбран настолько удачно, чтобы оно произошло наиболее безболезненно и с наименьшими жертвами для трудового люда. Это и было классическим образом осуществлено под ленинским руководством в октябре 1917 года.

В этом смысле исключительный интерес представляет письмо Ленина в ЦК партии от 24 октября 1917 года, в канун революции. Мы вновь и вновь поражаемся прозорливости Ильича, его уверенности в силе партии, рабочего класса, умению почувствовать и вовремя определить подготовленный ходом истории момент для решающего выступления. Величие и драматизм ситуации заключались и том, что следовало в одну ночь подвести итог многолетней, изнурительной борьбе. Именно в эту ночь, иначе будет поздно. Ленин тверд, решителен, настойчив, требователен.

Таково было веление исторического момента. В ЦК раздавались голоса: повременить, накопить еще силы, а потом уже выступить. Предлагали подождать съезда Советов, чтобы он провозгласил начало восстания. Но и противник мог накопить силы, перегруппировать их, нанести упреждающий удар. Поэтому Ленин и писал в тот день: «История не простит промедления революционерам, которые могли победить сегодня (и наверняка победят сегодня), рискуя терять много завтра, рискуя потерять все»\footnote{\textit{В. И. Ленин}. Полн. собр. соч., т. 34, стр. 436.}.

Ленин обосновывает право на вооруженное выступление, на насилие во имя высших интересов трудящихся и всего народа: «...народ вправе и обязан решать подобные вопросы не голосованиями, а силой...

Это доказала история всех революций, и безмерным было бы преступление революционеров, если бы они упустили момент, зная, что от них зависит спасение революции, предложение мира, спасение Питера, спасение от голода, передача земли крестьянам...

Промедление в выступлении смерти подобно»\footnote{Там же.}.

Глубочайшей заботой об интересах народа проникнуты эти слова! И 24 октября, как настаивал Ленин, власть была захвачена. Всероссийскому съезду Советов, собравшемуся 25 октября, было доложено, что власть находится в руках Петроградского Совета рабочих и солдатских депутатов в лице его Военно-революционного комитета.

Предсказание Ленина сбылось. История в течение нескольких месяцев напряженнейшей борьбы подвергла испытанию, проверке и полностью подтвердила прогноз, политику и тактику Ленина. Октябрьская революция — величайшее событие мировой истории. Она действительно повернула руль истории, положив начало всеобщему революционному обновлению мира.

Ленин был великий стратег и государственный деятель, он предвидел ход исторических событий. Это отчетливо видно на примере Брестского мира (1918 год). Первые три-четыре месяца победоносная революция в России «переживала период сравнительной и в значительной мере кажущейся самостоятельности и временной независимости от международных отношений»\footnote{\textit{В. И. Ленин}. Полн. собр. соч., т. 36, стр. 93.}. Империалисты в этот период были настолько заняты войной между собой, что им было не до Советской России, которая этим пользовалась. «Буржуазия, находясь в мертвой схватке борьбы друг с другом, была парализована в своем наступлении на Россию»\footnote{Там же, стр. 96.}.

По вскоре, как и предполагал Ленин, выяснилось, что ни Германия, ми Антанта понес не склонны предоставить Советской республике спокойно существовать. Германская армия перешла в наступление, захватила Украину, Белоруссию, Прибалтику. Одновременно к вторжению готовились и страны Антанты, признавшие в марте 1918 года на конференции Англии, Франции и Италии и Лондоне необходимость союзной интервенции в Восточной России.

Соотношение сил было совсем не в пользу Советской России. Более того, царская армия разложилась и начала распадаться еще до Октября. Провозглашенный в октябре лозунг о мире был воспринят вконец измотанными солдатскими массами с восхищением. Солдаты и офицеры расходились по домам. Большевики еще только приступили к созданию новой, социалистической армии. Среди офицеров сторонников Советов было крайне мало. Красная гвардия-малочисленна, не способна противостоять регулярной армии кайзера.

Ленин видел, что нужен мир, нужна передышка, чтобы организоваться, укрепить власть, дисциплину, начать налаживать транспорт, экономику, создать армию, выждать, пока изменятся международные условия, произойдет подъем революционного движения на Западе.

Путь к этому был лишь один — пойти на временные территориальные уступки империалистам, лавировать, отступить хотя бы «в полупорядке», как выражался Ленин, для того, чтобы сохранить все, что можно сохранить, прежде всего то, что выше всего, — начавшуюся социалистическую революцию и Советскую власть. Не считаясь с революционными эмоциями, отвергая точку зрения тех, которые «пытаются ответить относительно тактики революции на основании непосредственного чувства»\footnote{\textit{В. И. Ленин}. Полн. собр. соч., т. 36, стр. 92.}, пойти на унизительный, «похабный» мир.

Ни Ленин оказался в трудном положении, одно время даже в меньшинстве в ЦК партии. Возникла угроза раскола партии.

Отвергая мир, одни («левые коммунисты») требовали революционной войны против германского империализма. Московское областное бюро партии приняло резолюцию, требующую начать революционную войну и биться даже вплоть до утраты Советской власти. Это была опасная, авантюристическая линия «левизны». Как выявилось на IV чрезвычайном Всероссийском съезде Советов, она смыкалась с линией оппозиционных правых партий.

Другие поддерживали Троцкого, который выдвинул лозунг «ни мира, ни войны». Фактически это означало открыть ворота для немецкого вторжения.

Ленин доказывал, что в этот момент решается судьба революции: или перемирие с немцами, передышка, дающая выигрыш во времени, или гибель Советской власти. Он боролся в ЦК партии как лев, он заявлял даже, что если не будет принят мир, то он уйдет с поста председателя правительства и из состава ЦК.

По вине ведшего переговоры Троцкого Брестский мир был подписан только 3 марта 1918 года.

Ленин воевал за этот мир и отвоевал его. Борьба вокруг Брестского мира — ярчайший, классический пример того, как надо организовать отступление для передышки, чтобы перейти снова к наступлению. Уметь совершать такие маневры далеко не просто Это было настоящее отступление. Оно было необходимо для того, чтобы подготовить партию, страну для нового наступления. Искусство обеспечивать в необходимых случаях организованное отступление, которое даже труднее, чем вести наступательные бои.

Интересно еще и то, что Ленин не пытался делать прогнозы о возможных сроках передышки, ибо для этого не было достаточных данных. Он неоднократно высказывался по этому поводу. На VII съезде партии: «Не розовые краски я рисую насчет передышки; никто не знает, сколько будет продолжаться передышка, и я не знаю. Смешны потуги тех, которые стараются из меня выжать, сколько будет продолжаться передышка»\footnote{\textit{В. И. Ленин}. Полн. собр. соч., т. 36, стр. 33.}. Вновь на IV съезде Советов Ленин упорно отказывается дать прогноз, не имея для этого данных: «...Меня приглашают ответить на вопрос: неделю, две или больше будет передышка? Я утверждаю, что на всяком волостном сходе и на каждой фабрике человека, который от имени серьезной партии с подобным вопросом будет выступать к народу, его народ высмеет и прогонит вон, потому что на всяком волостном сходе поймут, что нельзя задавать вопросы о том, чего нельзя знать». И еще раз о том, что спрашивают, ««надолго ли эта передышка?», на это ответить нельзя, потому что не было международной объективной революционной ситуации»\footnote{Там же, стр. 116, 118.}.

Вот как ответственно и серьезно подходил Ленин к прогнозам и предсказаниям: при отсутствии необходимых для этого фактов и обстоятельств никакими силами предсказаний выжать из него было невозможно.

Прошло всего лишь около восьми месяцев. Германия потерпела поражение, в ноябре там вспыхнула революции. Постановка коренным образом изменилась, и по предложению Ленина ВЦИК издал декрет об аннулировании Брестского договора. «Похабный» мир был похоронен.

Жизнь еще раз подтвердила правильность избранной Лениным тактики в сложнейших и невыгодных условиях. Эта тактика была частью ленинского стратегического плана на закрепление Советской России в международных отношениях, постепенного укрепления ее международного положения, на принуждение империалистов к сосуществованию с новым строем, на использование в этих целях противоречий между ними.

И своих предвидениях Ленин опирался на опыт революционных масс, на опыт партии в целом и ее руководящего коллектива, на глубокий и всесторонний анализ жизни. Этим он и отличался от кабинетных «оракулов» разных мастей. Превыше всего ставил он творчество масс, считал его основой общественных процессов и неисчерпаемым источником революционной энергии, ибо, как он подчеркивал, «ум десятков миллионов творцов создает нечто неизмеримо более высокое, чем самое великое и гениальное предвидение»\footnote{\textit{В. И. Ленин}. Полн. собр. соч., т. 35, стр. 281.}.

\subsection*{\center * * *}

Ленинское научное предвидение было обращено в будущее. Именно Ленин выдвинул коренные вопросы общественного развития на долгие годы вперед. Поэтому ленинское учение продолжает оставаться живым, в нем черпают силы новые и новые поколения революционеров.

Не ставя задачей даже назвать все поставленные и решенные Лениным важнейшие вопросы, хотелось бы обратить внимание читателя лишь на некоторые из них.

Это — ленинское учение об империализме как высшей и последней стадии капитализма, о его развитии в направлении государственно-монополистического капитализма, приобретшего теперь столь реальные формы и содержание.

Это — ленинское учение о социалистической революции, о многообразии путей и форм перехода от капитализма к социализму, подтверждаемое ныне всем ходом мирового революционного процесса.

Это — ленинские планы построения социализма в Советском Союзе, сочетания национального и интернационального начал. Планы индустриализации страны, кооперирования крестьянства в условиях диктатуры пролетариата, осуществленные по ленинским заветам в Советском Союзе и осуществляемые в наше время большой группой стран, объединяющих свои усилия на пути к общей цели — коммунизму.

Это — ленинские идеи об объединении борьбы международного рабочего класса с национально-освободительным движением, его предвидение гибели колониальной системы, появления реальных возможностей перехода отдельных стран в соответствующих условиях к социализму, минуя капитализм.

Это — ленинские установки на длительное сосуществование и борьбу двух систем, конечная цель которой — всемирная победа нового строя.

И многое, многое другое.

Огромны, поистине беспредельны и трудно постижимы широта ленинской мысли, круг охваченных им проблем. Но все, о чем он думал, что творил, чему посвятил свою жизнь и борьбу, нее подчинено одной светлой цели — социалистической революции, коммунизму. Он отчетливо представлял очертании этого будущего, естество и душу коммунизма.

Об этом образно писал Максим Горький: «...Владимир Ильич Лепим так хороню знал историю прошлого, что мог и умел смотреть на настоящее из будущего... Он умел и мог делать это — мне кажется потому, что половиною великой души своей жил и будущем; железная, но гибкая логика его показывала ему отдаленное будущее в формах совершенно конкретных, реальных»\footnote{«В. И. Ленин и А. М. Горький. Письма, воспоминания, документы», 1961, стр. 287.}.

\newpage
	\section[Письмо Ленину]{\center \textbf{Письмо Ленину}}


Шел 1919 год...

Положение на Кавказе было исключительно тяжелым.

С каждым днем мы все острее чувствовали необходимость установления живой, непосредственной связи с Москвой, с Центральным Комитетом партии, с Лениным. Нельзя было решать крупнейшие вопросы политики только по своему уму-разуму. Нельзя было дальше работать без необходимой информации, без советов и указаний центра, да еще в таких сложных условиях. Мы, конечно, хорошо понимали, что наладить связь с Москвой будет делом нелегким. Нас разделяло не только расстояние — разделяла гражданская война, фронт контрреволюции. И все-таки связь надо было устанавливать несмотря ни на что.

В мае 1919 года, будучи членом бюро краевого комитета партии, я послал из бакинского подполья в ЦК партии и Ленину, как Председателю Совнаркома, письмо с информацией о положении в Закавказье и о наших задачах. Я боялся, что это письмо может не попасть к Ленину (кто там примет, в чьи руки оно попадет, куда направят — я не знал). Мне очень хотелось, чтобы Ленин его увидел, и казалось, что только таким путем ЦК может разобраться в наших делах, и мы получили бы директивы по волнующим нас вопросам. Поэтому л в адресе подчеркнул: ЦК РКП(б) и Председателю Совнаркома Ленину.

«Благодаря годичной оторванности Закавказья от центра, — писал я Ленину, — мы лишены Ваших директив и помощи в нашей борьбе, а Вы или вовсе не знакомы с настоящим положением вещей в Закавказье, или же неправильной информацией безответственных лиц введены в заблуждение и поэтому не гарантированы от серьезных ошибок во внешней политике по отношению к нашему краю. Пользуясь маленькой возможностью, информирую Вас...»\footnote{ЦПА ИМЛ, ф. 5, оп. 1, ед. хр. 475, л. 1.}

Мое письмо-доклад Ленину было довольно пространным. Я писал об обстановке, которая сложилась в Баку и Закавказье, об оккупационном режиме, о роли и политике азербайджанского национально-буржуазного правительства, о том, что в Азербайджане больше, нежели в Грузии или Армении, «горючего материала, больше социальных обостренных противоречий, больше классовой почвы для пролетарской революции, больше недовольства и ненависти к существующему правительству», имеется хорошо организованный рабочий класс, находящийся под влиянием и руководством коммунистов. Писал о том, что «Бакинская организация признала независимый Советский Азербайджан, с тесной Политической и экономической связью с Советской Россией», и что этот лозунг очень популярен и может сплотить вокруг себя широкие массы трудящихся азербайджанцев и поднять на восстание». Писал о нашей Рабочей конференции, как именовался тогда постоянный орган представителей промысловых и заводских комитетов; рассказывал о начале и исходе всеобщей забастовки бакинских рабочих, о росте влияния партии среди бакинских рабочих; писал о положении в Грузии и Армении; сообщал о хозяйничанье англичан вкупе с белогвардейцами в Петровске (ныне Махачкала), о разгроме ими Дагестанского комитета партии, о революционных настроениях дагестанских крестьян и организации повстанческих отрядов в Дагестане...

Я сообщал Ленину, что мы готовим вооруженное восстание, создаем опорные повстанческие пункты в крае и сможем выступить, как только узнаем, что части Красной Армии и Флота готовы нам помочь. Просил Ленина и ЦК партии направить в Баку опытных партийных работников (особенно из мусульман), а также прислать партийную литературу на русском, азербайджанском и армянском языках.

Это письмо, отпечатанное на полотне и зашитое в одежду, с большим риском для себя провезли из Баку через деникинский фронт в Москву двое наших молодых коммунистов — Шура Берцинская и Тигран Аскендарян. Баку в то время был фактически отрезан от центральных районов России. Северный Кавказ и в значительной мере Донбасс находились в руках Деникина, Колчак рвался к Волге с востока.

Когда Тигран уезжал в Москву, я дал ему, кроме того, строгий наказ: во что бы то ни стало привезти стенограмму VIII съезда партии, состоявшегося во второй половине марта 1919 года. Нам было известно об этом съезде очень мало: газеты из Советской России к нам не доходили; о том, что происходило на съезде, мы знали только по слухам да из сообщений буржуазной прессы.

Тиграну и Шуре было поручено передать письмо лично Ленину. Они успешно выполнили это поручение и уже в июле 1919 года вернулись — через Астрахань и Каспийское море — обратно в Баку с группой ответственных партийных работников, которые были нам в ту пору так необходимы.

\newpage
\section[Путь в Москву. У Кирова в Астрахани]{\center \textbf{Путь в Москву. \\ У Кирова в Астрахани}}


Партийные организации и революционное движение и Баку и середине 1919 года уже в значительной степени оправились после поражения, понесенного в 1918 году. Советская власть тогда пала в результате иностранных интервенций — германо-турецкой и английской, опиравшихся на местные контрреволюционные группировки.

Одно время наша Красная гвардии удерживала некоторую часть юго-западного побережья Каспия, где действовала Советская власть и была провозглашена Ленкоранская Советская республика. На нее возлагались большие надежды: красный флот мог бы высадить там десант и, опираясь на поддержку красногвардейских частей и крестьянства, войти в Баку, где рабочий класс ждал момента для выступления. Однако из-за измены в штабе флота после выхода из Волги и остановки в форте Александровском корабли были там настигнуты деникинским флотом и уничтожены. Тем самым перспектива похода Красного флота на Кавказ, в Азербайджан и на этой основе развитие революции и утверждение Советской власти в Закавказье на определенное время оказались исключенными. Нужно было определить новые перспективы, к чему и как нам готовиться дальше.

Назрели и приобретали все большее значение партийно-организационные вопросы. По ним не было единой точки зрения, и это еще более осложняло положение.

Все это вместе взятое требовало новых установок для кавказской краевой парторганизации, нужна была поездка в Москву для доклада о положении дел и получения необходимых указаний. К тому же за год с лишним, после поражения большевиков на Кавказе, ни один член крайкома не был в Центральном Комитете партии. Кавказским крайкомом партии было решено, что в Москву с докладом в ЦК от имени крайкома надо ехать мне. Признаться, сам я этого хотел и был рад, когда товарищи выдвинули мою кандидатуру.

Я знал, что путь будет тяжелым и рискованным, но предвкушал радость личной встречи с Лениным. Мне тогда было 24 года.

В Баку стало известно, что в ноябре должен собраться VII Всероссийский съезд Советов (на самом деле он состоялся 5–9 декабря 1919 года). 26 сентября президиум Рабочей конференции Баку и его районов выдал мне мандат, в котором говорилось:

«Выдан сей мандат от Президиума Центральной Рабочей Конференции г. Баку тов. Анастасу Микояну в том, что он делегируется на Всероссийский съезд Советов, предполагаемый быть в ноябре м-це»\footnote{ЦГАОР СССР, ф. 1235, оп. 6, ед. хр. 81, л. 77.}.

Путь в Москву был один: на рыбацкой лодке пять-шесть дней по Каспию в Астрахань. Этот путь уже был освоен организованной бакинскими большевиками флотилией лодок, возивших из Баку авиационный бензин для Советской России.

Этим путем я и воспользовался. 16 октября, после изнурительного плавания в штормовую погоду, миновав деникинские кордоны, я прибыл в Астрахань.

Меня доставили к Кирову на квартиру. Дома его не оказалось. Какая-то старушка, видимо хозяйка квартиры, на вопрос, где Киров, ответила: «А он с утра в Совете, речи произносит».

В то время в Астрахани было очень плохо с продовольствием. Потом, когда я ходил по улицам города, видел, как кое-где с рук продавались моченые яблоки, иногда попадалась вяленая вобла.

Кирову приходилось довольно туго. Ему надо было всюду поспеть, всем объяснить, разъяснить, а главное, поднять дух голодных людей.

Наконец, появился Киров. Это была моя первая встреча с ним. До этого мы были знакомы только по оживленной переписке, которая завязалась у нас около полугода назад.

Встретились мы как хорошие, давние друзья.

Я подробно рассказал ему все, что собирался говорить в ЦК партии. Особенно детально мы обсуждали конкретные вопросы помощи повстанцам в Кавказских горах, а также наметили меры по усилению вывоза бензина из Баку в Астрахань.

Человек очень живой, пытливый, умный, ясно и четко мыслящий, он мгновенно разобрался во всех тонкостях, и это было особенно приятно. Его положительное отношение к нашей позиции по тем вопросам, которые нас особенно волновали, подбодрило меня, вселило уверенность, что они будут успешно рассмотрены и в ЦК партии.

Киров приятно поразил меня своей работоспособностью, оперативностью, умением быстро схватывать суть дела и незамедлительно принимать решения по конкретным вопросам. Было видно, что все нити военной, государственной и партийной работы тянулись здесь прямо к нему, и он, опираясь на доверие товарищей, пользуясь высоким авторитетом, умело осуществлял руководство. В дни пребывания в Астрахани мы часто общались с ним, близко узнали друг друга и стали навсегда друзьями.

Киров тех дней запечатлелся в моей памяти исключительно собранным, подтянутым, необычайно цельным человеком, обладавшим к тому же очень твердым характером. И внешним своим обликом он располагал к себе людей. Невысокого роста, коренастый, симпатичный, он обладал каким-то особенным голосом и необыкновенным даром слова. Когда он выступал с трибуны, голос у него был мощный, приятный, покорявший слушателей. В личных беседах он не был многоречив. Но говорил ясно, четко, умел внимательно слушать собеседника, любил вставить острое словцо и сам был отличным рассказчиком.

Выяснилось, что Киров поддерживает постоянную связь с Лениным по телеграфу, сообщает о положении дел, запрашивает указания, передает в центр поступающую в Астрахань информацию с Кавказа.

В первый же день Сергей Миронович послал Ленину телеграмму о моем прибытии в Астрахань и предстоящей поездке в Москву. В большой телеграмме Киров сообщил Ленину обо всем услышанном от меня: о размахе восстания горцев, о мерах, которые в связи с этим собирается принять Деникин.

С Кировым мы решили, что я пробуду в Астрахани несколько дней: хотелось дождаться Федю Губанова, чтобы, как мы договаривались в Баку, ехать вместе в Москву.

Время пребывания в Астрахани не пропало у меня даром. Удалось осуществить ряд важных дел. В частности, 21–22 октября с пользой для дела я встретился с находившейся там большой группой коммунистов-армян, собиравшихся через Баку пробраться в Армению для нелегальной работы.

Обрадовался я встрече с прославленным чекистом Атарбекяном, который был правой рукой Кирова в борьбе с контрреволюцией и которого я знал еще по 1916 году в Эчмиадзине.

С тревогой ждал я прибытия Губанова. На четвертый-пятый день возникло, а затем усилилось подозрение, что он попал в лапы деникинцев. Впоследствии выяснилось, что его лодку настигли белогвардейцы, сам он был арестован и вскоре погиб.

Приходилось ехать одному. Регулярного сообщения с Москвой не было. Уехать можно было лишь с какой-либо оказией.

«Такая оказия есть, — сказал мне Киров. — Через несколько дней сюда должен прибыть со своим поездом член Реввоенсовета республики Смилга с группой военных работников. Смилга пробудет в Астрахани день-два, и ты вполне сможешь уехать с ним в Москву. Так ты попадешь туда раньше, чем любым другим способом».

Так все и произошло. 26 октября я уехал в Москву и поезде, которым возвращался Смилга. Придавая большое значение моей поездке в Москву, Киров предполагал, что для Ленина окажется важным доклад о положении дел на Кавказе непосредственно от участника происходивших там событий. Поэтому Киров телеграммами от 24 октября Стасовой и от 27 октября Ленину и Стасовой вновь сообщал о моем выезде в Москву.

\newpage
\section[На докладе у Ленина]{\center \textbf{На докладе у Ленина}}


2 ноября 1919 года я прибыл в Москву. Прямо С вокзала направился в Центральный Комитет партии, аппарат которого размещался тогда и здании на Воздвиженке (ныне проспект Калинина).

Миш направили и комнату, в которой работала Елена Дмитриевна Стасова. Когда я вошел в небольшой и, помнится, несколько темноватый кабинет, высокая белокурая женщина, как мне показалось, аристократической внешности, стоя у невзрачного стола, разговаривала с каким-то товарищем Она обернулась на стук закрываемой двери и, остановив собеседника движением руки, поинтересовалась, кого я ищу и по какому поводу. Я представился, сказал, что мне необходимо видеть секретаря ЦК Стасову, и на это услышал коротенькое: это я. Елена Дмитриевна уже знала о предстоящем моем приезде из телеграмм Кирова. Приветливо улыбнувшись, она попросила меня присесть и подождать окончания разговора с товарищем.

Через несколько минут я уже отвечал на вопросы Елены Дмитриевны, передал ей письмо Кирова, написанное им буквально перед самым моим отъездом из Астрахани. Я знал, что Киров высказывался за поддержку наших позиций по организационно-партийным вопросам. Расспросив, как я доехал, Стасова направила меня к Владимиру Ильичу, в Кремль. Сказала, что Ленин дал поручение: как только я появлюсь в городе, сразу же доставить меня к нему. В тот же день вечером он принял меня в своем кабинете, один.

Когда я открыл дверь в кабинет Ленина, он, приветственно, с прищуром улыбаясь, встал из-за письменного стола и вышел мне навстречу. Дружелюбно и просто пожил руку. Спросил, как я доехал. Ленин предложил мне сесть на стул, стоявший у письменного стола, и сам вернулся на свое место.

— Рассказывайте, рассказывайте...

Когда я начал говорить, Ленин сразу как-то преобразился: весь превратился по внимание, улыбка сошла с лица, выражение глаз стало серьезным, пытливым.

Я слышал, что Ленин человек простой, но не представлял его таким, каким увидел. Он сразу же создал атмосферу непринужденной деловой беседы. Поначалу и сильно волновался, но вскоре собрался и смело, без смущения стал ему докладывать.

Рассказал о больших успехах большевиков Азербайджана за полугодие, прошедшее с весны 1919 года, о сплочении бакинского пролетариата вокруг нашей партии, о сочетании подпольной партийной работы с легальными формами деятельности. О том, что есть в Баку постоянно действующая Рабочая конференция из делегатов предприятий — нечто вроде Совдепа, которая является сильным органом и фактически руководит рабочим движением в городе. В марте нынешнего года нам, коммунистам, удалось вытеснить эсеро-меньшевистское большинство из президиума Конференции и, пользуясь поддержкой рядовых членов Конференции, занять руководящее положение в ней и в ее президиуме.

Я говорил о том, что бакинские профсоюзы в начале 1919 года тоже находились в руках меньшевиков, организационно и идейно были очень слабы, а порой и беспомощны; но нам все-таки удалось занять ведущее положение в профсоюзах, а также в рабочих клубах районов Баку, превратив их в базы для развертывания массовой политической работы среди рабочих, в пункты связи и явок партийных организаций. Рассказал о том. как прошел общекавказский съезд профсоюзов, принявший резолюцию, предложенную коммунистами, отметив, что на этот съезд из Грузии прибыли представители только нескольких профсоюзов; союзы, которые находились под руководством меньшевиков, предвидя, что на съезде они окажутся в меньшинстве, не послали своих представителей, устроили свой съезд профсоюзов Грузии, отколовшись тем самым от общекавказской профсоюзной организации.

Я счел нужным сообщить Владимиру Ильичу, как нам удалось, вопреки усилиям меньшевиков, получить большинство на перевыборах правления Каспийского кооперативного объединения, взять управление им в свои руки. Этот кооператив стал опорным пунктом нашей партийной работы. Его предприятия использовались для устройства безработных коммунистов, для посылки людей в командировки по партийным делам под видом кооперативной работы, для организации связи, для устройства явок и пр.

Еще по пути в Москву я все время думал о предстоящей встрече с Лениным, о том, как лучше сделать доклад. Решил сначала принести одни только факты, потом их проанализировать, обобщить и сделать выводы. Так я и поступил. Владимир Ильич с жадностью слушал меня, когда речь шла о фактах. Но как только я пытался перейти к выводам, он меня вежливо прерывал, вставал, ходил по комнате, расспрашивал о дополнительных фактах. Подходил к карте: «А ну, где это, посмотрим, где Дагестан, где Чечня?» Смотрим. «Сколько там партизанских отрядов?» — вновь спрашивает Ленин. Отвечаю и опять начинаю обобщать. Ленин снова задает вопросы, уточняет. Тогда только я понял, что мне надо подробнее рассказывать о фактах и не стараться делать выводы: Ленин сам их сделает, и, конечно, лучше меня.

Я доложил Ленину о политическом положении в республиках Закавказья, в Дагестане, Чечне, Ингушетии и Кабарде. Отметил, что наиболее революционная ситуация сложилась в Азербайджане, где не только в Баку, но и в уездах развернулась борьба против буржуазного правительства, против помещиков, за захват земель. Коммунистам удалось ваять на себя руководство этим движением.

После трагической гибели Ленкоранской Советской республики наиболее сильно антиправительственное движение развернулось н Казахском уезде и в Карабахе. А в Дагестане, Чечне и Кабарде идет партизанская борьба против деникинских войск.

Усилиями грузинских меньшевиков и при поддержке муссаватистского правительства образовано правительство горских народов Кавказа. Однако фактически власти оно по имеет, решающей силой там являются красные партизанские отряды.

В Кизлярском районе и особенно в Новороссийском и Туапсийском районах действуют отряды так называемых «зеленоармейцев», состоявшие из недовольных крестьян, бывших северокавказских красноармейцев и дезертиров из деникинской армии. Они ведут боевые действия против Деникина. Краевой комитет партии принимает энергичные меры к установлению связи с ними, чтобы использовать их в борьбе с деникинской армией, для разрушения тыла противника.

Ленин задал вопрос: каково экономическое положение в буржуазных республиках Закавказья?

Я ответил, что относительно благополучно положение в Грузии. Азербайджан, владея огромными запасами уже добытой нефти, но лишившись российского рынка, переживает депрессию из-за невозможности вывезти и продать нефть.

Самое тяжелое положение в дашнакской Армении. При общей бедности населения, отсутствии промышленности и малоземелье экономическое положение здесь катастрофическое. Около 300 тыс. армян-беженцев, пришедших из Западной Армении вместе с отступавшими русскими войсками, очутились в ужасных условиях. Широко рекламируемая помощь Англии и Америки на деле ничтожна, народ голодает.

Ленин спросил, каковы взаимоотношения между националистическими правительствами Закавказья. Я ответил, что они грызутся между собой, их раздирают территориальные споры. Предпринятая летом на встрече представителей правительств попытка наладить сотрудничество результатов не дала.

В Грузии, вопреки меньшевистскому правительству, соотношение внутренних общественных сил развивается н нашу пользу. Большинство активной части крестьян и рабочих, отдельные подразделения регулярной армии не только хорошо к нам настроены, но и готовы восстать по призыву нашей партии. В некоторых районах и армейских частях даже не всегда удается предупредить преждевременные стихийные выступления. Так было, например, в Поти, Михайлове, Ланчхутах, Сухуме и др.

В Армении как экономическое положение, так и политические условия тяжелые, коммунисты подвергаются сильным преследованиям, они вынуждены вести работу только в подполье, так как никаких возможностей для легальной работы нет. В народных массах, среди недовольных существующим положением, очень сильна тяга к. России. В случае победы над Деникиным трудовое население Армении, включая и мусульман, поднимется на борьбу за Советскую власть.

Вдруг Ленин, хитро усмехнувшись, спрашивает: а как осуществляют демократию грузинские меньшевики?

Я ответил, что никакой демократией в меньшевистской Грузии и не пахнет Наоборот, там жестоко подавляют крестьянские восстании, возникшие в некоторых уездах. Арестовали и держат и тюрьмах без суда и следствии большую группу наших партийных товарищей. Арестован и сидит и Кутаисской тюрьме известный Ленину по швейцарской эмиграции Миха Цхакая. Конечно, нет и речи о свободе печати. Большевики не могут издавать ни легальных газет, ни брошюр, они работают подпольно.

Ленин очень горячо реагировал на мой рассказ и высказался примерно так: я этих меньшевиков хорошо знаю! От них другого нельзя и ожидать.

Я вспомнил, что в Грузии как-то начала выходить одна газета большевистского направления, но не прошло и двух недель, как меньшевики ее закрыли. Это было в дни, когда мы пытались создать единый фронт против Деникина. Летом, когда Деникин одержал ряд побед на фронте и появилась угроза его вторжения в Закавказье, сильно встревожились и мелкобуржуазные круги населения и даже правительства Закавказья. Желая создать затруднения Деникину на Кавказе и тем помочь Красной Армии, мы обратились с призывом к мелкобуржуазным партиям и национальным буржуазным правительствам с предложением о единых действиях против деникинской опасности, обещая им всяческую поддержку в этом деле. Наш призыв нашел отклик в широких массах трудящихся, под давлением которых грузинские меньшевики согласились вступить в переговоры и созвать общекавказский рабочий съезд для организации борьбы против деникинского нашествия. Этот период мы успешно использовали для легальных выступлений среди рабочих Тифлиса. Однако меньшевики увидели, что развернувшаяся кампания льет воду на большевистскую мельницу. Они добились у английских представителей гарантии, что те не допустят нашествия Деникина в Закавказье, и после этого прервали всякие переговоры с нами. Началась новая волна репрессий против коммунистов, разнузданной антибольшевистской пропаганды.

Рассказал я Ленину о том, что партийные организации Закавказья успешно готовятся к вооруженному восстанию и ждут указаний ЦК о времени его проведения. Сами мы не можем определить его наиболее целесообразного срока. ЦК партии знает положение на фронтах, а мы хотели бы приурочить восстание к моменту приближения Красной Армии к Кавказу.

Я сказал далее, что в Азербайджане сейчас кроме РКП(б) существуют две коммунистические организации. Это «Гуммет» (Энергия), объединяющая местных рабочих азербайджанской национальности, и «Адалет» (Справедливость), объединяющая рабочих — выходцев из Персии. Представители всех коммунистических организаций после некоторых споров пришли к выводу о том, что раздробленность рядов коммунистов не соответствует современному положению и задачам предстоящего восстания; «Гуммет» и «Адалет», созданные давно, совсем в других условиях, сыграли большую роль по вовлечению мусульманских рабочих в коммунистическое движение, а теперь надо иметь единую коммунистическую организацию для всех, независимо от национальной принадлежности.

Поскольку мы выдвинули лозунг создания независимого Советского Азербайджана, тесно связанного с Советской Россией, то такой единой организацией должна стать Коммунистическая партия Азербайджана, объединяющая всех коммунистов республики и входящая в состав РКП(б).

Я сообщил Ленину, что это единодушное мнение коммунистов Азербайджана находится в противоречии с принятым в июле этого года решением ЦК, согласно которому «Гуммет» признается самостоятельной Коммунистической партией Азербайджана с правами областного комитета партии.

Далее, я заявил, что коммунисты Армении и Кавказский крайком РКП(б) не признают находящийся в Москве и возглавляемый Г. Айкуни ЦК Компартии Армении. Он и его группа не имеют никакой связи с местными парторганизациями Армении.

Коммунистические» организации Армении избрали в условиях подполья комитет, получивший название Арменком, для руководства работой всех коммунистов Армении. Арменком входит в состав Кавказской краевой партийной организации и через него в РКП(б). Он признает руководство крайкома и не признает группу Айкуни. Эта группа, и спою очередь, признает только непосредственное руководство ЦК РКП(б) и отвергает руководство со стороны Кавказского крайкома партии, не желает входить в состав Кавказской краевой организации коммунистом и тем самым раскалывает ряды коммунистов Кавказа.

Рассказал Ленину, что, будучи в Астрахани, я встретился с группой армян-коммунистов, собиравшихся ехать на Кавказ, с которыми был Айкуни. Приняв участие в их собрании, я выступил против позиции Айкуни, и все присутствовавшие приняли предложенную мною резолюцию, а Айкуни остался в одиночестве и решил вернуться в Москву.

— Кавказские коммунисты не мыслят победу без единства коммунистов Кавказа, — говорил я.

По наивности я думал, что Ленин прямо в ходе беседы или сразу по ее окончании даст мне определенные ответы на поставленные вопросы, выскажет свое мнение и соображения по ним.

Ленин сказал, что все эти вопросы — о компартиях Азербайджана и Армении, об их отношении к крайкому, — равно как и другие поставленные мною вопросы, надо изучить и обсудить в ЦК, а потом уже принимать по ним решения. Для этого надо, чтобы я все изложил в письменном виде, что я позже и сделал.

Что же касается объединения коммунистов одной организации по территориально-производственному принципу, независимо от их национальной принадлежности, Ленин сказал, что это совершенно правильно. Правильно также, чтобы в организовавшихся на окраинах России самостоятельных государствах коммунистические организации работали в виде самостоятельных компартий, входящих в состав РКП(б).

Эта первая встреча и беседа с Владимиром Ильичей Лениным никогда не изгладится из моей памяти.

\subsection*{\center * * *}

После беседы с Лениным я вернулся на Воздвиженку, где мне дали направление в третий дом Советов (бывший Божедомский переулок, ныне Делегатская улица. Теперь в этом здании находятся Президиум Верховного Совета и Совет Министров РСФСР). В Москве я был впервые, ни трамваи, ни другой городской транспорт не работали, и, чтобы добраться до места, мне дали автомобиль.

На следующее утро мне удалось узнать, где живет семья Шаумяна. Оказалось, что она занимает квартиру во дворе этого же дома, и я зашел к ним. Екатерина Сергеевна Шаумян настояла на том, чтобы я перебрался к ним на квартиру. И вновь, как в Баку, дом Шаумяна стал для меня родным домом.

Наутро Стасова вновь приняла меня, как мы и договорились накануне. Подробно рассказал ей о положении дел на Кавказе и о тех вопросах, которые крайком ставит на рассмотрение ЦК. Затем, по совету Стасовой, я зашел к К. Т. Новгородцевой, заведовавшей тогда Орготделом ЦК, и также информировал ее обо всем. Она произвела на меня очень хорошее впечатление, как серьезный, вдумчивый партийный работник, простая в обращении и располагавшая к себе собеседника. Позже я узнал, что это была жена скончавшегося в начале года секретаря ЦК и Председателя ВЦИК Я. М. Свердлова.

\subsection*{\center * * *}

Много дней я переживал и обдумывал беседу с Лениным. Поражало, что Владимир Ильич подробнейшим образом интересовался всеми вопросами Кавказа, в том числе и деталями. Задавал многочисленные вопросы, уточнял, что происходит в Баку и в Азербайджане в целом, в Грузии, Армении, Дагестане, Чечне и у других народов Кавказа.

Казалось, до этого ли ему было. И внешнее и внутреннее положение Российской республики продолжало оставаться крайне тяжелым. В конце года, по представленным Ленину данным ЦСУ, Советская власть распространялась всего лишь на 75 млн. человек, т. е. примерно на половину населения царской России в 33 губерниях Республики. Это были губернии преимущественно потребляющие.

Из Архангельска и Мурманска английские войска продвигались к югу, на Вологду, белая армия Юденича стояла под Петроградом и, перейдя в наступление в конце сентября, в середине октября приблизилась к городу. На основном в то время деникинском фронте в конце октября произошел перелом, движение деникинских полчищ на север было приостановлено, они были отброшены от Орла. Была снята прямая угроза нашествия деникинцев на Москву. Однако положение в конце года все еще оставалось сложным. На Восточном фронте под давлением Колчака советские войска в середине октября были вынуждены отойти к Тоболу. Заводы, за исключением нескольких, работавших на оборону, стояли. Советская республика почти лишилась источников сырья, топлива, хлеба. Производство металла дошло до минимума. Связь с Советским Туркестаном была прервана казачьими войсками Дутова, в результате чего прекратилась поставка в Россию хлопка.

Как же велика была у Ленина вера в победу, если в такое труднейшее время он счел возможным более двух часов расспрашивать меня о том, что происходит на далеком Кавказе, какие вопросы волнуют местных коммунистов, как мы собираемся их решать!

Я расскажу далее, как Ленин и Политбюро не один раз еще вернутся к поставленным нами вопросам, окажут решающую помощь в освобождении Кавказа, в установлении там Советской власти.

\newpage
\section[Ленин на II Всероссийском съезде коммунистических организаций народов Востока: надо слиться в общей борьбе с пролетариями других стран]{\center \textbf{Ленин на II Всероссийском съезде коммунистических организаций народов Востока: надо слиться в общей борьбе с пролетариями других стран}}


В марте 1919 года был создан Коммунистический Интернационал, о необходимости которого Ленин спорил еще в Апрельских тезисах. Состоялся его I конгресс, положивший начало идейному и организационному объединению коммунистов всех стран под лозунгом диктатуры пролетариата. В этом было коренное отличие Коминтерна от международных объединений социал-демократов. Но было и еще одно существенное отличие: на конгрессе Коминтерна впервые в истории революционного движения объединились коммунисты страны, победившей в социалистической революции, пролетарии Европы и Америки с тружениками колониальных и зависимых стран. В этом плане важное значение имели и проводившиеся тогда съезды коммунистических организаций Востока.

Стало известно, что 22 ноября 1919 года в Москве открывается II Всероссийский съезд коммунистических организаций народов Востока. Мне захотелось в качестве гостя присутствовать на этом съезде, чтобы получить информацию и поучиться на опыте этих организаций решению национального вопроса.

Съезд собрался в Кремле, но не в Свердловском зале, а в помещении поменьше, которое вмещало около 80 делегатов. Там царила атмосфера интимности и идейной близости.

В день открытия съезда неожиданным и приятным для всех сюрпризом стало появление Ленина. Приветствуя делегатов, Ленин сразу же сказал, что самым существенным сейчас является отношение народов Востока к империализму и революционное движение этих народов. Он заявил, что «революционное движение народов Востока может сейчас получить успешное развитие, оно может получить разрешение не иначе, как в непосредственной связи с революционной борьбой нашей Советской республики против международного империализма»\footnote{\textit{В. И. Ленин}. Полн. собр. соч., т. 39, стр. 318.}.

Подробно рассказан о победах Красной Армии и переломе и ходе гражданской войны, он добавил, что «те победы, которые мы одерживаем теперь над Колчаком, Юденичем и Деникиным, эти победы означают наступление новой полосы и истории борьбы мирового империализма против стран и наций, поднявших борьбу за свое освобождение». И далее: «Я думаю, что то, что проделала Красная Армия, ее борьба и история победы будут иметь для всех народов Востока гигантское, всемирное значение»\footnote{Там же, стр. 319, 321.}.

Остановившись на причинах, приведших к этим победам, и на значении внутреннего разложения империализма, Ленин основное внимание сосредоточил на определении характера борьбы и задачах народов Востока, пути решения которых, как он выразился, «вы не найдете ни в одной коммунистической книжке, но решение которых вы найдете в общей борьбе, которую начала Россия»\footnote{Там же, стр. 330.}. В России удалось соединить борьбу рабочих с борьбой крестьян. То же самое, хотя и в других формах, надо осуществить в международном масштабе.

Важнейший теоретический вывод Ленина сводился к тому, что «социалистическая революция не будет только и главным образом борьбой революционных пролетариев в каждой стране против своей буржуазии, — нет, она будет борьбой всех угнетенных империализмом колоний и стран, всех зависимых стран против международного империализма»\footnote{Там же, стр. 327.}.

Вот почему, говорил Ленин, обращаясь к представителям народов Востока, — вам предстоит и истории развития мировой революции, которая будет, судя по началу, продолжаться много лет и потребует много трудов, — вам предстоит к революционной борьбе, и революционном движении сыграть большую роль и слиться в этой борьбе с нашей борьбой против международного империализма... слиться в общей борьбе с пролетариями других стран»\footnote{\textit{В. И. Ленин}. Полн. собр. соч., т. 39, стр. 328, 330.}.

Трудящиеся массы стран Востока должны знать, «что международный пролетариат является единственным союзником всех трудящихся и эксплуатируемых сотен миллионов народов Востока»\footnote{Там же, стр. 330.}.

Особый упор делал Ленин на поиски своеобразных форм этого союза передовых пролетариев всего мира с живущими часто в средневековых условиях трудящимися массами Востока. Он подчеркивал: «Здесь перед вами стоит задача, которая не стояла раньше перед коммунистами всего мира: опираясь на общекоммунистическую теорию и практику, вам нужно, применяясь к своеобразным условиям, которых нет в европейских странах, суметь применить эту теорию и практику к условиям, когда главной массой является крестьянство...»\footnote{Там же, стр. 329.}.

Немеркнущее, основополагающее значение имеют эти мысли Ленина об органической связи пролетарской революции с борьбой угнетенных народов против империализма. Выступая на съезде коммунистов, представляющих мусульманские народы России, Ленин прямо обращается ко всем народам Востока, открывая перед ними перспективы их революционной борьбы при поддержке Страны Советов и мирового пролетариата.

Сегодня мы видим, как победа первого социалистического государства над фашизмом способствовала росту революционных сил в странах Запада и Востока. Образовалась мировая система социализма. Крах старой колониальной системы и победа ранее угнетенных народов в национально-освободительной борьбе произошли именно в силу единства действий и взаимной помощи мировой системы социализма, пролетариата капиталистических стран и угнетенных народов, в тяжелых боях завоевывающих свою независимость, благодаря союзу этих трех великих революционных сил современного общественного развития.

Насколько свежи и актуальны ленинские мысли через полвека, в новой международной обстановке, видно из того, что они лежат в основе документов международного Совещания коммунистических и рабочих партий, состоявшегося и Москве в 1969 году. Насколько далеки от этих ленинских идей потуги тех, кто, называя себя марксистами-ленинцами, выдвигают новоиспеченную теорию победы мировой революции путем «окружения мирового города мировой деревней»!

Ленин подчеркивал возрастающее значение народов Востока в мировой революции. Он не только говорил о том, что Советская Россия окажет народам Востока помощи в революционном движении, но что развертывающаяся борьба народов Востока против империализма является большой помощью пролетариату Запада в его борьбе за победу революции в их странах.

Выступление Ленина произвело огромное впечатление на присутствующих. На меня — особенно большое, потому что я впервые слушал публичное выступление Ленина. Я отметил себе, что Ленин не коснулся конкретных вопросов, которые волновали делегатов и которые они должны были обсудить и принять по ним определенные решения. Накануне он проводил специальное совещание с делегатами, знал те вопросы, которые их волновали. Своим же выступлением Ленин как бы давал понять, что все конкретные проблемы можно решить только после того, как будет внесена полная ясность в поставленные жизнью новые вопросы революционной борьбы.

Я был еще на нескольких заседаниях этого съезда. Он продолжался почти 10 дней.

У меня было стремление поучиться у выступавших товарищей, перенять все ценное из их опыта разрешении национального вопроса у мусульманских народов Поволжья, Урала. Об их жизни я почти ничего не знал.

В теории национального вопроса российские коммунисты были вооружены, кап ни одна другая партия. Накануне и в ходе первой мировой войны Ленин всесторонне разработал национальный вопрос, который как раз к тому времени приобрел особенно актуальное значение. В частности, он разъяснил положение Программы партии 1903 года о праве наций на самоопределение, предусматривающее требование о самоопределении вплоть до отделения. Это, конечно, как подчеркивал и сам Ленин, вовсе не означало, что большевики рекомендуют или проповедуют отделение. Ленин писал, что пролетариат, «признавая равноправие и равное право на национальное государство... выше всего ценит и ставит союз пролетариев всех наций, оценивая под углом классовой борьбы рабочих всякое национальное требование, всякое национальное отделение»\footnote{\textit{В. И. Ленин}. Полн. собр. соч., т. 25, стр. 274–275.}. Но Ленин твердо стоял за то, чтобы населявшие Россию нации добились права на самоопределение вплоть до отделения от России и сами свободно решали, оставаться ли им в рамках единого государства или выделиться в самостоятельное государство.

Ленин был твердо за то, чтобы народы колоний получили право на самоопределение и отделение от метрополий. Это был лозунг свободы и независимости колоний, и он был куда весомее многих сентиментальных, но беспредметных речей об «уравнении прав» и «восстановлении гуманизма» по отношению к угнетенным народам.

Ведь европейская правая социал-демократия не выступала за свободу колоний и за право их отделения от метрополий — настолько она была привязана к империалистической политике буржуазии своих стран. Оппортунисты предпочитали разглагольствовать о «культурной автономии», как это делали австрийские социал-демократы. Ленинские работы в этой области были крепким ударом и по правой социал-демократии.

Как и многие другие товарищи, я, например, еще до революции прочитал статьи Ленина по национальному вопросу, которые были напечатаны в журнале «Просвещение» в 1913–1914 годах. В марте 1918 года Шаумян опубликовал в Баку в «Известиях Совета рабочих и солдатских депутатов», а также в «Бакинском рабочем» письмо Ленина к нему от 1913 года по национальному вопросу. Шаумян поместил рядом и свои комментарии к этому письму, подчеркивая, что «национальный вопрос является одним из важнейших вопросов нашей революции», а также, что «взгляды, защищаемые нашей партией в настоящее время, в точности высказывались Лениным еще в 1913 году»\footnote{\textit{С. Г. Шаумян}. Избр. произв., т. 2. Госполитиздат, 1958, с. 202.}. В письме Ленина серьезной критике подверглись некоторые взгляды самого Шаумяна по национальному вопросу, в частности по вопросу о государственном языке.

Помню еще, какое сильное впечатление на всех нас в Баку произвела телеграмма Ленина Шаумяну от 14 мая 1918 года, в которой, высоко оценивая деятельность Шаумяна как чрезвычайного комиссара по делам Кавказа, Ленин писал:

«Мы и восторге от вашей твердой и решительной политики. Сумейте соединить с ней осторожнейшую дипломатию, предписываемую, безусловно, теперешним труднейшим положением, — и мы победим.

Трудности необъятны. Пока нас спасают только противоречия и конфликты и борьба между империалистами. Умейте использовать эти конфликты: пока надо научиться дипломатии»\footnote{\textit{В. И. Ленин}. Полн. собр. соч., т. 50, стр. 73–74.}.

Призыв «научиться дипломатии» — это было что-то новое для нас, молодых коммунистов, в тот период революционной бури отметавших все «буржуазное».

Само слово «дипломатия» мы третировали как недостойное революционеров. И вдруг — «научиться дипломатии», использовать противоречия между империалистами! Хорошо помню, какое глубокое впечатление это тогда произвело. Но становилось с каждым днем очевиднее: став у руля государственного управления, Ленин не только учил других брать на вооружение дипломатию, но сам в совершенстве владел этим оружием, поставил его на службу интересам Советской страны и международного революционного движения, вносил в дипломатию неизвестные до того, присущие только социализму элементы и принципы.

И национальном вопросе жизнь выдвигала все новые и новые общие и частные проблемы. Дело постепенно шло к совершенно новому по своей форме и по своей сущности государственному образованию — союзу социалистических республик Все чаще то там, то здесь возникали споры вокруг поисков форм самоопределения, взаимоотношений между советскими республиками, понимания сущности многонационального государства, автономии отдельных народов, федерации и т. д.

По мере возникновении и проявления практических потребностей пролетарской революции, в ходе исторических событий Ленин непрестанно развивал свое понимание национального вопроса, уточнял программные требования и находил наиболее подходящие формы их разрешения.

В упомянутом выше письме 1913 года Шаумяну Ленин отрицательно высказался в отношении предложения Шаумяна о федерации как форме объединения населявших Россию наций, не желающих отделиться от нее. Видимо, он опасался, что такая форма может создать, в частности, помехи в развитии единой экономики. Однако после Февральской революции Ленин непосредственным образом ощутил небывалый и повсеместный рост национального сознания у народов, населявших Россию. Уже на I съезде Советов в июне 1917 года, в преддверии социалистической революции, Ленин, бичуя великодержавную шовинистическую политику эсеров и меньшевиков, провозгласил лозунг: «Пусть Россия будет союзом свободных республик»\footnote{\textit{В. И. Ленин}. Полн. собр. соч., т. 32, стр. 286.}.

Это было новое, очень важное, поворотное ленинское слово в решении национального вопроса народов России.

Здесь истоки и того, что Советская социалистическая Россия в начале 1918 года была объявлена федеративной республикой. Высказанное Лениным еще в 1917 году предвидение нашло свое воплощение в создании Союза Советских Социалистических Республик — новой, невиданной в истории формы самоопределения наций, их тесного единения и братского сотрудничества на основах равноправия.

У нас на Кавказе довольно длительное время дебатировался вопрос о правомерности провозглашения независимых национальных советских республик. Коммунисты Баку уже весной 1919 года выдвинули лозунг создания независимого Советского Азербайджана, который был бы теснейшим образом связан с РСФСР, а также считали правильным, как я уже рассказывал, образование Компартии Азербайджана, которая входила бы и гостив РКП(б). Такие же соображения были пригодными, как нам представлялось, и для Грузии, и для Армении. Другие же партийные организации, объединявшиеся Кавказским крайкомом, как и некоторые члены крайкома, были против этого, считая, что-де такое решение вопроса может подорвать межнациональную сплоченность трудящихся.

Я привел примеры, более мне знакомые, чем другие, с целью показать, каким жгучим был национальный вопрос в то время, какие и вокруг чего шли споры, поиски. Примерно такого же порядка вопросы, подчас в еще более острой форме, возникали и к других районах страны. Ото отражалось и на ходе съезда коммунистических организаций народов Востока.

Очень острой, можно сказать, беспощадной критике подверглась в прениях деятельность Центрального бюро коммунистических организаций народов Востока при ЦК РКП(б), неудовлетворительно работавшего во всех отношениях. Было видно, в каких тяжелых условиях шла работа в ряде губерний Поволжья и Урала, где в результате нашествия Колчака, чехословацкого корпуса, атамана Дутова и др. неоднократно сменялись власти.

Были очень резкие споры и столкновения взглядов выступавших. Чувствовалось, что налицо большой разнобой во взглядах при разрешении конкретных вопросов национальной политики, и дело дошло даже до создания группировок, занимавших непримиримые позиции. В частности, это касалось отношений между татарами и башкирами.

Я шел на съезд с лучшими представлениями о положении дел. Многое для меня явилось полной неожиданностью. Первые дни на съезде произвели на меня тяжелое впечатление. Но я был доволен тем, что многое узнал о фактическом положении дел в этих районах.

В числе других ораторов выступил и представитель московской группы азербайджанского «Гуммета» Исрафилбеков. Он, между прочим, сказал о моем приезде в Москву, а также о моем заявлении со ссылкой на мнение бакинских товарищей, что под флагом партии «Гуммет» не представится возможности успешно работать в Азербайджане.

Следует пояснить, что вскоре после приезда в Москву я имел встречи с Наримановым п другими активными работниками «Гуммета», которые тоже были в то время здесь. Их я информировал обстоятельно о положении дел и задачах, стоящих и Азербайджане. Особенно подробно аргументировал выработанную у нас в Баку точку зрения о необходимости создания единой Компартии Азербайджана. Вопреки моим ожиданиям, они согласились с этим. Это упрощало решение вопроса.

В целом съезд явился хорошей школой для молодых национальных кадров народов Востока, населявших нашу страну, содействовал сближению и дальнейшей разработке важных вопросов национального строительства и взаимоотношений между нациями и представлял особую ценность благодаря новым, весьма назревшим и глубоким мыслям Ленина об особенностях и ходе революционного и национально-освободительного движения в странах, отставших в своем развитии.

\subsection*{\center * * *}

Менее чем через год удалось собрать I съезд народов Востока, в котором участвовали уже и представители многих других стран. Решение провести этот съезд в Баку было принято в конце июня 1920 года Исполкомом Коминтерна и ЦК РКП(б). Было создано организационное бюро по подготовке съезда, в которое вошли Серго Орджоникидзе, Елена Стасова, Нариман Нариманов и я.

Это был съезд не только коммунистов, но и беспартийных, антиимпериалистически настроенных национальных деятелей, представителей широких масс трудящихся и их организаций. В нем приняли участие представители 37 национальностей, в том числе из Афганистана, Египта, Индии, Персии, Китая, Кореи, Сирии, Турции и других стран, треть из них не была коммунистами.

Съезд начал свою работу 1 сентября 1920 года. Накануне собрался Бакинский Совет, на котором было доложено о проведенной работе по подготовке к съезду. К вечеру заседание Совета подходило к концу, но так как ожидалось прибытие поезда с представителями Коминтерна и делегациями компартий ряда стран, то депутаты решили не расходиться и организованно встретить гостей. Ночью прибыл поезд. Все депутаты Совета, многочисленные рабочие делегации, несмотря на поздний час, пришли на торжественную встречу. Прямо с вокзала гостей привезли в театр, где возобновилось заседание Совета. Закончилось оно в пятом часу утра. Слушали приветственные выступления представителей Коминтерна. Собрание проходило с большим подъемом, говорили вдохновенно. Мы, бакинцы, гордились тем, что у нас, в пролетарской цитадели, через четыре месяца после победы Советской власти будет заседать съезд народом Востока

Съезд работал неделю и принял ряд важных решений, выразил солидарность с тезисами состоявшегося за месяц до итого II конгресса Коминтерна по национальному и колониальному вопросам. Были одобрены и вскоре опубликованы воззвание к народам Востока с призывом к борьбе против колонизаторов и воззвание к трудящимся Европы, Америки и Японии с призывом поддержать освободительное движение народов Востока. На съезде выступили представители компартий Европы и Америки: Бела Кун (Венгрия), Джон Рид (США), Томас Квелч (Англия), Россмер (Франция), представители многих народов Востока.

Ленин дал весьма высокую оценку этому съезду. В своей речи 15 октября 1920 года на совещании председателей уездных, волостных и сельских исполнительных комитетов Московской губернии, говоря о II конгрессе Коминтерна и съезде народов Востока в Баку, он отметил, что «это те международные съезды, которые сплотили коммунистов и показали, что во всех цивилизованных странах и во всех отсталых восточных странах большевистское знамя, программа большевизма, образ действий большевиков — есть то, что для рабочих всех цивилизованных стран, для крестьян всех отсталых колониальных стран является знаменем спасения, знаменем борьбы, что действительно Советская Россия за эти три года не только отбила тех, кто бросался, чтобы ее душить, но и завоевала себе сочувствие трудящихся во всем мире, что мы не только разбили наших врагов, но мы приобрели и приобретаем себе союзников не, по дням, а по часам»\footnote{\textit{В. И. Ленин}. Полн. Собр. соч., т. 41. стр. 357.}.

\newpage
\section[Ленин на VIII Всероссийской партийной конференции: мы стоим накануне крупнейшего перелома к лучшему]{\center \textbf{Ленин на VIII Всероссийской партийной конференции: мы стоим накануне крупнейшего перелома к лучшему}}


2 декабря 1919 года в Москве начала работать VIII партийная конференция. Она продолжалась три дня, в течение которых было проведено шесть заседаний.

Я был очень обрадован, получив возможность впервые участвовать в таком высоком партийном форуме. Все, что говорилось на конференции, было для меня ново, новыми были люди, дела, которые их волновали, вопросы, которые обсуждались. Обстановка на конференции была товарищеская, простая, и это способствовало откровенному обсуждению волнующих проблем.

Самое большое, захватывающее впечатление произвело на меня активное участие Ленина в работе конференции. Несмотря на исключительную занятость, в дни решающих боев на фронтах, перегруженный государственной и партийной работой, Ленин много времени и усилий уделял успешному проведению конференции. Он открыл конференцию кратким вступительным словом и вел все первое заседание; на втором заседании сделал доклад о политической работе ЦК и затем, после прений, выступил с заключительным словом; при обсуждении украинского вопроса Ленин дважды брал слово.

С самого начала конференции и до ее конца я весь превратился в слух, не пропускал ни одного выступлении, ни одного высказывания. Нет нужды объяснять, какое воздействие на меня и на всех нас произвели речи Ленина.

Делегаты говорили просто, подчас грубовато, но доходчиво, по существу, убежденно, и это производило на всех глубокое впечатление. Особенно на тех, кто, как и я, попал на такое совещание впервые и чувствовал себя взволнованно.

Ленин своим поведением, своим отношением к товарищам, своим чрезвычайным вниманием окидывал такое влияние на нею обстановку, что все выступавшие — и не только опытные, но и такие, как я, молодые коммунисты, — говорили, не опасаясь провала, без напряжения, с чувством и уверенностью, что тебя поймут. Ленин располагал к себе не только своей доступностью, но и всем существом своим. Ленин выступать умел так, что у слушателей создавалось впечатление, будто оратор не «держит речь», а запросто с тобой разговаривает.

По количеству делегатов VIII конференция была малочисленной: с решающим голосом нас было 45 делегатов и с совещательным — 73 человека. Это объяснялось тем, что далеко не во всех губерниях и краях закончилась гражданская война. К тому же некоторые районы только недавно получили свободу, не успели еще организоваться и послать делегатов на конференцию. Связь с Советским Туркестаном, например, была прервана войсками генерала Дутова.

Во вступительном слове Ленин сообщил, что из-за тяжелого военного положения, которое переживала страна еще несколько месяцев назад, конференция собралась с опозданием против срока, установленного решением VIII съезда партии, — созывать конференции один раз в три месяца.

«Мы созываем теперь конференцию... — говорил Владимир Ильич, — в такой момент, когда нам удалось достигнуть гигантского улучшения на фронтах и когда мы уверены, что стоим накануне крупнейшего перелома к лучшему как в международном положении, так и в отношении военном...»\footnote{\textit{В. И. Ленин}. Полн. собр. соч., т. 39, стр. 341.}

Какой прозорливостью были преисполнены эти слова Ленина, показали последующие исторические события. Через год кончилась гражданская война, и страна начала переходить к мирному строительству. Ярким проявлением этого явилось принятие и декабре следующего, 1920 года на VIII съезде Советов плана электрификации страны. Но тогда, в декабре 1919 года, всего этого еще нельзя было себе даже представить, особенно тем, кто прибыл прямо из глубокого подполья.

При выборе руководящих органов конференции произошел незначительный инцидент, который тем не менее произвел на меня большое впечатление, так как в нем ярко проявилось отношение Ленина к обязанностям коммуниста. Дело в том, что после оглашения списка лиц, намечаемых в мандатную комиссию конференции, Сапронов заявил о снятии своей кандидатуры. Скорее всего, Сапронов, представлявший оппозицию, оформившуюся позже в группу «демократического централизма», но уже нашумевшую на VIII съезде партии, претендовал на избрание не в мандатную комиссию, а в президиум или секретариат конференции. На это председательствовавший Ленин с горячностью заявил: «Президиум держится того мнения, что отказы недопустимы, а раз недопустимы отказы, то тем более недопустимы мотивировки отказов. Вопрос исчерпан. Значит, утверждается»\footnote{«Восьмая конференция РКП(б). Протоколы». М., 1961, стр. 4.}.

На этом заседании, после принятия порядка дня, были заслушаны два приветствия. Первым приветствовал конференцию Султан-Галиев от имени II Всероссийского съезда мусульманских коммунистических организаций народов Востока. Коммунисты мусульманских народов, говорил Султан-Галиев, понесут на Восток красное знамя, революционное знамя коммунизма, полученное из рук наших учителей — товарищей русских коммунистов, и разбудят спящий Восток.

Затем Ленин объявил: «Слово для приветствия от имени кавказских организаций имеет т. Микоян».

Мне выпала честь приветствовать конференцию от имени подпольного Кавказского краевого комитета РКП(б). Чтобы не излагать своими словами теперь, ровно через 50 лет, содержание того выступления, ограничусь короткими выдержками из протокола конференции. Я говорил: «...ходом событий целые два года весь Кавказ был оторван от центра... После двухгодичного опыта и после хозяйничанья буржуазии на Кавказе, у нас, наконец, намечается свой ход развития революционных действий, и мы на Кавказе стоим накануне Октябрьской кавказской революции... ость твердая надежда, что в ближайшие месяцы кавказские народные массы подымутся с красным знаменем и руках — со знаменем Коммунистической партии и установят Советскую власть по всему Кавказу... Заканчивая свое приветствие, я думаю и подчеркиваю, что партийная конференция... даст особые директивы нам для того, чтобы эта политика велась на местах не на наш собственный риск, а под вашим руководством и при вашем содействии»\footnote{«Восьмая конференция РКП(б). Протоколы», стр. 8–9.}.

И то, что обсуждалось в декабре 1919 года, начало осуществляться уже через полгода. Революция победила сперва в Азербайджане, в апреле — мае 1920 года, в ноябре — в Армении и в марте 1921 года — в Грузии. Вечернее, второе заседание конференции, как и все последующие, также отличалось большой деловитостью и насыщенностью повестки дня.

Заседание началось докладом Ленина о политической деятельности ЦК, затем Крестинский доложил об организационной работе ЦК, состоялись прения по этим докладам. После заключительного слова Ленина на этом же заседании был заслушан доклад Чичерина о международной политике.

Надо ли говорить, какое впечатление на всех и, конечно, на меня, новичка, произвел доклад Ленина, его четкий, острый анализ двухлетних итогов существования нового государства и вытекающих отсюда задач для дальнейшей практической деятельности.

Он говорил, что самые большие трудности у нас уже позади, что потерпела крах попытка Антанты удушить Советскую республику собственными войсками, потому что иностранные рабочие и крестьяне в военной форме отказались воевать против Советской власти.

На меня, как и на других делегатов, большое впечатление произвело то место выступления Ленина, в котором он говорил, что высадка французских войск на юге России кончилась рядом восстаний французских матросов. В настоящее время, несмотря на свирепую военную цензуру, которая, хотя войны и нет, продолжает существовать в якобы свободных странах — Англии и Франции, из доходящих до нас отдельных экземпляров газет видно, что сведения о восстании матросов на французских военных судах в Черном море попали во французскую прессу.

Вся коммунистическая французская и английская пресса сообщает о том, что несколько французских матросов осуждено и сослано на каторгу. Ленин сказал, что «имя тов. Жанны Лябурб, которую французы расстреляли в Одессе за большевистскую агитацию, стало лозунгом для французской социалистической рабочей печати не только коммунистического крыла...» Газета «Юманите» «имя Лябурб сделала лозунгом борьбы против французского империализма, за невмешательство в дела России»\footnote{\textit{В. И. Ленин}. Полн. собр. соч., т. 39, стр. 344, 345.}.

Жанна Лябурб, дочь деятеля Парижской коммуны, жившая в России, участвовала в революции как большевичка. Организовала в России французскую коммунистическую группу, с февраля 1919 года в Одессе она, работая в подполье, вела страстную пропаганду среди французских солдат и матросов против интервенции, за коммунизм. В марте, в возрасте 40 лет, она была арестована и расстреляна командованием французских оккупационных войск.

Это событие, видимо, произвело на Ленина глубокое впечатление. О Жанне Лябурб он также подробно и с любовью расскажет через несколько дней на VII съезде Советов и вновь подчеркнет, что «это имя стало известно всему французскому пролетариату и стало лозунгом борьбы, стало тем именем, вокруг которого все французские рабочие, без различия казавшихся столь трудно преодолимыми фракционных течений синдикализма, — все объединились для выступления против международного империализма»\footnote{Там же, стр. 391.}.

Ленин говорил и о том, что провалились попытки Антанты мобилизовать против нас маленькие соседние государства, находившиеся в зависимости от Антанты и враждебно относившиеся к большевикам. Буржуазные правительства этих государств поняли, что мы являемся для них если не союзниками, то соседями более надежными, чем империалисты, которые их грабят. Поэтому они заняли позицию нейтралитета и не поддались давлению Антанты.

Из победы над Колчаком и Деникиным, резюмировал Ленин, «мы имеем право сделать тот вывод, который является самым существенным для нас и во всей нашей деятельности должен нами руководить: исторически побеждает тот класс, который может вести за собой массу населения»\footnote{\textit{В. И. Ленин}. Полн. собр. соч., т. 39, стр. 350–351.}. Он объяснял, что крестьянство стало на сторону большевиков потому, что ему пришлось на практике испытать власть Колчака и Деникина, и оно сделало выбор в нашу пользу, окончательно стало на нашу сторону.

Уже в то время, когда еще шла гражданская война, Ленин прозорливо видел, что этот главный политический урок — переход на нашу сторону миллионного крестьянства «мы должны применить к тем задачам внутреннего строительства, которые теперь, по завершении нашей победы над Деникиным, будут стоять на очереди дня...»\footnote{Там же, стр. 354.}.

Ярко прозвучала в речи Ленина необходимость как можно больше сил уделить продовольственному вопросу: «это — основа всего»\footnote{Там же, стр. 358.}. На данном этапе, говорил он, разверстка хлеба должна лечь в основу нашей деятельности, это есть ссуда крестьянства «голодному рабочему».

Коротко коснулся Ленин партийных задач. Он сказал, что, с одной стороны, естественно, когда к правящей партии стремятся примазаться худшие элементы с целью использовать пребывание в ней в своих личных интересах, так как это партия правящая. С другой стороны, только передовая часть рабочего класса, его авангард в состоянии вести свою страну. Поэтому партия берет в свои ряды прежде всего тех, кто готов идти на фронт или ведет себя и тылу, прежде всего в труде на субботниках, как настоящий коммунист.

Говоря о перерегистрации членов партии, Ленин заявил: «Мы не боимся приступать к изъятию тех, кто не вполне надежен. Мы достигаем этого также тем, что доверяем члену партии, приходящему к нам в трудную минуту. Те члены партии, как показывает сегодняшний отчет Центрального Комитета, те тысячи и сотни тысяч, которые приходили к нам, когда Юденич стоял в нескольких верстах от Петрограда, а Деникин к северу от Орла, когда вся буржуазия уже ликовала, эти члены партии заслуживают нашего доверия. Такое расширение партии мы ценим»\footnote{\textit{В. И. Ленин}. Полн. собр. соч., т. 39, стр. 361.}. Ленин призвал к работе среди беспартийных рабочих и крестьян, потому что это вернейшие наши друзья, их поддержка нам обеспечена.

Отчет ЦК об организационной работе был предварительно опубликован в «Известиях ЦК РКП(б)», и поэтому Крестинский в своем докладе лишь вкратце остановился на основных организационных задачах партии.

После проведенной перерегистрации с целью очистки партии от нежелательных элементов почти вдвое сократилось число ее членов. Но, широко открыв двери рабочим, крестьянам и красноармейцам путем организации «партийной недели», партия обеспечила приток более 200 тыс. новых надежных коммунистов, вступивших в партию в тяжелые для Советской России дни. Всего в партии к VIII конференции было около 350 тыс. человек. Эти цифры, однако, не охватывали коммунистов многих краев и районов, где еще не закончилась гражданская война (Сибирь, Украина, Кавказ и др.).

В прениях по отчету ЦК выступил и я. Мне хотелось рассказать товарищам о том, что на Кавказе мы на своей шкуре испытали оба метода борьбы империалистов против нас, о которых говорил Ленин (как непосредственное вторжение войск империалистов, так и использование против нас маленьких государств, возникших на окраинах России).

Далее я счел нужным обратить внимание конференции на необходимость налаживания тесной связи с зафронтовыми и действующими в подполье организациями, которые ждут помощи и нуждаются в директивных указаниях ЦК. Я говорил: «Необходимо посылать своих работников не в момент завоевания нами той или другой области, а в момент, когда эта область изнывает под игом противника. Надо, чтобы наши работники к моменту завоевания той или другой области уже слились с массой, работали там, и не были бы чужим элементом, когда перейдет власть к ним в руки. За это время молено узнать, кто хороший работник, кто идет на опасность, на риск. Этого не было сделано по отношению ко многим окраинам. Это необходимо сделать по отношению к Кавказу и другим областям»\footnote{«Восьмая конференция РКП(б). Протоколы», стр. 39.}.

В прениях не было сделано ни одного критического замечания по политическому отчету ЦК. Даже оппозиционер Сапронов, выступая по докладу Ленина, сказал, что «мы привыкли в политических вопросах идти за Лениным». Но он занял особую позицию по организационному вопросу в духе линии оппозиционной группы «демократического централизма».

Ввиду этого Ленин, намеревавшийся было после прекращения прений отказаться от заключительного слова, счел нужным выступить по поводу предложения Сапронова об исключении из Проекта инструкции местным комитетам РКП о работе в деревне пункта об оказании совхозами и коммунами помощи окрестным крестьянам. Ленин заострил политическую сущность этого попроси: «Отношения коммун и советских хозяйств к окружающему крестьянству — это один из самых больных вопросов всей нашей политики... Сейчас мы идейно завоевали сибирского Крестьянина, освободив его от Колчака. Но это не будет прочно, если мы не сумеем поставить этого дела так, чтобы этому крестьянину оказывалась реальная помощь...»\footnote{\textit{В. И. Ленин}. Полн. собр. соч., т. 39, стр. 365.}. Возражал Ленин и против высказываний Сапронова, направленных на умаление роли партийных организаций в руководстве советской работой, на упразднение при волостных партийных комитетах организаторов по работе в деревне.

Конференции приняла важные решения в области партийного и советского строительства. Она подробно обсудила и утвердила новый Устав партии, проект которого был опубликован за несколько месяцев до этого в печати. Устав учитывал опыт партийного строительства в условиях первых двух лет Советской власти. Это был первый Устав легальной партии, правящей коммунистической партии, руководящей силы Советского государства. Новым в нем было, в частности, создание впервые института кандидатов в члены партии. После некоторых споров был установлен кандидатский стаж для рабочих и крестьян не менее двух месяцев, а для остальных не менее шести месяцев.

При обсуждении Устава партии я выступил в прениях. Я считал необходимым поднять вопрос о том, что проект Устава не отражает фактического положения в той организационной работе, которая была проведена на Украине, на Кавказе, в Прибалтике и других районах. В Уставе предусматривалось организационное построение губернских, уездных и волостных парторганизаций. А на Украине, у нас на Кавказе, в Туркестане и в некоторых других местах давно существовали объединения более высокого звена — краевые и областные комитеты. Целесообразность внесения в Устав раздела о краевых организациях я мотивировал так: «Они не только соответствуют определенным культурным, экономическим, этнографическим и другим условиям, но главным образом были и административными объединениями, которые отчасти остаются и до сих пор, отчасти перестроились. Краевой комитет объединяет все губернские, областные и окружные организации, находящиеся на территории Кавказа. Необходимость этого объединения показала сама жизнь... ЦК вряд ли может ориентироваться в тех сложных взаимоотношениях, которые существуют на Кавказе между отдельными губернскими, областными и другими организациями. Только путем создания краевого кавказского комитета, который сплачивает все партийные ряды на Кавказе, там может установиться одна линия поведения»\footnote{«Восьмая конференция РКП(б). Протоколы», стр. 144–145.}.

Мне представлялось существенным оттенить также вопрос о том, что работа партийных организаций, находящихся в оккупированных районах, по ту сторону фронта, протекала в подпольных условиях и полное применение нового Устава для них не всегда возможно. Поэтому я считал необходимым ввести в Устав примечание о том, что при работе в нелегальных условиях допускаются те или иные отступления от партийного Устава.

Несмотря на настойчивость, мне не удалось убедить конференцию.

Представляло интерес обсуждение вопроса о работе с новыми членами партии, которые большой массой были приняты в «партийную неделю», в дни величайшей опасности для Советской власти, когда новый слой рабочих и крестьян перед лицом угрозы деникинского наступления на Москву вступил в партию, а многие из них сразу лее пошли на фронт.

В ряде организаций большой процент, а в некоторых и половину составляли молодые коммунисты, вполне достойные быть в рядах партии, но не имеющие еще опыта партийной работы, необходимых знаний, партийной закалки, навыков партийной дисциплины. Конференция наметила ряд практических мероприятий по просвещению, вовлечению в советскую и партийную работу молодых товарищей, по выполнению ими партийных поручений, воспитанию дисциплины. Важное место в работе конференции занял вопрос о политике партии по отношению к трудовому крестьянству Украины, освобождающемуся от деникинских банд. Задачи, стоявшие перед партией в деле строительства Советской власти на Украине, были чрезвычайно сложными и трудными, так как условия украинской деревни были своеобразные: основную массу составляло среднее, сравнительно зажиточное крестьянство. В первые шесть месяцев существования советской государственности на Украине тамошние коммунисты допустили ошибки в отношении среднего крестьянства, оттолкнувшие его от Советской власти. Но захват Деникиным украинских районов и введенный им режим вызвали резкое недовольство крестьянства, оно поднялось против Деникина, и в нем усилились настроения в пользу Советской власти. Среднее крестьянство, на опыте проверившее две исключающие друг друга системы власти, стало склоняться и сторону Советов. На Украине развернулось широкое антиденикинское движение, хотя и повстанческих отрядах было немало анархистских, антисоветских тенденций. Многие из этих отрядов превратились в банды.

Теперь, когда Красная Армия успению продвигалась на Южном фронте, отбросив деникинские и петлюровские войска, когда на Украине начала восстанавливаться Советская власть, потребовалась ясная и определенная тактика и политика для ликвидации анархистско-противосоветского повстанчества, а главное, правильная политика, направленная на завоевание сочувственного отношения к нам среднего украинского крестьянства.

Ленин предупреждал, что отношения с основной массой крестьян «на Украине могут быть второй раз испорчены, если мы не сумеем провести своей политической линии в жизнь»\footnote{\textit{В. И. Ленин}. Полн. собр. соч., т. 39, стр. 364–365.}.

ЦК разработал и Ленин лично написал проект решения о политике коммунистов на Украине, который и был предметом обсуждения на конференции. Украинские товарищи прямо не спорили с ленинским проектом, не критиковали его, но в выступлениях некоторых из них чувствовалось непонимание ряда вопросов, а иногда отстаивание старых, ошибочных позиций. Это относилось главным образом к желанию сохранить в руках государства помещичьи земли и перенести центр тяжести в крестьянском вопросе на подавление кулака.

Ленин критиковал украинских товарищей за то, что они не всегда последовательно проводят линию ЦК партии. Высказываясь против «самостийности», Ленин подчеркивал, что имеет в виду самостийность не государственную, а партийную.

Главным в решении ЦК и в выступлениях Ленина по вопросу об отношении к крестьянству было стремление привлечь на советскую сторону не только бедняцкое, но и середняцкое крестьянство Украины.

Выступавший на съезде Раковский, выражая недоверие к крестьянству, отстаивал по существу старую политику сохранения помещичьих имений в виде совхозов. Он не понимал, что специфические условия Украины, необходимость достижения, как говорил Ленин, «блока с украинским крестьянством» требовали раздела между мелкими крестьянами Украины большей части бывших помещичьих имений. Ленин говорил: «По поводу речи т. Раковского приходится сказать, что когда он заявил, что советские хозяйства должны быть основанием нашего коммунистического строительства, то это неверно. Никоим образом так ставить дело мы не можем. Мы должны признать, что только весьма небольшую часть культурных хозяйств мы должны дать под советские хозяйства, иначе мы блока с мелким крестьянством не получим, а нам этот блок необходим... поэтому большую часть советских хозяйств нам нужно отдать для фактического раздела»\footnote{\textit{В. И. Ленин}. Полн. собр. соч., т. 39, стр. 171.}.

Было отчетливо видно, что в целом делегаты полностью одобряют политическую линию, развитую Лениным. В подтвержденной конференцией резолюции ЦК, составленной Лениным, записано, что «задачей Советской власти на Украине является завоевание к себе доверия со стороны не только крестьянской бедноты, но я широких слоев среднего крестьянства, которое своими подлинными интересами теснейшим образом связано с Советской властью»\footnote{«Восьмая конференция РКП(б). Протоколы. Стр. 190.}.

В резолюции о Советской власти на Украине особое значение имеет первый пункт, в котором указывается: «...ЦК считает необходимым еще раз подтвердить, что РКП стоит на точке зрения признания самостоятельности УССР»\footnote{Там же, стр. 189.}. Важным являлся также пункт резолюции о необходимости правильного подхода к украинскому языку и культуре. В нем говорится:

«...ЦК РКП вменяет в обязанность всем членам партии всеми средствами содействовать устранению всех препятствий к свободному развитию украинского языка и культуры. Поскольку на почве многовекового угнетения в среде отсталой части украинских масс наблюдаются националистические тенденции, члены РКП обязаны относиться к ним с величайшей терпимостью и осторожностью... всячески противодействуя попыткам искусственными средствами оттеснить украинский язык на второй план, стремясь, наоборот, превратить украинский язык в орудие коммунистического просвещения трудовых масс»\footnote{Там же, стр. 189–190.}.

Странным было поведение Троцкого ни партийной конференции. Он не выступил ни по одному вопросу в поддержку липни ЦК, не проронил буквально ни слова. И лишь когда Бубнов сказал, что одной из центральных задач является очищение Украины от банд, он бросил загадочную реплику: «Выгнать на украинском языке»\footnote{«Восьмая конференция РКП(б). Протоколы», стр. 103.}, на что Бубнов сразу же ответил: «На каком угодно языке — только выгоните»\footnote{Там же.}. Такая реплика Троцкого отражала его пренебрежение к вопросу о национальных языках и фактически шла вразрез с одним из важнейших пунктов резолюции о Советской власти на Украине.

Резолюция ЦК и выступления Ленина по украинскому вопросу сыграли решающую роль в разработке и проведении в дальнейшем Компартией Украины правильной политики как в национальном вопросе, так и по отношению к крестьянству.

VIII партийная конференция имела значение большее, чем обычная, хотя бы потому, что она приняла новый Устав партии, что обычно относится к функции съезда. Она всей своей работой как бы обобщила двухлетний опыт партии как партии правящей в условиях победы советского строя. Она наметила задачи партии на период завершения гражданской войны, определила пути дальнейшего укрепления диктатуры пролетариата и упрочения союза рабочего класса и трудового крестьянства.

Конференция проходила накануне очередного съезда Советов и обсудила поэтому основы проектов постановлений этого съезда в области советского и хозяйственного строительства, вопросов военных и внешнеполитических. Ряд предложений делегатов конференции — о работе Советов, об уточнении Конституции и др. — был передан фракции съезда Советов для дальнейшего обсуждения.

\newpage
\section[Ленин на VII Всероссийском съезде Советов: три победы над Антантой и три простейшие задачи, создающие возможность построить социалистическую республику]{\center \textbf{Ленин на VII Всероссийском съезде Советов: три победы над Антантой и три простейшие задачи, создающие возможность построить социалистическую республику}}


На следующий день после окончания VIII партийной конференции, 5 декабря 1919 года, в здании Большого театра открылся VII Всероссийский съезд Советом. Во вступительном слове Калинин сказал, что этот съезд является самым многочисленным из всех съездов и полнее, чем предыдущие съезды, представляет страну. Он привлек внимание делегатов к печальному событию — гибели в начале года Карла Либкнехта и Розы Люксембург в Берлине. Они состояли почетными членами ВЦИК. Калинин сказал и о великой потере нашей — кончине Якова Свердлова, Председателя ВЦИК. Он упомянул также имена других членов ВЦИК, которых мы потеряли. По его призыву съезд почтил память павших товарищей вставанием.

На съезде из 1366 делегатов подавляющее большинство составляли коммунисты. Было несколько десятков беспартийных и несколько депутатов с мест, избранных от мелких прокоммунистических партий и групп, возникших из осколков распадавшихся мелкобуржуазных партий.

Мне представилось крайне интересным, что президиум пригласил для участия к работе съезда по три — пять представителей разных партий и групп. Как правило, эти партии заявляли о признании принципов коммунизма и Советской власти, но в чем-то, одни в большем, другие в меньшем, расходились с нашей партией. Критерием для приглашения на съезд представителей была посылка членов этих партий на фронты гражданской войны. С совещательным голосом в работе съезда участвовали также представители оппозиционных партий — меньшевиков и эсеров.

Заседание съезда началось с горячих приветствий от иностранных коммунистических делегаций, представителей фронтов и Красного флота. Приветствия произносились с большим подъемом, вдохновенно и вызывали бурное одобрение переполненного зала Большого театра. Атмосфера, которая царила на съезде, производила глубокое впечатление на меня, впервые попавшего на заседание верховного органа Советской власти. Были также приветствия от Украины, от пролетариата порабощенной Латвии, от народов Востока и др. Выступали с приветствиями также представители оппозиционных партий, в числе которых был один из видных лидеров меньшевиков — пресловутый Дан, известный соглашатель и оборонец, сыгравший крайне отрицательную роль в 1917 году. Для меня было неожиданным его выступление о поддержке Советской власти. Пользуясь разными эффектными ораторскими приемами, он старался произвести впечатление, заклиная, будто их партия, несмотря на расхождения с большевиками, намерена выступать вместе с ними, когда речь идет о защите революции. Он призывал к единству революционных рядов и жаловался, что для этого, мол, не создаются необходимые условия.

Енукидзе сообщил о многочисленных телеграфных приветствиях с фронтов, от рабочих и крестьян. Съезд получил приветствие от революционной корейской делегации, которая ожидалась на съезде, но члены делегации заболели в пути и задержались в Симбирске. Это была очень теплая телеграмма, в ней говорилось о борьбе корейских революционеров против отечественной буржуазии и японских интервентов, о том, что революционеры мешают японцам вести борьбу против республики Советов и желают полного торжества Советской власти в России. Все мы с особым интересом слушали эти приветствия, ведь мы тогда не представляли себе, что происходило в далекой Корее. Уже к концу съезда представитель Корейской социалистической партии добрался до Москвы и выступил с речью на съезде.

Заслуживает внимания и то, что съезд принял специальную резолюцию, оглашенную секретарем ВЦИК Аванесовым, о Фридрихе Адлере, который был в те времена весьма популярной личностью в международной революционной жизни. Сын лидера австрийских социал-демократов, он пошел к председателю австрийского правительства и на месте застрелил его из пистолета. Из уважения к его отцу, Ф. Адлер не был расстрелян, а приговорен к вечной каторге. Но вскоре революционные события в Австрии открыли перед ним двери тюрьмы. Он стал тогда героем в глазах революционных масс, особенно в нашей стране, был избран почетным членом ВЦИК. В последующее же время, когда в Австрии сложилась революционная ситуация, когда Советы рабочих депутатов готовы были взять власть в свои руки, Адлер, пользуясь своим прошлым авторитетом, сделал все для того, чтобы запугать рабочих, не допустить перехода власти в их руки и сохранить буржуазное правительство. Он стал на путь борьбы с австрийскими коммунистами, что являлось прямой изменой делу пролетариата. И вот с этой мотивировкой он исключался из состава ВЦИК, как недостойный быть его почетным членом.

В повестке дня съезда было много вопросов: доклады ВЦИК и Совнаркома; военное положение; о Коммунистическом Интернационале; продовольственное положение; топливный вопрос; советское строительство; выборы ВЦИК.

По трем практическим вопросам: топливо, продовольствие и советское строительство — прения велись в трех разных секциях.

«Гвоздем» съезда было, конечно, выступление Ленина. С одобрения делегатов он выступил с объединенным отчетом — как от Совета Народных Комиссаров, так и от ВЦИК. Этот свой доклад Ленин построил в основном на положениях, намеченных им в докладе на партийной конференции.

Владимир Ильич аргументированно доказал, что за прошедшие два года гражданской войны Антанта потерпела три поражения в войне с нами, а мы тем самым одержали три большие и прочные победы.

Это — провал попытки Антанты силами своих войск разделаться с Советской Россией; войска пришлось отозвать, ибо солдаты Антанты оказались неспособными вести борьбу против революционной России. «Мы у нее отняли ее солдат... — сказал Ленин, — отняли у нее рабочих и крестьян, одетых в солдатские мундиры»\footnote{\textit{В. И. Ленин}. Полн. собр. соч., т. 39, стр. 391, 393.}. Это была первая победа.

Вторая победа — провал попытки Антанты действовать против нас силами малых буржуазных государств, возникших на окраинах России. Мы своей правильной политикой добились их нейтралитета.

Мы начали отвоевывать у Антанты в ее собственных странах и мелкую буржуазию и образованное мещанство, которые, будучи вначале настроенными против нас, теперь вместе с рабочими выступали против вмешательства Антанты в дела России, против блокады. Это ярче всего проявилось во Франции. Это наша третья победа.

Ленин предостерегал, однако, что из всего этого нельзя делать слишком поспешные выводы. Нет сомнения, что империалисты свои попытки возобновят. Но главное позади. «Нужно, конечно, помнить, что враг нас подкарауливает на каждом шагу и сделает еще массу попыток скинуть нас всеми путями, какие только смогут оказаться у него: насилием, обманом, подкупом, заговорами и т. д.»\footnote{Там же, стр. 407.}

Империалисты и их подголоски клевещут на нас, обвиняя в терроризме, говорил Ленин. Террор был нам навязан. Он был вызван нашествием вооруженных сил Антанты, вторгшихся в нашу страну и воевавших против нас. Разве это не террор, когда их дипломаты, пользуясь неприкосновенностью, организуют у нас белогвардейские восстания? Разве голодная блокада Советской страны не является террором? Всемирный капитал душил, душит и обрекает на голодную смерть рабочих и крестьян. Это же относится к положению средних классов в революционной России, которые также несут свою долю жертв. Это объясняется тем, что правительства Антанты поставили нас в неслыханно тяжелые условия. И только окончательная победа над ними создаст условия для улучшения положения средних классов.

Ленин отметил далее, что в ходе гражданской войны та часть среднего крестьянства, которая была недовольна нами, испытан на себе ужасы режима Колчака и Деникина, повернулась в сторону Советской власти, укрепив тем самым положение Сойотов в деревне. Ленин остановился на трех главных моментах внутреннего положения и на трудностях, без преодоления которых невозможна окончательная победа Советской власти.

Первое — это недостаток хлеба и голод в крупных рабочих центрах. В данных условиях, говорил Ленин, выход из положения только один — если крестьяне отдадут государству излишки хлеба в ссуду (ибо государство пока не способно в обмен дать товары), то рабочий класс, восстановив промышленность, вернет свой долг крестьянину сторицей.

Вторая трудность — отсутствие топлива. Наш транспорт не имеет возможности вывезти даже те запасы хлеба, которые образовались в местах заготовок, чтобы накормить рабочих Питера и Москвы. Ленин говорил: «...пройдитесь по рабочим кварталам Москвы и вы увидите страшный холод, страшные бедствия, которые сейчас обострились из-за вопроса о топливе»\footnote{\textit{В. И. Ленин}. Полн. собр. соч., т. 39, cтр. 409.}. В октябре кризис с топливом особо остро чувствовался. «...Это была катастрофа, это был голод для рабочих целого ряда заводов и фабрик Москвы, Петрограда и целого ряда других мест. Результаты этой катастрофы сказываются до сих пор... Мы вылезаем из этой катастрофы, но еще далеко не вылезли»\footnote{Там же.}.

«И третий бич на нас еще надвигается, — говорил Ленин, — вошь, сыпной тиф, который косит наши войска. И здесь, товарищи, нельзя представить себе того ужаса, который происходит в местах, пораженных сыпным тифом, когда население обессилено, ослаблено, нет материальных средств, всякая жизнь, всякая общественность исчезает. Тут мы говорим: «Товарищи, все внимание этому вопросу. Или вши победят социализм, или социализм победит вшей!»\footnote{Там же, стр. 410.}.

Приведя справку наркома здравоохранения Семашко, в которой, в частности, говорилось, что из Москвы на борьбу с тифом отправлено 150 врачей, Ленин твердо заявил, что и этот кризис мы также начинаем преодолевать.

Поражала нас острота постановки вопросов Лениным. Он чувствовал необходимость рассказать народу всю тяжелую правду ради того, чтобы это дошло до каждого человека, чтобы всех мобилизовать и скорее преодолеть создавшиеся трудности.

Ленин призывал применить приобретенный в гражданской войне опыт в борьбе с этими внутренними трудностями с полной уверенностью в победе.

В конце доклада Ленин огласил проект резолюции по вопросу о международной политике, утвержденный до этого партийной конференцией. Проект был принят единогласно. В этой написанной собственноручно Лениным резолюции одобрялась мирная политика Советского правительства и ему поручалось «систематически продолжать политику мира, принимая все необходимые для ее успеха меры». Резолюция была разослана правительствам Англии, Франции и Италии как мирное предложение Советского правительства. Но правительства этих стран, потерпевшие поражение в необъявленной войне — вооруженной интервенции в Россию, кичащиеся своей высокой культурой и цивилизацией, не снизошли даже до того, чтобы рассмотреть эти предложения и дать ответ молодому государству.

И сегодня мы свидетели подобных действий империализма. Ведя необъявленную войну во Вьетнаме и терпя в ней поражение за поражением, американские империалисты бесконечно затягивают рассмотрение по существу миролюбивых предложений свободолюбивого вьетнамского народа.

Доклад Ленина встретил горячее одобрение делегатов. Но были на съезде и выступления иного характера. В прениях одним из первых выступил известный лидер меньшевиков Мартов. В период рождения партии он был близок к Ленину, но начиная со II ее съезда стал противником Ленина и ленинизма. Я его, конечно, знал по литературе, но впервые увидел здесь, на съезде. Внешне он никакого впечатления не производил: низкого роста, худенький, тщедушный, в пенсне. Говорил тихим голосом, монотонно, не пользовался ораторскими приемами. Фразы у него были длинные, но логически стройные.

Он не спорил с Лениным и даже по-своему подтверждал правильность многих выдвинутых им тезисов в области международного положения страны. Эта речь отражала новую обстановку, определенные изменения в настроениях меньшевиков. Мартов был вынужден считаться с действительными победами Советской власти, он не мог оспаривать линию нашей внешней политики, которую развил Ленин в своем докладе, и ограничился лишь туманным требованием «открыть поскорее окно в Европу». Но вот в конце речи он перешел к внутренней политике, стал читать декларацию партии меньшевиков.

Декларация носила злобный антисоветский характер, хотя и преподносилась под видом поддержки Советской власти. Была полна лжи и клеветы на работу советских органов власти сверху донизу, якобы нарушавших на практике положения Конституции. В ней утверждалось, будто ЧК незаконно применяет террор, и выдвигалось требование предоставления свободы печати, союзов, собраний для всех граждан республики, в том числе для врагов революции.

Смысл устного выступления Мартова и декларации отличались, как небо и земля. Это было настолько разительно, что в ходе чтения декларации были слышны возмущенные реплики, из зала кричали, что эта декларация отражает прошлогоднюю линию меньшевиков, линию 1918 года, когда меньшевики входили в самарское правительство учредиловки и контактировали с Колчаком.

В прениях выступил также Вольский, бывший председатель Комитета Учредительного собрания. Он говорил от имени своей группы, отколовшейся от правых эсеров. Начал он с того, что партия эсеров вела вооруженную борьбу с Советской властью (в это время в зале поднялся шум, слышались крики «позор!»), но год тому назад, когда Красная Армия двинула свои силы против Колчака, он и его группа решили порвать с правыми эсерами, прекратить борьбу против Советской власти и направить свои силы против реакции, объединив свои усилия с усилиями Советской власти.

Высказываясь за Советскую власть, он говорил о необходимости пересмотра существа Советской Конституции, выдвигал обвинение, сводившееся к тому, что «конституционные рамки, установленные самими коммунистами, ими же не соблюдаются». Он ратовал за «трудовое избирательное равенство». На словах высказываясь за III Интернационал, он предлагал «включить в ряды Интернационала все революционное социалистическое движение в мировом масштабе». Тем самым он требовал принятия в Коминтерн таких лжесоциалистических партий, как эсеры и меньшевики, что в корне противоречило существу Коминтерна.

С приветствиями и в прениях выступили также представители партий и групп, отколовшихся от разных правых партий. Среди них такие, как «революционные коммунисты», украинские «боротьбисты», максималисты, социал-демократы, интернационалисты и др., даже «сектанты-коммунисты».

Меня это поразило. Я и не подозревал, как много сохранилось еще разных политических группировок, и все они вроде в своих приветствиях ратовали за Советскую власть! Не было ясно только, что они практически делают для поддержки Советской власти. Почему они сохраняют свои группировки?

Декларация меньшевиков, выступление представителя Бунда с критикой Конституции, требование свобод для оппозиционных партий вызвали искреннее возмущение делегатов съезда.

Внимательно выслушав своих политических противников, Ленин со свойственным ему полемическим задором начал свою заключительную речь. Он сразу взял под обстрел Мартова и оглашенную им декларацию. Он сказал, что 15 лет уже наблюдает за поведением меньшевиков и видит теперь, что после уроков Колчака и Деникина, больших побед Красной Армии меньшевики начали колебаться в сторону Советской власти, ближе к диктатуре пролетариата, что отразилось и в речи Мартова. Но декларация, прочитанная Мартовым, находится в полном противоречии с этим настроением. Она не отражает положения, она призывает возвратиться назад, к буржуазной демократии. Жизнь же идет своим чередом. Часть меньшевиков, чем дальше, тем больше, будет приближаться к Советской власти, хотя и с величайшим трудом, вопреки этой платформе, против своей воли.

Ленин говорил, что, вопреки утверждениям меньшевиков, у нас впервые в истории осуществлена самая полная, самая действительная демократия для трудящихся и Советская Конституция отражает и обеспечивает эту демократию.

Нас упрекают, что мы не соблюдаем Конституцию, «я утверждаю, что мы Конституцию соблюдаем строжайшим образом»\footnote{\textit{В. И. Ленин}. Полн. собр. соч., т. 39, стр. 423.}.

В это время из ложи, где сидела оппозиция, раздался голос: «Ого!» Ленин немедленно откликнулся на этот возглас: «И хотя из ложи, которая в прежние времена была ложей царской, а теперь является ложей оппозиции... я слышу иронический возглас «ого!», — тем не менее я сейчас это докажу»\footnote{Там же.}. Зал ответил хохотом и аплодисментами.

Ленин зачитал параграф 23 Конституции, где говорилось: «Руководствуясь интересами рабочего класса в целом, РСФСР лишает отдельных лиц и отдельные группы прав, которые используются ими в ущерб интересам социалистической революции», — и добавил, обращаясь к оппозиции: «Если вы хотите, чтобы мы соблюдали Конституцию, то хотите ли вы, чтобы соблюдался и параграф 23-й?»\footnote{Там же.} Это заявление также было встречено аплодисментами.

Ленин говорил очень откровенно, напрямик. Он подчеркнул, что большевики никогда не рассматривали свою деятельность как образец совершенства, не считают таким образцом и Конституцию, не случайно и на этом съезде стоит вопрос о внесении в нее некоторых изменений.

Но большевики никогда не давали и пустых обещаний о свободе и равенстве для всех, вне всяких классовых позиций. «И параграфом 23-м мы говорим, что для переходного периода не рисуем молочных рек и кисельных берегов. Мы говорим, что нам нужно продержаться не месяцы, а годы, чтобы закончить этот переходный период»\footnote{Там же, стр. 425.}. Ленин пояснил, что свободы и равенства не может быть для буржуазии и спекулянтов. Более того, поначалу мы не можем обещать даже равенства рабочих и крестьян, их положение в условиях диктатуры пролетариата — на период острейшей классовой борьбы — неодинаково.

Иное дело меньшевики. Им Ленин не менее прямо и открыто заявлял: «Вы... обещаете народу все, для того, чтобы этого не исполнять. А мы ничего такого не обещаем...»\footnote{\textit{В. И. Ленин}. Полн. собр. соч., т. ЗЭ, стр. 423–424.}

Снова и снова Ленин пояснял, что не может быть свободы, демократии, равенства вообще. «Свобода, — но для какого класса и для какого употребления? Равенство, — но кого с кем?»\footnote{Там же, стр. 424.} Только так надо ставить вопрос.

Обращаясь к Мартову и оппозиции, Ленин говорил: «Покажите нам образец страны, образец вашей прекрасной меньшевистской конституции. Может быть, найдете такой образец хотя бы в истории Самары, где была меньшевистская власть? Может быть, вы найдете его в Грузии, где сейчас меньшевистская власть?.. Покажите такой пример, и мы поучимся. Но вы и показать этого не можете...»\footnote{Там же.}

С тех пор прошло 50 лет. Эти же вопросы можно было бы задать с тем же основанием и современным правым социал-демократам, в том числе и тем, которые не один год стояли у власти, и тем, которые сейчас правят в некоторых странах. Ни одна из этих партий так и не выдвинула альтернативы социализму, не ликвидировала эксплуатации, не предоставила трудящимся свободы и равенства, не придерживалась и не придерживается истинной, социалистической демократии.

Ленин ответил далее на обвинения оппозиции в том, что в 1918–1919 годах не всегда соблюдались сроки выборов, нерегулярно созывались сессии ВЦИК, особенно в последнее время, когда велись ожесточенные бои против армий Колчака и Деникина: члены ВЦИК в большинстве своем были на фронтах, и это полностью оправдывается интересами революции — отрывать их на заседания в такой критический момент было бы в ущерб ходу борьбы. И какой жалкой выглядела на этом фоне реплика бундовки о том, что на фронт можно было послать других людей!

Ленин предложил подойти к этому с другим мерилом — сравнить с тем, что делалось во время войны в странах буржуазных, и он утверждал, что если опросить делегатов с мест, когда собирались Советы, то такой опрос докажет, «что в такое трудное время, как время войны, когда действие европейских конституций, веками установленных, вошедших в привычку западноевропейского человека, почти целиком было приостановлено, в это время Советская конституция в смысле участия народных масс в управлении и самостоятельном разрешении дел управления на съездах и в Советах и на перевыборах применялась на местах в таких размерах, как нигде в мире»\footnote{\textit{В. И. Ленин}. Полн. собр. соч., т. 39, стр. 419–420.}.

Это была наступательная и страстная ленинская речь. Он не собирался уступать ни одной позиции, не оставил без ответа ни одного упрека, ни одного критического замечания представителей оппозиционных партий.

Ленин высмеял тех людей, которые были в Учредительном собрании, а теперь начинают учить коммунистов, как поставить защиту революции, выступают против ЧК. Они были с Колчаком, с белогвардейцами, и эти же белогвардейцы их выкинули вон. Насколько мне помнится, Ленин, когда говорил об этом, был не на трибуне, а отошел от нее несколько в сторону. И тут-то он сделал такой общепонятный жест ногой, который показывал, как выбросили меньшевиков вон. Это вызвало дружный хохот делегатов съезда и аплодисменты.

Декларация, зачитанная Мартовым, говорил Ленин, повторяет буржуазную клевету о том, будто коммунисты представляют меньшинство рабочего класса. И это смехотворно. Когда такого рода клевету самых худших из хищников, зверей империализма, заявил Ленин, «повторяет здесь Мартов от имени Российской социал-демократической рабочей партии (смех), тогда я говорю себе, что надо быть начеку и знать, что тут ЧК необходима!»\footnote{Там же, стр. 422.}. Это снова вызвало аплодисменты всего зала.

Фактически это был последний бой Ленина от имени коммунистов с меньшевиками, остатками уже правых оппортунистических партий. В следующем году Мартов эмигрировал в Германию, где стал издавать свой злопыхательский «Социалистический вестник». За борьбу против Советской власти в начале 1922 года Дан был выслан за границу, как враг Советского государства. Меньшевистская и другие партии распадались на глазах у всех.

Это хорошо понимал Ленин, поэтому-то он и закончил свое выступление таким обращением: «Кто нам хочет помогать, того мы принимаем с величайшей радостью независимо от его прошлого, не считаясь ни с какими названиями. И мы знаем, что таких людей из других партий и беспартийных к нам идет все больше и больше, и этим обеспечивается наша победа»\footnote{\textit{В. И. Ленин}. Полн. собр. соч., т. 39, стр. 425.}.

Раздались громкие аплодисменты, крики «браво!», делегаты встали с мест. Хорошо чувствовалось общее удовлетворение от сознания величия и мастерства Ленина, владевшего сердцами посланников трудового народа и разбившего в пух и прах все аргументы оппозиции. Росло понимание того, что оппозиционные партии себя окончательно исчерпали и сходят с исторической сцены русской революции.

Закрывая съезд, Ленин подвел его итоги, отметил значение дискуссии с оппозиционными партиями о демократии и Советской власти и призвал в ближайшее время утроить усилия, как в страдную пору, добить Деникина и завоевать себе мир прочный и долгий, добиться «такого конца гражданской войны, который откроет нам возможность мирного социалистического строительства на долгое время»\footnote{Там же, стр. 436.}. С таким настроением и покидали делегаты свой съезд Советов.

Мое приподнятое настроение объяснялось еще и тем, что меня, молодого человека, которому было всего 24 года, избрали в состав верховного органа Советского государства — кандидатом в члены ВЦИК.

\newpage
\section[Центральный Комитет решает организационные вопросы нашей партии на Кавказе]{\center \textbf{Центральный Комитет решает организационные вопросы нашей партии на Кавказе}}


Наступил уже декабрь, а обсуждение поставленных мною перед ЦК вопросов о партийном строительстве на Кавказе оттягивалось.

Еще 14 ноября Политбюро под руководством Ленина постановило: ««Вопрос об организации нашей партии на Кавказе»...обсудить в соединенном заседании Политбюро и Оргбюро, вызвав на это заседание Сталина и допустив в качество докладчиков Микояна и Аванесова, а в качестве участника — Енукидзе»[71].

Время шло, я торопился вернуться на Кавказ. И вот 4 декабря, после окончания VIII партийной конференции, я увидел Ленина, который шел по коридору к себе в кабинет. Я подошел к нему и сказал, что кавказские вопросы не продвигаются вперед, а я тороплюсь с отъездом на Кавказ и прошу ускорить их рассмотрение. Ленин одобрительно отнесся к моей просьбе и сказал, что примет необходимые меры.

20 декабря меня пригласили в Кремль на заседание Политбюро ЦК. В нем участвовали три члена Политбюро, четыре члена Оргбюро, Аванесов от Наркомнаца и я. Сталина в Москве не было. Это было объединенное заседание Политбюро и Оргбюро, которое вел Ленин. Оно проходило и маленьком, скромно обставленном, уютном кабинете Ленина.

До заседания я успел ознакомиться с проектом резолюции, которую представил Каменев. Основные положения крайкома в ней были отражены.

Резолюция начиналась с принципиальной вводной части — совершенно правильной, имеющей теоретическое значение не только для Кавказа, но и для всех окраин России. Я чувствовал в ней руку Ленина.

Имелся пункт, предусматривающий, что все коммунистические организации Кавказа подчиняются Кавказскому краевому комитету РКП. Все коммунисты Азербайджана, работающие в Азербайджане, образуют Компартию Азербайджана, — так это и было определено коммунистами Азербайджана.

В Грузии не предусматривалось создания компартии. Это было понятно, поскольку грузинские товарищи тогда не ставили этого вопроса. В резолюции был еще ряд пунктов организационно-практического характера, связанных с работой Компартии Армении.

Каменев доложил, что этот проект не согласован со Сталиным, мнение которого до сих пор не получено. Я выступил коротко, поддержав проект резолюции в целом. Сказал, что проект достаточно ясно предусматривает подчинение всех коммунистических организаций Кавказа крайкому партии, тем самым подразумевается и Компартия Армении.

Ленин, высказавшись совсем коротко, заявил, что резолюция в основном правильная и следует ее принять за основу. О Компартии Армении Ленин предложил в решении записать особо, сформулировав этот пункт следующим образом:

«Признавая необходимым существование Коммунистической партии Армении, как территориальной организации в пределах существующей республики Армении, подчиненной краевому Кавказскому комитету РКП...»[72]

Затем Ленин предложил выработать детальную резолюцию и внести ее на утверждение Политбюро, что и было принято.

Я ушел с заседания, пораженный деловитостью и строгим порядком ведения его Лениным. Никаких речей, лишних разговоров. Высказывания были кратки и ясны.

Довольный решением вопроса, я не стал дожидаться выработки детальной резолюции. Я торопился вернуться в бакинское подполье до прихода Красной Армии, которая в то время стремительно рвалась на юг, громя белогвардейские войска Деникина.

За несколько дней до отъезда из Москвы я написал краткую записку Владимиру Ильичу, что, «разрешив организационный вопрос, Политбюро ЦК отложило вопрос о нашей политике на Кавказе»[73], и просил, если это возможно, принять меня еще раз на несколько минут, чтобы получить совет по отдельным неотложным вопросам, в частности о восстании в Грузии. Дело в том, что грузинские товарищи перед моим отъездом просили обязательно привезти им ответ Ленина. Через пару дней Стасова по поручению Ленина сообщила мне, что о восстании в Грузии и другие вопросы, касающиеся Кавказа, сейчас решать нельзя и нецелесообразно. Наши войска ведут стремительное наступление на юге. Меняется не только военное, но и политическое положение. Это изменит обстановку и на Кавказе. По мере продвижения Красной Армии все эти проблемы будут рассматриваться и решаться с учетом конкретно сложившейся обстановки.

Это было настолько правильным, что вызвало у меня чувство раскаяния — стало как-то стыдно, — оказалось, я зря побеспокоил своей запиской Владимира Ильича.

С другой стороны, я был удовлетворен, считая, что сделал все возможное для выполнения данных мне крайкомом партии поручений.

Я выехал через Ташкент и Красноводск в Баку. Тогда это был наиболее короткий путь, ибо Волгу у Астрахани уже сковало льдом. Для ускорения поездки я воспользовался специальным поездом, которым Фрунзе направлялся в Ташкент.

Много позже мне стало известно, что 3 января 1920 года состоялось постановление Политбюро по вопросу о партийных организациях на Кавказе.

Вводная часть этого постановления опубликована в 1968 году, во второй книге 3-го тома Истории КПСС.

Ввиду ее важного значения считаю целесообразным привести ее здесь:

«Борьба с местным шовинизмом и создание благоприятных условий пропаганды социалистической революции среди наций, находившихся под игом российского царизма, делают необходимым, чтобы в фактически образовавшихся «самостоятельных» государствах в пределах бывшей Российской империи коммунистические организации работали в виде самостоятельных коммунистических партий. Особенно важен подобный способ организации на Востоке, ввиду того что завоевание доверия народов, живущих в Азии или в преддвериях Азии, составляет одну из основных задач Советской России и РКП»[74].

В протоколе Политбюро от 3 января 1920 года записано: «Вопрос о партийных организациях на Кавказе. Утвердить предложенную т. Каменевым резолюцию с небольшими поправками»[75], а сама эта резолюция называется «Постановление ЦК РКП о коммунистической работе в Армении».

Первый пункт постановления начинается такой фразой: «Работой всех коммунистических организаций на территории Кавказа руководит Кавказский краевой Комитет РКП...», его решения одинаково обязательны для всех национальных коммунистических партий на Кавказе, а их финансирование происходит только через крайком, которому принадлежат все права областного комитета партии[76].

Один пункт, по существу, повторяет решение Политбюро ЦК РКП(б) от 20 декабря 1919 года, принятое в моем присутствии, о том, что в пределах существующей республики Армении коммунистическая работа ведется Коммунистической партией Армении. Несколько организационных пунктов постановления касаются только Компартии Армении. В постановлении также говорилось, что вне пределов республики Армении работа среди армянского пролетариата ведется отделами по работе среди национальных меньшинств, находившимися при местных комитетах национальных коммунистических партий или РКП(б).

В феврале 1920 года, в условиях подполья, состоялся первый съезд Коммунистической партии (большевиков) Азербайджана, провозгласивший создание партии и избравший ее руководство. На съезде была обсуждена программа борьбы за победу Советской власти

В мае 1920 года был образован ЦК Компартии Грузии и утвержден его состав из товарищей, работавших в Грузии.

В июне оформилось создание Коммунистической партии Армении и ее ЦК, работающих в пределах Армянской республики.

\subsection*{\center * * *}

В конце апреля 1920 года Красная Армия разбила деникинские части на Северном Кавказе. Освободив Северный Кавказ, XI армия подошла к границе с азербайджанским буржуазным государством.

Одновременно в Баку большевики уже полностью подготовились к вооруженному восстанию. Действия бакинских рабочих и Красной Армии были хорошо скоординированы. Командовал Кавказским фронтом Тухачевский, членом Реввоенсовета фронта был Орджоникидзе. XI армией командовал Левандовский, члены Реввоенсовета армии — Киров и Мехоношин.

Что касается меня, то, пробираясь из Закаспия в Баку, я 22 апреля добрался до Петровска (ныне Махачкала), в штаб XI армии, и был там назначен комиссаром передового отряда бронепоездов под командованием Ефремова, целью которого было ворваться в Баку с расчетом помешать буржуазному правительству уничтожить нефтяные промыслы.

27 апреля наши войска пересекли границу Азербайджана и вышли на помощь бакинским рабочим. Головной бронепоезд уже на заре 28 апреля прибыл на вокзал в Баку, где мирно, без кровопролития, при большом перевесе революционных сил победила власть Советов[77].

Ленин придавал этой победе очень серьезное значение. Уже 29 апреля, выступая на Всероссийском съезде рабочих стекло-фарфорового производства, Ленин говорил:

«...Вчера же нами была получена весть из Баку, которая указывает, что положение Советской России направляется к лучшему; мы знаем, что наша промышленность стоит без топлива, и вот мы получили весть, что бакинский пролетариат взял власть в свои руки и сверг азербайджанское правительство. Это означает, что мы имеем теперь такую экономическую базу, которая может оживить всю нашу промышленность»[78]. Съезд горячо приветствовал создание Азербайджанской Советской Республики.

Через несколько дней, 5 мая, Ленин от имени Совнаркома подписал приветственную телеграмму Советскому социалистическому правительству Азербайджана[79].

В свою очередь бакинский пролетариат, трудящиеся Азербайджана испытывали чувство горячей любви, преданности к своему вождю, о чем говорит, к примеру, телеграмма, посланная Ильичу Бакинским Советом:

«ОТ БАКИНСКОГО СОВЕТА

КРАСНОАРМЕЙСКИХ И МАТРОССКИХ

ДЕПУТАТОВ

11 июля 1920 г.

Москва, Кремль, ЛЕНИНУ

Первое торжественное заседание Бакинского Совета красноармейских и матросских депутатов избрало Вас, дорогой Ильич, своим почетным председателем. С гордостью сообщаем Вам, что из 650 депутатов 648 коммунистов, 2 сочувствующих. Надейтесь на нас как на стальную твердыню международного коммунизма на рубеже красного Севера и возрождающегося Востока.

Да здравствует наш дорогой Ильич!

По поручению президиума — Микоян»[80].

\subsection*{\center * * *}

В августе 1920 года состоялось решение ЦК партии о моем направлении на партийную работу в Нижний Новгород. На пути туда по окончании в Баку съезда народов Востока я попал в Москву как раз накануне IX партийной конференции. Она представляет существенный интерес для понимания внутрипартийной борьбы перед X съездом.

\newpage
\section[Ленин на сентябрьской партийной конференции 1920 года: если мы сплотим силы и напряжем их, победа будет за нами!]{\center \textbf{Ленин на сентябрьской партийной конференции 1920 года: если мы сплотим силы и напряжем их, победа будет за нами!}}


IX Всероссийская партийная конференция состоялась 22–25 сентября 1920 года. Ей предшествовал период огромных побед на фронтах, иностранные интервенты были изгнаны с территории Советской республики, разбиты полчища Колчака и Деникина. Но за несколько месяцев до конференции вновь разгорелась война. Активизировались остатки деникинских войск в Крыму, организованные Врангелем при помощи французского оружия и золота. Почти одновременно с этим империалистами Антанты было организовано наступление панской Польши на Страну Советов в целях захвата белорусских и украинских земель. Вновь пришлось послать на фронт десятки тысяч коммунистов, чтобы обеспечить скорейшую победу над белополяками и Врангелем.

Ленин открыл конференцию и выступил с политическим отчетом Центрального Комитета партии. Он коснулся главных событий международного положения, особенно изменений, вызванных нападением белополяков на Советскую Россию. Он отметил, что еще в начале года Советское правительство обратилось к Польше с предложением заключить мирный договор на выгодных для поляков условиях, с территориальными уступками с нашей стороны. Однако, видимо, империалисты и польские националисты под водительством Пилсудского неправильно поняли готовность Советского Союза к уступкам, истолковали это как слабость Советской России и решили начать военные действия, чтобы отхватить побольше территорию. Они захватили Киев и ряд других городов Украины. «Лишний раз подтвердилась истина, — говорил Ленин, — что буржуазная дипломатия не способна понять приемов нашей новой дипломатии открытых прямых заявлений»[81].

Красная Армия, собрав силы, разгромила белополяков на Украине, в Белоруссии и добралась до подступов к Варшаве. Но из-за растянутости коммуникаций, переутомления войск, плохого снабжения фронта и ряда других причин Красная Армия потерпела поражение в боях под Варшавой и отступила более чем на сотню верст назад. Однако она удержалась далеко впереди тех позиций, которые занимала до польского нападения.

Ленин сказал, что, несмотря на наше поражение под Варшавой, удар по польской армии — серьезное потрясение для самой Польши. В рабочем классе там началось брожение. Польская армия теперь разбавлена, пополнена теми рабочими и крестьянами старших возрастов, которые участвовали в предыдущей войне, в социальном составе она уже не та, что до понесенных ею поражений.

Антанта подталкивает Польшу к новой войне против нас. Однако теперь и польские мелкобуржуазные партии понимают, что война принесет лишь дальнейшее разорение, и поэтому предпочитают мир. Этот шанс мы и хотим использовать и снова предлагаем Польше выгодный для нее и невыгодный для нас мир, лишь бы избежать тяжелой для нас зимней кампании. «Если нам суждена зимняя кампания, мы победим, в этом нет сомнения, несмотря на истощение и усталость»[82].

Ленин сказал и о другом последствии контрудара, нанесенного по Польше, — воздействии на революционное движение Европы, особенно Англии. О печальной памяти «ультиматуме Керзона» он сказал: «Когда английское правительство предъявило нам ультиматум, то оказалось, что надо сперва спросить об этом английских рабочих»[83]. А рабочие выступили под лозунгом «Руки прочь от России» и стали организовывать «комитеты действий», хотя их организации и возглавлялись, как определил тогда Ленин, «злостными меньшевиками».

Он коснулся также задач ликвидации Врангеля. Красная Армия подтягивала силы к югу. За несколько дней до начала конференции Фрунзе был назначен командующим Южным фронтом и принят Лениным.

И мне помнится, в разгар польского наступления широко известный генерал Брусилов и еще несколько видных генералов выступили с обращением к бывшим царским офицерам: где бы они ни находились, поддержать Красную Армию в борьбе с белополяками. Это было выражением некоторого перелома в наиболее честных кругах старого русского офицерства и интеллигенции.

В области внутренней политики главным вопросом на конференции были задачи партийного строительства. С докладом по этому вопросу выступал Зиновьев.

На конференцию прибыли делегаты из разных концов страны, с разным опытом — из восточных районов Сибири, с Волги, Северного Кавказа, юга России. Были и такие, что только вышли из подполья, были делегаты от частей действующей армии.

Но отдельные выступления произвели на нас неприятное впечатление. Задавались, например, демагогические вопросы: «Будет ли соответствовать свобода критики свободе кушать персики?» — с намеком на то, что кого-то ЦК «за критику» послал работать из центра на юг. Спросила об этом Коллонтай, Ленин был вынужден отвечать и на подобные вопросы[84].

В выступлениях некоторых представителей рабочих центров России, которые вынесли на своих плечах всю горечь войны и разрухи, голода и нищеты, чувствовалось тяжелое настроение и недовольство внутренней жизнью. Очень острой критике подвергалось руководство, резко говорилось о «верхах» и «низах», об отсутствии демократии.

Но была иногда и самая настоящая демагогия. Наиболее разнузданно выступал представитель группы так называемого «демократического централизма» Сапронов. На деле эта группа пеклась, и притом весьма односторонне, лишь о демократии без конца и края и фактически исключала централизм. Отпор ей был дан еще на IX съезде партии, но группа не унималась.

Выступали и другие представители оппозиции: Лутовинов, Кутузов — руководящие работники профсоюзов. (Затем они влились в состав так называемой «рабочей оппозиции».)

Зиновьев, излагая линию ЦК и аргументируя необходимость изменения методов партийной работы, развития внутрипартийной демократии, как мне показалось, перехлестывал, например, в своих развязных нападках на «комиссаров, которые разъезжали в своих вагонах», обещал «самую развернутую» демократию. Видимо, он ставил целью обезоружить оппозицию, но делал это с перебором, что на меня лично производило неприятное впечатление. Я ехал с ним из Баку в одном поезде — поезде Коминтерна, которым возвращалась часть делегатов со съезда народов Востока. Меня поразило, что Зиновьев, который не только сам ехал в отдельном вагоне, в чем мы не видели ничего особенного, но и держал себя крайне отчужденно, не общался со своими весьма уважаемыми спутниками, теперь же с трибуны конференции вдруг выступает с резкими нападками на комиссаров, которые, мол, «ездят в своих вагонах».

Ленин внимательно слушал выступления. Их характер, их тон вызывали, видно, у него тревогу. Он счел необходимым вновь взять слово. Это было очень серьезное выступление. Ленин сказал, что некоторые речи ораторов в прениях выражают сильное переутомление, проявлявшееся в истеричности. Он заметил при этом: «Я бы не сказал, что тут демагогия. Переутомление физическое дошло до истеричности»[85].

В выступлениях Ленина проявилось глубокое понимание того, что происходит в стране, понимание причин, вызывавших недовольство, знание возможностей, которыми располагает партия в сложившихся условиях, умение выбрать из выступлений на конференции все здоровое, полезное. Например, в выступлениях представителей оппозиции было много демагогии, но Ленин пояснил, что нельзя все сказанное оппозиционерами считать демагогией. Он этим хотел облегчить положение оппозиционно настроенных делегатов, открыть перед ними возможность сплочения вокруг ЦК, вокруг тех предложений, которые принимались конференцией. Ибо Ленин отчетливо представлял, какое это не простое дело в условиях крайнего переутомления и отчаянной усталости, физического истощения и недовольства выступать с призывом к новому напряжению сил, к отрешенности от личных интересов, к откладыванию до конца войны уже назревших и перезревших вопросов. А ведь лейтмотивом ленинских выступлений на конференции было именно это: «...если мы сплотим силы и напряжем их, победа будет за нами»[86]. Пока идет война — «величайшее напряжение и никаких разговоров...»[87]

Ленин говорил на конференции о трудностях, которые партия переживала в силу условий, сложившихся за годы гражданской войны. Военная обстановка не позволяла органам партии, коммунистам часто собираться, обсуждать вопросы. Он предупреждал, что только мирные условия позволят как следует развернуть демократию и провести в жизнь, что мы намечаем. Это будет зависеть, например, от того, будем ли мы воевать с Польшей, вести зимнюю военную кампанию, или будем жить в мире.

В моменты, когда страна находилась в смертельной опасности, когда Колчак добрался до Волги, а Деникин до Орла, о какой демократии могла идти речь?

Но, выступая в прениях, Ленин подчеркнул, что уже на этой конференции требуется определить необходимые серьезные шаги в сторону развития внутрипартийной демократии, что Центральный Комитет это хорошо понимает. Ленин поддержал предложения Московской конференции по вопросу о демократии и внес несколько дополнений. В частности, он поддержал идею создать «контрольную комиссию», которая могла бы заняться разбором поступающих в ЦК жалоб с мест на ведомства, на центральные органы. Выборы контрольной комиссии, конечно, являлись делом предстоящего X съезда партии. Но, не дожидаясь съезда, было признано необходимым выбрать временную контрольную комиссию на этой конференции.

Мы должны, говорил Ленин, поставить вопросы демократии на повестку дня сейчас, в момент опасности военной, но мы их будем ставить иначе, когда войны не будет. То есть Ленин правильно разъяснял, несколько поправляя Зиновьева и делая упор на то, что отсутствие необходимой демократии в военное время и методы работы вызывались самой обстановкой. Нельзя одни и те же методы применять в период войны и в мирных условиях.

Не буду пересказывать подробно все, что говорил Ленин, остановлюсь только на одном, как мне кажется, весьма характерном эпизоде, показывающем, сколь неустойчивым еще было положение. Ленин отметил, что «общее положение республики... улучшилось настолько, что сейчас мы имеем возможность обсуждать с большим хладнокровием: мы не ставим теперь вопроса о преждевременном прекращении конференции, как ставили несколько раз в эпоху наступления Колчака и Деникина. Бывали партийные съезды, с которых, не дождавшись их окончания, уезжал целый ряд ответственных работников прямо на фронт»[88].

И что же! Едва закончилось его выступление, еще во время заседания, Ленину передают телеграмму с Западного фронта о новом наступлении белополяков, с предложением ускорить работу конференции или отъезд с нее делегатов-военных от этого фронта. Ленин тут же пишет на телеграмме: «и то и то принять тотчас»[89]. К этому можно добавить, что в дни состоявшегося через полгода X съезда партии, проходившего во время кронштадтского мятежа, положение для многих делегатов оказалось еще сложнее.

Прения показали значительный разнобой, граничивший с отсутствием единства. Однако Ленин своими выступлениями старался сплотить партию. Не в малой степени именно благодаря этому делегаты уезжали с лучшим настроением, чем приехали. Но причины, вызвавшие разногласия, не могли быть быстро устранены, и вскоре партии пришлось вновь столкнуться с оппозицией, навязавшей ей в конце года общепартийную дискуссию.

\subsection*{\center * * *}

После окончания партийной конференции Ленин участвовал в работе трех заседаний сессии ВЦИК. Он не выступал с речами, а сидел и внимательно слушал, делал записи для себя. Он внимательно выслушал доклад наркома просвещения Луначарского и содоклад Невского от комиссии ВЦИК.

Луначарского отличало особое умение строить доклад так, что он захватывал слушателей. Обсуждались новые для Советской Республики вопросы. Особенно интересным был для меня доклад наркома просвещения, жадно хотелось узнать, что нового вносится социализмом в дело народного образования.

Луначарский увлекательно говорил о задачах народного образования в молодой социалистической стране на фоне нищеты, острой нехватки школьных помещений, многие из которых были заняты военными учреждениями под лазареты или казармы. Школы не отапливались, не было дров, не было материалов для ремонта зданий. Гражданская война поглощала все. На десяток учеников выдавался один карандаш, одна тетрадь на весь год. И многие другие мрачные цифры со всей откровенностью приводил Луначарский.

Таковы последствия навязанной нам войны, говорил Луначарский. Скоро война кончится, и тогда удастся выделять больше материальных средств, чтобы достичь действительного расцвета образования от начального до высшего, ликвидации неграмотности взрослых.

Доклад Луначарского одновременно был проникнут оптимизмом. Он с восторгом говорил, как в таких условиях — при нехватке учителей, учебников — дети заполняют школы в количествах больших, чем это было при царе, когда не было такой нищеты. Дело просвещения зиждилось на энтузиазме учителей, в своей массе повернувшихся в сторону Советской власти, на поддержке рабочих и крестьян. С невероятной, невиданной ранее быстротой росло число лиц от мала до велика, которые садились за парту.

Для иллюстрации Луначарский приводил убедительные цифры по отдельным городам и деревням в сравнении с прошлым. Он показывал, какие принципиальные изменения внесены в постановку образования, какие ненужные предметы исключены — закон божий, мертвые языки — греческий и латынь. Школам придается трудовой характер. Открываются рабочие факультеты для подготовки в вузы, курсы для подготовки учителей и т. д.

Ленин внимательно слушал и что-то все время записывал. Приятно было наблюдать, как со скромностью простого труженика глава правительства и вождь партии напрягал свое внимание, весь сосредоточившись, делал заметки. Он не прерывал выступавших, не подавал реплик. Было видно, что Ленин вникает в излагаемые факты, изучает, обобщает их.

Впоследствии были опубликованы его записи по вопросам народного образования, сделанные на этой сессии. В них отмечались заслуживающие наибольшего внимания моменты из выступлений не только докладчика и содокладчика, но и ораторов, выступавших в прениях. Из этих записей видно, с каким знанием дела, как компетентно подходил Ленин к новым проблемам и нуждам образования в первой стране социализма.

После прекращения прений с мест раздались голоса — дать слово Ильичу, просьбы, чтобы он высказался. Но Ленин с места заявил, что он пришел сюда лишь для того, чтобы осведомиться о деятельности Наркомпроса, поэтому брать слово не хочет.

\subsection*{\center * * *}

Следует еще, пожалуй, напомнить, что на этой сессии ВЦИК было принято заявление о войне с Польшей, в котором с нашей стороны предлагался ряд уступок, ради того, чтобы добиться подписания перемирия, избежать затяжной зимней военной кампании. В документе говорилось, что эта война не отвечает интересам ни польского, ни советского народов. В документе содержалась оговорка, что, если эти предложения не будут приняты буржуазной Польшей до 5 октября, они теряют силу.

Польша заявила о своем согласии с нашими предложениями. Военные действия были прекращены 12 октября в Риге польское командование подписало предварительные условия мира. Мирный договор был подписан 18 марта 1921 года.

\newpage
\section[Ленин на VIII Всероссийском съезде Советов: план ГОЭЛРО — вторая Программа партии]{\center \textbf{Ленин на VIII Всероссийском съезде Советов: план ГОЭЛРО — вторая Программа партии}}



Минуло три года существования Советской власти. Советская республика, пройдя через тяжелейшие испытания, выстояла. Таков был главный итог пройденного пути. Наступила уже не передышка, а целая историческая полоса мирного развития. Полный веры в возможность успеха социалистического строительства, Ленин дал оценку новому положению Советской республики: «...мы имеем новую полосу, когда наше основное международное существование в сети капиталистических государств отвоевано»[90].

Новый период выдвигал и совершенно новые задачи в области внутренней и внешней политики партии и государства. Это четко видел и сформулировал Ленин, выделив главные звенья и стараясь именно на них сосредоточить усилия партии, рабочего класса, всех трудящихся.

Эти вопросы были поставлены в центр внимания VIII съезда Советов (декабрь 1920 года) и X съезда партии (март 1921 года) и решены ими.

Определяющими должны стать вопросы хозяйственного строительства. Надо приступить к созданию экономического фундамента социализма путем разработки первого долгосрочного народнохозяйственного плана. Таким планом явился план электрификации России — ГОЭЛРО, одобренный VIII съездом Советов.

Надо разработать пути восстановления сельского хозяйства. Здесь главное — в изменении отношения класса пролетариата к классу крестьянства. Эти вопросы также были предметом обсуждения на VIII съезде Советов и нашли свое конкретное завершение в выработке начал новой экономической политики X съездом партии.

Третью основную задачу в переходе к мирному строительству Ленин видел в более широком, чем в период войны, привлечении трудящихся к активному участию в хозяйственной и политической жизни, развитии партийной и советской демократии, укреплении руководящей роли партии и ее отношений с рабочим классом, авангардом которого она является.

Наконец, в области внешней политики — постепенная нормализация отношений с капиталистическими странами, развитие экономических связей. В этом плане Ленин придавал большое значение вопросу об иностранных концессиях как средству вынудить иностранный капитал оказать нам помощь прежде всего современной капиталистической техникой.

Эти вопросы приходилось решать в условиях политического и хозяйственного кризиса в стране, вызванного нехваткой продовольствия, топлива, сырья, разрухой в промышленности, на транспорте и в сельском хозяйстве, последствиями небывалой засухи, в условиях крайнего переутомления трудящихся из-за лишений, перенесенных в первой мировой и гражданской войнах. Все это вызвало в массах недовольство, шатания и колебания даже в рядах партии.

«...Мы наткнулись на большой, — я полагаю, на самый большой, — внутренний политический кризис Советской России. Этот внутренний кризис обнаружил недовольство не только значительной части крестьянства, но и рабочих»[91] —так остро и откровенно характеризовал Ленин обстановку, сложившуюся в начале 1921 года.

Положение еще более было осложнено навязанной Троцким дискуссией о профсоюзах, которая вдохновила и другие оппозиционные группы.

\subsection*{\center * * *}

VIII съезд Советов открылся 22 декабря 1920 года. Это был самый представительный с начала революции съезд Советов. В нем участвовали делегации от всех районов и окраин Советской страны. К тому времени образовалась Бухарская Советская Республика, Белорусская Советская Республика. Одна Грузия оставалась еще под властью меньшевиков. В Азербайджане за несколько месяцев, в Армении за месяц до этого победила Советская власть. На съезде выступали с приветствиями представители Азербайджана и Армении.

Делегатов и гостей было очень много, больше, чем постоянных мест в Большом театре. Сидели в крайней тесноте. Все проходы были заняты. Настолько было тесно, что решили после первого заседания съезда аннулировать гостевые билеты и оставить места только для делегатов, которых было 2537 человек[92].

Бурными аплодисментами был встречен Ленин, выступивший с отчетом ВЦИК и Советского правительства. Он сразу же оговорился, что не имеет в виду делать подробный доклад о всей деятельности правительства, о принятых законах и т. д. Он сказал, что хочет обобщить главные уроки за прошедшие годы революционной борьбы и сделать выводы о самых важных политических и хозяйственных задачах, которые следует поставить перед Советским правительством и трудящимися в хозяйственном строительстве.

Коротко говорил Ленин о международном положении республики, о подписании предварительного мирного договора с Польшей.

Он подчеркнул огромное значение быстрой, славной победы над хорошо вооруженным войском Врангеля. Отметил наличие договорных отношений с прибалтийскими буржуазными республиками, с соседними восточными странами. Все это в совокупности позволяет получить уже серьезную, длительную передышку.

Ленин заявил, что в целом одержана большая победа, по нужно быть сильными, ибо мы не имеем гарантии от новой войны. Отсюда — необходимо сохранить военную готовность, с тем чтобы, сокращая численность армии, поднять качество подготовки, ее боевую мощь. Коснулся значения торговых переговоров с Англией и необходимости развития внешней торговли.

Он особо остановился на значении концессий как с точки зрения международной политики в смысле заинтересовывания капиталистов концессиями у нас, так и, главным образом, в целях ускорения экономического развития страны. Говорил о концессиях продовольственных, лесных и горных; подчеркивал опасность, которую концессии несут с собой, но оговаривал, что эта опасность меньшая, чем если бы мы от них отказались и тем самым замедлили наш экономический рост.

Разработка Лениным вопроса о концессиях, специальный доклад на коммунистической фракции съезда Советов и заключительное слово по этому докладу были важной составной частью вырабатывавшейся им в этот период новой экономической политики.

Хозяйственный фронт выдвигается теперь как самый главный и основной. «Сейчас весь гвоздь политического момента, — говорил Ленин, — состоит в том, что мы переживаем как раз переломный, переходный период, некоторый зигзаг, — период, когда от войны мы переходим к строительству хозяйственному»[93].

Успехи гражданской войны надо видеть в том, что мы убедили крестьян, что руководство пролетариата спасает их от эксплуатации и насилия со стороны капиталистов и помещиков. Меньшевики, эсеры и все враги советского строя не уставали обвинять большевиков в принуждении крестьянства, выдавая это чуть ли не за сущность диктатуры пролетариата. Далеко не все и в партии понимали соотношение принуждения и убеждения, оппортунисты и тут вносили путаницу. Разъясняя сущность диктатуры пролетариата, умело соединяющей принуждение и убеждение, убедившей крестьянство, что руководство пролетариата есть единственное средство, которое спасает его от эксплуатации и насилия, Ленин резюмировал: «И только потому, что мы могли убедить в этом крестьянина, только поэтому наша политика принуждения, основанная на этом прочном и безусловном убеждении, имела такой гигантский успех»[94]. Теперь, говорил Ленин, надо убедить крестьян в том, что при таком руководстве пролетариата молено и нужно провести в жизнь великий хозяйственный план, восстановить сельское хозяйство, промышленность, а впоследствии постепенно реорганизовать и сами основы мелкого крестьянского хозяйства.

В связи с этим Ленин дал характеристику и оценку аграрного законопроекта Совнаркома об укреплении и развитии сельскохозяйственного производства и помощи крестьянскому хозяйству. Это было первым шагом в изменении политики партии по отношению к крестьянству. Ленин уже тогда предлагал и дальнейшие шаги, подводил к вопросу о замене продразверстки продналогом, однако в силу ряда причин разработка их задержалась, и дальнейшее развитие они получили на X съезде партии.

Ленин многократно подчеркивал в своем докладе необходимость обсуждать на съездах и на собраниях уже не только политические, а во все большей мере хозяйственные вопросы, чтобы на трибунах съездов, собраний побольше выступало агрономов, инженеров, чтобы изучался опыт районов, предприятий, отраслей в целях ускорения хозяйственного строительства. Он даже сказал: «Самая лучшая политика отныне — поменьше политики»[95], к чему сразу же, на следующий день, пытался придраться меньшевик Дан, а позднее оппозиционеры и фракционеры в партии.

Центральным в ленинском докладе было разъяснение целей и значения плана электрификации (ГОЭЛРО). Основой его явились директивы, утвержденные еще IX съездом партии. План за это время был тщательнейшим образом детализирован и уточнен при непосредственном участии Ленина. К работе над ним было привлечено около 200 виднейших ученых и специалистов. Теперь план в качестве особого вопроса стоял в повестке дня съезда Советов.

Ленин назвал его второй Программой партии. Программа партии дополнялась планом работ по воссозданию всего народного хозяйства на основе современной техники. Сердцевиной его была электрификация страны, без чего, как подчеркивалось в докладе, нельзя было перейти к действительному строительству. Именно здесь Ленин и сформулировал свою бессмертную мысль: «Коммунизм — это есть Советская власть плюс электрификация всей страны»[96].

План был рассчитан на 10–15 лет. За этот срок в России предстояло построить новые тепловые и гидроэлектростанции общей мощностью 1,5 млн. киловатт. Цифра, конечно, скромная по нынешним масштабам, но по тем временам для России это был гигантский шаг вперед.

Ленин утверждал, что только когда Россия будет электрифицирована, когда под промышленность, сельское хозяйство, транспорт удастся подвести современную техническую базу, только тогда победа революции будет окончательной. Таким образом, ленинский план ГОЭЛРО был первым в истории России комплексным планом развития всех ее производительных сил, подлинно революционного преобразования технически отсталого, аграрно-промышленного государства в передовую индустриальную социалистическую державу.

Здесь мне хочется сослаться на роман «Хождение по мукам» известного советского писателя Алексея Толстого, чьи произведения я весьма ценю. Он образно и правдиво описал обстановку, в которой VIII съезд Советов принимал ленинский план электрификации России. Я привожу отрывок из романа Алексея Толстого с некоторыми сокращениями.

«В пятиярусном зале Большого театра, в тумане, надышанном людьми, едва светились сотни лампочек красноватым накалом. Было холодно, как в погребе. На огромной сцене... с колосников свешивалась карта Европейской России, покрытая разноцветными кружками и окружностями, — они почти сплошь заполняли все пространство. Перед картой стоял маленький человек, в меховом пальто, без шапки это был инженер, а впоследствии советский академик, Глеб Кржижановский, соратник Ленина еще со времен петербургского «Союза борьбы», один из активных создателей плана ГОЭЛРО. — А. М.)... В руке он держал длинный кий и... указывал время от времени концом кия на тот или иной цветной кружок, загоравшийся тотчас столь ярким светом, что... становились видны напряженные, худые лица, с глазами, расширенными вниманием.

Поднимая кий, он указывал на будущие энергетические центры и описывал по карте окружности, в которых располагалась будущая новая цивилизация, и кружки, как звезды, ярко вспыхивали в сумраке огромной сцены. Чтобы так освещать на коротенькие мгновения карту, — понадобилось сосредоточить всю энергию московской электростанции, — даже в Кремле, в кабинетах народных комиссаров, были вывинчены все лампочки, кроме одной — в шестнадцать свечей».

Это описание не художественный вымысел, все исторически, можно сказать, документально верно. Я могу говорить об этом с уверенностью, так как сам был участником этого съезда Советов. Да, тогда так жила и Москва, и вся страна. Люди страдали от голода, не было топлива, электроэнергии производилось ничтожно мало.

В 1920 году, когда принимался план ГОЭЛРО, Россия произвела всего 500 млн. киловатт-часов электроэнергии. А в 1969 году было выработано 689 млрд. киловатт-часов, больше, чем в Федеративной Республике Германии, Англии и Франции, вместе взятых (594 млрд. киловатт-часов).

Прочитав описание Алексея Толстого, перестаешь, пожалуй, удивляться, что английский писатель-фантаст Герберт Уэллс, приезжавший в Россию в 1920 году и беседовавший с Лениным, но поверил и возможность осуществления программы электрификации Советской страны. В книге «Россия во мгле», вышедшей на русском языке два года спустя, он назвал план ГОЭЛРО утопией.

Уэллс воздал, однако, должное убежденности Ленина в том, что грандиозный план обязательно будет осуществлен: «Он видит... как подымается обновленная и счастливая, индустриализированная коммунистическая держава. И во время разговора со мной ему почти удалось убедить меня в реальности своего провидения».

Приехав к нам вторично, уже в 1934 году, Уэллс кардинально изменил свое мнение и не мог не отметить поразительных успехов Советского Союза за такой короткий исторический срок. Принятие съездом постановления «Об электрификации Республики» проходило в обстановке огромного подъема. Это ярко отразилось и в стиле принятого решения, где записано, что съезд расценивает «план электрификации России, как первый шаг великого хозяйственного начинания». Съезд поручил Совнаркому разработать постановление «о поголовной мобилизации всех, обладающих достаточной подготовкой, научной или практической, для пропаганды плана электрификации и преподавания необходимых знаний для его понимания». Решение заканчивается выражением уверенности, что все рабочие и трудящиеся крестьяне «не остановятся ни перед какими жертвами для осуществления плана электрификации России во что бы то ни стало и вопреки всем препятствиям»[97]. Мне остается добавить, что ленинский план ГОЭЛРО был выполнен в части электрификации уже к началу 1932 года.

Рожденный ленинским гением, этот план стал основой превращения Советской страны в страну сплошной электрификации.

\subsection*{\center * * *}

В прениях на съезде Советов от партии социал-демократов выступил известный лидер меньшевиков Дан, который заявил о признании огромных побед Советской власти на фронтах гражданской войны, но пытался подвергнуть критике внешнюю политику правительства, а также перспективы внутренней политики, начертанной Лениным. Выступил также представитель от меньшинства социалистов-революционеров (эсеров) Вольский. Дан обвинял Советское правительство в «тайной дипломатии», уверял, будто политика Советской власти непонятна западноевропейским пролетариям, а Коминтерн вносит разрозненность в их ряды. Он объявлял немарксистским призыв соединить борьбу пролетариата с национально-освободительной борьбой, брал под защиту социал-демократов Запада и меньшевиков Грузии. В области внутренней политики основное обвинение Дана сводилось к тому, будто власть рабочих острием своим направлена теперь против крестьян. Он договорился и до того, что главное, мол, «не в том, много или мало у нас будет электрических лампочек», и т. д. Ораторы имели время для выступления не более 15 минут. Однако когда выступал Дан, он дважды просил продлить время, и ему продлевали, выступал он более получаса. Съезд слушал его внимательно не потому, что на съезде оппозиция имела поддержку, а просто делегаты интересовались выдвигаемыми им аргументами.

И не случайно после заключительного слова Ленина, когда голосовалась резолюция меньшевиков, ни один голос не был подан за эту резолюцию. Все коммунисты и беспартийные были против.

Ленин оценил выпады меньшевиков против нашей внешней политики как пособничество империалистам, во вред Советской власти (Мартов, будучи уже за границей, выступал против Советской власти). Во внутренней политике это — смыкание с капиталистическими элементами, курс на восстановление капитализма. (Дан предлагал предоставить русским капиталистам те же права и льготы, которые предусматривались для иностранных концессионеров.)

Душой съезда был Ленин. Отношение к Ленину делегатов съезда, а большинство из них впервые видело Ленина, ибо многие до этого были в подполье, на окраинах, было исключительно хорошим. Открыто проявлялась искренняя любовь к нему, уверенность, что все сказанное им — это истинная правда.

Во всю силу проявилось умение Ленина вскрыть предстоящие трудности, вселить веру в возможность победы социализма, умение сосредоточить внимание на главных узловых вопросах, умение не только двигать дело вперед, но и требовать учиться на практическом опыте, проверять планы и решения в ходе их исполнения, внося в них необходимые изменения и уточнения. Значение съезда наряду с принятием плана электрификации состояло в том, что на нем Ленин доказал необходимость подведения новой экономической базы под военно-политический союз рабочих и основных масс крестьянства, а именно базы хозяйственного союза, укрепления и развития сельскохозяйственного производства. Ленин говорил на VIII съезде, что «без практического массового улучшения хозяйства мелкого крестьянства нам спасения нет: без этой базы невозможно никакое хозяйственное строительство, и какие бы то ни было великие планы — ничто»[98]. В осуществление этой задачи съезд принял решение «О мерах укрепления и развития крестьянского сельского хозяйства».

Съезд уделил также большое внимание важности согласования деятельности местных органов хозяйственных наркоматов, создав в этих целях губернские экономические совещания при губисполкомах на правах их комиссий. Этим самым усиливалась роль местных органов в управлении производством и их привлечение к выполнению общегосударственных задач. В аналогичном направлении съездом вынесено постановление «О советском строительстве», где даются указания о более тесной связи и установлении нормальных взаимоотношений центральных органов власти с местными Советами.

На заключительном заседании 29 декабря 1920 года съезд принял обращение «Ко всем трудящимся России», где была выдвинута решающая задача текущего момента— «приступить к мирному труду. Отдадим же этому труду все силы»[99]

\newpage
\section[Перед X съездом партии: партия больна, партию треплет лихорадка]{\center \textbf{Перед X съездом партии: партия больна, партию треплет лихорадка}}


Исторический поворот от войны к мирному строительству проходил, однако, нелегко. Внутреннее положение в стране было крайне тяжелое. Разруха царила в промышленности и на транспорте. Выплавка стали в 1920 году по сравнению с довоенным уровнем сократилась в 21 раз. На X съезде партии, например, приводились данные о том, что гвоздей было произведено по весу в полтора раза меньше, чем в царское время... касторки! Угольная промышленность Донбасса была разрушена, Баку давал мало нефти — отсюда жестокий топливный кризис. Многие предприятия, в том числе такие, как Путиловский завод в Петрограде, не работали. Сельское хозяйство находилось в состоянии крайнего упадка, что осложнялось опустошительной засухой 1920 года в Поволжье, где начался голод. Продукция сельского хозяйства в стране в том году составила менее двух третей довоенной. Повсеместно обострилось продовольственное положение. Весной 1921 года пришлось снова снизить нормы выдачи хлеба населению в крупнейших промышленных центрах, в том числе в Москве и Петрограде. Все это дополнялось трудностями, порожденными демобилизацией огромной армии (перевозка людей, устройство на работу, снабжение продовольствием, упадок деревни).

Многие рабочие разбежались по деревням, чтобы как-то прокормиться, продержаться, а многие оставшиеся занимались производством зажигалок и подобных изделий, продажей их на рынке. Численность рабочих в промышленности сократилась вдвое против 1913 года. Росло недовольство демобилизованных, городского и сельского населения. Во многих местах действовали террористические, бандитские отряды, мешавшие работе органов Советской власти. Перед X съездом партии вспыхнул контрреволюционный мятеж в Кронштадте. Мелкобуржуазные настроения захлестывали крестьянство. Началось брожение и в рабочем классе, возникали забастовки, колебания в его политических настроениях, в его отношении к своему авангарду и к Советам. Все это усугублялось тем, что и в партии положение было критическим.

В условиях гражданской войны и интервенции, естественно, все было подчинено интересам войны. Демократия, как внутрипартийная, так и общесоветская, в силу необходимости была свернута. И если это было оправданным во время войны, то фактическое продолжение этих методов, подчас вынужденное, уже при переходе к мирному строительству стало одной из причин, вызывавших недовольство в партии и народе.

Сильно изменился и состав партии. Количество членов партии возросло до 700 тыс. человек. В профсоюзах тогда было около 7 млн. человек. При этом лучшие и испытанные коммунисты были мобилизованы на фронт, а затем на хозяйственно-административную работу по управлению страной, что вызывало ослабление влияния старых партийных кадров и большевистских традиций на предприятиях и в рабочем классе в целом.

Увеличение числа членов партии благодаря призывам во время так называемых «партийных недель» привело к тому, что в ней оказалось много людей хотя и преданных партии, но недостаточно политически воспитанных, не закаленных борьбой, легко поддававшихся в тяжких условиях жизни мелкобуржуазному влиянию. К тому же в годы гражданской войны в партию в массовом порядке было принято немало бывших меньшевиков, бундовцев, боротьбистов и других выходцев из правых партий и группировок, которые перешли на сторону Советской власти и большевистской партии. В ее ряды проникли и люди, которые были не прочь примазаться к господствующей партии з погоне за привилегиями.

В связи с этим во многих партийных организациях возникало состояние разброда и недовольства, что явилось почвой, удобной для развития группировок со своими платформами, оттенками и оттеночками.

В такой во всех отношениях сложной обстановке и развертывалась предсъездовская дискуссия о профсоюзах. Внешне в ней как будто даже и не затрагивались основные вопросы, которые стали перед партией и народом. Но в действительности речь шла о методах подхода к массам, о взаимоотношениях между партией и рабочим классом, между пролетариатом и крестьянством, когда страна приступала к строительству социализма в мирных условиях. Речь шла о дальнейшей судьбе партии и всей Советской страны.

Дискуссия была навязана партии Троцким, который в этот переломный и очень сложный момент предпринял попытку, для того чтобы дать новый бой Ленину и ленинизму. Поводом, сулившим ему наиболее обещающие возможности, Троцкий считал вопрос о роли и задачах профсоюзов, ибо партия начала перестраивать их работу в соответствии с мирными условиями.

Ленин поначалу возмутился: он считал, что в такое время нельзя растрачивать силы партии на бесплодные споры о профсоюзах. Место и роль профсоюзов были определены Программой партии, принятой в 1919 году VIII съездом партии. В 1920 году IX съезд партии принял специальное решение «По вопросу о профессиональных союзах и их организации», в котором профсоюзы определялись как школа коммунизма, уточнялись некоторые формулировки Программы. В сентябре 1920 года ЦК обратился к губкомам и укомам партии со специальным письмом, в котором наметил меры помощи профсоюзам. В апреле этого же года состоялся III съезд профсоюзов, в котором участвовал Ленин. Его выступление определило основное содержание принятых съездом резолюций. Вопросы профсоюзов обсуждались на сентябрьском пленуме ЦК, и никаких существенных разногласий не возникало.

Именно в этом смысле Ленин обращает внимание партии на то, что «все разногласия у Троцкого выдуманные»[100], что таких разногласий нет. «Троцкий старался их указать, но не мог»[101]. Ленин настаивает на том, что «пора перейти от выдумывания и преувеличивания принципиальных разногласий к деловой работе»[102]. Ленин старается вести дискуссию, когда она становится неизбежной, «в вежливой и товарищеской форме...»[103], еще и еще раз призывает к переходу «от интеллигентских разговоров о пустых разногласиях к деловой работе»[104].

Но не получается деловой работы. Троцкому нужна была дискуссия, платформа, организованная фракция с дальним прицелом, с претензией на руководство партией.

Здесь, пожалуй, уместно напомнить характеристику, данную Лениным Троцкому в статье «О нарушении единства, прикрываемом криками о единстве» еще в 1914 году, уже после организации Троцким беспринципного антипартийного «Августовского блока» (1912 год):

«Старые участники марксистского движения в России хорошо знают фигуру Троцкого, и для них не стоит говорить о ней. Но молодое рабочее поколение не знает ее, и говорить приходится, ибо это — типичная фигура для всех тех пяти заграничных группок, которые фактически также колеблются между ликвидаторами и партией.

Во времена старой «Искры» (1901–1903) для этих колеблющихся и перебегающих от «экономистов» к «искровцам» и обратно была кличка: «тушинский перелет» (так звали в Смутное время на Руси воинов, перебегавших от одного лагеря к другому)...

«Тушинские перелеты» объявляют себя выше фракций на том единственном основании, что они «заимствуют» идеи сегодня одной, завтра другой фракции. Троцкий был ярым «искровцем», в 1901–1903 годах, и Рязанов назвал его роль на съезде 1903 года ролью «ленинской дубинки». В конце 1903 года Троцкий — ярый меньшевик, т. е. от искровцев перебежавший к «экономистам»; он провозглашает, что «между старой и новой «Искрой» лежит пропасть». В 1904–1905 году он отходит от меньшевиков и занимает колеблющееся положение, то сотрудничая с Мартыновым («экономистом»), то провозглашая несуразно-левую «перманентную революцию». В 1906–1907 году он подходит к большевикам и весной 1907 года заявляет себя солидарным с Розой Люксембург.

В эпоху распада, после долгих «нефракционных» колебаний, он опять идет вправо и в августе 1912 года входит в блок с ликвидаторами. Теперь опять отходит от них, повторяя, однако, по сути дела их же идейки.

Такие типы характерны, как обломки вчерашних исторических образований и формаций, когда массовое рабочее движение в России еще спало и любой группке «просторно» было изображать из себя течение, группу, фракцию, — одним словом, «державу», толкующую об объединении с другими.

Надо, чтобы молодое рабочее поколение хорошо знало, с кем оно имеет дело, когда с невероятными претензиями выступают люди, не желающие абсолютно считаться ни с партийными решениями, которые с 1908 года определили и установили отношение к ликвидаторству, ни с опытом современного рабочего движения в России, создавшего на деле единство большинства на почве полного признания указанных решений»[105].

И в этот период, в конце 1920 года, Троцкий вновь пользуется своими старыми методами борьбы с ленинизмом, так беспощадно раскритикованными Лениным. Он, как и прежде, выступает с претензией на диктат, на командование, не считаясь с партийными решениями, не подчиняясь партийной дисциплине.

А теперь можно рассказать и о самой дискуссии, которая вошла в историю партии под названием «профсоюзной».

\subsection*{\center * * *}

Прибыв на VIII съезд Советов в Москву, мы узнали, что в ЦК партии и в центральных профсоюзных кругах идет дискуссия и ожесточенная борьба разных точек зрения. Борьба завязалась на V Всероссийской конференции профсоюзов в первые дни ноября. Выступая на коммунистической фракции этой конференции, Троцкий бросил там крылатые слова о необходимости «перетряхнуть профсоюзы», требовал их «огосударствления», милитаризации, ссылаясь при этом на опыт руководства находившегося под его влиянием (он занимал в то время также должность наркома путей сообщения) профсоюза транспортных рабочих и его ЦК (Цектрана).

На самой конференции, после обсуждения, принимаются правильные, поддержанные впоследствии Лениным, тезисы Рудзутака о производственных задачах профсоюзов. Казалось бы, вопрос исчерпан...

Однако Троцкий, закусив удила, переносит борьбу в ЦК, ополчившись на первых порах на руководившего тогда профсоюзами Томского, «посмевшего» возразить против «перетряхивания». Центральный Комитет поручает особой комиссии с участием Троцкого в деловом порядке рассмотреть вопросы, вызывающие разногласия, и подготовить предложения для ЦК.

Троцкий, нарушая дисциплину, неожиданно и без всяких оснований отказывается работать в этой комиссии, что означало «продолжение борьбы и вынесение ее за пределы ЦК»[106], становится на путь создания фракции, раскалывает Цектран, делает все, чтобы столкнуть профсоюзы с партией, а внутри ЦК прилагает все силы для вербовки новых сторонников.

Во время работы VIII съезда Советов Троцкий выпускает брошюру-платформу, в которой заявляет, что партийному «съезду придется выбирать между двумя тенденциями».

Политический смысл такого заявления с точки зрения дальнего прицела ясен. Троцкий надеялся, что, пользуясь трудностями переломного периода, он вызовет кризис в партии, через ожесточенную внутрипартийную борьбу в масштабе всей страны ему удастся добиться того, что партия в итоге выберет его линию и отвергнет ленинскую. Трудно даже представить себе все пагубные политические последствия для партии и страны в случае такого оборота дела. Троцкий рвался к власти, строя свои расчеты на том, что он всего добьется, находясь у руля руководства армией и транспортом. В Секретариате ЦК сидели его сторонники. Все это вскружило голову Троцкому, и опять проявилось его подлинное нутро — нутро дезорганизатора партии, опасного врага большевизма, каким он был еще в дореволюционный период.

Преступный перед партией акт — выдвижение совершенно ошибочной платформы и создание на ее основе фракции — был вызовом Центральному Комитету со стороны Троцкого. Это было и попыткой отвлечь партию от самых насущных проблем устранения хозяйственной разрухи.

И действительно, как я уже рассказывал, съезд Советов был поглощен обсуждением на своих заседаниях, на секциях, на коммунистической фракции животрепещущих вопросов. А в его кулуарах муссировались сенсационные слухи об острой борьбе среди работников профдвижения, о расхождениях и группировках среди членов ЦК партии.

Начатая Троцким борьба оживила оппортунистов всех мастей, и вскоре ленинской линии в профсоюзном движении были противопоставлены уже семь различных платформ, из которых надо выделить три главные:

а) Троцкого и «буферную» Бухарина, которые затем объединились в одну платформу Троцкого — Бухарина. Вокруг этой платформы быстро сколотилась фракция, в которую вошла группа руководящих деятелей партии и государства. Она отражала извращенное понимание главных и текущих задач партии, была носителем бюрократических и «военизированных» методов работы в партии, профсоюзах, хозяйственных органах и Советах;

б) «рабочая оппозиция» во главе со Шляпниковым, отражавшая уклон в сторону анархо-синдикализма. Она настаивала на передаче управления народным хозяйством «съезду производителей», считала профсоюзы высшей формой организации рабочего класса. Профсоюзы должны были выбирать органы управления промышленностью, а партия и Советы лишались даже права отвода кандидатур. Ленин считал и этот синдикалистский уклон крайне опасным, антикоммунистическим. Несмотря на кажущуюся противоположность троцкизму («огосударствливание профсоюзов» у одних, «опрофсоюзивание государства» у других), они начали находить между собой общий язык и поддерживали друг друга;

в) группировка «демократического централизма» во главе с Сапроновым — нечто среднее между троцкистами и анархо-синдикалистами. Извращая принципы демократического централизма, они выступали за слияние Совнаркома с Президиумом ВЦИК, требовали коллегиального «демократического» управления в промышленности, принижали роль партии в руководстве работой Советов, вносили невообразимую путаницу в вопросы партийного строительства.

Менее значительными были платформы Игнатова, Рязанова, Ногина и др., но и они обостряли внутрипартийную борьбу, осложняли подготовку партии к X съезду.

\subsection*{\center * * *}

К концу работы VIII съезда Советов мы неожиданно поздно вечером узнали, что по окончании его состоится дискуссионное собрание коммунистической фракции съезда совместно с ВЦСПС и руководством московских профсоюзов. Стало известно также, что на днях состоялось совещание части делегатов съезда с работниками профсоюзов, на котором со своей платформой выступал Троцкий. Я, в частности, не был приглашен туда.

Большой театр, как уже говорилось, едва вмещал делегатов съезда. К ним теперь прибавилось несколько сот работников профсоюзов. Еще задолго до начала заседания зал был набит битком, многие стояли позади, в проходах, в уголках зала. С мест раздавались голоса, будто в зал проникли и беспартийные делегаты, что их надо удалить с фракционного заседания. Но была такая суматоха, что предпринять что-либо было уже невозможно.

Мне удалось заранее занять очень удобное место — в третьей ложе справа, если смотреть со сцены. Было хорошо видно все, что происходит в зале: как волновались собравшиеся, какой острый интерес проявляли они к тому, что предстояло выслушать.

Первым выступил Зиновьев — с критикой платформы Троцкого. Критиковал он также и бухаринскую «буферную» платформу, и неправильную позицию Шляпникова. Зиновьев был не глубоким, но искусным оратором, с изрядной долей демагогии. Ему удалось несколько раз вызвать одобрение и аплодисменты удачными шутками, едким высмеиванием противников, острыми ударами по ним.

Со стороны оппозиционеров раздавались возмущенные выкрики против докладчика. Но у меня уже в ходе доклада и наблюдений за реакцией делегатов сложилось убеждение, что большинство собрания поддерживает ленинскую линию, и это очень радовало.

Вторым выступил Троцкий. Как оратор он был сильнее Зиновьева. Своими остротами, отчеканенными фразами и крылатыми словечками, артистичностью ему удавалось добиться шумного одобрения и аплодисментов.

Ожидалось, что выступление Бухарина, представлявшего «буферную» платформу, будет направлено на примирение «крайних» позиций, будет не обострять, а смягчать разногласия. Вышло же наоборот. Бухарин обрушился с резкими нападками на ленинскую линию, не делая никаких выпадов против платформы Троцкого, не критикуя его неправильных положений. Он пошел даже дальше, взяв одновременно под защиту одно из синдикалистских предложений Шляпникова.

Выступление Шляпникова, изложившего свою платформу, было бледным, малоубедительным и не произвело особого впечатления на собрание. Резко демагогическим было выступление опытного оратора Сосновского, рьяного сторонника Троцкого.

У меня складывалось удручающее настроение. Речи нескольких оппозиционных ораторов против одного сторонника ленинской линии создавали ощущение, будто все большее количество колеблющихся перекидывалось на сторону оппозиции.

Возникали невеселые мысли. Чем кончится собрание? Поражением ленинской линии? Отражает ли это положение в партии? Какой оборот примет уже разгоравшаяся вовсю общепартийная дискуссия и как она повлияет на предстоящий через два месяца X съезд партии?

В это время вдруг, неожиданно для всех, на сцене появляется Ленин. Видимо, ему сообщили, какой неблагоприятный для партии оборот приняла здесь дискуссия.

В критический для партии момент, несмотря на болезнь, Ленин решил прийти и с ходу выступить на этом очень ответственном всероссийском собрании партийных, советских и профсоюзных работников.

Сидя близко к сцене, я хорошо видел, как болезненно выглядел Ленин. Он был крайне взволнован, возбужден. Когда он начал говорить, стало заметно, как грудь его колышется от учащенного сердцебиения. В зале наступила необычная для той бурной ночи тишина. Делегаты были явно озадачены, слушали Ленина с напряженным вниманием, чувствовалось, что волнение Ленина передалось и им.

Ленин начал с извинения за то, что вынужден выступить, не выслушав речей предшествовавших ему ораторов. Но он настолько нездоров, что не был в состоянии поступить иначе. Однако вчера он прочитал опубликованные в ходе дискуссии материалы, основным из которых является брошюра Троцкого. Вчитываясь в нее, сказал Ленин, «я удивляюсь, какое количество теоретических ошибок и вопиющих неправильностей сконцентрировано в ней»[107].

Разобрав по косточкам положения Троцкого и Бухарина, он разбил их в пух и прах, неопровержимо доказав их несостоятельность и теоретическую фальшивость.

По вине Троцкого, говорил Ленин, в ЦК создалась «каша и кутерьма; это в первый раз в истории нашей партии во время революции, и это опасно. Гвоздем было то, что получилось раздвоение, получилась «буферная» группа Бухарина, Преображенского и Серебрякова, которая больше всех навредила и напутала»[108]. Если бы я умел рисовать карикатуры, заметил Ленин, я бы Бухарина нарисовал с ведром керосина, подливающим его в огонь. Смыкание Бухарина с Троцким и его роль в дискуссии Ленин подытожил кратко и образно грибоедовским афоризмом: «Шел в комнату, попал в другую».

«Надо ли было, — спрашивал Ленин, — выносить все такие споры на широкую дискуссию? Заниматься этим безделием? Занимать нужные для нас недели перед партийным съездом?»[109] Ответы его на эти вопросы были недвусмысленно отрицательными.

Разобрав и высмеяв «теоретические выдумки» Троцкого о «сращивании» профсоюзов с хозорганами и «производственную демократию» Бухарина, Ленин заключил: «и принципиально, и теоретически, и практически один вывод — про тезисы Троцкого и позицию Бухарина: унеси ты мое горе!»[110]

Ленин упрекал оппозиционеров и в том, что они совершенно игнорируют текущий момент, его политическую сторону. А суть этого момента заключалась прежде всего в том, что окончание войны, переход к мирному строительству настоятельно требовали пересмотра политики по отношению к крестьянству. Здесь были самые большие трудности, игнорирование которых ставило под вопрос союз рабочего класса с крестьянством.

В тезисах и платформах оппозиции и не пахло пониманием этого. Более того, Троцкий и его сторонники голосовали в ЦК 12 января против принятой пленумом ЦК в тот день резолюции: «не преувеличивать разногласия, не отвлекать внимание и силы членов партии от насущных и неотложных практических задач»[111].

Ленин, как это известно теперь по документам, сразу же по окончании военных действий организовал в деловых комиссиях разработку предложений об отказе от продразверстки и переходе к продналогу, что стало позже стержнем новой экономической политики, заменившей экономическую политику военного коммунизма. Ленин был намерен уже на VIII съезде Советов провести решение об отмене денежных налогов (что и было сделано) и о введении продовольственного налога взамен продразверстки. Но вопрос о продналоге не был вынесен на обсуждение съезда, так как предложения не были еще подготовлены, и Ленин по хотел и таком важном вопросе проявлять торопливость. Он внимательно присматривался, изучал настроения крестьян, в частности беспартийных крестьян — участников съезда Советов, и поставил эти вопросы на X съезде партии.

А в это самое время Троцкий твердил о «перетряхивании» профсоюзов, требовал сделать это центральной задачей партии. Ленин с жаром разоблачал его бюрократический подход к массовым организациям, убеждал собравшихся в Большом театре, что это вопрос, связанный с самой сутью понимания диктатуры пролетариата. Это уже была политическая квалификация. Ленин говорил, что у нас не получается дружная работа «из-за расхождения по вопросу о методах подхода к массе, овладения массой, связи с массой. В этом вся суть»[112]. Однако в речи Ленина явно звучало желание не обострять ситуацию в партии, предпринять еще одну попытку наладить деловую работу, не поступаясь, естественно, принципами. С возмущением цитируя утверждение Троцкого, что партии надо будет «выбирать между двумя тенденциями», которое «совсем звучит странно», Ленин в этом первом своем выступлении в публичной дискуссии о профсоюзах еще не обвиняет прямо Троцкого и Бухарина в фракционной деятельности, о чем так убедительно скажет он в своих последующих речах и статьях.

По всему духу выступления чувствовалось, что Ленин уже знал о создании фракций с особой платформой, с определенным составом участников, со своей дисциплиной и о том, что оппозиционеры стали на путь раскола партии. В статье «Кризис партии» Ленин определит дату создания Троцким фракции: 25 декабря. Но здесь Ленин говорит еще о платформах и лишь намекает, что они перерастают во фракции, старается по возможности предотвратить или, по крайней мере, не подталкивать их окончательное оформление. Он называет и зачитывает платформу, которая уже принята профсоюзами, — это тезисы Рудзутака «Производственные задачи профсоюзов».

Ленин остро критикует самого себя за то, что он «проглядел» эти тезисы и только за день до выступления с ними ознакомился. «Вот это — платформа, — восклицал Ленин, — она во сто раз лучше и того, что написал т. Троцкий, много раз обдумав, и того, что написал т. Бухарин (резолюцию пленума 7 декабря), совершенно не обдумав»[113].

Подытоживая, он сказал, что вынесение разногласий на широкую партийную дискуссию и на партийный съезд есть величайшая ошибка. Заявив, что тезисы Троцкого — политически вредная вещь, Ленин, под бурные продолжительные аплодисменты закончил: «И наш партийный съезд, я уверен, эту политику осудит и отвергнет»[114].

Надо было присутствовать в зале, чтобы до конца прочувствовать, как Ленин, находившийся в нелегком положении, так как не слышал ни докладчика, ни оппозиционных ораторов, с нарастающей силой, шаг за шагом овладевал умами и чувствами большинства собравшихся. Это выражалось в напряженном внимании, с которым его слушали, в аплодисментах, которые раздавались в наиболее острые моменты его речи. Он не пользовался эффектными ораторскими приемами, говорил просто, понятно. Покорял железной логикой своих суждений. Покорял искренностью своего выступления.

Сердца его сторонников наполнялись радостью, гордостью за своего вождя. Поздно ночью мы расходились, удовлетворенные, исполненные уверенности, что и нам на местах на партийных собраниях, на губернских конференциях удастся, опираясь на ленинские высказывания, нанести поражение оппозиционным фракциям.

\subsection*{\center * * *}

Большую поддержку ленинской линии в начале января 1921 года, спустя несколько дней после описанной мной дискуссии 30 декабря, оказала петроградская партийная организация. Это произошло на четырехтысячном собрании коммунистов, где от оппозиции выступал Бухарин и где всего лишь 20 человек голосовало против ленинской линии. На районных партийных собраниях Петрограда против оппозиционеров голосовало 95–98 % присутствовавших. Это было первое тяжелое поражение, нанесенное питерцами всем оппозиционным фракциям. Петроградская организация выпустила обращение к партии, направленное против платформы Троцкого. В Петрограде было издано отдельной брошюрой и выступление Ленина на дискуссии 30 декабря. Ее мы получили на местах, и она очень помогла партийным кадрам разобраться в сути дискуссии о профсоюзах. Из Центрального Комитета поступали, однако, в подавляющих количествах троцкистские материалы, поскольку секретарями ЦК были сторонники Троцкого.

Ленин высоко оценил обращение петроградцев к партии, указав, что им открывается новый этап внутрипартийной борьбы, что это «большой шаг вперед к оздоровлению»[115], «переход от борьбы фракций, образуемых сверху, к вмешательству организаций снизу»[116].

Иное положение было в Москве. Здесь чувствовался идейный разброд. Московский комитет занимал промежуточную, «буферную» позицию, в нем действовали все оппозиционные группы, троцкисты, шляпниковцы и сапроновцы пробрались в руководство ряда районных комитетов. Дело дошло до такого, по выражению Ленина, курьеза, что большинством в один голос МК принял решение, осуждающее обращение питерцев к партии, оценивающее его как «опасную» попытку создать особый центр по подготовке съезда партии. Это выступление МК, которое было опубликовано в «Правде» одновременно с питерским обращением (пролежавшим в портфеле редакции 10 дней, так как редактором центрального органа был в то время Бухарин), оказалось фактической поддержкой Троцкого. Оно подлило еще порцию керосина в огонь дискуссии, подбодрило сторонников оппозиционных групп.

Перелом в ходе дискуссии был создан 14 января опубликованием ленинской платформы в виде проекта постановления X съезда партии «Роль и задачи профессиональных союзов», подписанной Лениным и еще семью членами ЦК и кандидатом Петровским, а также Лозовским (как членом комиссии ЦК для подготовки этого решения) и вошедшей в историю нашей партии как «платформа десяти».

В ней подтверждались прежние принципиальные решения партии о профсоюзах, как школе управления, школе хозяйствования, школе коммунизм;»; определялось место профсоюзов в общей системе диктатуры пролетариата; формулировались новые задачи профсоюзов в обстановке окончания гражданской войны и перехода к социалистическому строительству. Подчеркивалось, что основой деятельности профсоюзов должен стать метод воспитания и убеждения, широкое развертывание демократических начал.

Появление ленинской платформы подтверждало, однако, тревожные разговоры о сложности положения внутри Центрального Комитета. Ее подписало лишь восемь членов ЦК из 19, т. е. меньшинство. Остальные или придерживались других платформ, или колебались. Соотношение сил в ЦК, следовательно, мало еще отличалось от самого критического момента, возникшего на пленуме ЦК 7 декабря, когда, пользуясь словами Ленина, впервые при голосовании «наша группа чекистов остается в меньшинстве»[117] и восемью голосами против семи одобряется резолюция, предложенная Бухариным (за спиной которого стоял Троцкий) с требованием «производственной демократии».

Вскоре после обнародования «платформы десяти» оформляют свои платформы Бухарин, Сапронов, Игнатов. Особенное возмущение Ленина вызвала бухаринская платформа. Бухарин теперь уже «буферил» между Троцким и синдикалистами, выступая за обязательность принятия кандидатур на хозяйственные посты, выдвигаемых профсоюзами, вставал одной ногой на позиции синдикализма, совершал поворот, по характеристике Ленина, «не столько исторический, сколько истерический», прокладывая дорогу к объединению всех антипартийных групп.

«Платформа десяти» сразу же оказала оздоровляющее влияние на московскую партийную организацию, где не могло долго продолжаться противостояние сил разброда и распада силам сплочения. Вскоре на местах мы узнали, что уже 17 января состоялось расширенное заседание Московского комитета совместно с представителями районов и уездов. На этом заседании выступили представители всех платформ, включая и местную синдикалистского толка платформу «игнатовцев».

Ленинскую платформу защищал Сталин против тезисов Троцкого, Бухарина и других оппозиционеров.

Результаты голосования оказались сокрушительными для Троцкого и других оппозиционных групп. За тезисы Ленина было подано 76 голосов, Троцкого — 27, Игнатова— 25, Сапронова—11, Бухарина — 5, Шляпникова — 4, Ногина — 0, Рязанова — 0. При переголосовании двух основных платформ за Ленина было подано 84 голоса, за Троцкого — 27. Московская организация восстановила свою честь, приняв обращение к партии о единодушной поддержке ленинской платформы. Наряду с позицией петроградцев решение москвичей приобретало общепартийное значение, означало утверждение реальных возможностей преодоления «болезни партии» или, как кратко формулировал Ленин: «К партсъезду вылечить»[118].

Сформулировав такую задачу, Ленин развивает кипучую деятельность, неослабно следит за ходом дискуссии, активнейшим образом в ней участвует, организует настоящее наступление здорового большинства партии на всех противников ее единства.

21 января в «Правде» публикуется статья Ленина «Кризис партии». В ней он подробнейшим образом анализирует каждый этап развернувшейся внутрипартийной борьбы, дает политическую оценку платформ и фракций, предупреждает партию о тенденции к объединению антипартийных групп на беспринципной основе за пределами коммунизма, под лозунгом «кто больше обещает «прав» беспартийным, соединяйтесь по случаю партсъезда РКП!..»[119] Ленин доказывает, что оппозицию объединяет отречение от руководящей роли партии, и призывает партию к борьбе с идейным разбродом.

«Надо иметь мужество смотреть прямо в лицо горькой истине, — писал Ленин. — Партия больна. Партию треплет лихорадка». Ленин закончил статью мыслями о том, что «болезнью нашей партии, несомненно, постараются воспользоваться и капиталисты Антанты для нового нашествия, и эсеры для устройства заговоров и восстаний»[120], но что партия полна решимости преодолеть болезнь и, сплотившись, она не ослабнет, а окрепнет к съезду и после него.

Эта статья представляет огромный интерес и сегодня как мощный удар по троцкизму и оппортунизму разных мастей. В ней с предельной точностью и ясностью конспективно рассказано о возникновении и ходе дискуссии, обо всех перипетиях ее развития, дана исчерпывающая политическая оценка одного из самых драматических моментов во всей истории нашей партии.

Кстати говоря, в ленинской статье есть очень интересное отступление от хода изложения, которое ярко характеризует Ленина как человека, как высоко принципиального и самокритичного партийного деятеля. На дискуссии 30 декабря Ленин критиковал Троцкого за его понимание нашего государства как государства рабочего. Из такой абстракции Троцкий делал неправильные выводы по отношению к крестьянству. Ленин тогда сказал: «Не совсем рабочее, в том-то и штука. Тут и заключается одна из основных ошибок т. Троцкого... У нас государство на деле не рабочее, а рабоче-крестьянское...» И кроме того, «государство у нас рабочее с бюрократическим извращением»[121].

И вот через три недели, возвращаясь к этому вопросу, Ленин пишет, что должен исправить свою ошибку, ибо, читая отчет об этом собрании, увидел, что был неправ. «Мне надо было сказать: «Рабочее государство есть абстракция. А на деле мы имеем рабочее государство, во-1-х, с той особенностью, что в стране преобладает не рабочее, а крестьянское население; и, во-2-х, рабочее государство с бюрократическим извращением». Читатель, который захочет прочесть всю мою речь, увидит, что от этой поправки ни ход моей аргументации, ни мои выводы не изменяются»[122], т. е., иными словами, и при такой поправке критика позиций Троцкого остается в силе.

23 и 24 января Ленин выступает с докладом и заключительным словом о роли и задачах профсоюзов на коммунистической фракции съезда горнорабочих и в эти же дни пишет брошюру «Еще раз о профсоюзах, о текущем моменте и об ошибках тт. Троцкого и Бухарина» размером в 2,5 печатных листа, которую заканчивает 25 января. Это были кульминационные дни дискуссии.

Всех нас, сторонников Ленина, весьма обрадовали результаты дискуссии на комфракции съезда горнорабочих. Там против Ленина выступали Троцкий, Шляпников. При голосовании за ленинскую платформу было подано 137 голосов, Троцкого — всего 8, Шляпникова — 61 голос.

Решающее, пожалуй, значение в наступившем в ходе дискуссии переломе приобрела ленинская брошюра «Еще раз о профсоюзах». После прочтения статьи «Кризис партии» и этой брошюры мне стало предельно ясно то, о чем говорил Ленин еще 30 декабря и что развил и разъяснял здесь.

Прежде всего это бесплодность дискуссии о профсоюзах и ее ненужность.

Вредность дискуссии заключалась в том, что она отвлекала партию от решения на съезде действительно серьезных проблем, лежавших в основе кризиса, который переживала страна, и связанных с судьбой диктатуры пролетариата. Между строк можно было без труда понять, что, навязав партии дискуссию о профсоюзах, Троцкий ставил перед собой цель — изменить руководство большевистской партии, ее политическую линию.

Указав на опасность, грозящую от существования внутри партии фракций, которая могла привести партию к расколу, Ленин, мне казалось, сознательно не ставил точки над «и», не делал вытекавших из положения крайних выводов, всемерно стараясь сохранить возможность дружной работы после съезда с: теми, кто начал борьбу и поднял дискуссию против партийной линии.

В брошюре Ленин дополняет свое выступление 30 декабря в Большом театре критическим разбором выступлений ораторов, которых он тогда не слышал, но которые он изучил теперь по стенограмме, изданной ВЦСПС. В результате этой критики камня на камне не осталось от платформ Троцкого, Бухарина, Шляпникова (которого Ленин не критиковал в первом своем выступлении, так как не знал ни его речи, ни зачитанной им тогда платформы, да и считал, видимо, необходимым сосредоточить в тот момент критику на Троцком и Бухарине). Эта брошюра — блестящий образец марксистской полемики по принципиальным, политическим вопросам, волновавшим партию. Ленинская брошюра и сейчас читается с захватывающим интересом.

Ленин остро критиковал и высмеивал бухаринские утверждения. В выступлении Бухарина на дискуссии 30 декабря и в платформе его фракции говорилось, что верно и положение «платформы десяти» — «профсоюзы это школа коммунизма» и одновременно правильно положение Троцкого — профсоюзы — это «административно-технический аппарат управления производством». Бухарин считал необходимым «соединение этих обоих положений». Раскрывая теоретическую несостоятельность позиции Бухарина, Ленин обвинял Бухарина в подмене диалектики марксизма эклектизмом, широко распространенным в модных, реакционных философских системах. Ленин писал: ««И то, и другое», «с одной стороны, с другой стороны» — вот теоретическая позиция Бухарина. Это и есть эклектицизм. Диалектика требует всестороннего учета соотношений в их конкретном развитии, а не выдергивания кусочка одного, кусочка другого»[123]. Ленин утверждал, что Бухарин стоит на точке зрения логики формальной или схоластической, а не логики диалектической или марксистской. «Диалектическая логика, — подчеркивал Ленин, — учит, что «абстрактной истины нет, истина всегда конкретна», как любил говорить вслед за Гегелем, покойный Плеханов». Я понимал из этих ленинских слов, что Бухарин, хотя и теоретически эрудированный человек, но не владеет марксистской диалектикой.

\subsection*{\center * * *}

Меня тогда обрадовала эта ссылка на Плеханова и поразила в тогдашних условиях столь высокая оценка его философских трудов и то, что Ленин тут же сказал о значении философских трудов Плеханова для марксизма, для марксистского воспитания членов партии.

Ведь все мы знали, что после II съезда партии Плеханов перешел в лагерь меньшевизма, а в годы войны и Февральской революции вконец опозорил себя, заняв крайне правый фланг социал-шовинизма. Ленин писал: «...В скобках уместным, мне кажется, заметить для молодых членов партии, что нельзя стать сознательным, настоящим коммунистом без того, чтобы изучать — именно изучать — все, написанное Плехановым по философии, ибо это лучшее во всей международной литературе марксизма...»[124]

А в примечании Ленин указал, что «рабочему государству, по-моему, следует требовать от профессоров философии, чтобы они знали изложение марксистской философии Плехановым и умели передать учащимся это знание»[125].

Думается, что это указание Ленина и сегодня сохраняет свою силу. Нечего греха таить, есть у нас такие профессора философии, которые забыли о нем и его не выполняют. Очень жаль. Чем раньше это будет исправлено, тем будет лучше для марксистско-ленинского образования коммунистов.

В юности и на меня большое впечатление произвели философские труды Плеханова. Мне было приятно, что в новых условиях Ленин так умело отделил политика — меньшевика Плеханова от Плеханова-философа тех времен, когда он был марксистом.

В связи с этим мне вспомнилось выступление С. Шаумяна на заседании Бакинского Совета в 1918 году через несколько дней после смерти Плеханова, где он противопоставил двух Плехановых, сказав: «И если... он изменил своим взглядам, на которых он воспитывал целые поколения, то это можно объяснить, по его собственной теории, теми изменениями, которые произошли с его личностью...»[126]

Следует заметить, что Ленин, по свойственной ему скромности, естественно, далее намека не сделал на свои философские труды, представляющие еще большую ценность для формирования коммунистического мировоззрения членов партии.

Это умение дать многогранную оценку фактам, событиям, политическим деятелям — характерная черта Ленина. Как тонкий политик, Ленин умел разбираться в людях, видеть все стороны их характера, выделить положительные и отрицательные черты в человеке, заметить их взаимодействие и изменение по мере развития. Рассматривал он их, конечно, с точки зрения интересов социализма, и у него всегда это было главным отправным пунктом в оценке политических действий, событий, людей и их роли в них.

В этой связи привлекает внимание необычная, с обычной точки зрения, статья Ленина в газете «Правда» за декабрь 1918 года — «Памяти тов. Прошьяна». В то время я был лишен возможности знать эту статью, так как находился в закаспийской тюрьме при белогвардейцах. Значительно позже я с ней познакомился. В числе других левых социалистов-революционеров Прошьян с декабря 1917 по март 1918 года входил в состав Совнаркома в качестве наркома почт и телеграфов. Меня поразили не столько облик этого интересного человека, сколько умение Ленина так тонко разобраться в духовном складе этого революционера, в сложном переплете противоречий его мировоззрения, объяснить причины его взлетов и падений.

Ленин пишет: «Прошьян выделялся сразу глубокой преданностью революции и социализму. Не про всех левых эсеров можно было сказать, что они социалисты, даже, пожалуй, про большинство из них сказать этого было нельзя. Но про Прошьяна это надо было сказать, ибо, несмотря на верность его идеологии русских народников, идеологии несоциалистической, в Прошьяне виден был глубоко убежденный социалист. По-своему, не через марксизм, не от идей классовой борьбы пролетариата, этот человек сделался социалистом»[127].

Глубоко преданный революции и социализму, он, оказавшись волей судеб в штабе социалистического государства, повседневно сталкиваясь с практическим разрешением многочисленных новых, трудных, больших и малых проблем начального строительства социализма, вызывавших немало споров и разногласий, изо дня в день видел, как большевики во главе с Лениным решают коренные вопросы революции, как строят государство нового типа, закладывают основы социалистического строя. Находясь под прямым влиянием Ленина в течение менее чем четырех месяцев (но каких месяцев!), он лояльно сотрудничал и вел за собой некоторых других социалистов-революционеров, «становился, — как пишет Ленин, — решительно на сторону большевиков-коммунистов против своих коллег, левых социалистов-революционеров, когда они выражали точку зрения мелких хозяйчиков и относились отрицательно к коммунистическим мероприятиям в области сельского хозяйства»[128].

В ходе совместной работы они так сблизились, что до Брестского мира, «казалось, — писал Ленин, — что разногласий между нами сколько-нибудь существенных уже не осталось». Прошьян прямо стал говорить Ленину о необходимости слияния левых эсеров с большевиками.

Ленин отнесся сдержанно к этому предложению Прошьяна, назвав его преждевременным. Однако он отнюдь не отрицал сближения на практической почве.

«Полное расхождение принес Брестский мир, а из расхождения, при революционной последовательности и убежденности Прошьяна, не могла не произойти прямая, даже военная борьба... Но пример Прошьяна, — писал Ленин, — показал мне... — как разногласия в общих основах миросозерцания с неизбежностью проявили себя при трудном повороте истории».

После Бреста Прошьян вышел из состава правительства, а летом принимал участие в левоэсеровском мятеже в Москве, захватил телеграф (всего на один час!) — оповещал Россию о свержении большевиков[129].

Статью Ленин заканчивает следующим:

«А все же таки Прошьяну довелось до июля 1918 года больше сделать для укрепления Советской власти, чем с июля 1918 года для се подрыва. И в международной обстановке, создавшейся после немецкой революции, новое — более прочное, чем прежде, — сближение Прошьяна с коммунизмом было бы неизбежно, если бы этому сближению не помешала преждевременная смерть»[130].

После прочтения этой статьи сознанию читателя открывается еще одна обаятельная черта образа вождя мировой революции.

Но хочется подчеркнуть снова важнейшую особенность оценок Ленина: они всегда принципиальны, партийны. Если речь шла о судьбах революции и социализма в борьбе с идейными и политическими противниками, Ленин был непримирим, резок, беспощаден, отстаивая дело партии. Достаточно вспомнить письма Ленина от 18 и 19 октября 1917 года, когда он узнал о чудовищной измене Каменева и Зиновьева, раскрывших перед врагами секретное решение ЦК о подготовке вооруженного восстания. «Уважающая себя партия не может терпеть штрейкбрехерства и штрейкбрехеров в своей среде»[131],— пишет он.

Интересно, как Ленин отзывался об этой черте своего характера. В споре с левым эсером Камковым о Брестском мире, в ответ на его резкое провокационное выступление против большевиков, Ленин говорил: «И уж, конечно, товарищи, не мне возражать против резких слов. Никогда я против этого возражать не буду. Только, чтобы быть резким, нужно на это иметь право, а право на резкость дает то, чтобы слово не расходилось с делом. И вот такое маленькое условие, которое многие интеллигенты не ценят, а рабочие и крестьяне... улавливали, и у них сходится и слово и дело»[132].

\subsection*{\center * * *}

Поражало то, что стержнем всех ленинских выступлений против оппозиции было доказательство отрыва оппортунистов от жизни, от «текущего момента», как говорил Ленин, подмена главных вопросов не только второстепенными, а просто надуманными. Главным же вопросом, как еще и еще раз подчеркивал Ленин, он считал происходящее в связи с окончанием войны и переходом к мирной жизни изменение отношения рабочего класса к крестьянству. Он все клонил в эту сторону: вот здесь надо думать, искать, здесь можно и поспорить с пользой для дела.

Самое важное — без торопливости присмотреться и решить возникающие вопросы: что и как изменить, что привнести нового. Никто другой в ходе дискуссии не касался этого, не понимал так, как Ленин, потребностей новой ситуации, в которой оказалась революция в этот «переходный период в переходном периоде» (выражение Ленина). Именно это выдвигало целую кучу политических и экономических проблем, во сто крат более важных, чем надуманные, теоретически ошибочные, немарксистские тезисы и платформы по профсоюзным вопросам. Дискуссия поглотила всю подготовку к съезду. Бесконечные собрания, споры во всех звеньях почти парализовали деловую работу партийных, советских, профсоюзных органов, проникли и в армию, взбудоражили и потрясли всю партию и народ. Оппозиционеры сеяли в массах недоверие к ЦК, возлагали вину за хозяйственные трудности на ленинское руководство, пытаясь его опорочить.

Свою брошюру Лепим заканчивал уверенными выводами: «За один месяц и Питер, и Москва, и ряд провинциальных городов ужо показали, что партия откликнулась на дискуссию и подавляющим большинством отвергла ошибочную линию т. Троцкого. Если на «верхах» и в «периферии», в комитетах, в учреждениях, несомненно, были колебания, то масса рядовых членов партии, рабочая партийная масса в большинстве именно подавляющем высказывалась против этой ошибочной линии»[133].

Уральское бюро ЦК, партийное руководство Екатеринбурга качнулось вначале к троцкизму. Немалую роль сыграл здесь сторонник Троцкого секретарь ЦК Крестинский, имевший тесные связи с верхушкой этих организаций, где он работал до революции. Затем, под влиянием перечисленных выше событий, состоялись повторные дискуссии, на которых, несмотря на участие в них самого Троцкого, верх одержала платформа Ленина и его сторонников.

Сложным оказалось положение и в Баку, где в результате непонимания смысла дискуссии городская конференция в своем большинстве проголосовала за платформу Троцкого. Ошибка эта была вскоре исправлена на съезде Компартии Азербайджана, в феврале 1921 года, принявшем большинством голосов ленинскую линию. С докладом здесь выступил Серго Орджоникидзе, блестяще разъяснивший суть дискуссии и вскрывший антипартийный характер позиций Троцкого и Шляпникова. Когда на X съезде партии я стыдил тех бакинских делегатов, которые голосовали за тезисы Троцкого, они чувствовали себя очень плохо, тяжело переживали эту свою грубую ошибку и на съезде стали на позиции Ленина.

Крайне тяжелая ситуация сложилась в начале дискуссии и в нижегородской партийной организации, где я работал с октября 1920 года. В руководстве и активе основных районных партийных организаций, где была сосредоточена основная масса рабочих, — в Сормове и Канавине — господствовали сторонники Шляпникова. В другом рабочем районе города — Береговом (водники и портовики) в большинстве оказались сторонники Троцкого. И только в парторганизации небольшого Выксунского рабочего района губернии большинство было на стороне «платформы десяти». Этой платформы придерживались еще и в городском районе, парторганизации которого состояли преимущественно из служащих учреждений. Уездные организации представляли собой весьма пеструю картину.

В этих условиях в губкоме достигли договоренности о том, что выборы делегатов будут производиться, как и во всей партии, пропорционально полученному каждой платформой числу голосов.

Для нас, сторонников Ленина, было ясно, что если не склонить на нашу сторону парторганизацию Сормовского или Канавинского района, то на губернской конференции мы окажемся в меньшинстве.

Секретарем городского райкома работал Иконников, ныне проживающий в Москве. Он держался правильной партийной линии, в его районе дело обстояло благополучно. И я, выступая здесь, знал положение в районе. До революции Иконников вел партийную работу в Сормове и имел широкие связи с рабочими. Поэтому мы сговорились, что он на несколько дней переберется в Сормово, не обращая на себя внимания районного руководства, побывает на квартирах знакомых ему старых большевиков-рабочих и разъяснит им положение в партии, смысл дискуссии. Он должен был сказать им прямо, что Троцкий противопоставляет себя Ленину. Хотя троцкистов в Сормове нет, но здесь заправляют шляпниковцы, а это одного поля ягоды. Если победит платформа Троцкого или Шляпникова и потерпит поражение ленинская платформа, то будет нанесен непоправимый удар руководству партии, и Ленин окажется в немыслимо трудных условиях работы. Вот какой вопрос решается в этой по видимости профсоюзной дискуссии. Поэтому каждый рабочий-большевик должен побеседовать в этом духе с близкими товарищами по партии, чтобы они на дискуссионном собрании района голосовали за Ленина и его платформу, против Шляпникова и ею анархо-синдикалистской платформы. Как информировал Иконников, его беседы с рабочими давали весьма благоприятные результаты.

Настал день собрания в Сормове. Сторонники Шляпникова не сомневались в своей победе и пели себя весьма самоуверенно. Мы были полны надежд и сомнений, очень волновались, хотя и не подавали виду.

Выступая вторым докладчиком от сторонников «платформы десяти», я старался говорить попроще и подоходчивей. оголяя суть разногласий. Для большей убедительности я опирался на примеры из местной жизни Сормова, ставя своей задачей разбить критику партии оппортунистами и разоблачить их лживые демагогические обещания быстро и резко улучшить жизнь рабочих, если победит их платформа.

Припоминаются два момента из моего выступления, которые произвели определенное впечатление на рабочих и встретили одобрение большинства присутствовавших. Я говорил примерно так:

Сторонники платформы Шляпникова хотят здесь доказать, будто тяжелые условия, в которых работает и живет рабочий класс, есть результат неправильной линии и плохого руководства со стороны партии и Ленина. Но кому из вас неясно, что корни наших трудностей надо искать в последствиях длившихся семь лет разорительных империалистической и гражданской войн и вызванной ими хозяйственной разрухи? Под руководством Ленина война кончилась сокрушительной для врагов победой Советской власти. Разбиты и изгнаны из пределов России могущественные империалистические интервенты, которые были куда сильнее нас. Разбита вдребезги внутренняя контрреволюция помещиков и капиталистов. Какое же может быть сомнение в том, что в условиях завоеванного мира мы под испытанным уже ленинским руководством с таким же успехом преодолеем переживаемые трудности и построим социализм?

Напрасно оппортунисты по всей стране из кожи лезут вон, пытаясь убедить большевистскую партию, будто выход из трудностей возможен при победе платформ Троцкого или Шляпникова и при смене руководства партией и страной. Сормовские рабочие не поддадутся этому обману и станут стеной за Ленина, за его политическую платформу. Это встретило бурное одобрение.

И другое. Я коснулся демагогических обещаний, записанных в платформе «рабочей оппозиции». Зачитал пункт, в котором они утверждали, что они, как только победят, проведут трамвай к рабочим поселкам. Позвольте, говорил я, спросить сторонников этой платформы: Челышева — самого влиятельного члена губкома партии, подпись которого в платформе красуется вслед за подписью Шляпникова, и Чернова — директора сормовского промышленного гиганта, чья подпись также украшает эту платформу... Кто же мешал им до сих пор, стоя у власти, в губернии и на заводе, провести трамвай из города до Сормова? Ведь это было в пределах их прав и обязанностей. Может быть, Ленин им помешал? Тогда докажите это здесь на фактах и документах. Доказать этого вы не можете, ибо Ленин не меньше, а больше вашего заботится о нуждах народа. В ужасную войну, в условиях разрухи, голода и холода, где нам было найти средства и хотя бы рельсы для трамвая, когда металла не хватало даже для производства винтовок? Теперь, когда мы вступаем в полосу мирного развития, партия приступает к разработке мер по восстановлению хозяйства и новому строительству. Тогда станет возможным и провести трамвай и сделать многое другое для поднятия благополучия трудящихся. Это мы сделаем успешно под испытанным руководством Ленина, за которым без колебаний пойдут большевики сормовской организации, проголосовав единодушно за «платформу десяти». Это также встретило шумное одобрение собрания. Оппозиционные лидеры были в полной растерянности.

После нескольких выступлений рабочих было внесено и принято предложение о прекращении прений. Платформа «рабочей оппозиции», которая голосовалась первой, получила около четверти голосов. «Платформа десяти» — примерно три четверти. Другие платформы здесь не обсуждались и не голосовались.

Председательствующий предложил делегатов на губернскую конференцию избирать пропорционально количеству голосов, полученных каждой платформой, как это делалось и на других партийных собраниях. Вдруг из задних рядов поднимается внушительная фигура рабочего, который говорит твердо и четко: мы большинством приняли платформу Ленина. На губернскую конференцию от Сормова надо выбирать только тех, которые поддерживают волю большинства. Нам не гоже, чтобы на конференции одни сормовичи говорили за платформу Ленина, а другие за платформу Шляпникова. Голосованием это предложение было принято. Избраны были только делегаты — сторонники «платформы десяти». Получилось так, что ни один из членов Сормовского райкома не был избран делегатом, что ошарашило и без того растерянных лидеров оппозиции.

Дискуссионное собрание в Канавинском районе прошло до этого, и я также выступал вторым докладчиком. Но там у нас не было хороших связей с рабочими, мы не сумели провести подготовительной работы в низах и оказались в заметном меньшинстве, получили лишь несколько делегатов — сторонников «платформы десяти». Тем не менее в итоге на губернской конференции значительным большинством победила ленинская линия.

И так почти повсюду. К съезду из семи антиленинских платформ выжили лишь две: объединенная платформа Троцкого — Бухарина и анархо-синдикалистская платформа Шляпникона. Но и они уже были лишены былой силы, доживали последние дни, их сторонники оказались неспособными дать сколь-либо серьезный бой на съезде. Число их сторонников на съезде составляло около 15 %. Партия пошла за Лениным и его соратниками. План Ленина «вылечить до съезда» партию оказался осуществленным.

Это обеспечило возможность повернуть съезд и партию в целом к главному вопросу — вопросу перехода к новой экономической политике. Это было главным и основным, и именно этого добивался Ленин. Съезд подвел итоги внутрипартийной борьбы, дал исчерпывающие оценки антипартийным фракциям и поставил на них крест, принял политические и организационные меры, призванные предотвратить возможности подобных кризисов, воспрепятствовать появлению фракций, укрепить единство партии на основе ленинских норм и дальнейшего совершенствования принципа демократического централизма. Но об этом в следующей главе.

\newpage
\section[Ленин на X съезде партии: наладить отношения рабочих и крестьян. Восстановление единства партии]{\center \textbf{Ленин на X съезде партии: наладить отношения рабочих и крестьян. Восстановление единства партии}}


X съезд партии проходил с 8 по 16 марта 1921 года. К этому времени уже была одержана полная победа на фронтах гражданской войны. Ленин в своей речи при открытии съезда сказал: «...мы в первый раз собираемся на съезд при таких условиях, когда вражеских войск, поддерживаемых капиталистами и империалистами всего мира, на территории Советской республики нет... Три с половиной года неслыханно тяжелой борьбы, но отсутствие вражеских армий на нашей территории, — это мы завоевали!»[134]

Вопросы, которые разделяли партию в предсъездовской дискуссии, раскололи и Центральный Комитет. Тогдашние секретари ЦК (Крестинский, Серебряков, Преображенский) были на стороне Троцкого — Бухарина. Заполучив в свои руки аппарат ЦК, троцкисты начали проводить свою кадровую политику, беспрепятственно наводняли местные организации, в том числе и нижегородскую, где я в то время работал, троцкистскими брошюрами и другими фракционными материалами.

Но к моменту созыва съезда вся партия уже как бы самоопределилась. В подавляющем большинстве губернские партийные конференции (за исключением, кажется, Самарской) отдали явное предпочтение ленинской линии. Исход партийных конференций предопределил, что на съезде подавляющее большинство составили сторонники «платформы десяти», т. е. ленинской платформы. Не оставалось сомнений, что съезд встанет на ленинскую позицию, против троцкистов и «рабочей оппозиции». Съезду предстояло закрепить положение, которое сложилось в партии, т. с. поддержать ленинский курс и отклонить вес другие платформы, а главное — принять меры к восстановлению единства партии, поставить неодолимые препятствия тем, кто снова захотел бы подорвать это единство.

В краткой речи при открытии съезда Ленин охарактеризовал главные трудности момента и узловые задачи перехода от военного коммунизма к мирному хозяйственному строительству. Эти задачи касаются основ самих отношений между классами, которые в нашем обществе остались, и этот вопрос должен быть одним из главных, которые предстоит разобрать и разрешить. Другая главная задача съезда — чтобы в партии «и не только формально работа была более сплоченной, более дружной, чем прежде, чтобы не было ни малейших следов фракционности...»[135]. Эти основные мысли Ленин неоднократно, в разных аспектах, повторит и разовьет в своих многочисленных выступлениях на съезде.

В докладе о политической деятельности Центрального Комитета Ленин, как это для пего было обычно, выделил прежде всего важнейшие события истекшего года и подал их так, чтоб предоставить больше материала для размышлений — о причинах именно такого хода революции, о значении выявившихся ошибок и их уроках для будущего, о накопленном опыте, о выводах, которые подлежит сделать из пережитого, и путях решения вставших перед партией задач.

Многие трудности и ошибки он отнес за счет длительной гражданской войны и ее тяжелых последствий, крайне затруднивших сам переход к мирной хозяйственной деятельности, в результате чего оказался невозможным быстрый и сколько-нибудь планомерный переход от войны к миру. Недостаточно оказались учтенными, как утверждал, подкрепляя это фактами, Ленин, не только технические, но и экономические, и социальные трудности, связанные с демобилизацией армии. Но мы, подчеркивал, с другой стороны, Ленин, не имели достаточных данных и опыта для этого, «и только по окончании войны мы увидели всю ту степень разорения и нищеты, которые надолго осуждают нас на простое только излечение ран»[136], А в результате, когда сотни тысяч демобилизованных не могут найти приложения своему труду, обнищавшие, разоренные, привыкшие воевать и чуть ли не смотрящие на войну, как на ремесло, мы оказываемся втянутыми в новый вид гражданской войны, которая находит свое выражение в бандитизме[137].

Подобного же порядка ошибка была допущена и в области продовольственной. Не сумев оценить всей опасности надвигавшегося к весне продовольственного кризиса, мы поддались естественному стремлению увеличить выдачу продовольствия голодающим рабочим. Не имея опорного пункта для правильного учета, не могли сразу установить должной меры. Мы должны были провести, самокритично заявлял Ленин, умеренное увеличение выдачи пайков, чтобы накопить необходимый запас на черный день, который наступил к весне. Аналогичная ошибка была допущена и в области топлива.

Вывод: «Нужна гораздо более длительная подготовка, более медленный темп, — вот урок, который мы за этот год пережили...»[138].

Не ставя задачей пересказывать часть доклада, посвященного хозяйственным и социальным трудностям того времени и содержащего много интересных подробностей, мне хотелось бы лишь подчеркнуть, что, на мой взгляд, он имеет далеко не только исторический интерес. Начальные шаги в области мирного строительства, откровенное раскрытие трудностей и ошибок, способы их преодоления — весь этот опыт разоренной и отсталой тогда страны, приступавшей к широкому строительству новой жизни, может быть полезен и даже очень полезен, например, странам, только освободившимся от колониального гнета и подчас в чрезвычайных обстоятельствах прокладывающим путь к экономической независимости и социальному прогрессу. Надо хорошенько знать этот опыт и многочисленным нашим товарищам, имеющим контакты с развивающимися странами.

Касаясь дискуссии о профсоюзах, Ленин заявил, что «эта роскошь была действительно совершенно непозволительной и что, допустив такую дискуссию, мы, несомненно, сделали ошибку... выпятили на первое место вопрос, который по объективным условиям не может стоять на первом месте...»[139], и что тем самым отвлекли внимание партии от насущного и грозного вопроса — наступившего кризиса в стране.

Здесь оправдалась пословица, продолжал докладчик, «нет худа без добра». К сожалению, добавил он под смех всего зала, «худа оказалось немного много, а добра немного мало».

Ленин заметил, что за одни платформы голосовали по преимуществу «верхи партии». Прямо он их не назвал, но было ясно, что при этом имеются в виду платформы троцкистов и «демократического централизма». Платформа же, крайне неудачно называемая «рабочей оппозицией», представляет собой явно синдикалистский уклон.

Важнейший вывод — партия проявила себя настолько зрелой, что, видя некоторое шатание «верхов», мобилизовалась и громаднейшим большинством наиболее значительных партийных организаций быстро пришла к правильному решению. Судя по всему, именно это Ленин и считал той долей «добра», которую в итоге все же дала дискуссия.

Подробно остановился Ленин на кронштадтском мятеже, который был поднят под лозунгом «Советы без коммунистов», «За свободную торговлю» и т. п. Мятеж свидетельствовал не только о мелкобуржуазном, эсеро-меньшевистском влиянии. С одной стороны, он сливался с белогвардейщиной, а с другой, отражал недовольство крестьянства, а также части беспартийных рабочих. Учитывая шатания среди части членов партии, Ленин указал, что существует угроза мелкобуржуазной контрреволюции, которая, «несомненно, более опасна, чем Деникин, Юденич и Колчак вместе взятые, потому что мы имеем дело со страной, где пролетариат составляет меньшинство, мы имеем дело со страной, в которой разорение обнаружилось на крестьянской собственности»[140].

Ленин видел выход только в том, чтобы суметь перейти к мерам экономического обеспечения крестьянства, в том, чтобы обеспечить ему возможность свободы хозяйничать.

В числе этих мер, прежде всего, переход от продразверстки к натуральному налогу. Но, пояснял Ленин, мы мыслим его теперь не так, как при издании закона о продналоге 30 октября 1918 года (который в силу гражданской войны не удалось осуществить), а иначе: «...мы смотрим на него не только с точки зрения обеспечения государства, а также с точки зрения обеспечения мелких земледельческих хозяйств»[141].

Говоря о борьбе с бюрократизмом, Ленин всячески подчеркивал значение этой «язвы» и упирал на то, что борьба с бюрократизмом абсолютно необходима и так же сложна, как борьба с мелкобуржуазной стихией. Однако, с другой стороны, Ленин фактически предупреждал, что руководство этой борьбой должно находиться в руках партии. Нельзя позволить перехватить и извратить лозунг борьбы с бюрократией ни «рабочей оппозиции», ни другим силам. «Когда в Москве были беспартийные собрания, — говорил он, — ясно было, что из демократии, свободы они делают лозунг, ведущий к свержению Советской власти»[142].

В другом месте Ленин пояснял, что нельзя, как это делают представители «рабочей оппозиции», выдвигать лозунг немедленно «покончить с бюрократизмом», ибо это задача длительная, сложная и трудная, а подобный лозунг, и особенно фактическая невозможность его осуществления, могут быть на руку лишь антипартийным силам.

С такой искренностью и прямотой были поставлены Лениным все узловые вопросы переживаемого момента, так ясно, без прикрас, была обрисована тяжелая обстановка в стране, такие важные выводы были сделаны для дальнейшей работы партии, что подавляющее большинство делегатов было вполне удовлетворено политическим отчетом ЦК.

Поэтому начавшиеся было прения быстро прекратились. Выступили фактически только представители «демократического централизма», да большей частью сторонники «рабочей оппозиции». Прения оказались на редкость бесцветными, выступавшие говорили, как выразился потом напрямик Ленин, «удивительно бессодержательно»[143].

После доклада Ленина был еще доклад секретаря ЦК, троцкиста Крестинского, об организационной деятельности ЦК, ничем не примечательный. Один из ораторов в прениях довольно метко заметил, что если Ленин говорил о ряде ошибок в работе ЦК, то Крестинский, докладывая об организационной деятельности ЦК, даже не упомянул ни об одной из них. Но и сравнение-то это неподобающее: великий революционер, с одной стороны, и партийный администратор бюрократического толка — с другой. Нашел кого с кем сравнивать!

В заключительном слове Ленин подверг резкой критике как выступление Коллонтай на съезде, так и особенно выпущенную ею к съезду брошюру с изложением линии «рабочей оппозиции», а также тезисы Шляпникова. Он еще раз дал отпор поползновениям передать управление экономикой «съезду производителей» или его уполномоченным. Подчеркивая политическую опасность этого злополучного «кредо», Ленин указал, что «теперь, при свете кронштадтских событий, этот тезис тем более странно слышать»[144]. Он едко высмеял утверждение в брошюре Коллонтай о том, будто «рабочая оппозиция» — это классово спаянная, классово сознательная и классово выдержанная часть нашего промышленного пролетариата.

Под хохот всего зала Ленин иронически заметил: «Ну, слава богу, так и будем знать, что тов. Коллонтай и т. Шляпников — «классово спаянные, классово сознательные»»[145].

Очень возмущен был Ленин демагогическим выступлением Шляпникова, который, в частности, привел пример, что где-то гноят картошку, и с пафосом спрашивал съезд, почему наркома продовольствия Цюрупу не предают суду.

Ленин парировал: «А я задаю вопрос: почему не предают суду Шляпникова за такие выступления?»[146] Суть выcказывания Ленина по этому поводу сводилась к тому, что если Шляпников всерьез думает об отдаче под суд Цюрупы, то почему он вовремя не обратился с этим к партийной контрольной комиссии, а пришел на ответственнейшее собрание и бросает здесь слова, в которые сам не верит, которые он доказать не может, это — демагогия, на которой базируются анархо-махновские и кронштадтские элементы.

Но достойно пристального внимания то, что острейшую критику в адрес Шляпникова и Коллонтай Ленин завершил заявлением о готовности ЦК прислушаться к критике бюрократизма со стороны оппозиции. Ленин призвал к деловой работе тех, кто всю энергию расходовал на упреки в адрес ЦК, на внутрипартийную борьбу, и заявлял о готовности ЦК пойти на определенные уступки и сработаться с теми сторонниками «рабочей оппозиции», которые связаны с рабочими массами и хотят на деле помочь ЦК в борьбе с бюрократизмом, в развитии внутрипартийной и советской демократии, отказываются от фракционной деятельности.

Заключительное слово Ленин произносил взволнованно, не скрывая желания раз и навсегда положить конец дезорганизаторской фракционной деятельности оппозиции, льющей в конечном счете воду на мельницу мелкобуржуазной контрреволюции. «Не надо теперь оппозиции, товарищи, — восклицал он, — не то время! Либо — тут, либо — там, с винтовкой, а не с оппозицией. Это вытекает из объективного положения, не пеняйте»[147]. Эта его речь произвела сильное отрезвляющее впечатление даже на многих сторонников оппозиции.

\subsection*{\center * * *}

После первых двух дней работы съезда меня неожиданно вызвали в Кремль на какое-то совещание к Ленину. Это было после вечернего заседания съезда. 9 марта совещание состоялось — не в кабинете Ленина, а в каком-то другом помещении, на первом этаже. Собралось человек пятнадцать или несколько больше. Расселись в небольшой комнате, Ленин сидел за письменным столом. Представителей Петрограда здесь, как и до этого, на съезде не было в связи с тем, что все они были заняты подавлением кронштадтского мятежа.

Ленин начал с объяснения цели совещания. Сказал, что исход партийных конференций на местах и состав делегатов доказывают, что «платформа десяти» будет принята съездом. Это уже не проблема. Однако этого крайне недостаточно. Есть опасность, что, приняв правильное решение и отвергнув фракционную деятельность, съезд может избрать ЦК, примерно, в том же составе, что и на предыдущем съезде. В результате в ЦК попадет много сторонников Троцкого и может случиться так, что они лишат Центральный Комитет возможности деловой работы, будут затевать споры по любому вопросу, углублять разногласия, и вскоре ЦК снова окажется в положении фактического раскола. Нельзя повторять то, что произошло после IX съезда, нетерпимо то, что имело место в последние месяцы. Надо партию гарантировать от этого. Необходимо добиться, чтобы состав ЦК был трудоспособным, сплоченным, был обеспечен большинством сторонников платформы, которую примет съезд, и чтобы приверженцы других платформ, хотя и были бы представлены в ЦК, но составляли бы незначительное меньшинство.

Главная гарантия заключается в том, чтобы в ЦК подавляющее большинство — две трети голосов — принадлежало твердым сторонникам «платформы десяти». Остальные места дать некоторым представителям платформы Троцкого, «рабочей оппозиции», а также «демократического централизма», поскольку в числе сторонников оппозиции много видных партийных деятелей, полезных в работе и популярных в стране. Нужно избрать в состав ЦК действительно пролетарские элементы.

Количество членов и кандидатов в члены ЦК целесообразно в этой связи несколько увеличить, скажем, членов ЦК — на шесть человек, кандидатов — на три.

Это предложение встретило общее одобрение присутствующих. Затем Ленин перечислил именно тех членов Центрального Комитета, которых не следует переизбрать в новый состав ЦК. Из числа сторонников троцкистской платформы он назвал Крестинского, Преображенского, Серебрякова, Смирнова И. Н., из «рабочей оппозиции» — Коллонтай. Ленин высказался за избрание в ЦК из сторонников «рабочей оппозиции» Шляпникова и Кутузова членами ЦК и Киселева, как помнится, также Медведева и Лутовинова — кандидатами.

Ленин высказался за включение в ЦК от фракции «демократического централизма» кандидатуры Сапронова. Все названные им кандидатуры вызвали оживленные споры. Ленин воспользовался этим для разъяснений. Он сказал примерно так: нам нужно не просто отвергнуть оппозиционные платформы, нам необходимо обеспечить единство партии. У оппозиции есть верхушка бюрократическая, но есть и значительная масса рабочих, искренних людей, которые в силу трудностей, переживаемых нашей партией и страной, поддаются колебаниям. И нужно создать условия, которые не отталкивали бы их от нашей партии. Это будет способствовать тому, чтобы они стали на правильные позиции. Для этого необходимо вести работу по сплочению партии. Предлагалось ввести в состав ЦК Шляпникова, Кутузова и Сапронова, зная об их недостатках, но учитывая тот факт, что они отражают хотя и отсталые, но действительно имеющие место настроения в партии и в народе. Когда в Центральном Комитете будет наше твердое большинство, то присутствие нескольких оппозиционеров не создаст угрозы единству ЦК. С другой же стороны, наличие их в составе ЦК даст возможность, опираясь на силу чекистской дисциплины, требовать и от них выполнения решений ЦК, подчинения этим решениям. К тому же через них можно будет знать о настроениях и мнениях, которые имеются в определенных слоях населения, рабочего класса, учитывать это в работе Центрального Комитета, взять здоровое зерно и не дать неподходящим настроениям разрастаться дальше, Таким образом, их избрание не только не ослабит борьбу с оппозицией, но облегчит борьбу с попытками сохранить фракционность.

На съезде, — говорил Ленин, необходимо принять постановление, которое осуждало бы фракционность, распустило все существующие фракции, исключало условия для создания и восстановления фракций, установило несовместимость фракционной деятельности с пребыванием в партии. Более того, нужно предусмотреть возможность исключительных мер, чтобы даже член Центрального Комитета, если он стал на путь фракционности, мог бы быть исключен из ЦК и даже из партии. Конечно, к этому следует прибегать в самом крайнем случае с соблюдением такого порядка: если две трети всех членов и кандидатов в члены ЦК и ЦКК, собранных вместе, будет голосовать за исключение, только тогда член ЦК может быть исключен из ЦК и из партии за фракционную борьбу. Такое решение будет дополнительной уздой. При наличии такого решения тем более не остается никакой опасности в том, что Шляпников и некоторые другие сторонники оппозиции будут набраны в состав ЦК в качестве меньшинства.

Ленин детально охарактеризовал сущность «рабочей оппозиции, что отражено в сохранившемся конспекте его выступления[148]. Но об этом подробнее я скажу, вспоминая доклад Ленина по этому вопросу на самом съезде.

В отношении анархо-синдикалистского уклона, говорил Ленин, нам придется принять специальное решение, которое не просто бы отвергало этот уклон, а, учитывая трудное положение в стране, разъясняло бы партии его опасность и вредность, антимарксистский, анархический характер.

Все это было настолько убедительным, что никем никаких сомнений не высказывалось.

Но как добиться, чтобы съезд избрал подходящий состав ЦК, продолжал Ленин. Без должной подготовки и организации этого не достичь. На съезде будет трудно аргументировать, трудно обсудить, обеспечить поддержку всех подходящих и отклонение нежелательных кандидатур, если не провести предварительную организационную работу. Поэтому сегодня необходимо договориться о том, чтобы каждому поработать в своих делегациях, разъяснить необходимость поддержки выдвигаемых нами кандидатур и отказать в поддержке нежелательных кандидатур. А главное — надо созвать совещание всех тех делегатов съезда, которые избраны по «платформе десяти». На таком совещании обсудить вопрос о составе ЦК, отвести ряд кандидатур сторонников Троцкого, а также определить, кого из членов оппозиции персонально следует избрать.

Ленин говорил, что следует написать извещение о созыве такого совещания. Извещение должно служить и приглашением и пропуском на совещание. Оно должно быть коротким, напечатанным типографским способом. Но я не считаю, сказал Ленин, возможным для этого дела использовать государственную типографию. На всякий случай я привел одного старого товарища, коммуниста-подпольщика, у которого есть шрифт и есть ручной станок, на котором он обеспечит за ночь печатание извещения в достаточном количестве.

Сталин подал реплику, что созыв совещания делегатов, избранных по «платформе десяти», будет использован троцкистами и другими оппозиционерами для обвинения нас во фракционности и наши действия могут быть неправильно истолкованы съездом.

Ленин, добродушно улыбаясь, шутливо сказал Сталину: что я слышу от старого заядлого фракционера?! Даже он сомневается в необходимости созыва совещания делегатов, стоящих на «платформе десяти»! Вы должны знать, что Троцкий давно собирает сторонников своей платформы, да и сейчас, пока мы с вами разговариваем здесь, наверное, собрал свою фракцию. То же самое делают и Шляпников и Сапронов. Зачем закрывать глаза на хотя и неприятный, но явный факт существования фракций в партии? Именно созыв такого совещания сторонников «платформы десяти» обеспечит условия, которые исключили бы всякую фракционность в нашей партии в дальнейшем.

Сталин не стал спорить, и никто другой не оспаривал справедливости ленинского предложения. Все согласились, что надо созвать такое совещание.

Ленин попросил карандаш, чтобы написать текст извещения. Томский предложил свой карандаш. Ленин стал писать, и карандаш поломался. Тогда Ленин с юмором сказал: вот что значит профсоюзы — не имеют даже приличного карандаша. Правильно Троцкий предлагает «перетряхнуть» профсоюзы! Все, конечно, засмеялись. Тогда Каменев дал Ленину свой карандаш, и он стал писать. Начал так: «Делегаты съезда, избранные по платформе «Десяти»...»

Один из товарищей заметил, что неправильно так писать, потому что некоторые губернские конференции выбирали делегатов не по платформам, а персонально, и среди них большинство сторонников «десяти», а не сторонников оппозиции. Если так написать, тогда они не будут приглашены на совещание.

Ленин с удивлением заметил: что за странные большевики и большевистские конференции, которые при наличии предсъездовской дискуссии с фракционными платформами и решения ЦК о выборах по платформам выбирают на съезд делегатов без учета их отношения к платформам?

Но факт остается фактом. Тогда Ленин предложил внести изменение. Извещение адресовать не избранным по «платформе десяти», а сторонникам платформы. Он добавил, что надо тщательно проверить, чтобы среди приглашенных не оказались троцкисты и другие фракционеры. Поэтому у дверей нужно поставить товарищей Рудзутака и В. Шмидта, которые хорошо знают местные кадры, чтобы не пропустить кого не следует.

Участникам данного совещания, говорил далее Ленин, надо знать, что происходит во фракциях съезда, изучать делегатов, их мнение и пр. для того, чтобы учесть все это на нашем совещании...

На этом совещании в числе других был и Артем (Сергеев)[149], с которым ранее я не был знаком. Выходили с совещания мы с ним вместе и разговорились. Ленина я давно знаю, говорил он, и вот какая странность получается. Подчас я не согласен с ним по некоторым вопросам. Вот по вопросу сращивания профсоюзов с государством у меня также были колебания. Но я на опыте убедился, что всегда, когда я расходился с Лениным, он в конце концов оказывался прав, а я нет. Поэтому при обсуждении того или иного вопроса я высказываю то, что думаю, но когда дело доходит до голосования, то голосую за предложения Ленина, зная, что и на этот раз он будет прав.

...На следующий день нам передали известное количество приглашений на совещание сторонников «платформы десяти», и мы должны были раздать их тем товарищам, которых мы знали как избранных по ней или как ее сторонников.

Через день состоялось совещание сторонников «платформы десяти». У двери стояли и пропускали приглашенных Рудзутак и В. Шмидт. Совещание началось выступлением Ленина о том, что съезд в своем подавляющем большинстве стоит на правильной принципиальной платформе, в партии налицо серьезные успехи и необходимо повести дело так, чтобы и после съезда обеспечить проведение выработанной им линии. Главное — не допустить, чтобы повторилось положение, когда ЦК оказался расколот и была навязана дискуссия, которая парализовала партийные организации, внесла в партию угрозу раскола. Поэтому надо избрать в Центральный Комитет таких людей, которые по совести и с чувством ответственности проводили бы линию, определенную съездом.

Перейдя к вопросам персональным, он начал с трех секретарей Центрального Комитета и сказал, что их кандидатуры надо отвести из состава ЦК, потому что они являются носителями бюрократического стиля работы (ни один из них впоследствии при голосовании на съезде в состав ЦК не прошел). Ленин отвел кандидатуры еще нескольких троцкистов, избранных в ЦК предыдущим съездом. При обсуждении кандидатур в состав ЦК раздалось много голосов против Шляпникова, Кутузова, Сапронова, намечаемых членами ЦК, а также против некоторых сторонников оппозиции, выдвигавшихся кандидатами в члены ЦК. Ленин высказался за кандидатуры Шляпникова, Кутузова, Сапронова, Киселева, привел примерно те же аргументы, что и на предшествовавшем узком совещании. Однако совещание не поддержало кандидатуру Сапронова (он и не был избран съездом). Следует отметить, что в список, предложенный совещанием сторонников «платформы десяти», была включена кандидатура Осинского, одного из лидеров фракции «демократического централизма».

Надо еще пожалуй, заметить, что в ходе предварительного согласования кандидатур при переговорах с представителями фракций и с учетом изменений, которые произошли по отдельным кандидатурам, в представленный сторонниками «платформы десяти» список кандидатов в члены Центральной Контрольной Комиссии был включен лидер нижегородской «рабочей оппозиции» Челышев, хотя первоначально Ленин считал, что в ЦКК ни одного оппозиционера вводить не следует. После успешного совещания сторонников «платформы десяти», в котором участвовало внушительное большинство делегатов съезда, Владимир Ильич, казалось бы, имел все основания быть вполне спокойным за дальнейший ход съезда. Не оставалось сомнений, что принятие съездом всех намеченных Лениным и встретивших уже одобрение большинства делегатов решений полностью обеспечено. Это было большое, можно даже сказать, великое для партии дело. Но Ленин на этом не успокоился. Терпеливо взялся он за кропотливую работу по убеждению тех групп делегатов, которые были не согласны с большинством. На частных совещаниях отдельных групп он убеждал, откровенно разъяснял спорные вопросы, выслушивал возражения и отвечал на них, добивался, чтобы оппозиционеры и уклонисты отошли от неправильных позиций, поняли существо своих ошибок и присоединились к точке зрения большинства делегатов.

Ленин хотел до конца использовать партийный, товарищеский метод убеждения, великим мастером которого он был, использовать его вовсю еще до предстоявшего голосования таких острых вопросов, как особое постановление о единстве партии, включавшее крайнюю меру, впервые предлагавшуюся к применению в партийной практике, осуждение анархо-синдикалистского уклона, выборы руководящих органов. Он добивался подведения прочной базы убеждения даже при заведомо ясном исходе голосования. При этом интересовали Ленина не заядлые фракционные лидеры, а прежде всего давние коммунисты. особенно из рабочих, делегаты как с решающим, так и с совещательным голосом, которые были связаны с массами и отражали их настроения.

Так Ленин организовал частное совещание сторонников «рабочей оппозиции», на котором он несколько раз выступал, выслушивал участников совещания, приводил новые аргументы, вновь и вновь повторял старые, всяческими способами доказывал правоту своих позиций. Настойчиво добивался, чтобы до сознания уклонистов дошла гибельность фракционности для рабочей партии, правящей государством, жизненная необходимость единства ее рядов, без чего невозможно завершить дело социалистической революции. Ленин выражал готовность включить в выборные органы партии лучших, преданных партии людей из их группы, готовность использовать все здоровое, что было в их критике, для дальнейшего развития демократии, для борьбы с бюрократизмом, в целях улучшения жизни рабочих. Он снова и снова призывал к дружной работе всю партию после X съезда, без чего нельзя было бы справиться со сложнейшей обстановкой в стране.

Такого же порядка совещание Ленин провел и со сторонниками фракции «демократического централизма». Но и этим он не ограничился.

Было созвано многолюдное совещание сторонников «платформы десяти» и других платформ из числа рабочих— старых большевиков, как делегатов съезда, так и гостей. На нем Ленин призывал до конца понять, из какого серьезного положения партия должна окончательно выбраться. Предлагал распустить на съезде все фракции, запретив впредь создание таковых. Избрать состав руководящих органов с включением известного числа инакомыслящих, но чтобы было обеспечено в них твердое большинство убежденных сторонников принимаемых съездом решений. Только это может исключить возможность повторения такого рода дискуссии, какая была навязана партии перед съездом, когда ЦК фактически раскололся и поставил партию на грань раскола.

Ленин особенно подробно разъяснял теоретическую несостоятельность и политическую опасность анархо-синдикалистского уклона, дальнейшая пропаганда которого должна быть определена как несовместимая с принадлежностью к партии коммунистов.

На всех этих частных совещаниях, или, как тогда говорили, «хождениях Ленина по трем-четырем комнатам», он, заостряя, выпячивая суть спорных вопросов и последствия этих споров, не допускал резкостей, личных выпадов, выражений, которые могли кого-либо обидеть или оттолкнуть. Говорил в духе товарищества, как с единомышленниками в главном и основном, проявлял терпимость и уважение к собеседникам. Трудно даже себе представить, сколько сил и энергии потратил Ленин, сколько умения и терпения он проявил в этой утомительной, кропотливой, но весьма полезной работе по убеждению заблуждавшихся товарищей. Многих из них действительно удалось убедить, у других — посеять зерна сомнения в их правоте, устранить ожесточенность, перераставшую подчас уже чуть ли не во враждебность. Все это решающим образом содействовало созданию товарищеской атмосферы в дальнейшей работе съезда и во всей партии.

Затем Ленин соединил эти разбредшиеся по комнатам группы в большое частное совещание в кремлевском зале Свердлова, на котором, как говорил Ленин в заключительном слове по докладу о единстве партии и анархо-синдикалистском уклоне, была группа «рабочей оппозиции» в полном составе и ряд виднейших товарищей, представителей всех оттенков[150] и были достигнуты «соглашения групп между собою»[151], в том числе и по списку кандидатов в руководящие органы (в частности, и по вопросу о том, чтобы представителя «демократического централизма» Сапронова избирать по списку кандидатов в члены ЦК, а не членов ЦК).

Перед выборами представители оппозиции робко пытались уговорить делегатов выбирать в ЦК лишь по личным качествам, не учитывая приверженности к той или иной платформе. Однако такое предложение не встретило поддержки.

14 марта вечером состоялось закрытое заседание съезда, посвященное выборам. Во многом опять повторилась картина, известная нам по описанным выше частным совещаниям. Ленин снова выступал по поводу кандидатур Кутузова, Шляпникова и Каменева.

Из членов ЦК, избиравшихся на VIII и IX съездах, значительным влиянием пользовался Иван Никитич Смирнов. У губкомовских партийных кадров о нем было мнение как о способном и эрудированном партийном руководителе, как об обаятельном человеке. Возглавлял он тогда Сибирское бюро ЦК. К удивлению и огорчению многих, в дискуссии он проявил себя как ярый сторонник Троцкого и повел за собой многих руководящих работников Сибири. Было опасение, что при тайном голосовании он получит столько голосов, что будет избран, и тем самым в ЦК окажется еще один лишний троцкист. Поэтому Ленин счел нужным лично выступить с отводом его кандидатуры в члены ЦК. (Даже после отвода Лениным за Смирнова голосовало более одной четверти делегатов, он был избран кандидатом в члены ЦК, а на следующем, XI съезде — снова в члены ЦК.)

Несколько военных делегатов с Северного Кавказа неожиданно стали с мест выкрикивать возражения против кандидатуры Серго. Сидели эти делегаты в последних рядах и шумели на весь зал. Один из них поднялся на трибуну и стал говорить, что, мол, Орджоникидзе кричит на всех, командует, не считается с местными работниками, а потому не может быть в составе ЦК. Это демагогическое выступление оказало влияние на настроение делегатов съезда, многие из которых к тому же не знали Орджоникидзе. Выступавший был политработник Врачев из Грозного. Он хорошо запомнился, потому что щека его была перевязана марлей, видимо, по причине зубной боли.

Позже, на XIII партийной конференции, Врачев признал, что выступал он с отводом Серго на X съезде партии не по своей собственной инициативе, а под большим давлением фракции Троцкого, к которой он, Врачев, в те годы примыкал. Он рассказал, что во время работы X съезда участвовал в тайном совещании троцкистов — делегатов съезда, где и было принято решение об отводе из состава ЦК ряда товарищей — сторонников Ленина, в том числе и Орджоникидзе. (Впоследствии Врачев отошел от троцкизма, признал свои ошибки и занял правильную партийную позицию.)

В защиту Орджоникидзе выступил Сталин. Говорил он в спокойном тоне, тихим голосом. Привел биографические данные, рассказал о работе Серго в подполье, на фронтах гражданской войны и рекомендовал избрать его в ЦК. Было видно, однако, что ему не удалось убедить делегатов. Они продолжали шуметь.

Тогда выступил Ленин. Он сделал примерно следующее заявление: я знаю товарища Серго давно, еще со времен подполья, как преданного, активного, бесстрашного революционера. Хорошо показал он себя в эмиграции. Сыграл выдающуюся роль в подготовке Пражской конференции партии в 1912 году, был тогда избран членом ЦК. Вел активную работу в Петрограде в период подготовки и проведения Октябрьской революции.

В гражданской войне он показал себя храбрым, способным организатором. Но в критике выступавших товарищей есть одно правильное замечание по адресу товарища Серго. Это то, что он кричит на всех. Это верно. Громко говорит, но вы, наверное, не знаете, в чем дело? Он и со мной, когда разговаривает, так же кричит. Потому что он глуховат на левое ухо, как-то особенно мягко и тепло сказал Ленин. Потому и кричит, думает, что его не слышат. Но не следует этот недостаток принимать всерьез при выборах в ЦК...

Это вызвало улыбки и даже добродушный смех у делегатов. Стало ясно, что Ленин отлично поддержал кандидатуру Орджоникидзе, отбив атаки на него оппозиции (хотя в скобках надо заметить, Ленин никогда не спускал Серго отдельные срывы, у него случавшиеся). А ведь были серьезные опасения, что окажется много голосов против. После выступления Ленина, при тайном голосовании, Орджоникидзе получил подавляющее большинство голосов. Как и Дзержинский, он получил 438 голосов из 479 (один только Ленин был избран единогласно) — больше, чем ряд других тогдашних видных членов ЦК.

И сейчас еще свежо в памяти, как поразила меня тогда наблюдательность Владимира Ильича. Я тоже замечал, что Серго слышит хуже других, но не знал, что он глуховат именно на левое ухо.

\subsection*{\center * * *}

Свой доклад на съезде «О замене разверстки натуральным налогом» — так скромно именовался тогда этот коренной вопрос — вопрос о нэпе — Ленин начал с утверждения, что он является «прежде всего и больше всего вопросом политическим, ибо суть этого вопроса состоит в отношении рабочего класса к крестьянству»[152]. Отношения между этими двумя классами — борьба их или соглашение между ними — определяют судьбы всей пашей революции.

Проанализировав состояние крестьянства, Ленин констатировал, что одним из результатов военного коммунизма было то, что «деревня нивелировалась, выравнилась, т. е. резкое выделение в сторону кулака и в сторону беспосевщика сгладилось. Все стало ровнее, крестьянство стало в общем в положение середняка»[153].

Однако военный коммунизм и продовольственная разверстка нарушили отношения между рабочим классом и крестьянством, дальше такое положение терпеть нельзя. «Крестьянство формой отношений, которая у нас с ним установилась, недовольно... оно этой формы отношений не хочет и дальше так существовать не будет. Это бесспорно. Эта воля его выразилась определенно»[154].

С этим необходимо считаться, и, как трезвые политики, мы должны нашу политику по отношению к крестьянству пересмотреть.

Крестьянство не прочь восстановить частный капитал и свободную торговлю целиком. Однако в конечном смете это не в его интересах, так как не может привести ни к чему другому, кроме как к реставрации власти помещиков и капиталистов. Вот это-то и надо ему доказать. Рабочий класс может пойти лишь на частичную свободу торговли. Но здесь, именно в этом, возможность соглашения между рабочим классом и крестьянством, несмотря на глубокую рознь экономических интересов пролетариата и мелкого земледельца.

Зная экономические требования крестьянства, говорил Ленин, есть основания утверждать, что мелкого земледельца можно удовлетворить двумя вещами: «Во-первых, нужна известная свобода оборота, свобода для частного мелкого хозяина, а во-вторых, нужно достать товары и продукты»[155].

Как это было свойственно Ленину, обозревая экономическую политику военного времени, он сделал упор на уроках, вытекающих из пройденного этапа. Общий его вывод: в условиях войны, в которые мы были поставлены, основы этой политики были правильны, хотя при ее применении были допущены ошибки и преувеличения. С обычной прямотой Ленин резко говорит об этих ошибках: «...мы зашли дальше, чем это теоретически и политически было необходимо...»[156], и затем: «мы слишком далеко зашли по пути национализации торговли и промышленности, по пути закрытия местного оборота... мы меры не соблюли, не знали, как ее соблюсти»[157].

Но объективности ради Ленин отметил, что тут была налицо и «вынужденная необходимость» — в условиях неслыханно тяжелой войны ничего не оставалось делать, как действовать по-военному и в области экономической. «И чудом было, — говорил Ленин, — что такую войну выдержала разоренная страна, и это чудо не с небес свалилось, а оно выросло из экономических интересов рабочего класса и крестьянства, которые создали это чудо своим массовым подъемом»[158].

Ленин упомянул, что в Москве есть целый слой интеллигентов, злорадствующих в связи с бедственным положением страны и пытающихся создавать вокруг этого «общественное мнение». Они острят на тему о том, что от коммунизма остается лишь загадочная картинка, вроде человека, у которого внизу костыли, а вместо лица — сплошная перевязка. В связи с этим Ленин сказал прямо и очень серьезно: «Россия из войны вышла в таком положении, что ее состояние больше всего похоже на состояние человека, которого избили до полусмерти: семь лет колотили ее, и тут, дай бог, с костылями двигаться! Вот мы в каком положении! Тут думать, что мы можем вылезти без костылей, — значит ничего не понимать!»[159] И далее, он обращает внимание еще на одну сторону ситуации: «Надо признать такой факт, как переутомление и изнеможение масс... Это изнеможение, это состояние — близкое к полной невозможности работать»[160].

Поразительно остро и правдиво описывает Ленин трудное состояние страны и народа, из какого нам надо было выходить.

В чем же видит он выход, предупреждая, что никакого простого решения быть не может, что придется кое-чем поступиться, лишь бы сохранить завоеванное, удержать пролетарскую власть?

Это прежде всего переход от продовольственной разверстки к продовольственному налогу. Общий размер налога должен быть существенно меньше разверстки— приходится думать о снабжении продовольствием не только города, но и деревни.

Однако эта мера окажется действенной лишь одновременно с допущением свободы торговли в рамках местного хозяйственного оборота. Это порождает серьезные опасности. И об опасностях Ильич говорил не менее откровенно. Не надо закрывать глаза, предупреждал он, на то, что замена разверстки налогом означает, что «кулачество из данного строя будет вырастать еще больше, чем до сих пор»[161]. Свобода торговли не может означать ничего другого, как в известной степени движения назад, к капитализму, поэтому в практическом осуществлении эта политика чрезвычайно трудна. Вся задача заключается в том, чтобы, сделав уступку крестьянству, избежать крупных уступок капитализму. Все сводится к вопросу о пределах местного хозяйственного оборота. Но сейчас на эти вопросы нельзя дать ответа — надо дальше посмотреть наш практический опыт, чтобы убедиться, как все это лучше сделать. Скоро Центральный Комитет напишет в партийные организации письмо, ориентировал Ленин, где лучше, чем я говорю сейчас, будет сказано: «ничего не ломайте, не смешите, не мудрите наспех, поступайте так, чтобы максимально удовлетворить среднее крестьянство, не нарушая интересов пролетариата»[162]. Теперь же, на съезде, «мы должны решить этот вопрос принципиально, оповестить об этом крестьянство, потому что посев на носу»[163].

Но если объявить о некоторой, даже ограниченной местными рамками, свободе торговли, то нужно дать товары. Иначе никакой заинтересованности в производстве не появится. Отсюда вынужденная необходимость на время пересмотреть импортные планы, ориентированные на закупку средств производства. «Прежняя наша программа была теоретически правильна, но практически несостоятельна»[164]. Придется вместо намеченного ранее покупать на имеющееся золото продовольствие и другие предметы потребления, уголь».

При этом Ленин сделал и такое заявление, которое нас тогда очень удивило: «Золото-то у нас есть, но золото продать нельзя...»[165] Такова была еще одна сложность того времени, когда в условиях жесточайшей блокады даже на золото было трудно получить из-за границы нужные нам товары. Однако, заметил Ленин, и в этих условиях у нас есть известная экономическая возможность достать товары и кое-что уже делается. Удалось купить некоторое количество продовольствия, а также угля, что может оживить индустрию Петрограда и текстильную промышленность.

Ленин рассказал также о некоторых планах получения займов от иностранных капиталистов и предоставления им концессий, что он связывал воедино с новой экономической политикой. «Тип экономических отношений, — говорил он, — который вверху имеет вид блока с иностранным капитализмом, даст возможность для пролетарской государственной власти свободного оборота с крестьянством внизу»[166].

Вопрос о концессиях Ленин считал чрезвычайно важным и неоднократно к нему возвращался. Он говорил о нем на заседании фракции коммунистов — делегатов VIII съезда Советов, на котором я присутствовал и о чем уже упоминал кратко в предыдущей главе. Он снова разъяснял и развивал вопрос о концессиях и на X съезде— в отчете о политической деятельности ЦК и в докладе о продналоге, в заключительном слове по этому докладу и еще раз в речи при закрытии съезда.

Вследствие крайнего обнищания и экономической отсталости страны Ленин имел в виду через концессии привлечь иностранный капитал в целях восстановления и развития промышленности и таким способом ускорить процесс экономического подъема Советской республики. Он предполагал, если иностранные капиталисты на это пойдут, дать им концессии в первую очередь в нефтяной, горнодобывающей, лесной промышленности. Расчет был на то, что концессионеры привезут новое, современное оборудование для предприятий, продовольствие и другие предметы потребления и облегчат тем самым тяжкую участь хотя бы части голодавших рабочих.

Ради этого он был готов пойти на серьезные уступки коммерческого характера, которые могли бы дать большие прибыли, достаточно привлекательные для капиталистов: «...и тут не жалко сотнями миллионов, а то и миллиардами поступиться из наших необъятных богатств, из наших богатых источников сырья, лишь бы получить помощь крупного передового капитализма. Мы потом с лихвой себе вернем»[167],— говорил Ленин.

Владимир Ильич предостерегал от выдвижения при переговорах о концессиях особых «коммунистических принципов» для деятельности капиталистов, чтобы их не отпугнуть.

Более того, он лично и тщательнейшим образом подготовил 10 важнейших положений примерного концессионного договора, взяв за основу все лучшее и наиболее подходящее для нас из практики капиталистических стран, а также из опыта русского капитализма. Принцип был ясный: лучше переплатить, но ограничить деятельность иностранного капитала коммерческой областью, не давая никаких прав на вмешательство во внутренние дела. Эти условия он зачитал в докладе о концессиях менее чем через месяц после X съезда — на коммунистической фракции ВСЦПС[168]. Самое замечательное в этих основных принципах концессионного договора, разработанных Лениным и принятых Совнаркомом 29 марта 1921 года, было доведенное до совершенства сочетание правовой защиты деловых коммерческих интересов концессионера, с одной стороны, и материальных и социальных интересов рабочих через договор концессионера с профессиональными союзами.

Никаких уступок, которые носили бы политический характер, Ленин не допускал. Основой основ для него было, «чтобы иностранный капитал был привлечен без власти»[169]. Он принципиально исключал возможность создания таких условий, чтобы концессионер получил хотя бы малейшую возможность вмешиваться в политику, так или иначе влиять на государственную или местную власть.

В речи, записанной 25 апреля 1921 года на граммофонную пластинку, Ленин, кроме того, указывал: «Концессия есть своего рода арендный договор. Капиталист становится арендатором части государственной собственности, по договору, на определенный срок, но не становится собственником. Собственность остается за государством»[170].

Эта сторона вопроса, ленинские выступления и его типовые условия концессионного договора имеют и сегодня принципиально важное и, прямо скажем, поучительное значение для развивающихся стран, борющихся за свою экономическую независимость. Естественно, что освободившиеся страны, даже те из них, которые стали на путь социалистической ориентации, вынуждены привлекать иностранный частный и государственный капитал для подъема своей отсталой, по вине колонизаторов, экономики.

И если для Страны Советов ленинские труды и высказывания по вопросам иностранных концессий и займов представляют теперь уже главным образом исторический интерес, то для развивающихся стран они актуальны и злободневны, так как могут помочь им формировать политику в деле привлечения иностранного капитала вообще и в форме концессий в частности на ленинском принципе «без власти».

И дальнейшем, в ходе работы съезда, по другому пункту повестки дня в отсутствие Ленина вновь возник вопрос о концессиях. Надо сказать, что все выступавшие делегаты поддержали политику в отношении концессий.

Но одно дело принципиальное согласие, другое — когда речь заходила о конкретных концессиях, тут понять необходимость этого оказывалось куда труднее. Так было и с бакинскими концессиями. Когда в дни съезда в беседе с бакинскими делегатами состоялся обмен мнениями, для меня выяснилось их внутреннее несогласие с передачей части бакинских промыслов в концессию иностранным капиталистам. И сам я не был убежден, что в этом случае нельзя обойтись без концессий. Тем более что в Баку, в отличие от Грозного, не было никаких разрушений нефтяных промыслов, многие нефтяные скважины были законсервированы фирмами ввиду отсутствия рынка сбыта. Казалось, что в Баку было много не только квалифицированных рабочих, но и инженеров и техников, которые во главе с талантливым и знающим инженером Серебровским успешно восстанавливали нефтяную промышленность. Представляюсь несомненным, что сдача концессий вызовет отрицательную реакцию рабочих Баку, лишь недавно нагнавших оттуда иноземных владельцев. Основываясь на этом, я сказал в прениях по упомянутому докладу, что поскольку принципиальная политика ЦК в вопросах концессий одобрена съездом (по докладу Ленина), то «задача съезда не состоит в том, чтобы решать, в каком месте какую концессию сдавать... Почему решать вопрос о Баку, а не решать об Урале?»[171] Исходя из этого, и предложил, одобрив общую резолюцию о концессионной политике, вопрос о концессиях в Баку снять.

Упоминавшийся уже мною делегат съезда от Кавказа Врачев рассказывал, что после моего выступления по вопросу о сдаче в концессию нефтепромыслов еще во время съезда Ленин вызвал его и спрашивал о том, как воспримут грозненские рабочие решение о сдаче грозненских нефтепромыслов в концессию. Врачев ответил, что ему очень трудно предвидеть, как рабочие отнесутся к этому, так как разговора на эту тему не было, и как рабочие воспримут это — сказать не может.

В результате прений в проект резолюции, где говорилось о том, что объектами концессий может явиться часть грозненских и бакинских промыслов, было добавлено: «поскольку на это последует согласие советской азербайджанской власти»[172]. Резолюция была принята съездом. Вместе с другими за нее голосовал и я.

В речи при закрытии съезда Ленин счел необходимым вернуться к этому вопросу и подробно разъяснить решение о концессиях в Баку и Грозном. «Вопрос этот, — сказал он, — был только вскользь затронут на съезде. Мне не удалось на этом заседании присутствовать, но мне говорили, что у некоторой части товарищей осталось чувство недовольства или недоверия. Я думаю, что оснований для этого быть не должно. ЦК разбирал вопрос о сдаче концессий в Грозном и Баку основательно. Несколько раз создавались специальные комиссии, требовались специальные отчеты от заинтересованных ведомств. Разногласия были, было несколько голосований, после последнего голосования ни одна группа членов ЦК и ни один член ЦК не пожелали воспользоваться своим бесспорным правом апелляции к съезду. Я думаю, что новый ЦК имеет полное право, формальное и по существу, опираясь на решение съезда, решить этот крупный вопрос. Без концессий мы не можем рассчитывать на помощь высокооборудованной современной капиталистической техники. Не пользуясь ею, мы не имеем возможности основу нашего крупного производства в таких областях, как добыча нефти, имеющей исключительное значение для всего мирового хозяйства, поставить правильно... Дело сводится к тому, чтобы улучшить наше экономическое положение, усилить техническое оборудование нашей республики, увеличить количество продуктов, количество предметов продовольствия и потребления для наших рабочих. В этом отношении всякое облегчение имеет гигантское значение для нас. Вот почему сдачи в концессию части Грозного и Баку мы не боимся; сдав в концессию одну четверть Грозного и одну четверть Баку, мы используем эту сдачу — если удастся ее осуществить, — чтобы на остальных трех четвертях догнать передовую технику передового капитализма»[173].

Еще резче говорил по этому поводу Ленин в своем докладе 11 апреля 1921 года на специальном заседании коммунистической фракции ВЦСПС в связи с тем, что после партийного съезда Шляпников и Рязанов вели демагогическую агитацию против концессий, а часть профсоюзных работников проявляла колебания. Следует при этом отметить, что Шляпников до съезда в ЦК и на самом X съезде не высказывался против концессий.

В этом докладе Ленин, напомнив, что на съезде пришлось проводить специальное решение, подтверждающее декрет о концессиях и специально распространяющее его и на отдачу концессий в Баку и Грозном, заявил: «Некоторые бакинские товарищи не хотели примириться с мыслью о том, что и для Баку, может быть в особенности для Баку, концессии необходимы и что в Баку большую часть бакинских промыслов желательно отдать под концессии. Доводы были чрезвычайно разнообразны, начиная с того довода, что мы, мол, сами «исследуем», зачем нам иностранцев звать, продолжая тем, что старые, испытанные в борьбе с капиталистами рабочие не потерпят того, чтобы идти назад под ярмо капиталиста, и т. п.

Я не берусь сейчас судить, насколько в таких доводах была общая принципиальность или, так сказать, бакинский «патриотизм», бакинское местничество. Про себя я должен сказать, что я с этим взглядом боролся самым решительным образом...»[174]

15 марта, заканчивая доклад о замене разверстки натуральным налогом, Ленин еще раз подчеркнул, что главное, центральное звено сейчас — замена разверстки натуральным налогом. Он предупредил, что разработка в деталях и толкованиях мероприятий по замене разверстки налогом и оживлению местного товарооборота потребует нескольких месяцев. «А сейчас нам надо иметь в виду основное: нам нужно, чтобы о принятом вечером же было оповещено по радио во все концы мира, что съезд правительственной партии в основном заменяет разверстку налогом, давая этим целый ряд стимулов мелкому земледельцу расширять хозяйство, увеличивать засев, что съезд, вступая на этот путь, исправляет систему отношений между пролетариатом и крестьянством и выражает уверенность, что этим путем будет достигнуто прочное отношение между пролетариатом и крестьянством»[175]. Съезд встретил этот призыв вождя бурными аплодисментами.

Обсуждение доклада носило весьма специфический характер и имело даже некоторые странности. С одной стороны, в принципиальном отношении вопрос, кажется, был настолько ясен, что после выступления всего четырех ораторов съезд прения прекратил (а в прениях на съезде по другим вопросам было более 250 выступлений!). Никто из ораторов не выступил против предлагавшегося решения или сказанного в докладе. Настораживал, однако, поток записок докладчику, о чем я скажу ниже.

Меня поразило тогда, что ни Троцкий, ни Бухарин, ни Шляпников не сочли нужным выступить хотя бы с кратким словом об их отношении к коренному решению, которое съезд должен был принять.

Скажут: молчание — знак согласия. Но как раз в этом случае молчания было явно недостаточно. Они молчали так, как будто их не касалось изменение отношения пролетариата к крестьянству. (Этого нельзя сказать о других вопросах — Бухарин выступил на съезде 4 раза, Шляпников — 8, Троцкий — даже 10, хотя он и отсутствовал на большинстве заседаний.)

Из оппозиционеров выступил только один Преображенский, который выразил полное согласие с внесенным Лениным проектом ЦК. Затем он затронул ряд важных вопросов, которые возникали в области финансовой политики и денежной системы в связи с допущением местного оборота.

Впрочем, Троцкий мельком коснулся отношения крестьянству в своем содокладе о профсоюзах. Но коснулся хитро, путано, двусмысленно, так, что это никак не освобождало его от выступления по докладу Ленина о продналоге. Троцкий защищался от справедливого обвинения, брошенного ему сторонниками «платформы десяти», а именно: Троцкий не понимает, что политическая суть дискуссии о профсоюзах сводится вовсе не к их производственной роли, а к отношению авангарда рабочих к своему классу и об отношении этого класса к крестьянству. Отвергая это обвинение, Троцкий ссылался на свою годичной давности записку в ЦК с предложением, которое, по его словам, якобы «почти буква в букву совпадает с тем предложением о замене разверстки продовольственным налогом, которое вы теперь будете обсуждать и принимать»[176].

Это была попытка ввести съезд в заблуждение. Записка его предлагала, с одной стороны, свободу торговли без ограничений, а с другой, наряду с послаблением кулачеству, предлагалось углубить меры военного коммунизма в сельском хозяйстве, распространив разверстку также и на обработку земли. Нельзя не учитывать, что и время тогда было другое — гражданская война еще продолжалась.

Но даже если бы в этой ссылке Троцкого и было бы что-либо похожее на правду, это тем более обязывало его выступить на съезде и помочь делегатам, а следовательно и партии, понять необходимость коренного поворота в политике, привести новые аргументы, факты, доказательства, которые убеждали бы в правильности предложения ЦК.

К этому, казалось, обязывало Троцкого и еще одно обстоятельство: его близкий соратник по фракционной борьбе Сосновский, также выступая не прямо, а по другому поводу, проявил по меньшей мере полное непонимание новой политики, заявив на съезде, будто «мы вступаем сейчас в полосу капитуляции (подчеркнуто мною. — А. М.) перед мелкой буржуазией, что эта полоса, может быть, на следующем съезде приведет к капитуляции перед той же мелкобуржуазной стихией, — в зависимости от того, как будет развертываться революция в Европе»[177], и прямо противопоставлял меры, предлагавшиеся в свое время Троцким, нынешним мерам ЦК.

Троцкий не пожелал отмежеваться от этого антиленинского заявления, хотя и числился сторонником мер, внесенных на съезд ЦК. Нет, молчание Троцкого при обсуждении доклада Ленина о продналоге не было знаком согласия. Был бы согласен, он не упустил бы случая проявить свой ораторский талант. Он же не пожелал защищать проект ЦК, хотя проект этот и был одобрен в Центральном Комитете единодушно, как об этом заявил Ленин.

Возможно, Троцкий иначе понимал этот вопрос (отношение к крестьянству никогда не было сильной его стороной, он и раньше ошибался в этом), но отмалчивался, дабы не быть побитым вторично, и по этому, главному вопросу съезда. Поведение Троцкого в последующие годы подтверждает предположение о том, что его молчание означало просто желание зарезервировать свои позиции, не раскрывать их до более «подходящего» времени.

В заключительном слове Ленин особо коснулся содоклада наркома продовольствия Цюрупы. Выступление Цюрупы сводилось главным образом к дополнениям, а также предостережениям по ряду вопросов.

Цюрупа расходился с докладчиком в вопросе о роли кооперации в новых условиях. Дело в том, что в проекте резолюции, внесенном ЦК, и в докладе Ильича предлагалось отменить решение предыдущего, IX съезда партии, по которому деятельность кооперации строилась применительно к продразверстке, которая заменялась теперь натуральным налогом. Цюрупа возражал против предложения Ленина, чтобы заготовки хлеба Накомпрод проводил через кооперацию, на основе договора с Центросоюзом, Цюрупа не доверял кооперации, как организации враждебной, считал, что заготовки надо по-прежнему вести только через Наркомпрод. Ленин доказал ошибочность возражений Цюрупы.

Ленин получил от делегатов «гору записок», как он выразился. Это тоже было формой участия в обсуждении вопроса, хотя и своеобразной. Записки делились на две основные группы. Одна — вопросы техники, практики перехода от продразверстки к продналогу. Другая — в основном содержала опасения, что от поворота в экономической политике больше выиграет мелкая буржуазия, экономически враждебная коммунизму, а не крупная промышленность, представляющая из себя основу перехода к социализму.

Поэтому Ленин сосредоточил основные аргументы на доказательстве того, что при правильной постановке дела выиграть должен социализм: «Кто из нас использует эти средства лучше, от этого зависят и результаты. Я думаю, что если рабочий класс, имея в своих руках важнейшие отрасли крупной промышленности, сосредоточит внимание на наиболее важных из них, то он выиграет больше, чем выиграет мелкая промышленность, хотя бы пропорционально она и вырастала быстрее»[178].

Ленин признал, что не смог даже прочитать всех поданных ему записок, с остальными придется ознакомиться и учесть их при подготовке конкретных вопросов перехода к новой политике. Сейчас же «основному соображению— во что бы то ни стало увеличить количество продуктов — приходится на время подчинить все»[179], — вновь и вновь повторял Ильич.

Обсуждение вопроса о натуральном налоге на съезде особенно отчетливо показало, как тонко чувствовал Ленин пульс жизни, как глубоко понимал психологию и думы масс рабочих и крестьян. С какой страстью и настойчивостью, неутомимо убеждал он партию повернуть экономическую политику, чтобы поскорее вывести страну из этого отчаянного положения на путь оживления экономики, улучшения жизни трудящихся масс, как непременных условий экономического и политического укрепления диктатуры пролетариата, укрепления союза рабочих и крестьян на новой экономической базе, в интересах победы социализма.

На этой стадии разъяснения партии и народу принципиальной необходимости новой политики в отношении к крестьянству и начальных шагов в этом направлении Ленин, в силу отсутствия еще практического опыта, не торопился дать теперь же ответы на все возникающие вопросы. Он не торопился раскрывать все карты. Он намечал лишь первые, главные шаги той экономической политики, которую затем на основе опыта мест, мнения и пожеланий масс развернул в стройную систему. Она вошла в историю нашей страны и международного революционного движения как нэп.

После съезда Ленин скрупулезно изучил его материалы, включая и упоминавшуюся уже «гору записок», сделал еще несколько выступлений на различных собраниях и совещаниях, собрал новые материалы и в течение месяца, к 21 апреля, закончил весьма обстоятельную и содержательную брошюру — «О продовольственном налоге (Значение новой политики и ее условия)». По настоянию Ленина брошюра была выпущена накануне майской 1921 года Всероссийской партконференции и переведена на немецкий, английский и французский языки. Брошюра эта сыграла большую роль для полного понимания в партии и широких массах сути совершенного поворота, а равно и в деле рассеивания всяких недоумений и сомнений в правильности новой экономической политики.

В начальный период новой экономической политики возникали, складывались многие новые принципы хозяйствования, разработанные применительно к периоду перехода от капитализма к социализму. Получали развитие некоторые принципы (материальное стимулирование, сочетание моральных и материальных факторов, хозрасчет), которые в своем дальнейшем развитии оказались пригодными на более длительный период и так или иначе находят применение и в наши дни, в период зрелого социализма и развернутого строительства коммунистического общества.

\subsection*{\center * * *}

Как же была завершена на съезде профсоюзная дискуссия? Доклад Зиновьева и содоклады Троцкого и Шляпникова, а также и прения по ним были неинтересны для большинства делегатов, в том числе и для меня. Нового ничего не было. В течение более двух месяцев мы столько слушали, читали и сами говорили на эту тему, что нам уже чертовски надоело все это снова слушать. было ощущение, какое бывает от пережевывании уже разжеванного. Вопрос был исчерпывающе обсужден и решен на собраниях и партийных конференциях. Предстояло лишь узаконить именем съезда позорное поражение платформы Троцкого — Бухарина, Шляпникова и других и зафиксировать победу ленинской линии.

Жалкое впечатление производил Троцкий (как, впрочем, и Шляпников). Это был уже совсем не тот Троцкий, что месяца два тому назад, когда он рвался к победе. Будучи разбит наголову, он пытался теперь прятать концы в воду, вилял и изворачивался, отрицал, что он начал и раздул дискуссию, пытался многое свалить с больной головы на здоровую. По всему было видно, что он, делая хорошую мину при плохой игре, стремился теперь преуменьшить вред дискуссии для партии, стереть в памяти партии след от причиненного ей вреда, замалчивал свои позорные поступки и применявшиеся им антипартийные методы. Троцкий отмечал, что Ленин и дискуссии оставался архиосторожным. Сам же Троцкий говорил только намеками, в сбивчивых выражениях пытался доказать, будто резолюцию, внесенную «платформой десяти», принимать не следует, но тут же заявлял, что разногласия не так уж велики, что к ним надо отнестись с большим спокойствием и т. д.

Трудно сказать, собирался ли Ленин заранее выступать по этому вопросу повестки дня. Ясно лишь, что выступления Троцкого обойти он не мог. Он и начал свою речь с заявления, что Троцкий сегодня «особенно вежливо полемизировал со мной и упрекал или называл меня архиосторожным. Я должен его поблагодарить за этот комплимент и выразить сожаление, что лишен возможности вернуть его обратно. Напротив, мне придется говорить о моем неосторожном друге, чтобы выразить подход к той ошибке, из-за которой я так много лишнего времени потерял и из-за которой приходится теперь продолжать прения по вопросу о профсоюзах, по переходя к вопросам более актуальным»[180].

Стало ясно, что Троцкий своей речью пытался смазать значение дискуссии, принизить значение затрагивавшихся основных вопросов, исказить принципиальные положения, высказанные Лениным.

Это заставило Ленина, «как ни скучно повторять», снова остановиться на ходе дискуссии, заявить, что он «старался сорвать дискуссию в том виде, как она пошла», напомнить об исходном пункте — лозунге «перетряхивания», квалифицировать отказ Троцкого участвовать в профкомиссии, что было срывом дисциплины ЦК, и получилось «трясение партии и озлобление». Пришлось дать оценку непростительного публичного выступления Троцкого о том, что съезд должен выбирать между двумя тенденциями («Эти слова составляют политическую ошибку, из-за которой мы боремся»), сказать, что спор шел об отношении авангарда рабочего класса к пролетариату, о шатаниях Бухарина и «чехарде платформ», завершив все это выводом, что «в истории партии мы таких шатаний забывать не должны»[181] — этими словами завершил Ленин свое выступление.

Главный удар в этом выступлении Ленин сосредоточил на Троцком. По вопросу же об анархо-синдикалистском уклоне Шляпникова Ленин потребовал специального пункта в повестке дня и готовился выступить по нему с докладом.

Ленин выразил уверенность в том, что общими силами партия сможет упрочить, сплотить ряды, единой выйдет на тяжелую борьбу, которая ей предстоит, и победит.

Съезд закрыл дискуссию о профсоюзах, приняв ленинскую концепцию профсоюзов как школы хозяйничания, школы коммунизма, основного звена, при помощи которого партия связана с рабочим классом в целом. Приняв внесенную «платформой десяти» с Лениным во главе резолюцию «О роли и задачах профсоюзов», съезд отверг как троцкистское требование огосударствления профсоюзов, так и мелкобуржуазные претензии «рабочей оппозиции» на передачу «непосредственным производителям» в лице профсоюзов функций государства в управлении производством.

Учитывая ход предсъездовской дискуссии и самого съезда, Ленин внес в повестку дня вопрос, первоначально в ней не предусматривавшийся, о единстве партии и анархо-синдикалистском уклоне, выступил с докладом и составил проекты резолюций по этому вопросу.

Троцкистская платформа была отвергнута подавляющим большинством партии. Со всей определенностью выяснилось, что массы за ней не пошли. Другое дело — «рабочая оппозиция», которая имела связи с рабочими ряда районов и крупных предприятий и пользовалась там определенным влиянием, так как отражала реальные, хотя и отсталые, нездоровые, настроения определенной части рабочих. Анархо-синдикалистское направление уклона представляло особую опасность в период оживления мелкобуржуазной стихии, крайним выражением чего явился кронштадтский мятеж, глубоко отразившийся на всей работе съезда. Серьезность положения внутри партии видна из следующих слов Ленина в заключении по докладу о единстве партии: «Бывали ли на прежних съездах при самых резких разногласиях такие моменты, которые одной стороной касаются раскола? Не бывали. Есть ли они сейчас? Есть»[182]. Борьба Ленина с «рабочей оппозицией» на съезде произвела большое впечатление. Кроме специального доклада он касался этого вопроса практически во всех своих выступлениях, начиная с политического отчета ЦК, подходил к нему с самых различных сторон. Ленин стремился исчерпывающе, со всей глубиной раскрыть порочность теоретических предпосылок этой платформы, политическую ее вредность.

Как известно, основным тезисом «рабочей оппозиции» было: «Организация управления народным хозяйством принадлежит Всероссийскому съезду производителей, объединенных в профессиональные и производственные союзы, которые избирают центральный орган, управляющий всем народным хозяйством Республики». «Рабочая оппозиция» требовала организации рабочих комитетов для управления предприятиями «на демократических началах», без вмешательства партии и государства. Шляпников и его группа настаивали также на установлении порядка, при котором рабочие и служащие фабрик, заводов, шахт, рудников, предприятий, учреждений и служб транспорта и связи и всех видов сельского хозяйства являлись бы непосредственными распорядителями находящегося в их ведении имущества[183].

Ленин прежде всего высмеял несостоятельность тезиса о «съезде производителей», и, поскольку Шляпников в своем выступлении ссылался на то, что это выражение взято у Энгельса, Ленин сказал: Маркс и Энгельс говорили об объединении производителей, имея в виду тот период, когда не будет классов, о коммунистическом обществе. И ясно, что в момент самой бешеной борьбы классов говорить: «Всероссийский съезд производителей»— это и есть синдикалистский и анархистский уклон, который надо решительно отвергнуть.

Съезд признал пропаганду взглядов «рабочей оппозиции» несовместимой с принадлежностью к партии. Тем самым анархо-синдикалистский уклон уже в первые годы деятельности нашей партии, как партии правящей, потерял право на существование и не мог нанести ущерба в развитии партии и страны. Борьба Ленина и партии с этим уклоном имеет международное значение.

Но надо со всей силой подчеркнуть и другую сторону вопроса. Настаивая на прямой и точной характеристике анархо-синдикализма, его несовместимости с коммунизмом, Ленин отмечал, что главное, чем руководствуется партия, — это «желание предостеречь и поставить вопрос во всей полноте и принципиально»[184].

Желая отсеять здоровое от нездорового и облегчить переход на партийные позиции, если не всех, то большинства сторонников «рабочей оппозиции» (а к ней принадлежали главным образом рабочие), Ленин делал все возможное для сближения их с партией, для облегчения им такого перехода. Он, например, заявил, что, отвергая неправильные предложения сторонников «рабочей оппозиции», партия принимает все правильное, что содержится в их критике совершенных ошибок, в указании на бюрократические безобразия аппарата, необходимости развития демократии, отмечает в резолюции их заслуги в этом и оказывает высочайшее товарищеское доверие, включая их представителей в состав ЦК партии. Ленин обратился к ним с таким призывом: «Вы утверждаете, что мы мало боремся с бюрократизмом, — идите помогать нам, идите ближе, помогайте бороться ...»[185] В этом он видел целесообразную, действенную помощь уклонившимся от генеральной линии членам партии.

Сторонники «рабочей оппозиции», например Медведев, Шляпников, выступили с крайне резкой критикой предложенных Лениным резолюций о единстве партии и об анархо-синдикалистском уклоне, вплоть до обвинений в клевете. Ленин и в этой обстановке, сохраняя принципиальную ясность в определении и осуждении взглядов оппозиции, пытался найти с ними общий язык. Он убеждал, что его резко-критическое определение анархо-синдикалистского уклона соответствует фактам, что уклон не есть еще готовое течение, это еще то, что можно поправить, и он согласен поискать другие слова, точно выражающие эту мысль, внести смягчение. Сейчас другого слова он не имеет, но готов обсудить этот вопрос. Результаты обсуждения можно будет передать в президиум для редактирования с участием Шляпникова, для возможного смягчения формулировок.

Все это Ленин делал ради того, чтобы создать на съезде у членов распущенных фракционных группировок настроение товарищеского доверия, облегчить им и их сторонникам возможность встать после съезда на искренний и честный путь исправления своих ошибок. Такова была принципиальная ленинская линия, которую он точно выскажет в декабре того же года в письме Осинскому: «...надо не видеть «интригу» или «противовес» в инакомыслящих или инакоподходящих к делу, а ценить самостоятельных людей»[186].

Когда со стороны некоторых делегатов — сторонников «рабочей оппозиции» раздались голоса, что они откажутся от вхождения в ЦК и подадут в отставку со своих постов, если будет принята резолюция о единстве и об анархо-синдикалистском уклоне, Ленин выступил и призвал товарищей не становиться на путь «отставников», а работать вместе со всей партией. Более того, он предложил принять короткое решение съезда, призывающее их подчиниться партийной дисциплине и обязывающее оставаться на порученных им постах. Таким образом, Ленин не только не делал вывода, что носителей уклона, после того как съезд отверг их позиции, необходимо снимать с постов, наоборот, в целях единства партии, привлечения уклонистов на партийные позиции, считал необходимым специальным партийным решением запретить Шляпникову и другим представителям «рабочей оппозиции» уходить в отставку.

В результате критики Шляпников в своем выступлении был вынужден поправиться и сказать, что, говоря о производителях, он имел в виду работников промышленности и транспорта, а не частный капитал и крестьян.

Он кроме прочего сказал, что предложение о передаче руководства промышленностью и транспортом «союзам производителей» опирается якобы на параграф 5 экономической части Программы партии.

Ленин, видя, что оппозиция использует неточную формулировку из Программы партии и дает свое особое толкование этому пункту Программы, используя это в качестве повода для прикрытия своего отхода от партийных позиций на путь анархо-синдикализма, предложил дать законное, аутентичное толкование этому пункту Программы в резолюции съезда. Поскольку Программа принята съездом, съезд и может давать толкование, ибо только авторское толкование законно, а автором Программы был съезд. Предложение о толковании этого пункта Программы партии нашло свое отражение в резолюции съезда о синдикалистском и анархистском уклоне.

Ленин — подлинно великий человек, и часто это величие проявлялось в мелочах, которые имели, однако, далеко не мелкое значение.

Из числа сторонников «рабочей оппозиции» кандидатом в члены ЦК был избран Киселев, председатель ЦК Союза горнорабочих, сам из рабочих. Когда на съезде обсуждался ленинский проект резолюции о единстве партии, в котором в качестве чрезвычайной меры на случай возникновения опасности раскола предусматривалась особая процедура исключения за фракционную деятельность членов Центрального Комитета и ЦК и даже из партии, Ленин применил образное выражение: «пулеметы поставить», то есть наглухо закрыть все возможности возникновения фракционности. Киселев выступил с заявлением, что поскольку Ленин применил это выражение, то значит и он, Киселев, тоже должен будет «сидеть за пулеметом, и придется, вероятно, стрелять» из него по своим же товарищам, которые примыкали к «рабочей оппозиции». Он заявил, что не может «согласиться на такую роль пулеметчика» и потому снимает свою кандидатуру и надеется, что съезд примет во внимание его заявление и не заставит его стрелять в своих товарищей из пулемета.

После выступления Киселева снова слово взял Ленин и сказал:

«Товарищи, я очень сожалею, что я употребил слово «пулемет», и даю торжественное обещание вперед и образно таких слов не употреблять, ибо они зря людей пугают и после этого нельзя понять, что они хотят. (Аплодисменты). Никто ни из какого пулемета ни в кого стрелять не собирается, и мы абсолютно уверены, что ни т. Киселеву, ни кому другому стрелять не придется»[187].

Съезд бурными аплодисментами встретил эти слова Ленина.

Параграф 7 резолюции съезда о единстве по решению съезда не был тогда опубликован. Решение о его опубликовании было принято только в 1924 году, на XIII партконференции, в связи с тем, что еще в декабре 1923 года Троцкий навязал партии новую опаснейшую дискуссию.

Суть этого параграфа сводится к тому, что в случае проявления фракционности виновные могут быть исключены из партии, члены ЦК — переведены в кандидаты и даже исключены из партии. Такое решение должно быть принято особым собранием членов и кандидатов ЦК и ЦКК, большинством в две трети голосов.

По этому же пункту резолюции о единстве партии Рязанов выступил с поправкой, которая на первый взгляд, казалось бы, шла в развитие предложения Ленина. А именно: «Осуждая самым решительным образом всякую фракционность, съезд в то же время высказывается так же решительно против выборов на съезд по платформам».

Исходя из высокой принципиальности и чувства ответственности, чтобы не выйти за пределы необходимого, чтобы обеспечить единство партии, с одной стороны, а с другой — сохранить в партии возможность в исключительных случаях создать наиболее подходящую обстановку для обсуждения и решения трудных вопросов, Ленин выступил против этого предложения Рязанова, Разве можно ручаться, что не появится таких вопросов, как, скажем, заключение Брестского мира? — спрашивал он делегатов и разъяснял: «...в вашей резолюции написано: никаких выборов по платформам. Я думаю, что запретить этого мы не в силах. Если наша резолюция об единстве и, конечно, развитие революции нас сплотит, то выборы по платформам не повторятся. Урок, который мы получили на этом съезде, не забудется. Если же обстоятельства вызовут коренные разногласия, можно ли запретить вынесение их на суд всей партии? Нельзя! Это чрезмерное пожелание, которое невыполнимо и которое я предлагаю отвергнуть»[188].

Съезд согласился с Лениным и отверг поправку Рязанова.

Принятие резолюции о единстве партии имело огромное значение и для того периода и для всей последующей истории нашей партии. Она как бы завершала формулирование основных ленинских норм партийной жизни правящей Коммунистической партии. Съезд распустил все без изъятия, образовавшиеся на той или иной платформе группы и поручил всем организациям строжайше следить за недопущением каких-либо фракций.

В этом свете хотелось бы обратить внимание на прозорливые указания Ленина и X съезда нашей партии, подчеркнувших еще почти 50 лет назад такого рода положения: «Использование врагами пролетариата всяких уклонений от строго выдержанной коммунистической линии с наибольшей наглядностью показало себя на примере кронштадтского мятежа, когда буржуазная контрреволюция и белогвардейцы во всех странах мира сразу выявили свою готовность принять лозунги даже советского строя, лишь бы свергнуть диктатуру пролетариата в России».

Резолюция выявляла новейшие тактические приемы врагов Советской власти: «Эти враги, убедившись в безнадежности контрреволюции под открыто белогвардейским флагом, напрягают теперь все усилия, чтобы, используя разногласия внутри РКП, двинуть контрреволюцию так или иначе путем передачи власти политическим группировкам, наиболее близким по внешности к признанию Советской власти». Интересно еще и указание резолюции на то, что «пропаганда должна выяснить также опыт предшествующих революций, когда контрреволюция поддерживала наиболее близкие к крайней революционной партии мелкобуржуазные группировки, чтобы поколебать и свергнуть революционную диктатуру»[189].

С резолюцией о единстве тесным образом было связано и принятое съездом решение о проведении генеральной чистки партии, вызванное проникновением в партию мелкобуржуазных и других случайных элементов. Съезд ориентировал парторганизации на то, чтобы не гнаться за увеличением численности партии, сосредоточиться на улучшении ее качественного состава, повышении активности, развитии внутрипартийной демократии и укреплении ее единства.

Вопросам демократии в решениях съезда вообще было уделено первостепенное внимание. В резолюции по вопросам партийного строительства два раздела специально были посвящены развитию внутрипартийной демократии, резолюция о роли и задачах профсоюзов базировалась на усилении их роли в управлении экономикой и выработке производственных планов.

Враги партии пытались представить дело так, будто единство ее было достигнуто благодаря сужению демократии, установлению в ней «тоталитарных порядков». На самом деле все обстояло как раз наоборот: меры по пресечению фракционности сопровождались широким развитием внутрипартийной демократии, что логически вытекало из окончания войны и перехода к мирному строительству. Однако позже оппозиция — сперва Шляпников, а затем и Троцкий — всячески старалась легализовать свою фракционную деятельность против ленинского курса именно путем противопоставления решений X съезда о демократии его же решению о недопущении фракций.

В конце 1923 года троцкисты сделали новую попытку раздуть внутрипартийную борьбу, в частности они выступали против решений X съезда о единстве партии. В декабре 1923 года, выступая на Кубано-Черноморской партконференции, подвергнув резкой критике заявления троцкистов, я говорил:

«Оппозиция теперь хочет повторения ошибок, пользуясь отсутствием Ленина... Наша партия не допустит группировок и раскола. Партия должна сказать, что дискуссия это вещь хорошая, демократия еще лучше, а единство лучше этих двух вещей, вместе взятых.

Представители оппозиции говорят, что решения X съезда о демократии нужно проводить, а решение этого же съезда о фракциях временно забыть. Ни партия, нм ЦК на это не пойдут. Мы должны провести в жизнь решения X съезда и о проведении демократии и о недопущении образования группировок и фракций».

\subsection*{\center * * *}

Раз уж я так подробно рассказал о X съезде партии, придавая этому съезду первостепенное значение, нельзя не сказать и о том, что съезд уделил серьезное внимание национальному вопросу, приняв важные решения по докладу Сталина. Существенной стороной этих решений была разработка взаимоотношений с крестьянством угнетавшихся ранее национальных окраин. Съезд в очень резкой форме осудил великодержавно-шовинистические проявления и записал, что «ликвидация националистических и, в первую голову, колонизаторских шатаний в коммунизме является одной из важнейших задач партии на окраинах»[190]. С другой стороны, съезд осудил и замечавшийся в ряде республик уклон местных национальных кадров к буржуазному национализму, выражавшийся прежде всего в чрезмерном выпячивании национальных особенностей. Съезд ориентировал партию на то, что основной задачей является ликвидация унаследованного от капитализма фактического неравенства наций в их политическом, экономическом и культурном развитии.

При обсуждении национального вопроса Владимир Ильич на съезде не присутствовал.

В прениях в числе других выступил и я, высказав свое некоторое неудовлетворение докладом. Заявил, что в трудах Ленина, в партийных и государственных документах «национальный вопрос теоретически вполне освещен». Теперь мы должны заняться вопросами его практического разрешения и с учетом опыта определить формы советского строительства на окраинах, классовую структуру народов, их населяющих. В докладе же об этом не сказано. Я оспаривал утверждение Сталина о том, что Азербайджан относится к окраинам, «не ушедшим дальше первобытных форм полупатриархального, полуфеодального быта».

В заключительном слове Сталин сказал, что, очевидно, я смешиваю Баку с Азербайджаном, что я неправ в отношении Азербайджана и что Баку вырос не недр Азербайджана, а надстроен сверху усилиями нефтяных монополий, поэтому он относит Азербайджан в целом к группе отсталых окраин, «которым необходимо применить своеобразные методы втягивания этих окраин в русло советского хозяйства».

Это меня не переубедило. Следует отметить, что фактически в отношении Азербайджана никаких «своеобразных методов втягивания» в русло советского хозяйства не было применено. Там также была введена общепринятая структура советских органов власти.

\subsection*{\center * * *}

Ленин все дни съезда — а он продолжался 9 дней — вел невероятно активную, буквально титаническую работу.

Доклад политической деятельности ЦК и заключительное слово по нему. Большое выступление в прениях о профсоюзах. Доклад о замене продразверстки натуральным налогом и заключительное слово по этому докладу. Доклад о единстве партии и анархо-синдикалистском уклоне и заключительное слово. Два выступления по топливному вопросу, четыре выступления при обсуждении резолюции «О единстве партии». Пять выступлений при обсуждении кандидатур в состав ЦК. Речи при открытии и при закрытии съезда. А всего — 20 выступлений — видимо, больше, чем на любом другом съезде. И это, не считая частных совещаний и двух заседаний съезда по военному вопросу, в которых Ленин также принял участие (эти заседания не стенографировались). Три проекта резолюций съезда написаны собственноручно Лениным. К тому же текущая работа, да кронштадтский мятеж, за ходом ликвидации которого Ленин неустанно следит, вносит предложение послать туда большую группу делегатов съезда во главе с Ворошиловым.

Очень яркой была заключительная речь Ленина на съезде. Кратко, но очень точно говорил он об итогах съезда, проходившего «в момент чрезвычайно важный для судеб нашей революции», когда «элементы распада или разложения, мелкобуржуазная и анархическая стихия поднимают свою голову»[191]. Твердо рассчитывая на сплоченность партии, заявил Ленин, мы уверены, что она из тяжелого опыта революции вышла достаточно закаленной, чтобы всем тяжелым испытаниям и новым трудностям противостоять. Два коренных вопроса были на съезде: об отношении авангарда пролетариата к его массе и об отношении пролетариата к крестьянству, и оба эти вопроса разрешены съездом.

Ленин затронул вопрос о «способе докладов», как он это назвал, на местах о ходе съезда, о разногласиях на нем и, главное, об обстановке в стране.

Как бы опасаясь, что некоторые делегаты, не имея той убежденности, которая была у него самого, возвратясь на места, расскажут все в таком оголенном виде, как это говорилось на съезде, и это произведет впечатление еще большей неуверенности и даже безнадежности, Ленин в заключительном слове говорил о том как должны делегаты по возвращении, на местах, рассказать, что было на съезде. Необходимость в этой вызывалась еще и тем, что более трех четвертей делегатов съезда впервые присутствовали на партийном съезде. Ленин призывал проявлять необходимую осторожность, не преувеличивать опасности положения, не вызывать никакой путаницы и паники, а, наоборот, создавать уверенность в силе, сплоченности для преодоления трудностей, которые стояли тогда перед партией, — он зачитал проект директивы президиума отправляющимся на места делегатам.

Ленин подробнейшим образом рассказал о раздуваемой буржуазной прессой клеветнической кампании вокруг положения в Советской России и на этом фоне подчеркнул: «Я думаю, что большинство присутствующих членов съезда поняло то, какую меру нам надо отвести нашим разногласиям»[192], и закончил призывом к съезду, к партии: «...сплотившись на этом съезде, мы действительно выйдем из наших разногласий абсолютно едиными и с партией, более закаленной, которая пойдет все к более и более решительным международным победам!»[193]

Жизнь подтвердила, что Ленин оказался прав и на этот раз.

\newpage
\section[Ленин на X Всероссийской партконференции — май 1921 года: нэп всерьез и надолго]{\center \textbf{Ленин на X Всероссийской партконференции — май 1921 года: нэп всерьез и надолго}}


Прошло только два месяца после X съезда партии, и значительно раньше установленного срока была созвана X партконференция.

Сразу же при открытии конференции Ленин заявил, что у ЦК нет каких-либо особых, новых вопросов, которые следует ставить на ее обсуждение. Главная причина созыва конференции — это поступившие в ЦК письма и сообщения, в частности от Осинского, недавно объездившего ряд губерний. На местах многие не понимают значения поворота партии в отношениях пролетариата и крестьянства, имеется много сомнений, недоуменных вопросов, мимо которых нельзя пройти. Выступления товарищей с мест помогут прояснить, что это за вопросы, какие сомнения. Тогда можно будет сделать необходимые разъяснения, чтобы не осталось никаких недомолвок.

С другой стороны, ЦК на конференции сможет проинформироваться об опыте мест, узнать данные практики, чтобы сделать необходимые выводы и использовать их при разработке законоположений и инструкций, связанных с осуществлением продналога, с развитием местного оборота в нужных пределах.

Чтобы придать обсуждению вопросов о продналоге и свободе местного оборота сугубо практическое направление, по решению ЦК, кроме Ленина по этому вопросу на конференции выступило еще несколько докладчиков от разных центральных хозяйственных ведомств, осуществлявших новую политику.

Заместитель председателя ВСНХ Милютин сделал доклад о порядке взимания налога и о роли в условиях щи экономической политики нашей промышленности, как крупной, так и в особенности мелкой; о новой продовольственной политике доложил член коллегии Наркомпрода Свидерский; о работе кооперации по организации товарообмена между городом и деревней на основании договора с Наркомпродом выступил председатель Центросоюза Хинчук; по вопросу о финансах в условиях перехода к новой экономической политике от Министерства финансов выступил Преображенский; Радек информировал о предстоящем III конгрессе Коминтерна; Молотов, избранный на X съезде партии членом ЦК, а затем его секретарем, выступил по вопросу о ближайших задачах партийной работы; председатель ВЦСПС Томский информировал о работе IV съезда профсоюзов. Ленин выступил также по информации о профсоюзном съезде с критикой неправильного поведения Томского.

Ленин участвовал в составлении проектов решений конференции. Сохранилась записка Ленина от 24 мая 1921 года, связанная с проектом решения конференции о партийном строительстве.

Главной задачей намечаемого совещания заведующих орготделами губкомов он считал подготовку и сбор материалов по вопросам:

«1) о более тесной связи организаций партии с беспартийными массами.

2) о более широком и систематическом привлечении к работе беспартийных работников.

3) о реорганизации РаКри в направлении борьбы с бюрократизмом и волокитой, улучшения положения рабочих и крестьян и привлечения беспартийных к советской работе.

По всем этим вопросам совещание должно не только собрать и изучить весь материал практического опыта, но и выработать практические предложения для ЦК»[194].

Эти предложения Ленина нашли отражение в принятых конференцией решениях.

\subsection*{\center * * *}

В докладе о продовольственном налоге Ленин коснулся того нового вопроса, который особенно сильно меня заинтересовал. Ссылаясь на сообщение Осинского, он сказал, что на местах задают вопрос: будет ли новая экономическая политика проводиться «всерьез и надолго» или нет.

Подтверждая высказывание Осинского, Ленин сказал: ««Всерьез и надолго» — это действительно надо твердо зарубить себе на носу и запомнить хорошенько...»[195]

Дальше, полемизируя с Осинским насчет того, на какой срок рассчитана новая экономическая политика, Ленин заявил, что названный Осинским 25-летний срок является пессимистическим. В таком деле точный срок трудно рассчитать, во всяком случае «я не такой пессимист», говорил Ленин и отказался назвать конкретный срок.

Мы были озадачены тем, что Осинский отводил четверть века на ликвидацию частного капитала в торговле, в мелкой промышленности и переделку мелкого товарного сельскохозяйственного производства. Так долго ждать?! В ходе этого выступления, как иногда это Ленин делал, он предпринял неожиданный экскурс — дал неизвестную мне периодизацию коммунистической борьбы в России: первый период (курс) — с 70-х годов прошлого столетия до 1903 года; второй — с 1903 по 1917 год; третий — с 1917 по 1921 год. Четыре года последнего периода, говорил Ленин, по своему содержанию больше, чем первые 40 лет. Сравнивая работу партии с четырехлетним курсом высших наук, он определил так: «...мы держим переходный экзамен с третьего на четвертый курс; еще не выдержали его, но по всем признакам выдержим»[196].

Затем в ходе доклада Ленин вновь вернулся к вопросу о продолжительности нэпа. Излагая особо его волнующий вопрос, Ленин имел обыкновение по нескольку раз к нему возвращаться, чтобы оттенить какую-то его сторону или дать более точное определение.

Я был очень обрадован тем, что Ленин, вначале отказавшийся назвать конкретный срок нэпа, в ходе дальнейшего изложения, видимо, почувствовал необходимость, в противовес Осинскому, для правильной ориентации партии вновь вернуться к этому вопросу и определеннее высказаться по нему, противопоставить пессимистическому взгляду Осинского свой реалистический подход. Он сказал: «Если взять пример Осинского и говорить о годах, то мне думается, что тут надо 10 лет положить, потому что придется выдержать переходный экзамен с 3 на 4 курс, а затем нам надо хорошо пройти весь четвертый курс: тогда мы будем, действительно, непобедимы»[197].

Это был срок, более близкий по времени, более близкий нашим сердцам. Осинский брал свой срок с потолка, без марксистского экономического обоснования; Ленин хотя и не сказал, на чем он основывает 10-летний срок, но, поразмыслив, нетрудно было догадаться, из чего он исходил.

Ленин не раз говорил, что уступки капитализму вызваны отсутствием у Советской России развитой крупной промышленности, которая давала бы крестьянам достаточно товаров потребления и сельскохозяйственных машин.

Называя 10-летний срок, Ленин, несомненно, имел в виду срок, установленный VIII съездом Советов для полного осуществления плана ГОЭЛРО — плана электрификации страны, плана развития на ее базе крупной индустрии в 10–15 лет.

Ход событий свидетельствует, что по основному показателю — строительству электростанций — план ГОЭЛРО страна выполнила менее чем за 11 лет. В этот же период был в основном ликвидирован частный капитал в торговле и промышленности, ликвидирован класс кулачества и основная масса крестьян перешла на те или другие простейшие или сложные формы производственного кооперирования.

И в этом конкретном примере проявился свойственный Ленину дар предвидения.

\subsection*{\center * * *}

С поразительной смелостью и откровенностью Ленин раскрыл тяжелое положение, в котором находимся мы, наш рабочий класс, и вновь убедительно доказал, что все это будет преодолено, будет обеспечена победа социализма. Например, он говорил:

«В силу печальных условий нашей действительности, пролетарии вынуждены прибегать к способам заработка не пролетарским, не связанным с крупной промышленностью, а мелкобуржуазным, спекулятивным и путем ли хищений, путем ли частного производства на общественной фабрике добывать продукты себе и эти продукты обменивать на земледельческие, — в этом наша главная экономическая опасность, главная опасность всего существования советского строя»[198].

Эти слова Ленин произносил, когда последствия засухи и отсутствие запасов заставили сократить и без того малые хлебные пайки. Когда, выражаясь словами Ленина, такие пролетарские центры, как Москва и Петроград, «обречены временно на полуголодное, четверть-голодное существование»[199].

Пользуясь таким положением, меньшевики и эсеры продолжают кричать на весь мир, что поскольку пролетариат России деклассирован, то нужно отказаться от диктатуры пролетариата.

Мы им отвечаем: «...что даже тогда, когда пролетариату приходится переживать период деклассированности... он, несмотря на эти минусы, свою задачу завоевания и удержания власти осуществить может»[200].

Далее, говоря о значении новой политики соглашения пролетариата с крестьянством, Ленин сделал предельно ясное и очень интересное разъяснение сути вопроса. Это было важно потому, что в партии у многих по этому вопросу не было ясного представления.

Ленин сказал, что формула соглашения рабочего класса с крестьянством сама по себе настолько неопределенна, что враги Советской власти часто пускают ее в ход против нас. Они вкладывают в нее свое содержание, выгодное их классовым интересам, содержание, которое направлено на реставрацию капитализма.

С точки же зрения пролетариата соглашение между рабочим классом и крестьянством лишь тогда «является допустимым, правильным и принципиально возможным, когда оно поддерживает диктатуру рабочего класса и является одной из мер, направленных к уничтожению классов»[201].

Провозглашая новую экономическую политику, мы идем на компромисс, на известные экономические уступки крестьянству. Ленин всегда, как реальный коммунистический политик, допускал компромиссы и уступки, диктуемые зигзагообразным ходом развития. При этом его больше заботил не материальный объем тех или иных уступок. Решающим для него было — будут ли эти уступки укреплять диктатуру пролетариата и в конечном счете служить победе социализма. Нет почему здесь вновь он считает необходимым повторить:

«Главным вопросом для нас остается — и в течение долгого ряда лет неминуемо останется — правильное установление отношений между этими двумя классами, правильное с точки зрения уничтожения классов»[202].

Эта гениальная идея Ленина, взятая на вооружение Коммунистической партией, на протяжении долгих лет и трудной борьбы выдержала испытание истории.

\subsection*{\center * * *}

Наибольшее впечатление в заключительном слове Ленина на меня произвела та его часть, в которой он откликался на речь в прениях Варейкиса. Это был один из видных партийных губернских секретарей, он пользовался любым поводом, чтобы выступить на партийных конференциях. В практических вопросах его интересно было слушать, но этого нельзя было сказать, когда он пытался «теоретизировать», чувствовалось, что он рисуется и любуется собою, не замечая того, что частенько садится в лужу.

Мы, секретари губкомов, тогда получали и читали меньшевистский ежемесячный журнал «Социалистический вестник», который издавал Мартов в Берлине, начиненный ярко антисоветскими статьями. Прикрываясь марксистской фразеологией, меньшевики грубо искажали нашу действительность.

Варейкис сослался в своей речи на одну из статей, в которой Мартов утверждал, будто ленинская политика экономических уступок крестьянству является несостоятельной и кончится крахом, если не будут одновременно сделаны политические уступки крестьянству. Провозглашаемая политика только экономических уступок крестьянству — это, по Мартову, та же политика «зубатовщины», проводившаяся царским правительством в отношении пролетариата с целью обмана его.

Ссылаясь на эту статью Мартова, Варейкис вдруг предложил в какой-то мере, хотя бы частично, наряду с экономическими сделать и политические уступки крестьянству. Не совсем такие, как предлагает Мартов, но все же сделать какие-то политические уступки, иначе политика нэпа может принять характер «зубатовщины».

Поняв, как запутался оратор, в какую чудовищную путаницу он попал, Ленин сказал: «То, что говорил т. Варейкис, было в сущности не коммунистично и по содержанию мысли напоминало меньшевизм»[203].

Ленин разъяснил, что суть «зубатовщины» заключалась в том, чтобы мелкими экономическими уступками угнетенным классам поддержать власть класса угнетателей. Это не могло заставить пролетариат отказаться от захвата власти для уничтожения классов.

Пролетариат теперь стоит у власти, руководит крестьянством, ведя линию на ликвидацию классов, а не на мелкого производителя. Если мы с этого пути сойдем, мы попадем в мелкобуржуазный лагерь эсеров и меньшевиков, этих самых опасных врагов пролетариата. Ленин привел высказывание Милюкова, что теперь в России успешнее всего могут вести борьбу с большевиками «социалистические партии». Поэтому «эсерам и меньшевикам — «честь и место». Это буквальное выражение Милюкова, показывающее, что он умнее Мартова и Чернова по одному тому, что он представитель крупной буржуазии (хотя бы лично он и не был так умен, как Чернов и Мартов). И Милюков прав.

Милюков понимает, говорил Ленин, «что для перехода назад к капитализму необходимая ступенька — эсеровщина и меньшевизм. Буржуазии такая ступенька нужна, а кто этого не понимает, тот глупец»[204].

Очень интересна мысль Ленина, что в тех условиях крупная буржуазия не могла одержать победу на арене политической борьбы. Она была настолько скомпрометирована в глазах народных масс, что не могла рассчитывать на победу в открытой политической борьбе.

Мелкобуржуазным же партиям, если бы им удалось, прикрываясь флагом социализма, использовав недовольство крестьянства и городских мелкобуржуазных элементов, при шатаниях в самом рабочем классе кое-чего добиться в борьбе с большевизмом, то это было бы только ступенькой для Милюкова.

Эти партии болтались между двумя главными классами — пролетариатом и буржуазией, и поэтому не имели иной перспективы, как стать ступенькой для господства буржуазии. Вот почему нельзя идти ни на какие политические уступки крестьянству, выдвигаемые Мартовым. Мы крестьянство не обманывали. Мы протягиваем ему руку союза.

Идем на серьезные экономические уступки, чтобы обеспечить крестьянству выход из хозяйственной разрухи и улучшить жизнь. Мы прямо говорим:

«... Мы на всякие уступки пойдем в пределах того, что поддерживает и укрепляет власть пролетариата, который неуклонно, несмотря на трудности и препятствия, идет к уничтожению классов и коммунизму»[205].

Ленин считал своим долгом много раз настойчиво подчеркнуть, что вся наша политика вообще и в отношении и крестьянству в частности, все осуществляемые нами мероприятия исходят из одной главной задачи вести страну к коммунизму.

Обращает на себя внимание то, как Ленин прямо связывает оценку ума политического деятеля с тем, как он служит интересам своего класса. В данном случае он Милюкова считает более умным и прозорливым, хотя Мартов и Чернов, может, и были сами по себе умнее Милюкова.

Эту же мысль Ленин как-то уже высказывал в связи с кронштадтским мятежом, когда Мартов навязывал скомпрометировавший себя лозунг Учредительного собрания, а Милюков в противовес этому предлагал идти под флагом Советской власти против большевиков, или, как тогда формулировался им лозунг: «Советы без коммунистов».

Ленин говорил: Милюков правильно понимает, что стоит сдвинуть Советскую власть, свалить большевиков и это будет переходной ступенью к господству капиталистов.

«Мы говорим, что есть две дорожки: за Мартовым и Черновым, а через них — к Милюкову, или за коммунистами»[206].

Ленин считал смешными утверждения, будто мы обманываем крестьянство. Мы честно разъясняем ему нашу политику, наши цели и говорим, что будем руководствоваться интересами пролетариата, не допустим реставрации капитализма и обеспечим путь к коммунизму. Если крестьянство считает себя более усталым, то мы и делаем ему больше экономических уступок для обеспечения от реставрации капитализма.

Мы прямо говорим, в каких пределах эти уступки будут допускаться и в чем они будут состоять.

Говоря о том, что до сих пор на деле большая тяжесть падала на пролетариат, чем на крестьянство, Ленин заявил: «За три с половиной года диктатуры пролетариата больше бедствий взял на себя пролетариат, чем крестьянство. Это — совершенно очевидная и бесспорная истина»[207].

Должен признаться, я даже был рад, что Варейкис своими путаными выступлениями дал повод Ленину сделать эти интереснейшие теоретические и политические высказывания.

\subsection*{\center * * *}

После окончания работы партконференции 30–31 мая 1921 года проходила третья сессия ВЦИК.

Я был на этой сессии как член ВЦИК.

В ходе работы сессии Ленин, озабоченный тем, как обеспечить правильное проведение новой экономической политики на местах, посвятил свое выступление внесенным Советским правительством на рассмотрение сессии ВЦИК Наказу по осуществлению новой хозяйственной политики и постановлению о местных экономических совещаниях (ЭКОСО). Разработкой этих документов руководил непосредственно Ленин. Идея создания ЭКОСО принадлежала ему. Он всегда придавал большое значение сочетанию централизованного руководства органов республики с местной инициативой.

Создание ЭКОСО имело главной целью обеспечение в экономической области сочетания централизма с демократией, координации деятельности представителей центральных ведомств на местах, ликвидацию разнобоя в их работе. В своем выступлении Ленин говорил о важности налаживания регулярной письменной отчетности о проделанной работе как мест, так и центральных ведомств, показа достигнутых результатов и практического опыта. Ленин подчеркивал, что эта отчетность не должна быть бюрократически раздута, но главное, существенное должно быть в ней так отражено, чтобы местные работники, особенно беспартийные, могли с ней знакомиться.

\newpage
\section[Ленин на III конгрессе Коминтерна: завоевание на нашу сторону большинства пролетариата — вот главнейшая задача]{\center \textbf{Ленин на III конгрессе Коминтерна: завоевание на нашу сторону большинства пролетариата — вот главнейшая задача}}



III конгресс Коминтерна проходил спустя почти три года после окончания первой мировой войны. Открылся он в июне 1921 года. Я был тогда секретарем губкома нижегородской партийной организации и входил в состав русской делегации.

Ко времени III конгресса соотношение сил революционного движения и сил буржуазии в Европе существенно изменились. После крупных массовых движений, накал которых относился к 1919 году и к революционным событиям весны 1921 года в Германии, капитализму удалось добиться упрочения своих позиций. Буржуазия опиралась на предательскую политику социал-демократических партий II Интернационала. Создалась ситуация неустойчивого равновесия сил.

В связи с новой обстановкой III конгрессу предстояло обсудить новые задачи Коминтерна, выработать новую тактику.

Возникшие в Коминтерне разногласия потребовали предварительного обсуждения их в комиссиях Исполкома, чтобы договориться о единой точке зрения.

Ленин неоднократно встречался с иностранными делегациями, помогая им выбрать правильную ориентацию. В результате большинство представленных партий приняли за основу предложение российской делегации. Исключение составили германская, австрийская и итальянская партии, которые в противовес предложению нашей делегации выдвинули свою «теорию наступления», считая, что она должна быть положена в основу тезисов Коминтерна о тактике на ближайший период. (Правда, в конечном итоге конгресс принял тезисы единогласно.)

Конгресс открылся 22 июня в Большом театре. Ленин был тогда поглощен проведением новой экономической политики в стране, но это не помешало ему принять самое активное участие в работе конгресса и добиться того, что конгресс в один из серьезнейших периодов мирового революционного движения пошел по верному пути.

Ленин выступал в прениях по итальянскому вопросу 28 июня, затем 1 июля с речью в защиту тактики Коммунистического Интернационала и 5 июля сделал Доклад о тактике Российской Коммунистической партии.

В своем выступлении со всей силой убеждения и полемического мастерства Ленин обрушился на тех, которые поддались соблазну левых фраз, заявив: «Нам, русским, эти левые фразы уже до тошноты надоели»[208].

Высмеивая так называемую наступательную политику левацких элементов, Ленин утверждал, что политика наступления правильна в каждый определенный момент, в каждой данной стране. Но вести такую политику, когда для этого нет необходимых условий, проповедуя всегда и всюду только наступление, смешно, вредно, опасно.

Воздав должное героизму германских рабочих в революционном выступлении начала года, подавленному буржуазией при помощи социал-демократов, Ленин указал, что надо считаться со сложившейся теперь обстановкой в Европе, нельзя оставаться в плену левацких фраз о наступательной тактике.

Точнее, всем партиям Коминтерна необходимо сплотиться вокруг задачи распространения своего влияния в большинстве рабочего класса, суметь повести его за нами. Проделав эту работу, закрепив свое влияние в профессиональных союзах, можно создать условия для наступательной тактики.

Ленин обосновал и разъяснил тактику, предлагаемую нашей партией применительно к периоду, когда в мире «наступило равновесие, — разумеется, весьма неустойчивое»[209]. Первый урок, который должен быть сделан, — это глубокое изучение конкретного развития революции в передовых капиталистических странах, и основной тактикой должно стать «завоевание большинства пролетариата»[210].

Весьма примечательно, что в те годы, когда империализм безраздельно господствовал в колониях, подавляя малейшее движение угнетенных народов, Ленин прозорливо оценивал значение национально-освободительного движения, говоря, что «в грядущих решающих сражениях мировой революции движение большинства населения земного шара, первоначально направленное на национальное освобождение, обратится против капитализма и империализма»[211].

Основную часть своего доклада Ленин посвятил анализу внутреннего политического положения в нашей стране и вытекающей отсюда тактики для нашей партии. Он говорил, что в России стало практически необходимо «определить отношение держащего в своих руках государственную власть пролетариата к последнему капиталистическому классу, к глубочайшей основе капитализма, к мелкой собственности, к мелкому производителю»[212].

Осветив хозяйственную обстановку, сложившуюся к весне 1921 года, Ленин говорил: «Высший принцип диктатуры — это поддержание союза пролетариата с крестьянством...

Единственное средство, которое мы для этого нашли, был переход к натуральному налогу»[213].

Ленин раскрыл сущность и значение нэпа, и делегатам стало ясно, что нэп обеспечит спасение русской революции и имеет огромное международное значение.

Спустя месяц после окончания конгресса Ленин в «Письме к немецким коммунистам» дал исчерпывающую характеристику значения III конгресса. Он писал, что конгресс явился крупным шагом движения вперед Коммунистического Интернационала, что на этом конгрессе надо было «определить, как именно работать дальше, в отношении тактическом и в отношении организационном. Этот третий шаг мы и сделали»[214]. И тут же приводит существеннейшее место тактической резолюции III конгресса: «Завоевание на нашу сторону большинства пролетариата — вот «главнейшая задача»...»[215]

Преодолению назревавших разногласий и успешному проведению своей работы III конгресс Коминтерна был целиком обязан Ленину.

\newpage
\section[Ленин на IX Всероссийском съезде Советов: наши недостатки — продолжение наших достоинств]{\center \textbf{Ленин на IX Всероссийском съезде Советов: наши недостатки — продолжение наших достоинств}}


В конце декабря 1921 года состоялся IX Всероссийский съезд Советов. Ленин выступал на съезде с отчетным докладом ВЦИК и Совнаркома. Будучи крайне озабочен положением крестьянства, Ленин участвовал в совещаниях делегатов крестьян, которые происходили под председательством Калинина, и трижды на них выступил. Ленин внимательно слушал все выступления, записывал главное из них, выступал с разъяснениями по интересующим крестьян вопросам. Иногда он не был готов сразу ответить и говорил, что он записывает вопросы, получит нужные справки, разберется, примет нужные меры и даст ответ.

Встреченный бурной овацией и возгласами делегатов съезда: «Ура! Да здравствует наш вождь товарищ Ленин!», «Да здравствует вождь мирового пролетариата Ленин!», Ленин, приступив к своему отчетному докладу, сказал:

«Первый раз приходится мне давать такой отчет в обстановке, когда прошел целый год, и ни одного, по крайней мере крупного, нашествия на нашу Советскую власть со стороны русских и иностранных капиталистов не было».

Ленин воздал должное той поддержке и сочувствию трудящихся масс во всем мире, которые были «последним, самым решающим источником, решающей причиной того, что все направленные против нас нашествия кончились крахом...».

«Мы, — говорил Ленин, — оказались правы в самом основном. Мы оказались правы в своих предвидениях и в своих расчетах»[216].

Мы слишком хорошо знаем, какие бедствия несет за собой для рабочих и крестьян война. Поэтому мы должны самым осторожным и осмотрительным образом относиться к этому вопросу.

Затем Ленин заявил, что Советское правительство единственное, которое открыто признало преступность политики старой России в отношении других национальностей, населявших нашу страну, и что даже самая маленькая, невооруженная национальность «должна быть спокойна за то, что ничего, кроме мирных намерений, у нас по отношению к ней нет»[217], и что наше желание — во что бы то ни стало поддержать мир со всеми национальностями, входившими ранее в Российскую Империю.

Тут же Ленин счел необходимым предупредить те авантюристические круги, которые, вопреки мирным договорам, хотят попытаться нарушить нашу мирную работу. «Мы не допустим этого ни в коем случае и станем, как один человек, чтобы отстоять свое существование»[218].

Одним из главных результатов отчетного 1921 года, говорил Ленин, во внешнем положении страны является то, что, вопреки заявлениям буржуазных правительств, у нас существуют известные отношения с капиталистическим миром, поддерживается целый ряд довольно правильных торговых отношений и представительств с рядом стран; в этот год в деле торгового оборота с заграницей «мы чрезвычайно шагнули вперед»[219]. Импорт из-за границы в 1921 году составил 50 млн. пудов, т. е. в 3 раза больше, чем за три предыдущих года, вместе взятых. Наш экспорт за один год составил 11,5 млн. пудов вместо 2,5 млн. пудов за три предшествующих года.

Показывая динамику развития торговли с заграницей, Ленин привел еще цифры: в первом квартале отчетного года привоз был около 3 млн. пудов, во втором— 8 млн., в третьем — 24 млн. пудов.

«Цифры эти ничтожно малы, но все же они постепенно увеличиваются. И мы видим, как они увеличиваются в 1921 г., который был годом неслыханной тяжести»[220].

Это увеличение торговли происходит даже вопреки желанию буржуазных правительств, класса капиталистов, которые ненавидят нас ввиду того, что есть большая сила, чем эта. «Эта сила — общие экономические всемирные отношения, которые заставляют их вступить на этот путь сношения с нами»[221].

Переходя к внутреннему положению, Ленин сказал, что главный вопрос за отчетный год заключался в переходе к новой экономической политике. В том, «чтобы сделать первые шаги на этом пути, научиться их делать, приспособить к этому наше законодательство, нашу администрацию...»[222].

Вновь разъясняя правильность новой экономической политики, существом которой является союз рабочего класса с крестьянством, Ленин заявил, что мы только стали на этот путь и делаем в этом отношении шаг вперед, учимся, как налаживать новые отношения с крестьянством, через торговлю. Кое-чему научились, но еще плохо учимся. Призвал всех осознать эту нашу слабость, чтобы скорее ее ликвидировать и справиться с теми экономическими трудностями, которые еще тяжелыми гирями висят на ногах.

«Прекрасно вам известно, товарищи, какой неслыханной тяжестью обрушился на нас голод в 1921 году»[223], — заявил Ленин и доложил съезду о тех чрезвычайных мерах, которые принимает правительство, чтобы помочь крестьянству в снабжении семенами для посевов и дать хлеб для голодающих.

В частности, он заявил, что достигнуто соглашение о закупке за границей (включая американскую помощь) в течение первых трех месяцев на 60 млн. золотых руб. продовольствия для голодающих и семян. Это, конечно, мало, но принесет облегчение в преодолении отчаянной нужды. Мы надеемся добиться расширения посевов весной с большим успехом, чем было этой осенью, и привел цифры: в голодающих губерниях засеяно 75 % озимого клина, в губерниях, частично пострадавших от неурожая, — 102, в губерниях производящих —123, в губерниях потребляющих —126 %.

«Это, во всяком случае, показывает, что как ни дьявольски тяжелы были наши условия, а помощь крестьянству в деле расширения посевов и борьбы с голодом мы все-таки кое-какую оказали»[224].

Осенняя посевная кампания, несмотря на голод, была проведена успешно не только благодаря помощи государства, но главным образом благодаря тому, что крестьянство поняло и высоко оценило замену разверстки продналогом.

Затем Ленин сообщил, что нам необходимо не менее 230 млн. пудов хлеба, а располагаем мы пока 215 млн. пудов. Неизвестно, сколько еще удастся купить дополнительно за границей. Поэтому безусловной необходимостью является собрать продналог полностью.

Топливный кризис продолжал тормозить развитие экономики, несмотря на то что были достигнуты серьезные успехи, особенно во второй половине года, в добыче угля, торфа и нефти. Ленин особенно подчеркивал успехи донецких шахтеров, где выработка забойщика достигла довоенной нормы, хотя общая добыча угля в 1921 году составила всего 350 млн. пудов против 272 млн. пудов в 1920 году, вместо 1 млрд. 700 млн. пудов довоенное время.

Производство чугуна составило в ноябре отчетного года 270 тыс. пудов против 70 тыс. пудов ежемесячно в первой половине года. Увеличилась добыча торфа: в 1921 году 139 млн. пудов вместо 93 млн. пудов в 1920 году, тем самым далеко обогнав довоенные нормы добычи торфа. Привел также подробные цифры первых успехов строительства новых электростанций. Как бы подводя итоги, Ленин сказал: «Мы нисколько не скрываем от себя, что цифры, которые я приводил, показывают уровень мизерный, нищенский, но мы все же мы можем доказать этими цифрами, что, как ни был тяжел 1921 г., какие исключительные тяжести ни обрушивались на рабочий и крестьянский класс, все же мы поднимаемся, мы стоим на верном пути и, напрягая все силы, мы можем надеяться, что подъем станет еще выше»[225].

Большой интерес у меня вызвала ссылка Ленина на французскую поговорку, которая гласит, что «недостатки у человека являются как бы продолжением его достоинств»[226].

Ленин объяснял, что, если достоинства продолжаются больше, чем надо, обнаруживаются не тогда, когда надо, и не там, где надо, они становятся его недостатками. «Величайшим достоинством было то, что в области политической и военной мы сделали всемирно-исторический шаг, который вошел в мировую историю, как смена двух эпох... Энтузиазм, натиск, героизм, который оставался и останется навсегда памятником того, что делает революция и что она могла сделать, помогли решить эти задачи. Вот чем мы достигли нашего политического и военного успеха, и это достоинство становится теперь самым опасным нашим недостатком»[227].

Ленин настойчиво разъяснял необходимость решительного отказа от методов военного коммунизма, говорил о неотложной необходимости перестройки в новых условиях методов работы партийных, профсоюзных и советских кадров в соответствии с требованиями новой экономической политики. Он выражал недовольство тем, что мало сделано в этой области. Он призывал настойчиво учиться работать по-новому.

Затем Ленин сделал очень важный вывод в отношении дальнейшей работы ВЧК в условиях новой экономической политики. Воздав должное героической работе этого стража революции по борьбе с внутренними и внешними врагами советского строя, Ленин поставил вопрос об изменении функций ВЧК, сказав, что «та обстановка, которая у нас создалась, повелительно требует ограничить это учреждение сферой чисто политической...». Новая экономическая политика, говорил Ленин, «требует большей революционной законности»[228].

Ленин подчеркнул, что задача объединения рабочих И крестьян в прочном экономическом союзе поставлена нами правильно, «и эта задача не только русская, но и мировая».

Под бурные, долго не смолкающие аплодисменты Ленин закончил доклад словами:

«Мы эту задачу решим и союз рабочих и крестьян создадим настолько прочным, что никакие силы на земле его не расторгнут»[229].

Кому из нас тогда могло прийти в голову, что это — последний съезд Советов, на котором был Ленин.

Будучи больным, осенью 1922 г. Ленин принимает самое живое участие в решении всех важнейших вопросов, которые стоят перед партией. Одним из них был вопрос о том, какие взаимоотношения должны быть между независимыми советскими республиками и какие начала положить в основу их объединения.

На рассмотрение Ленина был представлен план Сталина об «автономизации», т. е. включении в состав РСФСР независимых союзных республик на правах автономии.

«Я, кажется, сильно виноват перед рабочими России за то, что не вмешался достаточно энергично и достаточно резко в пресловутый вопрос об автономизации»[230], — писал Ленин в декабре 1922 г.

Ленин отверг план автономизации. Он подчеркивал, что надо не уничтожать независимость республик, а создать «еще новый этаж» — Союз равноправных республик.

ЦК партии полностью поддержал его предложение. Ленин был душой подготовки исторического решения об объединении независимых равноправных советских республик к новое государство — Союз Советских Социалистических Республик.

\newpage
\section[Перед XI съездом партии]{\center \textbf{Перед XI съездом партии}}


В начале января 1922 года меня вызвали из Нижнего Новгорода в ЦК, не сообщив о причине. Когда я прибыл в Москву, то в ЦК сказали, что меня хочет видеть Сталин, и что мне следует пойти к нему на квартиру в Кремле. Он жил в здании, на месте которого теперь стоит Дворец съездов. Я зашел во второй этаж. Сталин занимал две комнаты с узким коридором.

Сталин принял меня приветливо. Сказал, что он вызвал меня по поручению Ленина. Речь идет о работе по подготовке к очередному, XI съезду партии. Он начал с вопроса о том, как обстоит с этой точки зрения положение в нижегородской губернской организации. Я коротко сообщил, что после длительной борьбы партийная организация оздоровилась, сплотилась вокруг ленинских идей и уверенно идет к предстоящему съезду. Отдельные группки и лица, связанные со Шляпниковым, есть, но они потеряли всякое доверие и связи с массами.

Видимо, Сталин задал этот вопрос, чтобы как-то начать беседу, ибо ЦК хорошо знал положение у нас. Поэтому он сразу перешел к делу.

Сталин сказал: условия, в которых идет подготовка к XI съезду, коренным образом отличаются от тех, которые были накануне X съезда. На горизонте не видно никаких разногласий и открытых группировок или политических платформ. Он не придает большого значения тому, что Шляпников с узкой группой своих сторонников, оторванный от масс, скрытно ведет групповую работу. Теперь уже он и его группа не представляют серьезной опасности, и если что-нибудь антипартийное возникнет, то быстро провалится.

Главная опасность может идти от Троцкого и его сторонников. Но пока они ведут себя тихо. Никаких заметных разногласий, могущих отразиться на партии, нет. Конечно, от Троцкого можно всего ожидать. До съезда остается еще два месяца. Он может выкинуть какой-нибудь политический трюк, но, по всему судя, это теперь маловероятно. Надо полагать, что скорее всего он, проученный горьким опытом X съезда, изберет другую тактику: пойдет на съезд без разногласий, платформ, демонстрируя полное единство. Этим он рассчитывает усыпить бдительность партии, восстановить свой авторитет, провести побольше своих сторонников в Центральный Комитет. При отсутствии платформ и разногласий делегаты будут давать свои голоса за кандидатов в центральные органы партии по соображениям только персональных достоинств их, предавая забвению прошлые принципиальные разногласия. И если в таких условиях в ЦК будет избрано относительно много бывших троцкистов, то это будет представлять опасность для работы ЦК после того, как съезд разойдется. Потом Троцкий может поднять голову, вызвать разногласия в ЦК и, опираясь на своих сторонников, затруднять работу ЦК, мешать тому, чтобы партия под руководством Ленина сосредоточилась целиком на положительной работе, может начать борьбу против Центрального Комитета.

Рассуждая так, мы озабочены тем, какие делегаты прибудут на съезд, много ли будет бывших троцкистов. Ведь на губернских конференциях теперь будут выбирать только по персональным качествам кандидатов. А среди ответственных работников много троцкистов, они пользуются доверием в своих организациях. Таких особенно много в сибирских губерниях. Мы опасаемся, что из Сибири может прибыть много делегатов-троцкистов. Поэтому Ленин поручил мне все это сказать вам, и, если вы согласны с этим, не поедете ли в Ново-Николаевск (Новосибирск) к Лашевичу[231] — передать от имени Ленина все, что я вам сказал.

Я без колебаний заявил, что понимаю, согласен с этим и готов поехать в Сибирь с этим поручением, заехав до этого на один день в Нижний Новгород. Сталин согласился. Он предупредил меня особо, что об этом можно рассказать только лично Лашевичу, члену Сиббюро, председателю Сибревкома: он поймет правильно, сообщит, кому найдет нужным, и сделает практические выводы, чтобы поменьше троцкистов оказалось среди сибирских делегатов.

Сталин добавил, что мне надо ехать как бы по личным, семейным делам. Эту беседу в свое время я не записывал, но содержание ее хорошо сохранилось в моей памяти.

В беседе Сталин был спокоен, произвел хорошее впечатление на меня. Я собрался было уходить, как вдруг тихо открывается дверь (это было вечером, уже темнело) — и входит Ленин, в пальто и кепке. Поздоровался и, улыбаясь, смотря на Сталина и на меня, в шутку, с присущим ему одному прищуром глаз, сказал: «Вы свои кавказские разногласия обсуждаете?» Сталин ответил, что он передал мне все, что было условлено, и что я согласен во всем и поеду через день к Лашевичу. Я был смущен этой неожиданной встречей с Лениным и поторопился уйти, попрощавшись с Лениным и Сталиным.

Прибыв в Ново-Николаевск, попал в непривычный для меня 45-градусный мороз. Дыхание мое задерживалось. По дороге от вокзала до Сибревкома я шел пешком, через каждые 5 минут заходя в лавчонки, чтобы перевести дух, согреться. Лашевича на месте не оказалось — он был в отъезде. Устроившись в общежитии, четыре дня я ждал, пока приехал Лашевич. Он обрадовался привезенному ему поручению от Ленина и сказал: «Хорошо, что вы приехали и привезли эти указания. Мы, как провинциалы, ничего подобного не думали, и наверняка немало бывших троцкистов было бы выбрано на съезд. Но теперь мы это учтем, передайте в Москве, чтобы Ленин не беспокоился за Сибирь».

Вернувшись в Москву, я зашел к Сталину и доложил об исполнении поручения.

\newpage
\section[Ленин на XI съезде партии: теперь цель — перегруппировка сил]{\center \textbf{Ленин на XI съезде партии: теперь цель — перегруппировка сил}}


XI съезд партии состоялся 27 марта — 2 апреля 1922 года, ровно через год после X съезда, поворотного в жизни нашей партии и страны.

На X съезде партии Ленин говорил о необходимости чистки рядов партии от примазавшихся, от тех, кто пришел в партию ради своих личных выгод, а не для служения революции, от тех пришельцев из других партий, которые не сбросили с себя багаж мелкобуржуазной идеологии.

Ленин был против гонки за большой численностью рядов партии, он добивался, чтобы в партии были только настоящие революционеры, преданные до конца тему делу, честно выполняющие высокий долг коммуниста. За время между X и XI съездами была проведена чистка партии, в результате которой было исключено или механически выбыло более 25 % состава партии.

Делегатов с решающим голосом на XI съезде было всего 522 человека, с совещательным — 146. Съезд проходил в Свердловском зале в Кремле. Мы сидели на венских стульях, которые легко можно было придвинуть друг к другу, и президиум занимал мало места. Сидели плотно, умещались все.

Съезд открыл Ленин кратким вступительным словом Он отметил, что в течение истекшего года империалистические государства на нас не нападали и мы имели возможность все свои силы посвятить задачам социалистического строительства и сделали в этой области только первые шаги. Но, говорил Ленин, «бедствия, которые обрушились на нас в этом году, были едва ли еще не более тяжелыми, чем в предыдущие годы»[232].

Ленин говорил о голоде и разорении — прямых последствиях империалистической войны, интервенции и гражданской войны. И мы вновь и вновь поражались тому, как Ленин умел без прикрас показать тяжесть действительности, одновременно внушая уверенность, что все трудности мы можем успешно преодолеть и преодолеем, если сохраним и укрепим единство нашей партии.

Отметив рост во многих странах силы и влияния компартий, Ленин говорил, что если мы в сотрудничестве с ними «сумеем трезво оценить свое положение и не побоимся сознать свои ошибки, то из всех этих трудностей мы выйдем победителями»[233].

Ленин предоставил слово Енукидзе, который от имени собрания представителей делегаций предложил избрать состав президиума из 19-человек, что было единодушно принято. Затем Томский от имени того же собрания представителей предложил избрать трех секретарей съезда: Енукидзе, Микояна и Кирова. И это было принято.

После продолжительного обсуждения и утверждения порядка дня на том же заседании с политическим докладом ЦК выступил Ленин.

Ленин начал довольно своеобразно: «Позвольте мне политический доклад ЦК начать не с начала года, а с его конца. Сейчас наиболее злободневным вопросом политики является Генуя»[234].

Ленин рассказал, как тщательно из лучших наших дипломатов подбирался состав делегации. «...А у нас теперь, — заметил Ленин, — советских дипломатов порядочное количество, не так, как в начале существования Советской республики»[235].

Подробно обрисовав сложную борьбу вокруг Генуэзской конференции в различных странах и коалициях, противоречивое отношение к ней политических партий, Ленин сказал, что в Геную мы идем, «как купцы. Нам надо торговать, и им надо торговать. Нам хочется, чтобы мы торговали в нашу выгоду, а им хочется, чтобы было в их выгоду. Как развернется борьба, это будет зависеть, хотя и в небольшой степени, от искусства наших дипломатов»[236].

Трудно определить исход этой конференции, но через Генуэзскую конференцию, если объективный ход событий будет для нас благоприятным, то мы целей своих достигнем.

Торговля между нами и капиталистическими странами— объективная хозяйственная необходимость, она проложит себе дорогу и неизбежно будет развиваться.

Ленин чувствовал пульс не только нашей жизни, но и буржуазных стран. Он умел отделить объективную необходимость и субъективные отклонения политики от необходимости, он предвидел, что в конце концов одержит верх объективная необходимость.

Вот уже почти полвека из года в год расширяются наши экономические связи с капиталистическим миром, растет объем торговли и количество стран, имеющих с нами экономические связи.

Сделав обзор главных вопросов политики за истекший год и задач предстоящего года, Ленин сказал: «Мне сдается (или, по крайней мере, такова моя привычка), что в политическом докладе ЦК нам надо вести речь не просто о том, что было за отчетный год, но и о том, какие за отчетный год получились политические уроки — основные, коренные, чтобы свою политику на ближайший год определить верно, чтобы кое-чему за год научиться»[237].

Весь отчетный год прошел под знаком новой экономической политики, которая научила нас многому. Назрела необходимость нового подхода к социалистической экономике, которая должна сомкнуться с экономикой крестьянской. Мы еще не создали экономической смычки, мы к ней только подходим и уже имеем некоторый опыт. Несомненно, мы создадим смычку социалистической экономики с крестьянской экономикой, и это решит в последнем счете «и судьбу нэпа и судьбу коммунистической власти в России»[238].

Золотыми буквами вписываются в историю нашей партии мудрые ленинские слова:

«Сомкнуться с крестьянской массой, с рядовым трудовым крестьянством, и начать двигаться вперед неизмеримо, бесконечно медленнее, чем мы мечтали, но зато так, что действительно будет двигаться вся масса с нами. Тогда и ускорение этого движения в свое время наступит такое, о котором мы сейчас и мечтать не можем. Это, по-моему, первый основной политический урок новой экономической политики»[239].

Это сказано вскоре после голодного 1921 года, в период ужасающей разрухи. И это не было выражением мечты и фантазии, а, говоря современным языком, явилось научно запрограммированным, математически точным прогнозом хода социалистического строительства в нашей стране, нашей борьбы за коммунизм.

Действительно, около 10 лет мы терпеливо укрепляли смычку с крестьянством, медленно двигались вместе с его основной массой вперед. Терпеливая, политически и практически плодотворная политика нашей партии подготовила уже в 30-е годы такое ускоренное движение вперед, которое привело к победе коллективного ведения сельского хозяйства, а в наши дни — к торжеству социалистических производственных отношений в стране.

В нашей экономической политике мы обязаны были доказать крестьянину на практике, что умеем не хуже капиталистов налаживать торговлю с деревней. Нам это было очень трудно доказать за первый год нэпа ввиду тех объективных трудностей, которые навалились на нас в связи с засухой, тяжелыми последствиями разорительных войн, а также ввиду того, что мы не научились еще торговать.

Бичуя неправильную практику работы многих коммунистов— хозяйственных руководителей, не желающих учиться торговать, Ленин требовал от них осознать, что во многих случаях, хотя они и стоят во главе учреждений, но дело ведут не они, а буржуазные элементы, засевшие в советском аппарате. Ленин критиковал бюрократические методы работы таких товарищей, ибо бюрократизм — это опаснейший враг советского строя.

Мы, говорил Ленин, хозяйничать не умеем, это за год доказано, а мы должны научиться этому делу. Тут предстоит «последний и решительный бой». Ленин предупреждал: «Либо мы этот экзамен соревнования с частным капиталом выдержим, либо это будет полный провал»[240].

Ленинские предупреждения вселяли тревогу, заставляли всех подтянуться, звали партию к учебе, к напряженному труду. Но тут же, в следующих фразах, Ленин выражал уверенность в грядущей победе: «Чтобы выдержать этот экзамен, для этого мы имеем политическую власть и целую кучу всяких экономических и других ресурсов — все, чего хотите, кроме уменья». Но «мы и эту трудность победим, несмотря на то что она гораздо больше, чем предыдущая трудность, потому что она лежит в нас самих»[241].

Большое внимание уделил Ленин и выступавшие в прениях товарищи вопросу о госкапитализме. Ленин говорил о Бухарине, Преображенском, Осинском, Шляпникове, которые, не вымолвив ни одного слова против политики нэпа на X съезде, стали критиковать ее теперь. Он высмеял их высказывания о характере и природе госкапитализма.

Для понимания этого вопроса они читают книги и учебники, когда-то написанные о госкапитализме. Но ведь в этих книгах ни слова не сказано, да и не могло быть сказано, о госкапитализме при диктатуре пролетариата.

«Даже Маркс не догадался написать ни одного слова по этому поводу и умер, не оставив ни одной точной цитаты и неопровержимых указаний. Поэтому нам сейчас приходится выкарабкиваться самим»[242].

Анализируя сегодня все, что было сказано в те годы Лениным о госкапитализме, приходишь к выводу, что, хотя госкапиталистическая форма хозяйства и не получила у нас широкого развития, вопрос этот надо было тогда всесторонне теоретически осмыслить. Это было необходимо тем более потому, что Бухарин, Преображенский, Ларин и другие вносили в вопрос о госкапитализме несусветную путаницу.

Госкапитализм при Советской власти «это — тот капитализм, который мы сумеем ограничить, пределы которого мы сумеем установить, этот государственный капитализм связан с государством, а государство это — рабочие, это — передовая часть рабочих, это — авангард, это — мы.

...И уже от нас зависит, каков будет этот государственный капитализм»[243]. У нас есть все необходимое, чтобы подчинить этот госкапитализм себе, а не быть подчиненными.

Перейдя к вопросу об отступлении в экономической политике, Ленин заявил, что мы год отступали. «Мы должны теперь сказать от имени партии: достаточно! Та цель, которая отступлением преследовалась, достигнута. Этот период кончается или кончился. Теперь цель выдвигается другая — перегруппировка сил. Мы пришли в новое место, отступление в общем и целом мы все-таки произвели в сравнительном порядке»[244].

Ленин счел нужным подчеркнуть, что группа Шляпникова сеяла панику, мешала тому, чтобы отступление провести дисциплинированно.

Ленин говорил, что нам надо было «в порядке отступать и сугубо держать дисциплину. Тот, кто сколько-нибудь вносит ноты паники или нарушение дисциплины, погубил бы революцию, потому что нет ничего труднее, как отступление с людьми, которые привыкли завоевывать, которые пропитаны революционными воззрениями и идеалами и в душе всякое отступление считают вроде того, что гнусностью»[245].

Ленин отлично знал закономерности, действующие и при наступлении, и при отступлении, решающую роль организованности и гибельность паники, все объективные и субъективные факторы, действующие в классовых боях и военных действиях. Все это партия вспоминала и учитывала не однажды и прежде всего в ходе Великой Отечественной войны, особенно в ее первые месяцы.

«Самая опасная штука при отступлении — это паника, — говорил Ленин. — Ежели вся армия (тут я говорю в переносном смысле) отступает, тут такого настроения, которое бывает, когда все идут вперед, быть не может. Тут уже на каждом шагу вы встретите настроение до известной степени подавленное... И в этом громадная опасность: отступать после победоносного великого наступления страшно трудно; тут имеются совершенно иные отношения; там дисциплину если и не поддерживаешь, все сами собой прут и летят вперед; тут и дисциплина должна быть сознательней и в сто раз нужнее, потому что, когда вся армия отступает, ей не ясно, она не видит, где остановиться, а видит лишь отступление, — тут иногда достаточно и немногих панических голосов, чтобы все побежали»[246].

Эти глубокие мысли Ленина помогают понять и осмыслить многие трудные минуты нашей революции— и в годы гражданской войны, и в период Великой Отечественной войны. Вскрытые им закономерности психологии масс типичны и для военных действий, и для действий в области социальной борьбы. Это не были абстрактные рассуждения, а вывод, сделанный из реального анализа гражданской войны и года экономического отступления.

Очень запомнились ленинские разоблачения Устрялова, эмигранта, бывшего колчаковского министра. Устрялов представлял в те дни часть русской эмигрантской буржуазии, которая якобы отказывалась от борьбы с Советской властью. Более того, устряловцы объявили себя сторонниками Советов, исходя из того будто, что нэп — это не тактика большевистской партии, а эволюция советского строя, т. е. начало перерождении нашей страны в буржуазное государство[247].

Ленин со всей серьезностью относился к устряловским настроениям. «Враг говорит классовую правду, указывая на ту опасность, которая перед нами стоит. Враг стремится к тому, чтобы это стало неизбежным... Это — основная и действительная опасность»[248]. На нас с оружием в руках пока сегодня не наступают. Тем не менее борьба с капиталистическим обществом стала в сто раз более ожесточенной и опасной, потому что мы не всегда ясно видим, где против нас враг и кто наш друг.

Ленин настойчиво внедрял в наше сознание мысль о том, что в переживаемые дни идет «отчаянная, бешеная, если не последняя, то близкая к тому, борьба не на живот, а на смерть между капитализмом и коммунизмом»[249].

Ленин ставил первейшую задачу — завоевать доверие крестьянина в новых условиях и укрепить смычку. Строить коммунизм руками одних коммунистов невозможно, надо привлечь к этому делу массы беспартийных. Надо добиться того, чтобы крестьянин сказал: «Как ни труден, как ни тяжел, как ни мучителен голод, а я вижу, что власть хотя и непривычная, хотя и необычайная, но от нее получается практическая, реально ощутимая польза»[250].

Вновь перечитывая выступления Ленина, поражаешься, какую убежденность нужно иметь в своей правоте, чтобы так беспощадно критиковать коммунистов, советский аппарат с единственной целью — поскорее исправить недостатки, преодолеть бескультурье, бюрократизм, неумение по-новому хозяйствовать. И как это сочеталось с такой глубокой верой в величие Советского государства!

Ленин говорил: «Ведь, мы живем в XX веке, и единственный народ, который вышел из реакционной войны революционным путем... — это русский народ, и вывела его русская революция. И то, что завоевано русской революцией, — неотъемлемо. Этого никакая сила не может взять, как никакая сила в мире не может взять назад того, что Советское государство было создано. Это — всемирно-историческая победа. Сотни лет государства строились по буржуазному типу, и впервые была найдена форма государства не буржуазного»[251].

И как не вспомнить ленинские слова насчет недостатков нашего госаппарата. «Может быть, наш аппарат и плох, но говорят, что первая паровая машина, которая была изобретена, была тоже плоха, и даже неизвестно, работала ли она. Но не в этом дело, а дело в том, что изобретение было сделано. Пускай первая паровая машина по своей форме и была непригодна, но зато теперь мы имеем паровоз. Пусть наш государственный аппарат из рук вон плох, но все-таки он создан, величайшее историческое изобретение сделано, и государство пролетарского типа создано, — и поэтому... во всем мире все рабочие тяготеют к Советскому государству. Вот те великие завоевания, которые нами достигнуты и которые являются неотъемлемыми»[252].

\subsection*{\center * * *}

В те годы на всех, пожалуй, съездах, конференциях и пленумах всеобщее внимание привлекал своими экстравагантными выступлениями Ларин. И на XI съезде партии Ленин высмеял Ларина, который не имел представления, что такое новая экономическая политика и как к ней нужно относиться. Ларин являлся политическим путаником, который выступал остроумно, с шутками, со всякими фантастическими планами и предложениями. Полемизируя с ним, Ленин очень интересно высказался о роли фантазии. Высмеяв Ларина, Ленин говорил, что он «человек очень способный и обладает большой фантазией. Эта способность чрезвычайно ценна. Напрасно думают, что она нужна только поэту. Это глупый предрассудок! Даже в математике она нужна, даже открытие дифференциального и интегрального исчислений невозможно было бы без фантазии. Фантазия есть качество величайшей ценности, но у тов. Ларина ее маленький избыток. Например, я бы сказал так, что, если бы весь запас фантазии Ларина разделить поровну на все число членов РКП, тогда бы получилось очень хорошо»[253]. Можно себе представить, какой смех и аплодисменты вызвала такая острая ленинская характеристика Ларина!

Преображенский, выступая с критикой руководства экономикой со стороны ЦК партии, предложил создать Экономбюро, т. е. параллельно с Политбюро и Оргбюро новую организацию, которая должна решать вопросы экономической политики. Ленин подверг резкой критике это предложение, принятие которого внесло бы мешанину и кашу в руководство страной, «все, что мы говорим о разделении партийной и советской работы, пойдет насмарку». Экономические вопросы неотделимы от политических, ибо «политика — это концентрированная экономика».

Ленин подверг справедливой критике Шляпникова и его группу за поведение после X съезда и неправильное выступление на XI съезде.

Весь истекший год Шляпников продолжал вести фракционную борьбу. Хотя фракция его распалась, но осколки ее остались как руководимая им группировка; выступления против решений ЦК вынуждали в соответствии с резолюцией X съезда о единстве партии ставить вопрос об исключении Шляпникова из ЦК. Такое решение не было принято из-за нехватки одного голоса до законных двух третей, могущих решить вопрос об исключении. Шляпникову было вынесено строгое предупреждение.

Накануне съезда в феврале Шляпников собрал группу своих сторонников из 22 человек и направил заявление Исполкому Коминтерна (жалоба с огульными обвинениями партии).

Ленин отметил примечательное явление: платформу Шляпникова перед X съездом подписало 38 человек, а его новую платформу только 9 человек из них, остальные от Шляпникова отошли.

Ленин не оспаривал права Шляпникова, как и любого коммуниста, подавать заявление в Коминтерн, но напомнил, что большинство членов ЦК было против позиции и действий Шляпникова. Шляпников продолжает фракционную работу, ведет дело к расколу. Мы не можем этого дольше терпеть.

На съезде партии острому обсуждению подверглась деятельность ЦКК в части, где дело касалось борьбы с остатками фракционности, за единство рядов партии. Речь шла главным образом о конфликтах, которые создавала «рабочая оппозиция» во главе со Шляпниковым.

В докладах и речах много говорилось о раскольнической деятельности оппозиции не только в Москве, но и на местах, особенно в промышленных центрах страны. Сторонники «рабочей оппозиции», в частности, говорили, что в Нижегородской губернии, где я был секретарем губкома, идет преследование сторонников «рабочей оппозиции», что в ЦК и ЦКК поступили законные жалобы, однако они не дали никаких результатов.

Действительно, Шляпников неоднократно обращался к Ленину с жалобами на якобы преследование его сторонников в нижегородской организации. И в течение менее года по решению ЦК и ЦКК 5 раз выезжали из Москвы руководящие товарищи для разбора конфликтов, искусственно создаваемых местной оппозиционной группой сторонников Шляпникова.

Вызвано это было тем, что лидер нижегородской оппозиции Челышев, состоявший членом ЦКК, вел двурушническую политику, тщательно скрывая продолжение своей фракционной деятельности и после X съезда партии. Он проводил всю работу по указке Шляпникова. Челышев держал связь с местными группками, инспирировал их на фракционные выступления, обещая при этом поддержку, и на посылку жалоб в центр о якобы неблагополучном положении в нижегородской партийной организации.

Первым поводом для жалобы в ЦК был конфликт губкома с руководящими работниками губернского совета профсоюзов. Руководство губпрофсовета в своем большинстве стояло на стороне «рабочей оппозиции».

Им удалось так организовать дело, что фракция 6-го губсъезда профсоюзов отвела из списка в состав президиума губпрофсовета двух кандидатов, включив вместо них двух сторонников «рабочей оппозиции», чем обеспечивалось оппозиционное большинство в руководящем органе профсоюзов.

Помню, мы на заседании губкома решили тогда согласиться с вводом в состав президиума одного из предложенных фракцией кандидатов — Козина, как местного работника, и предложили фракции пересмотреть свое решение в отношении московского работника Ищенко, ближайшего соратника Шляпникова, мотивируя тем, что он, постановлением ЦК партии, выведен из состава Цектрана и что в важной работе губпрофсовета нужна политическая выдержка, которой у Ищенко нет.

Они возражали, не соглашались. Губком, пользуясь своими уставными правами, принял решение, обязывающее фракцию голосовать на съезде за список с исправлением, внесенным губкомом.

На новом заседании фракции губсъезда было сообщено решение губкома. После бурных прений фракция отклонила решение губкома и постановила голосовать за собственный список без изменений.

Мы вновь собрались в губкоме с участием представителей фракции и подтвердили свое прежнее решение. Бюро губкома поручило мне пойти на заседание фракции съезда, разъяснить смысл решения губкома и обеспечить его выполнение.

Но оппозиционеры «поработали» основательно, а наши сторонники проморгали, как это было видно из реплик и выкриков участников собрания во время моего выступления. Тогда я достал из кармана Устав РКП(б) и зачитал пункт, который определял права и обязанности фракции во внепартийных организациях. Особо подчеркнул, что после вторичного решения губкома фракция обязана подчиниться, и, ссылаясь на Устав партии, потребовал от товарищей в порядке партийной дисциплины голосовать на съезде за список, санкционированный губкомом.

С мест раздавались выкрики о неправильном поведении губкома, о подавлении демократии, грозили пожаловаться в Центральный Комитет партии на поведение губкома. Я разъяснил, что их законное право послать жалобу в ЦК и изложить в ней свое мнение, но все коммунисты обязаны в соответствии с Уставом голосовать так, как решил губком. Несогласие и подача жалобы никого не освобождает от обязанности строго выполнять решение губкома. Фракция постановила подчиниться решению губкома. Группа сторонников Шляпникова заявила, что пошлет жалобу в Москву с просьбой прислать комиссию для разбора на месте.

К тому же еще до этого возник конфликт со сторонниками Шляпникова в Выксунском промышленном районе. Там местная фракция союза металлистов опротестовала решение райкома партии о переводе с профсоюзной на другую работу председателя райкома союза металлистов, который все время конфликтовал с райкомом партии. Фракция послала жалобу в Москву, прося присылки представителя ЦК для расследования.

Все жалобы в Москве продвигались через Шляпникова, и в июне 1921 года в Нижний для их проверки прибыл член ЦКК т. Литвин-Седой, старый коммунист-рабочий, участник декабрьского вооруженного восстания 1905 года в Москве.

Он побывал в Береговом и Выксунском районах, беседовал с коммунистами, с членами фракции губпрофсовета — авторами жалобы, участвовал в заседании пленума губкома. На нем Литвин-Седой заявил, что, произведя проверку, он считает действия губкома правильными, что жалобы о якобы разгоне сторонников Шляпникова («рабочей оппозиции») оказались необоснованными. Он пришел к выводу, что борьба губкома с остатками фракционных группировок велась в соответствии с решением X съезда партии. Заявил, что об этом напишет подробный доклад в ЦКК. Вскоре, 7 июля 1921 года, на заседании Политбюро ЦК, с участием Ленина, председатель ЦКК Сольц докладывал о результатах проведенного ЦКК расследования жалоб, поступивших из Нижнего от сторонников Шляпникова. Политбюро постановило: «Доклад принять к сведению».

Дважды приезжал в Нижний и Сольц — в августе и ноябре 1921 года. Оба его приезда были связаны с жалобами в ЦК и ЦКК о том, что в Нижнем исключали при чистке партии тех, кто прежде принадлежал к «рабочей оппозиции», и что якобы имело место групповое исключение.

Сольц оба раза беседовал с авторами жалоб, был в районах, разбирался в решениях комиссий по проверке и на заседании губкома информировал нас о положительной оценке результатов проверки работы комиссий.

Приезжал к нам и председатель Центральной комиссии по проверке, пересмотру и чистке партии Залуцкий в связи с тем, что директор Сормовского завода Чернов, подписавший перед X съездом партии платформу «рабочей оппозиции», будучи в Москве, жаловался в моем присутствии на заседании Оргбюро ЦК, где стоял вопрос о положении в Сормове. Чернов в докладе изложил, что Нижгубком и Нижгубисполком не руководят Сормовским райкомом партии, всячески дискредитировал секретаря этого райкома.

Оргбюро постановило поручить Залуцкому выехать в Нижний Новгород и принять участие в работе предстоящей губпартконференции, совместно с губкомом разобраться с жалобой Чернова и результаты доложить в ЦК.

Залуцкий проверил работу губкомиссии по чистке партии, рассмотрел имевшиеся жалобы сторонников оппозиции, побеседовав с ними. Проверил жалобу Чернова, обсудив вопрос с губкомом, признал ее необоснованной. Участвуя в работе губконференции, Залуцкий выступил с докладом о чистке партии, одобрив при этом деятельность нашей губернской комиссии. Возвратившись в Москву, Залуцкий представил положительную оценку как работы губкома, так и хода XII губпартконференции, которая продемонстрировала сплоченность и идейное единство организации, несмотря на жалобы отдельных сторонников Шляпникова.

За две недели до XI съезда партии шляпниковцам удалось втайне от райкома партии провести работу среди своих бывших сторонников в первичной партийной организации Канавинского железнодорожного участка и добиться того, что на собрании большинство высказалось за поддержку «заявления 22-х».

Губком сразу обсудил это событие и немедленно довел его до сведения ЦК. В связи с этим в Нижний прибыл член ЦК партии Томский. Было созвано новое собрание этой организации, где выступал Томский и призывал пересмотреть принятое решение. Однако собрание большинством голосов отказалось отменить прежнее решение, но одновременно постановило, что будет подчиняться тем решениям, которые примет предстоящий XI съезд партии.

Характерно, что это был единственный случай среди партячеек Нижегородской губернии, когда оппозиционерам удалось в это время получить поддержку.

Что касается обоснованности жалоб на действия губкома, то во всех случаях доклады приезжавших для проверки из центра руководящих товарищей содержали положительную оценку деятельности губкома и были одобрены ЦК и ЦКК. Они осуждали поведение сторонников бывшей так называемой «рабочей оппозиции» — ставленников Шляпникова, — фактически продолжавших фракционную борьбу.

Можно только удивляться выдержке и долготерпению Ленина, который ни одну жалобу не оставлял без ответа, хотя, может, чувствовал, что они необоснованны. Ленин считал необходимым объективно проверять все жалобы с помощью безупречных деятелей партии.

Следует отметить, что борьба за ликвидацию фракционной деятельности в соответствии с решением X съезда шла не только в Нижегородской губернии, где «рабочая оппозиция» накануне X съезда была сильна. В Пермской губернии парторганизацию крупнейшего завода возглавлял заядлый сторонник Шляпникова И. Мясников, который долго и упорно вел борьбу против партии, склоняя организацию на путь раскола, и выпустил антипартийную брошюру. Характерным лозунгом Мясникова было: свободу всем — от анархистов до монархистов включительно!

Ленин дорожил каждым коммунистом из рабочей среды (а Мясников был из старых рабочих), терпеливо объяснял и объяснял им политику партии, стремясь сделать все, чтобы вывести их на правильную дорогу, прибегая к уставным мерам принуждения.

Ленин написал специально письмо лично Мясникову, разъяснял, что в практике работы действительно много недостатков, которые надо устранять. Ленин призывал Мясникова переключить свои силы и опыт на устранение этих недостатков партийными методами и отказаться от совершенно неправильного, антисоветского лозунга. Мясников, который, кстати, тоже подписал «заявление 22-х» сторонников Шляпникова в Коминтерн, не поддался уговорам Ленина, прямо встал на путь раскола, за что был исключен из партии. А потом, через некоторое время, сбежал из России за границу.

\subsection*{\center * * *}

На XI съезде я выступал в прениях по докладу Зиновьева об укреплении партии. Касаясь мер, которые он предлагал для укрепления партии, я сказал, что предложенная им формулировка «была расплывчатая, неточная, а между тем здесь требуется уточнение», и предложил сделать это на заседании секции.

Далее я говорил: «Тов. Зиновьев преувеличивал, когда говорил, что почти что во всех губерниях склока растет. Получилась картина, как будто положение в партии ухудшается. Но я знаю, что, начиная с прошлого года партия на местах растет и оздоровляется. Ряд губерний, которые я знаю, в прошлом году были раздираемы склоками, группировками. Они сейчас оздоровели. Таковы: Нижегородская, Владимирская, Тульская, Самарская губернии».

Заметив, что докладчик правильно отметил наличие уклона «рабочей оппозиции», я сказал, что «нужно отметить еще другой уклон — это уклон в рядах партийной интеллигенции и советских работников, оторвавшихся от масс. Они потеряли сейчас перспективу, уверенность в силах рабочего класса, они на каждом собрании ноют и плачут. Это один из признаков неоформленного ликвидаторства...»[254] Оно «опасно, поскольку разлагает психологию ответственных работников, так как внедряет в их сознание уныние»[255].

Далее я говорил о необходимости борьбы с уклоном к национализму.

Касаясь размежевания работы Советов и партии, я сказал: «То, что говорил об этом т. Ленин, совершенно правильно. Правильно и постановление VIII съезда партии. Но не всюду это правильно проводится. По моему мнению, т. Троцкий и т. Зиновьев переборщили, говоря об отстранении партии от советской работы.

Я не удивляюсь, когда т. Троцкий это говорил. Он не имеет представления о губкомах, он военный человек. Но Зиновьев работает в Петроградском губкоме и думает, что все остальные губкомы плохи. Он говорит, что вопросы о посевной кампании губкомами в полторы минуты решаются.

В губкомах никогда так вопросов не решают»[256].

Затем я критиковал предложение докладчика о том, что хороших хозяйственников следует переводить из губкомов на хозяйственную работу. Это приведет к тому, говорил я, что «потом будут говорить губкому, чтобы он не вмешивался в хозяйственную работу, потому что он в ней не компетентен. Это приведет к фактическому отстранению партии от советской работы»[257].

Я заявлял, что губкомы, как органы партии рабочего класса, не могут не руководить хозяйственной работой. «Конечно, мы не должны вмешиваться в администрирование, в мелочи, но политику хозяйственной работы мы должны предопределять. Мы наглядно должны знать, как осуществляется новая экономическая политика»[258].

Я предлагал «предупредить от другого переборщения те губкомы, которые, увлекаясь советской работой, мелочной опекой над ней, очень часто отстраняются от партийной работы... то агитационная, пропагандистская и иная партийная работа слабеет». Далее я выражал неудовлетворение тем, что докладчик «недостаточно остановился на положении и работе фабрично-заводских ячеек при новой экономической политике»[259]. Коснулся также вопроса работы среди молодежи, в особенности в деревне, а также культмассовой работы среди крестьян, и упрекнул докладчика и том, что и эти вопросы, несмотря на их актуальность, недостаточно освещены.

\subsection*{\center * * *}

Работа съезда подходила к концу. Шло обсуждение и голосование проектов резолюций и постановлений. Вечером 2 апреля во время обсуждения резолюции «О печати и пропаганде», о которой доложил от Комиссии съезда работник аппарата ЦК Яковлев, выступил с места пресловутый Рязанов. Он стоял около президиума и внес предложение запретить печатание платных рекламных объявлений в партийной печати и сослался при этом на декрет от ноября 1917 года о введении государственной монополии на объявления, предусматривающей сосредоточение публикации всех платных объявлений только в изданиях Советского правительства в Петрограде и в изданиях местных Советов.

Рязанов настаивал на том, что и в условиях новой экономической политики не надо отступать от декрета и запретить печатать платные объявления во всей партийной печати, в том числе и в «Правде». По своей натуре Рязанов не мог обойтись без демагогии, он утверждал, что чтение объявлений Мосторга в «Правде» и в «Московском рабочем», например об открытии какого-то ресторана с великолепными закусками или об устройстве ночного клуба, подействует на рабочих провокационным образом. Рязанов предлагал избавить партийную печать от этой пакости. Когда, дескать, у партийной прессы будет отнят этот источник дохода, парторганизации подтянутся и, опираясь на массы, организуют широкое распространение газет, как это мы делали раньше.

Такая постановка вопроса оказалась для всех неожиданной, делегаты не были подготовлены, мы сразу не успели продумать и понять истинное существо вопроса. Кроме того, все были уже утомлены от напряженной работы в течение семи дней.

Выступление Рязанова произвело впечатление не только на делегатов с мест, но, судя по выражению их лиц, и на членов президиума. Никто из них не выступил против, не подал реплики. Создалось впечатление, что все согласны с Рязановым.

Один Яковлев выступил против предложения Рязанова. Он сказал, что на местах пресса партийно-советская, но это и без него было известно делегатам. Исключительно партийных органов печати мало. Касаясь существа поправки Рязанова, он отметил, что мы сможем ее принять, когда придет для этого время. Возвращение же теперь к старому декрету нанесет удар целому ряду местных газет.

Выступление Яковлева убедило меня только в том, что, запретив печатание объявлений, мы нанесем ущерб местным газетам. Но Яковлев ни слова не сказал о том, что запрет может вызвать финансовые затруднения и у «Правды». Я понял так, что эта большая газета полностью обеспечивает покрытие своих расходов выручкой от продажи тиража.

Поднявшись, я сказал с места: следует отклонить предложение Рязанова о запрещении печатать объявления во всей партийной прессе, ограничившись запретом только «Правды». Никто ничего другого не предложил и даже не попросил слова по этому вопросу. Председательствующий поставил мое предложение на голосование и оно было принято.

Я сидел недалеко от Марии Ильиничны Ульяновой, работавшей секретарем редакции «Правды», и видел, с каким напряжением следила она за ходом обсуждения вопроса. В силу присущей ей скромности и стеснительности, она обычно в таких местах не выступала. После голосования вопроса я заметил, как она была возбуждена. Она тихо встала и вышла ил зала через ближайшую дверь. Работа съезда продолжалась. Мне и в голову не приходило, что могло ее так взволновать и заставить внезапно выйти из зала заседания.

Только позже мне все стало ясно. Мария Ильинична знала точку зрения Ленина по вопросу о платных объявлениях и видела, что принятое делегатами решение ей противоречит. Болея за интересы «Правды», она пошла к Ленину. К самому концу работы съезда дверь за президиумом бесшумно открылась, вошел Ленин и сел на свободный стул у стола президиума. Нам показалось, что он пришел для выступления с заключительной речью. Ленин написан небольшую записку и передал ее председательствующему. Записка сохранилась:

«т. Каменев!

Говорят, съезд провел отмену объявлений в Правде?

Нельзя ли исправить, ибо это ошибка явная?»[260]

Записку можно понять в том смысле, чтобы инициативу постановки вопроса об отмене принятого решения взял на себя президиум съезда, Ленин хотел избежать необходимости оказывать личное давление на делегатов.

Но Каменев не обратился к съезду с таким предложением.

Тогда Ленин взял слово и выступил «к порядку». Он просил у съезда в отступление от существующего правила, согласно которому после принятия решения всякое вмешательство в вопрос не практикуется, дать ему 4–5 минут, чтобы высказаться против решения, которое было ошибочно принято. Когда Ленин сказал, что принятое решение защищает Рязанов, то Рязанов, прервав его, громко крикнул с места: это не верно.

Наглый отказ Рязанова от ответственности за это решение вызвал возмущение делегатов. Поднялся шум, начались выкрики, что он говорит неправду. В ответ Рязанов подал новую реплику: за принятие съездом решения отвечают делегаты съезда. Принятое решение было внесено делегатом с места. Продолжая речь, Ленин сказал примерно так: ну и прекрасно, что хоть одно несуразное решение обошлось без Рязанова. Принятое решение ошибочно и выглядит смешно.

Ленин говорил, что мы допустили частную торговлю, сами начали торговать пока очень плохо, учимся у купцов, как это получше делать. И теперь, как бы стесняясь происшедшего, вы запрещаете «Правде» публиковать коммерческие объявления, хотя бы и о ресторанах, принимая наивный вид, как будто вы чистенькие.

А вы подумали, спрашивал Ленин, откуда покрывать расходы на издание газеты «Правда», если лишить ее объявлений? Вот мы сейчас начинаем переговоры в Генуе, рассчитывая получить от капиталистов золото. Неужели вы думаете, что часть этого золота, если мы его получим, можно будет передать «Правде» на покрытие ее расходов? Или думаете, что можно взять для этого средства из государственного золотого фонда? Это невозможно. Может быть, вы думаете получить средства за счет принудительного сбора налогов с крестьян? Это чертовски трудно и неприемлемо.

Остается одно — разрешить «Правде» печатать платные объявления

Заканчивая речь, Ленин предложил отменить принятое решение о лишении «Правды» объявлений. Как практический результат это даст ей равные с «Известиями» материальные возможности.

Выступление Ленина вызвало у меня чувство недовольства собой. Зачем мне надо было наспех выступать, подходя к вопросу узко практически? Дело оказалось куда серьезнее.

Делегаты с огромным вниманием слушали Ленина. Было ясно, что он всех нас убедил в ошибочности принятого решения.

Тут Рязанов вновь сорвался с места и громко крикнул: пусть тот, кто внес принятое предложение, соберется с мужеством и назовет свое имя!

Меня возмутил тон этого предложения и удивило, — почему Рязанов так ставит вопрос, отлично зная, как и все делегаты съезда, что поправка к его предложению была внесена мной.

Как и еще несколько товарищей, я сидел около трибуны на ступеньках, ведущих из зала в президиум. Оттуда был виден и весь зал, и президиум, а главное — ни одно слово выступавших не могло проскочить мимо наших ушей.

Возмущенный, я встал и с места заявил: Рязанов предложил лишить все наши партийные газеты, включая «Правду», права печатать объявления. Я же внес поправку не принимать предложения Рязанова в отношении всей партийной прессы, а принять решение только в отношении «Правды».

Говоря, я смотрел прямо в глаза Ленину. Мои слова для делегатов не были новостью. Ленин слушал меня внимательно, и у него появилась чуть заметная добродушная улыбка. Увидев ее, я сел и успокоился.

Согласовав с Лениным, Каменев не поставил на голосование предложение об отмене принятого решении, а внес предложение: вопрос об объявлениях, как практический, передать на решение ЦК.

Затем он сказал, что прежде чем поставить это предложение на голосование, он даст слово тому, кто хочет защищать старое постановление. Никто из делегатов не пожелал выступить в защиту принятого решения об объявлениях, кроме Рязанова. Выступив развязно, он, по существу, повторил свою демагогическую речь, произнесенную им при защите своего предложения.

Сразу после его выступления председательствующий поставил на голосование предложение о передаче вопроса о печатании объявлений в «Правде» на решение ЦК партии, что и было принято.

\subsection*{\center * * *}

В речи при закрытии XI съезда партии Ленин отметил, что при сравнении этого съезда с предыдущими бросается в глаза «большая сплоченность, большее единодушие, большее организационное единство»[261]. По вопросам о профсоюзах, о нэпе разногласий в нашей партии не оказалось в сколько-нибудь заметном размере

Съезд еще раз показал неправоту наших врагов, утверждавших, будто наша партия впадает в старчество, теряет гибкость ума, гибкость всего своего организма. Когда нужно было, мы наступали на врага с беззаветной смелостью, подняв нашу революцию на невиданную в мире высоту. Мы доказали, что отступали за истекший год в общем и целом в революционном порядке.

Пролетариату развитых стран мира «не удастся решить своей задачи без того, чтобы соединить уменье беззаветно бороться и наступать с уменьем отступать в революционном порядке...

Теперь мы постановили признать отступление законченным.

Это значит, что по-новому ставится вся задача нашей политики»[262].

В заключение Ленин заявил, что «весь гвоздь теперь в том, чтобы авангард не побоялся поработать над самим собой, переделать самого себя, признать открыто свою недостаточную подготовленность, недостаточное уменье. Весь гвоздь в том, чтобы двигаться теперь вперед несравненно более широкой и мощной массой, не иначе как вместе с крестьянством, доказывая ему делом, практикой, опытом, что мы учимся и научимся ему помогать, его вести вперед»[263].

\subsection*{\center * * *}

XI съезд был последним съездом партии, в котором лично участвовал Ленин. Но он незримо присутствовал на XII съезде, когда обсуждались его последние статьи и письма о реорганизации Рабкрина, о кооперировании сельского хозяйства, о национальном и культурном строительстве. Ленин был с нами и на XIII съезде, когда обсуждалось его «Письмо к съезду».

Ленин жил и живет с партией — всегда и во всем. И как не было Ленина без партии, так нет и никогда не будет партии без Ленина.

\newpage
\section[Ленин на пленуме ЦК партии в мае 1922 года]{\center \textbf{Ленин на пленуме ЦК партии в мае 1922 года}}


Не может изгладиться из памяти заседание Пленума Центрального Комитета партии в середине мая 1922 года, на котором председательствовал Ленин. Я участвовал в этом пленуме, поскольку был избран на XI съезде кандидатом и члены ЦК.

Пленум проходил в зале заседаний Совнаркома. Члены и кандидаты в члены ЦК сидели за длинным столом, а Ленин занимал председательское место и был виден всем. Ленин в руках держал свои карманные часы и строго следил за соблюдением регламента выступавшими. Обстановка была строго деловая, никаких посторонних разговоров между членами и кандидатами в члены ЦК, которые могли бы помешать ведению собрания, не было. Докладчику Ленин давал время, как помнится, 3 минуты, в особых случаях — 7 минут, как, например, Рудзутаку, который докладывал о Генуэзской конференции, выступавшим в прениях — 1–2 минуты. В большинстве своем докладчики были опытные и вполне укладывались в регламент и научились короткими фразами выражать суть вопроса и понимать друг друга без длительных объяснений.

Орджоникидзе и Киров на этом пленуме отсутствовали. Обстоятельства работы на Кавказе не позволили в тот момент отлучиться. Отсутствовали и некоторые другие члены ЦК и кандидаты в члены ЦК. Потому на пленуме было всего 18 членов ЦК из 27 и 4 кандидата в члены ЦК из 19.

Первым вопросом обсуждался вопрос о Генуэзской конференции. Доклад делал Рудзутак, который входил в состав делегации на Генуэзскую конференцию. Он приехал с отчетом ЦК. Рудзутак критически отнесся к некоторым сторонам работы делегации, считая, что Чичерин и Красин идут в переговорах на большие уступки, чем это следует. Критиковал их за это. Потом изложил вновь поставленные вопросы для дальнейших переговоров перед делегацией в Генуе.

Видно было, что Ленин был полностью в курсе дела, поскольку практическое руководство делегацией проводил он (подготовка директив и пр.). Он не стал подчеркивать отдельные ошибки из того, что говорил Рудзутак, указав, что до настоящего времени генуэзская делегация правильно выполняла свою задачу. Он предложил дать Чичерину директиву по вопросам, которые были поставлены перед ЦК в докладе Рудзутака. Во-первых, считать недопустимым перенесение работ комиссии из предполагаемой Гааги в Россию. Если это трудно отклонить, то согласиться на устройство этой работы в Сочи или в Ливадии (Крым). Во-вторых, всемерно протестовать против запрещения сепаратных договоров отдельных держав с Советской страной, добавляя при этом, что мы не подчинимся этому ни в коем случае.

Вторым вопросом, по которому докладывали Бухарин и Зиновьев, был проект директив представителям нашей партии в Коминтерне, в связи с заседаниями представителей трех Интернационалов (II, II1/2 и III Интернационалы), создавшими «комиссию девятки».

В то время широко обсуждалась на страницах зарубежной рабочей печати идея создания единого рабочего фронта, выдвинутая Коминтерном, и в связи с этим предложение о подготовке Всемирного рабочего конгресса, который должен «организовать оборонительную борьбу рабочего класса против международного капитала»[264].

Зиновьев прочитал короткий проект директив, видимо заранее согласовав его с Лениным. Суть директив сводилась к тому, что делегация Коминтерна должна в переговорах на заседаниях «девятки» требовать быстрейшего созыва Всемирного рабочего съезда — ультимативно, в определенный короткий срок, который должен быть сразу же определен по согласованию сторон. В случае саботажа со стороны II Интернационала Коминтерн немедленно отзовет своих представителей из «девятки» и продолжит агитацию за единый пролетарский фронт.

В качестве уступки представителям II Интернационала допускалось, что представитель РКП заявит о готовности вычеркнуть из общей платформы обращение о защите Советской России, сосредоточит всю свою борьбу на защите 8-часового рабочего дня, против безработицы и т. п. В то же время, если представители двух Интернационалов изменят мнение и изъявят готовность не на словах, а на деле созвать Всемирный рабочий конгресс, то мы готовы будем снова послать наших делегатов в «девятку». Затем поручалось представителям нашей партии в Исполкоме Коминтерна обеспечить издание брошюры об опыте большевистской партии для иллюстрации тактики единого фронта. Это было принято без дискуссии.

После пленума предстояла сессии ВЦИК, поэтому на пленуме обсуждались некоторые вопросы ее повестки. Первым вопросом предусматривалось созвать совещание членов ВЦИК для рассмотрения предложений о районировании России и представить доклад в Политбюро. Политбюро обсудит состав комиссии по районированию. Следует ли ставить вопрос на сессии ВЦИК о Красной Армии и в каком объеме, было поручено решить Политбюро.

Почему-то в памяти у меня запечатлелось обсуждение вопроса о кампании по извлечению ценностей. Докладчиком по этому вопросу выступал Троцкий. Речь шла о том, чтобы извлекать из государственных учреждений, музеев разные ценности, не представлявшие исторического или научного интереса, но могущие быть проданными за валюту за границу. Я не знал, почему Троцкий этим вопросом занимался и был докладчиком по нему. Бесспорным для всех было предложение о том, чтобы признать необходимым максимально ускорить реализацию ценностей, уже собранных. Для осуществления этой директивы было решено образовать комиссию, предложив комиссии представлять в ЦК раза в месяц информацию о ходе реализации ценностей.

Спор вызвало одно предложение Троцкого, которое заключалось в том, чтобы 5 % от вырученной суммы передавать в распоряжение Реввоенсовета республики.

Сталин, который сидел рядом с Троцким, выступил только по этому вопросу, как обычно, тихо, спокойно говорил. Он высказался за отклонение этого предложения Троцкого о 5-процентном отчислении.

Троцкий, который был вспыльчив, вскипятился, стал горячо спорить со Сталиным и доказывать правильность своего предложения. Сталин говорил не больше одной минуты. Троцкий же — около 2–3 минут. Тогда Ленин, показав часы, сказал: «Предлагаю соблюдать регламент». Троцкий подчинился. Сталин вновь попросил слова и, заняв не более минуты, привел аргументы в пользу предложения о том, какую сумму из выручки давать Красной Армии. Он сказал, что не надо определять твердый процент на все время, а надо решать этот вопрос по мере необходимости, смотря на какие цели это требуется. Вся же сумма выручки должна поступать в распоряжение правительства.

Троцкий, еще более разгорячившись, вновь взял было слово. После нескольких фраз, им сказанных, было видно, что спор еще более разгорается и занимает слишком долго время пленума. Видимо, по существу Сталин был более прав, да и по настроению членов ЦК чувствовалось, что предложение Сталина было более правильное, а предложение Троцкого — «ведомственное», неправильное. Ленин поручил рассмотреть этот вопрос в двухдневный срок комиссии и внести предложение в Политбюро, а в случае разногласий в Политбюро обсудить этот вопрос на пленуме ЦК.

Это поразительное свойство Ленина — не допускать возникновения излишних разногласий, бесплодных прений и столкновения мнений в вопросах, которые в спокойной, деловой обстановке можно решить без какого-либо ущерба, — наиболее ярко выразилось в решении этого вопроса.

Предложение Ленина не означало принятия предложения ни Троцкого, ни Сталина, так что это не могло задеть их самолюбия. С другой же стороны, дело передавалось в руки комиссии, которая в деловом порядке подготовила бы правильное предложение.

Обсуждался также вопрос о так называемом «Московском комбинированном кусте», поставленный опять Троцким. Речь шла о некоторых сторонах деятельности объединения нескольких совхозов и предприятий, организованных военными, над которым осуществлял шефство Троцкий. Он оспаривал какие-то решения советских органов и выдвигал какие-то требования, которые не нужно было ставить на пленуме ЦК. Было решено рассмотреть это дело в советском порядке, не связывая советские органы ранее принятыми решениями ЦК. Поручено было Секретариату послать документы по этому вопросу членам ЦК, обязав последних ознакомиться со всем делом. Это и было, по существу, отклонением предложения Троцкого.

Мне запомнился и такой интересный случай. Я сидел рядом с Дзержинским. Дзержинский был заядлым курильщиком. Недалеко от того места, где мы сидели, была голландская печь. Дзержинскому хотелось слушать выступления и не покидать заседания, но он не мог выдержать без курения. И он, как это обычно бывает с ребятами, решил перехитрить Ленина, подошел к голландской печи так, чтобы Ленин не видел, покурил в отдушину и вернулся на место, довольный тем, что Ленин не мог видеть, что он курил, да и остальные не обратили на это внимания.

Несмотря на трения, которые имели место при обсуждении вопросов, Ленин ни разу голоса не повышал, не делал резких замечаний, избегал даже иногда высказываться по существу, чтобы создать более спокойную обстановку, деловую и серьезную. Это, конечно, на меня, как на новичка в ЦК, производило внушительное впечатление.

И тогда, и позже я убедился, насколько важно для политического руководителя ценить время, лишнего не говорить, уметь коротко изложить суть дела. И если вопрос нельзя решить на данном заседании, не допуская пустопорожних прений, он передавался в деловую комиссию, назначался короткий срок для работы и представления проектов решений. Только этим и можно, пожалуй, объяснить, как много вопросов самых разных в течение дня Ленин рассматривал, разбирался в них. И теперь, когда читаешь протоколы заседаний, записки и другие материалы Ленина, поражаешься огромному объему работ, многообразию и глубокому содержанию вопросов, рассматриваемых и разрешаемых им, будь то в плоскости теоретических, политических вопросов, будь то в плоскости практических дел, даже самых узких, не имеющих большого значения.

\newpage
\section[Ленин на IV конгрессе Коминтерна]{\center \textbf{Ленин на IV конгрессе Коминтерна}}


Я участвовал и в IV конгрессе Коминтерна в составе русской делегации теперь уже не от нижегородской парторганизации, а от северокавказской, где я работал секретарем Юго-Восточного бюро ЦК партии (в Ростове-на-Дону).

Конгресс открылся и провел свое первое заседание в Петрограде 5 ноября 1922 года, накануне пятой годовщины Советской власти. Это было торжественное заседание в цитадели Октябрьской революции — в знак глубокого уважения к, питерским пролетариям, давшим прочную победу первой в мире социалистической революции.

На этом заседании были приняты обращения к политзаключенным в тюрьмах, к рабочим и крестьянам Италии, к трудящимся России, к Красной Армии и Флоту, к рабочим Петрограда, Заседание завершилось заслушиванием речей представителей иностранных делегаций. Делегации возложили венки на Марсовом поле, воздав должное памяти павших героев революции. Наша, русская делегация, также возложила венок.

Питерский Совет специально просил Ленина приехать на эти торжества празднования пятилетнего юбилея победы Октябрьской революции. Ленин не смог приехать и прислал теплую приветственную телеграмму конгрессу Коминтерна и Петроградскому Совету, в которой писал: «Крайне сожалею, что не смогу быть на первом заседании конгресса и должен ограничиться письменным приветствием.

Несмотря на гигантские трудности, стоящие на пути коммунистических партий, Коминтерн растет и крепнет. Главная задача по-прежнему состоит в завоевании большинства рабочих. И эту задачу мы, несмотря ни на что, выполним»[265].

Затем Ленин, передавая петроградским рабочим лучшие пожелания и горячий привет, писал:

«Питерские рабочие должны быть и на хозяйственном фронте одними из первых. С радостью слышим мы о начале хозяйственного возрождения Петрограда. На ваше приглашение побывать в Петрограде надеюсь ответить приездом в скором будущем».

Отмечая праздник первого пятилетия Советской власти в России, которая более прочна, чем когда бы то ни было, Ленин закончил следующим:

«Советская Россия считает величайшей гордостью помогать рабочим всего мира и их трудной за низвержение капитализма. Победа будет за нами.

Да здравствует Коммунистический Интернационал!»[266]

Тогда мне представился счастливый случай впервые побывать в городе — колыбели пролетарской революции. Я находился под огромным впечатлением невиданной ранее мною красоты города: прямые, казалось, бесконечно длинные проспекты, красивые стройные многоэтажные дома, монументальные сооружения дворцов, соборов, памятников, многочисленные мосты через Неву и каналы, прорезавшие город.

Красота Петрограда настолько пленила, что я не замечал запущенности домов, улиц, заколоченных магазинов. Не замечалось и то, что народ, заполнявший улицы в эти праздничные дни, был одет плохо, имел вид исхудалый, утомленный. Бросались в глаза — революционный подъем, бодрость духа, блеск в глазах, отражение на лицах петроградцев чувства гордости от того, что они зачинатели и победители первой пролетарской революции; радость того, что многочисленные иностранные делегации коммунистов из многих стран мира прибыли к нам и на своем первом заседании приняли теплое обращение к питерским рабочим в дни праздника пятилетней годовщины Великого Октября. По возвращении делегатов конгресса в Москву мы участвовали 7 ноября в Большом театре в совместном торжественном заседании Моссовета, фабрично-заводских организаций и IV конгресса Коминтерна, посвященном пятилетию Октябрьской революции.

Деловые заседания конгресса начались в Кремле, в Андреевском зале (теперь он реконструирован, увеличен и является залом заседаний Верховного Совета СССР).

Ленин уже был больной, но он продолжал работать. Выступая 13 ноября на этом конгрессе, он начал с заявления, чтобы делегаты не ждали от него главного доклада повестки дня по теме «Пять лет российской революции и перспективы мировой революции». Он сказал, что по болезни не может сделать такого большого доклада и поэтому из первого раздела доклада о пяти годах российской революции он берет только тему о нэпе, как о важнейшем вопросе.

Нужно сказать, что у многих иностранных коммунистов не было правильного понимания сущности новой экономической политики. В то время иностранная буржуазия, меньшевики и эсеры представляли дело так, что большевики возвращаются к капитализму, что политика социализма провалилась и вводится вновь капитализм, т. е. то, чего им хотелось бы. Словом, рисовали, что произошла катастрофа для коммунистов. А многие иностранные коммунисты еще плохо понимали нэп и поддавались этой буржуазной пропаганде. В своей речи Ленин сослался на то, что еще в 1918 году, до развертывания гражданской войны, он говорил о пяти экономических укладах в нашей стране, о том, что главную роль играет мелкотоварный сектор, который непрерывно порождает капитализм. Он напомнил, что держался в 1918 году того мнения, что по отношению к тогдашнему хозяйственному состоянию Советской республики государственный капитализм представлял собой шаг вперед. Если бы тогда у нас установился госкапитализм, это было бы громадным успехом и вернейшей гарантией, что у нас упрочится и победит социализм. Он отметил, что в отношении социалистической страны может казаться странным утверждение, что несоциалистический элемент — госкапитализм расценивается выше, чем социализм, но эта форма была бы тогда более благоприятной. Ленин сказал, что он особенно подчеркивает это сейчас в своем выступлении, потому что только исходя из этого можно объяснить, что представляет собой новая экономическая политика. И он объяснял, что мы сейчас отступаем, но и вроде наступаем. Это наше отступление потому, что мы в своем экономическом наступлении слишком далеко продвинулись вперед. Но мы так далеко зашли, что у нас было куда отступать. Причем позиции, на которые мы отступали, намечались еще в 1918 году. И теперь, по истечении полутора лет, мы уже в состоянии подвести итоги.

Ленин не случайно сделал вопрос о нэпе основным содержанием доклада на конгрессе. В то время это был не только решающий вопрос русской революции, но и один ми важнейших вопросов будущих пролетарских революций и других странах. Важнейшей стороной новой экономической политики было то, что она являлась правильным путем обеспечения перехода к социализму.

Ленин уделял особое внимание доказательству того, что нэп с самого начала был необходимым шагом в направлении к социализму и что вызванный интервенцией и гражданской войной военный коммунизм был вынужденным явлением.

Ленин и в этой речи и во всех других никогда не замалчивал трудностей.

Как я уже говорил, огромное впечатление производила большая откровенность Ленина. Он рассказал о первых успехах новой экономической политики, объяснил, что мы уже хлеб имеем, пайки мы обеспечим, дела у крестьян поправились, легкая промышленность стала восстанавливаться. Но положение еще очень грудное. Для развития тяжелой промышленности нужны большие капиталы, которые нам никто не дает. А без тяжелой промышленности мы пропадем, если ее не поднимем. С особенной гордостью и радостью Ленин сообщил, что удалось накопить 20 с лишним млн. рублей золотом, которые будут использованы для поднятия тяжелой промышленности.

Далее он откровенно стал говорить о том, что мы наделали не мало глупостей и еще наделаем их впереди. Он добавил, что «никто не может судить об этом лучше и видеть это нагляднее, чем я». И он объяснил, что мы их наделали и делаем потому, что мы отсталая страна, наш народ не просвещенный, госаппарат засорен враждебными элементами и что нам не помогает ни одно государство. Конечно, в таких условиях естественны ошибки.

Он смело критиковал ошибки, совершенные партией, Советской властью, не ради критики вообще, а для того, чтобы вскрыть их и исправить, чтобы на ошибках учиться, чтобы их не повторять, чтобы их было меньше, чтобы правильно наметить задачи дальнейшего пути. Он говорил, что не ошибается тот, кто ничего не делает. Смелость Ленина, вскрывавшего наши ошибки и трудности, была неповторима.

Но, говорил Ленин, империалисты прочитают мое выступление и скажут, что вот Ленин сам признал, что большевики делают глупости. Но пусть они знают, говорил Ленин далее, что у них глупости невероятно большие и совсем другого рода.

И он привел такое сравнение, воспользовавшись тем местом из романа Тургенева «Рудин», в котором женоненавистник Пигасов, отрицавший способность женщин к логическому мышлению, объяснял, какая разница между ошибкой мужчины и женщины, говорил: мужчина может, например, сказать — дважды два не четыре, а пять или три с половиной, а женщина скажет, что дважды два — стеариновая свечка. Ленин сказал: я приведу для сравнения слова одного знаменитого русского писателя в несколько измененном виде: когда мужчина ошибется в счете, женщина скажет явную несуразность. Вот, заметил Ленин, на такой манер и можно сравнить наши ошибки, совершаемые большевиками, с ошибками, совершаемыми империалистами! Когда большевик делает глупости, то он говорит: «Дважды два — пять», а если его противники, т. е. капиталисты и «герои» II Интернационала, делают глупости, то у них выходит: «Дважды два — стеариновая свечка», т. е. полный абсурд. Вы посмотрите, говорил Владимир Ильич, что они делали? Они, хорошо не подумав, обещали Колчаку помощь, признали его даже верховным правителем России, а что из этого получилось? Полнейшее фиаско.

Ленин привел еще пример с Версальским миром. Слушая его, можно было добавить о том, о чем он уже упоминал в других выступлениях: и то, как империалисты со своими войсками вторглись в Россию и, не добившись ничего, ушли, оставшись «с носом», и о других больших глупостях, совершенных империалистами.

А об иностранных коммунистических партиях на IV конгрессе он говорил, что главная задача иностранных товарищей — учиться строить настоящую революционную партию, партию нового типа: постигнуть организацию, построение и метод революционной работы, приспосабливать опыт русской партии к собственным условиям, чтобы это была партия революционная и могла оказать влияние на беспартийные массы, вести их за собою в борьбе за пролетарскую революцию.

И такие партии были созданы, и в создании их принимал участие Ленин. Ленин участвовал в подготовке и в проведении четырех конгрессов Коминтерна. Здесь он проявил свои гениальные способности как организатор и вождь мирового коммунистического движения. На IV конгрессе была принята резолюция по русскому вопросу, в которой конгресс призвал трудящихся всех стран оказать всяческую поддержку Советской России, как авангарду мирового пролетариата.

Трудно было предположить, что речь Ленина на IV конгрессе станет предпоследним его публичным выступлением. Через семь дней после этого он выступил и Московском Совете с речью, в которой сказал широко известную теперь фразу: «...Владивосток далеко, но ведь это город-то нашенский»[267]. Это свое выступление Ленин закончил облетевшими весь мир пророческими словами: «Социализм уже теперь не есть вопрос отдаленного будущего... Мы социализм протащили в повседневную жизнь... Из России нэповской будет Россия социалистическая»[268]. Это было последнее публичное выступление Ленина. Но он продолжал вести разработку политики партии по социалистическому строительству, оставив в наследие написанные им в этот период важнейшие труды о кооперативном плане, о реорганизации Рабкрина в целях улучшения работы госаппарата, об укреплении единства партии, о союзе рабочего класса с крестьянством, о величайшей экономии в хозяйстве, чтобы малейшие сбережения сохранялись для развития крупной машинной индустрии.

Заветы Ленина в этих и в других его трудах служат путеводной звездой для нашей партии, для советского народа, для братских социалистических стран, для коммунистов и революционеров всего мира.

Ярким подтверждением этого является состоявшееся в Москве в июне 1969 года международное Совещание коммунистических и рабочих партий, в котором участвовали представители 75 стран. В их числе и молодые коммунистические и рабочие партии, которые под влиянием старых коммунистических партий, созданных при Ленине, возникали в освободившихся странах. И разве это Совещание не является триумфом ленинских идей, разве это не есть самый лучший способ отметить юбилей Ленина? Отметить его труды, поблагодарить за его учение, которое освещает наш путь вперед. Совещание имеет великое историческое значение. Главное — это идея единства мирового коммунистического движения в борьбе с империализмом, за что Ленин всегда боролся.

Ход Совещания, сто решения, выступления на нем — это победа марксистско-ленинских идей, светом которых освещаются современные события мировой жизни и дальнейшая программа действий. В работе Совещания претворены ленинские идеи.

Такое же большое, поистине историческое значение имеет торжественное заседание в Кремле по случаю 100-летнего юбилея Ленина с участием виднейших деятелей мирового коммунистического, рабочего и национально-освободительного движения более чем из 100 стран мира.

\newpage
\section[Новая дискуссия, навязанная Троцким]{\center \textbf{Новая дискуссия, навязанная Троцким}}


Опасения, высказанные Лениным накануне XI партийного съезда, что Троцкий может попытаться вновь начать открытую борьбу против ЦК, нашли подтверждение менее чем через два года.

В 1923 году, когда Ленин был тяжело болен и лишен возможности принимать какое-либо участие в партийной жизни, Троцкий навязал партии новую дискуссию. Он начал атаковать партийное руководство и партийный аппарат, заигрывая с учащейся молодежью, объявив ее — ни больше, ни меньше — барометром партии. Он противопоставлял и натравливал эту молодежь на рабочую часть партии и на ее основные кадры. Он явно добивался слома ленинского партийного аппарата и устранения ленинских партийных кадров. Его конечная цель как и прежде состояла в захвате руководства партией и изменении ее политики.

Сперва в стенах ЦК, а затем в открытой, навязанной им дискуссии Троцкий добивался ревизии ленинской генеральной линии, вопреки Ленину отрицал возможность построения социалистического общества в СССР без победы мировой революции.

Находясь тогда в Ростове, я от приезжавших из Москвы товарищей узнал, что там в вузах и некоторых учреждениях идет горячая дискуссия, в ходе которой происходят резкие оголтелые нападки троцкистов на партию, на ее линию.

Никакой информации от ЦК партии о начавшейся дискуссии мы не имели. Из сообщений товарищей я не мог составить ясного представления, какие конкретные аргументы выдвигает оппозиция против ЦК партии и каково вообще идеологическое содержание ее новой платформы.

Вот почему я выехал в Москву. Там я провел первый же день на собраниях в МГУ. С утра и до вечера, с небольшим перерывом, там происходили очень шумные дискуссионные выступления. Сидел я почти в последних рядах, намерения выступать у меня не было — хотелось только побольше послушать и разобраться самому, о чем идет спор.

Выступали главным образом разные оппозиционные ораторы. Как помнится, особенно разнузданными были демагогические выступления Рязанова и Мдивани. Ораторы они были довольно искусные и сумели добиться шумного одобрения и аплодисментов присутствующих.

Сторонников линии ЦК среди выступавших было очень мало, и выступали они не на высоком уровне. Фактически толковой защиты линии партии не было, а нападки на нее носили резкий, явно антипартийный характер. Я был удручен атмосферой, царившей на этих собраниях.

На следующий день я решил выступать на дискуссионном собрании медицинского факультета МГУ. У меня была цель — ударить по оппозиционерам, а заодно проверить и себя — смогу ли я найти нужные аргументы, чтобы убедить аудиторию в правоте линии партии и отбить атаку на нее со стороны оппозиции.

Мне удалось убедить собравшихся в неправильности установок оппозиции, склонить большинство на сторону ЦК партии.

После этого собрания я зашел на квартиру к Сталину и рассказал ему обо всем, что видел и слышал в МГУ, а также о том, что, по полученным мною из бесед сведениям, во многих вузах и ряде учреждений троцкисты на собраниях одерживают верх.

В результате, с возмущением заявил ему, создается впечатление, что в столице нет Московского комитета партии и все пущено на самотек.

Я был возбужден и выразил неудовольствие и поведением ЦК, который, по-моему, тоже самоустранился от фактически уже начавшейся в столице дискуссии и тем облегчает возможность троцкистам запутать неопытных товарищей в надежде на легкую победу.

Я спросил Сталина, почему ЦК до сих пор молчит, когда собирается выступить и как.

Помню, что я был удивлен, с каким невозмутимым спокойствием меня выслушал Сталин. Он сказал, что для волнений нет особых оснований. Все члены Политбюро решили выступить единым фронтом прочий Троцкого и его сторонников и почти уже закончили подготовку проекта постановления Политбюро ЦК и Президиума ЦКК «О партстроительстве», в котором разоблачается антиленинская идеологическая позиции Троцкого и разъясняются задачи партии в области внутрипартийной политики в ленинском духе.

Опубликование этого постановления, сказал мне Сталин, нанесет сокрушительный удар троцкизму и объединит вокруг ЦК самые широкие партийные массы.

Успокоенный этим заявлением, я уехал в Ростов. Однако, не дожидаясь опубликования постановления ЦК и ЦКК, созвал узкий актив партработников, на котором рассказал, что происходит в Москве, и предложил им быть бдительными и готовыми к открытой дискуссии.

Вскоре такая дискуссия началась, причем выявилось значительное число сторонников Троцкого и у час — в Ростове, Краснодаре и Владикавказе.

Но все рабочие организации были на стороне ЦК. Поэтому ни в одном районе края троцкистам не удалось добиться победы. То же самое происходило и по всей стране. Троцкий был полностью разбит, как говорится, в открытом бою. И хотя эта дискуссии отняла у нас много сил и времени, она оказалась и полезной в том смысле, что троцкизм был окончательно разоблачен и осужден партийными организациями.

Ленинские идеи восторжествовали и на этот раз, доказав всю свою жизненность и глубину.

Троцкий был настолько морально изолирован партией, что, оставаясь еще тогда членом Политбюро, не отважился даже явиться на XIII партконференцию, состоявшуюся в середине января 1924 года, на которой были подведены итоги прошедшей дискуссии. Еще до начала работы этой конференции он уехал из Москвы в Сухуми «на отдых», — или, как тогда говорили, «ждать у моря погоды».

\newpage
\section[Перестало биться сердце Ильича...]{\center \textbf{Перестало биться сердце Ильича...}}


Партийные массы и весь народ очень волновало состояние здоровья Ленина. Выступая в декабре 1923 года на Донском областном съезде Советов в Ростове, в ответ на многочисленные вопросы делегатов я пространно информировал их о ходе болезни и о состоянии здоровья Ленина. Сохранилась стенографическая запись этого выступления.

Я сказал о том, что «со слов лечащих врачей болезнь Владимира Ильича очень тяжелая и одно время его положение было очень опасно», но все-таки врачи «высказались, что он поправится. За это время он очень быстро физически окреп и ходит. Приезжал даже в Москву, правда, против воли врачей. Заходил в свой кабинет, стал рыться в книгах, выбрал и взял некоторые с собой». Далее напомнил, что весной состояние Ленина «было очень тяжелое и бюллетени издавались каждый день, а сейчас бюллетеней нет, потому что его здоровье медленно восстанавливается». Упомянул и о том, что были случаи, когда неправильно информировали о здоровье Ленина — то очень пессимистически, то чрезмерно оптимистически, — что создавало впечатление «что есть какая-то неправда». Поэтому «ЦК постановил сообщать о его здоровье только правду». Я сообщил, что следить за здоровьем Ленина «приглашены лучшие врачи, и страна сделала все, что можно было». Свое выступление я закончил словами: «Мы надеемся, что Владимир Ильич преодолеет свою болезнь, и в ближайшие годы он будет в нашей среде и будет руководить нами так, как и раньше»[269].

Оптимизм, которым было пронизано это мое заявление, не был тогда проявлением какого-то моего сугубо личного, субъективного представления о действительном состоянии здоровья Владимира Ильича.

Точно такого же мнения придерживались в то время и лечащие Ленина врачи, а также члены Центрального Комитета партии.

В правительственном сообщении от 22 января 1924 года о смерти В. И. Ленина было сказано:

«Ничто не указывало на близость смертельного исхода. За последнее время в состоянии здоровья Владимира Ильича наступило значительное улучшение. Все заставляло думать, что его здоровье будет и дальше восстанавливаться. Совершенно неожиданно вчера в состоянии здоровья Владимира Ильича наступило резкое ухудшение. Несколько часов спустя Владимира Ильича не стало».

Этот трагический день — 21 января 1924 года — навсегда запечатлелся в нашей памяти как день величайшего, подлинно всенародного, а для каждого из нас — и личного горя.

...Только недавно закончилась XIII партийная конференция. Подходила к концу работа XI Всероссийского съезда Советов, вскоре после которого должно было состояться открытие II съезда Советов Союза CCP.

Помню, 21 января, во второй половине дня я зашел на квартиру к Сталину, чтобы посоветоваться с ним но ряду вопросов, связанных с нашими северокавказскими делами...

Не прошло и 30–40 минут нашей беседы, как вдруг, неожиданно, в комнату ворвался крайне взволнованный и возбужденный Бухарин и не сказал, а как-то выкрикнул, что из Горок позвонила Мария Ильинична и сказала: «Только что, в 6 часов 50 минут, скончался Ленин».

Это так неожиданно. Мы были потрясены. Наступило минутное замешательство и молчание.

Мы все мгновенно оделись, на аэросанях поехали в Горки.

...В небольшой комнате на кровати лежал спокойный, как бы только что уснувший, Ленин...

Вскоре стали подъезжать другие члены Центрального Комитета.

Трудно описать эти тягостные минуты, когда все мы осиротело столпились около Ильича...

Собравшиеся члены Политбюро условились срочно подготовить и опубликовать правительственное сообщение о смерти Ленина, а также принять специальное обращение Центрального Комитета партии.

Было решено созвать экстренный пленум ЦК партии.

...Подавленные великим горем, многие с невысохшими еще от слез глазами, собрались мы на этот пленум в Кремле ночью с 21 на 22 января.

Говорить было трудно, слова как-то «не находились».

Мы избрали комиссию по организации похорон Владимира Ильича под председательством Дзержинского.

Приняли обращение «К партии. Ко всем трудящимся».

В обращении отмечались исторические заслуги Ленина перед нашей партией и страной, перед всем прогрессивным человечеством.

«Никогда еще после Маркса, — говорилось в обращении, — история великого освободительного движения пролетариата не выдвигала такой гигантской фигуры, как наш покойный вождь, учитель, друг. Все, что есть в пролетариате поистине великого и героического — бесстрашный ум, железная, несгибаемая, упорная, все преодолевающая воля, священная ненависть, ненависть до смерти к рабству и угнетению, революционная страсть, которая двигает горами, безграничная вера в творческие силы масс, громадный организационный гений, — все это нашло свое великолепное воплощение в Ленине, имя которого стало символом нового мира от запада до востока, от юга до севера...

Но его физическая смерть не есть смерть его дела. Ленин живет в душе каждого члена нашей партии. Каждый член нашей партии есть частичка Ленина.

Вся наша коммунистическая семья есть коллективное воплощение Ленина... Смерть нашего учителя — этот тяжелый удар сплотит еще сильнее лягни ряды. Дружной боевой цепью идем мы в поход против капитала, и никакие силы в мире не помешают нашей окончательной победе.

Эта победа будет самым лучшим памятником товарищу Ленину, тому, которого, как лучшего друга, массы звали своим «Ильичем»».

23 января члены ЦК, наркомы и представители рабочих делегаций несли на руках гроб с телом Ленина из Горок до ближайшей станции Герасимово, где ожидал специальный поезд.

В час дня траурный поезд прибыл на Павелецкий вокзал Москвы.

По улицам, сплошь заполненным народом, гроб с телом Ленина был перевезен и установлен в Колон-ион зале Дома Союзов.

Пять дней и ночей, бесконечным потоком, в торжественном молчании, нарушаемом лишь сдержанными рыданиями, шли мимо гроба рабочие и крестьяне, воины Красной Армии, интеллигенция, молодежь, делегации трудящихся зарубежных стран, люди самых разных национальностей, — шли, чтобы проститься со своим великим вождем, учителем и другом, шли, чтобы в эти скорбные минуты прощания запечатлеть навсегда в своей памяти дорогие черты Ильича.

Я хорошо помню эти суровые январские дни 1924 года.

Мороз доходил до 30 градусов, а люди все шли и шли... Часами они стояли на морозе, грелись у костров, горящих на улицах, чтобы потом попасть на две-три минуты к своему Ильичу и сказать ему свое последнее «прощай».

Свыше 900 тыс. человек прошло за эти дни около гроба Ленина.

... 26 января открылся II Всесоюзный съезд Советов. Первое заседание было посвящено памяти Владимира Ильича. Съезд принял постановление о переименовании Петрограда в Ленинград.

... 27 января, в 9 часов 20 минут, гроб с телом Ленина был перенесен из Дома Союзов на Красную площадь и установлен на специальном возвышении.

И вновь сотни тысяч людей в торжественном молчании, колонна за колонной прошли, прощаясь с великим Лениным...

В четыре часа под залпы орудий гроб был внесен в Мавзолей.

В этот момент по всей стране на пять минут прекратилась работа на предприятиях, остановились поезда и пароходы, остановились автомобили, прервались занятия в школах и высших учебных заведениях... Гудели гудки фабрик, заводов, паровозов... Народ прощался с Лениным...

\subsection*{\center * * *}

21 января 1924 года, в 6 часов 50 минут, перестало биться сердце Владимира Ильича Ленина.

Но Ленин продолжает жить и в наши дни, он будет жить в веках. Как сказал Маяковский:

\quad \quad \quad \quad \quad \quad \quad \quad Ленин

\quad \quad \quad \quad \quad \quad \quad \quad \quad \quad \quad и теперь

\quad \quad \quad \quad \quad \quad \quad \quad \quad \quad \quad \quad \quad \quad живее всех живых.

\quad \quad \quad \quad \quad \quad \quad \quad Наше знанье —

\quad \quad \quad \quad \quad \quad \quad \quad \quad \quad \quad \quad \quad \quad сила

\quad \quad \quad \quad \quad \quad \quad \quad \quad \quad \quad \quad \quad \quad \quad \quad и оружие.

Своими великими теоретическими открытиями он обогатил революционное учение Маркса — Энгельса.

Достоянием мирового пролетариата стал опыт ленинской революционной теории и практики:

создания Коммунистической партии и ее руководства боями пролетариата в разных исторических условиях;

ленинской стратегии и тактики, приведших к победе Октябрьской социалистической революции;

создания и руководства строительством первого социалистического государства в мире.

Образовав III Интернационал, Ленин обеспечил развитие и укрепление коммунистических партий во многих странах, что способствовало дальнейшему сплочению мирового коммунистического движения.

Ленинские идеи, победа Великой Октябрьской социалистической революции расшатали основы мировой капиталистической системы, пробудили к жизни национально-освободительное движение угнетенных народов, соединившееся с борьбой пролетариата в один революционный поток.

Воплощением в жизнь ленинских идей является возникновение и развитие мировой системы социализма.

Народы нашей страны под руководством ленинской партии успешно построили социалистическое общество, создают экономическую и культурную базу коммунизма, строят коммунизм.

Оглядываясь на прошлое, все больше и больше убеждаешься, какую великую, неодолимую историческую силу представляет бессмертное учение Ленина.

Особенно полно чувствуешь и понимаешь это, когда читаешь тезисы ЦК КПСС «К 100-летию со дня рождения Владимира Ильича Ленина».

Широко и глубоко раскрыта в них суть марксизма-ленинизма, его связь с современностью, великие задачи, которые встают перед нами в связи с знаменательной для всего человечества ленинской датой.






	
	\newpage
	\tableofcontents
	
	\thispagestyle{empty} % 
	
	\newpage
	
	\setcounter{secnumdepth}{0}
	
	\phantomsection
	
		\section*{Описание}
	
	{\bf Название:} Мысли и воспоминания о Ленине
	
{\bf Автор:} Микоян Анастас Иванович

{\bf Редактор:} \textit{Г. П. Шкаренкова}

{\bf Художник:} \textit{В. В. Сорокин}

{\bf Художественный редактор:} \textit{Г. Ф. Семиреченко}

{\bf Технический редактор:} \textit{Ю. А. Мухин}
	
{\bf Издательство:} М: Политиздат, 1970.
	
		{\bf Аннотация:} Имя А. И. Микояна хорошо известно в нашей стране и за рубежом. В настоящее время он — член Президиума Верховного Совета СССР, член ЦК КПСС.
		
		А. И. Микоян — один из тех, на чью долю выпало встречаться с В. И. Лениным, работать под его руководством.
		
		Воспоминания Микояна охватывают 1919–1924 годы. В них рассказывается о встречах и беседах с Владимиром Ильичем, о выступлениях Ленина на советских, партийных и других съездах, конференциях, совещаниях, участником которых был автор. Большое место в воспоминаниях отведено раскрытию роли Ленина в разрешении национального вопроса, в выработке принципов нэпа, в борьбе за единство партии.
		
		


		\thispagestyle{empty} % выключаем отображение номера для этой страницы


\newpage

\setcounter{secnumdepth}{0}  

\phantomsection	

\ \\
\ \\
\ \\
\ \\

\hangindent=2cm \hangafter=0  	\textbf{ИНФОРМАЦИЯ ОБ ОЦИФРОВКЕ}

\hangindent=2cm \hangafter=0  	\textbf{Источник}:  \href{https://rutracker.org/forum/viewtopic.php?t=5855912}{https://rutracker.org/forum/viewtopic.php?t=5855912} 

\hangindent=2cm \hangafter=0  	\textbf{Версия}: 1.0 - от 02 октября 2026

\hangindent=2cm \hangafter=0  	\textbf{Оцифровка}: lord199 (\href{mailto:lord199@mail.ru}{lord199@mail.ru})



\hangindent=2cm \hangafter=0  Если вы нашли какие-то опечатки и другие неточности, просьба связаться по электронному адресу.




\thispagestyle{empty} % выключаем отображение номера для этой страницы
	
\end{document}


